\documentclass[12pt]{article}
\def\unitnumber{1}
\def\pdftitle{Linear Algebra: Lecture Notes before the Midterm}
\def\runningtitle{Lecture notes}
\usepackage[T1]{fontenc}
\usepackage{lmodern,microtype}
\usepackage[a4paper,margin=15mm,headheight=14pt,headsep=5mm]{geometry}
\usepackage{amsmath,amssymb,amsthm,mathtools}
\usepackage{enumitem,booktabs,needspace}
\usepackage{pifont}
\usepackage{tikz}
\usetikzlibrary{arrows.meta}
\usepackage{fancyhdr}
\usepackage{xr-hyper}
\usepackage[hidelinks]{hyperref}
\makeatletter
\newcommand{\notetarget}[1]{\Hy@raisedlink{\hypertarget{#1}{}}}
\makeatother
\hypersetup{pdftitle={\pdftitle},pdfauthor={Chang-Ye Tu}}
\pagestyle{fancy}
\fancyhf{}
\fancyhead[L]{\small Linear Algebra}
\fancyhead[R]{\small\runningtitle}
\fancyfoot[C]{\small\thepage}
\renewcommand{\headrulewidth}{0.3pt}
\setlength{\parindent}{0pt}
\setlength{\parskip}{4pt plus 1pt minus 1pt}
\setlength{\emergencystretch}{1.5em}
\setlength{\abovedisplayskip}{7pt plus 2pt minus 2pt}
\setlength{\belowdisplayskip}{7pt plus 2pt minus 2pt}
\setlist[enumerate]{label=(\alph*),leftmargin=*,itemsep=2pt,topsep=3pt}
\setlist[itemize]{leftmargin=*,itemsep=2pt,topsep=3pt}
\newcommand{\R}{\mathbb R}
\newcommand{\contradiction}{\mathord{\text{\large\ding{56}}}}
\newcommand{\norm}[1]{\left\lVert#1\right\rVert}
\newcommand{\ip}[2]{\left\langle#1,#2\right\rangle}
\newcommand{\one}{\mathbf 1}
\DeclareMathOperator{\tr}{tr}
\DeclareMathOperator{\row}{row}
\DeclareMathOperator{\col}{col}
\DeclareMathOperator{\diag}{diag}
\DeclareMathOperator{\rank}{rank}
\DeclareMathOperator{\Span}{span}
\DeclareMathOperator{\nullity}{nullity}
\DeclareMathOperator{\im}{im}
\DeclareMathOperator{\id}{id}
\theoremstyle{definition}
\newtheorem{theorem}{Theorem}
\renewcommand{\thetheorem}{\unitnumber.\arabic{theorem}}
\newtheorem{proposition}[theorem]{Proposition}
\newtheorem{lemma}[theorem]{Lemma}
\newtheorem{corollary}[theorem]{Corollary}
\newtheorem{definition}[theorem]{Definition}
\newtheorem{example}[theorem]{Example}
\newtheorem{remark}[theorem]{Remark}
\newtheorem{problem}{Exercise}
\renewcommand{\theproblem}{\unitnumber.\arabic{problem}}
\newenvironment{solution}{\par\noindent\textit{Solution.}\ }{\par\smallskip}
\AddToHook{env/problem/before}{\needspace{8\baselineskip}}
\AddToHook{env/problem/after}{\nopagebreak[4]}
\def\theHtheorem{\unitnumber.\arabic{theorem}}
\def\theHproblem{\unitnumber.\arabic{problem}}
\begin{document}
{\LARGE\bfseries Linear Algebra\par}
\medskip
{\Large Lecture Notes before the Midterm\par}
\bigskip
\clearpage
\def\unitnumber{1}
\setcounter{theorem}{0}\setcounter{problem}{0}
\section*{Chapter 1. Matrices}

\subsection*{Matrices and products}

\begin{definition}\label{r1-def-matrix}
Let \(m,n\) be positive integers. An \(m\times n\) \textbf{matrix} \(A=(a_{ij})\) is a
rectangular array of real numbers, its \textbf{entries}; \(a_{ij}\) stands in row \(i\) and
column \(j\) (\(1\leqslant i\leqslant m\), \(1\leqslant j\leqslant n\)). The set of all
\(m\times n\) matrices is \(\R^{m\times n}\).
\begin{itemize}
\item \(A=B\) means that \(A\) and \(B\) have the same size and \(a_{ij}=b_{ij}\) for all
\(i,j\).
\item \(\R^n=\R^{n\times1}\) is the set of \textbf{column vectors} \(x=(x_i)\), and
\(\R^{1\times n}\) the set of \textbf{row vectors}. A \(1\times1\) matrix is identified with
its entry.
\item \(\col_j(A)\in\R^m\) is the \(j\)th column and
\(\row_i(A)=(a_{i1},\dots,a_{in})\in\R^{1\times n}\) the \(i\)th row of \(A\). We write
\(a_j=\col_j(A)\), so \(A=[a_1\ \cdots\ a_n]\).
\item Sums and scalar multiples are taken entrywise:
\((A+B)_{ij}=a_{ij}+b_{ij}\), \((cA)_{ij}=ca_{ij}\) (\(c\in\R\)); \(-A=(-1)A\) and
\(A-B=A+(-B)\). They obey the laws of numbers, such as \(A+B=B+A\) and
\(c(A+B)=cA+cB\), because these hold entry by entry.
\item The \textbf{zero matrix} \(0\) has all entries \(0\). The \textbf{identity matrix} is
\(I_n=(\delta_{ij})\in\R^{n\times n}\), where \(\delta_{ij}=1\) if \(i=j\) and \(\delta_{ij}=0\)
if \(i\ne j\); its columns \(e_1,\dots,e_n\) are the \textbf{unit vectors} of \(\R^n\).
\end{itemize}
\end{definition}

In statistics, measurements of \(k\) variables on \(n\) units are stored in an \(n\times k\)
\textbf{data matrix} \((x_{ij})\): row \(i\) holds the measurements on unit \(i\), column \(j\)
holds the \(n\) values of variable \(j\), and \(x_{ij}\) is the value of variable \(j\) on unit
\(i\).

\begin{definition}\label{r1-def-product}
Let \(A\in\R^{m\times n}\) and \(B\in\R^{n\times p}\). The \textbf{product}
\(AB\in\R^{m\times p}\) is defined by
\[
(AB)_{ij}=\sum_{r=1}^n a_{ir}b_{rj}=a_{i1}b_{1j}+a_{i2}b_{2j}+\dots+a_{in}b_{nj}
\qquad(1\leqslant i\leqslant m,\ 1\leqslant j\leqslant p)
\]
It is defined only when \(A\) has as many columns as \(B\) has rows.
\end{definition}

\begin{theorem}\label{r1-thm-rules}
Whenever the sizes allow the operations,
\[
A(BC)=(AB)C,\quad A(B+C)=AB+AC,\quad (A+B)C=AC+BC,\quad A(cB)=c(AB)=(cA)B,
\]
\[
I_mA=A=AI_n,\qquad 0A=0,\qquad A0=0\qquad(A\in\R^{m\times n})
\]
\end{theorem}
Each rule is checked entry by entry; for instance
\([A(BC)]_{ij}=\sum_r a_{ir}\sum_s b_{rs}c_{sj}=\sum_s\bigl(\sum_r a_{ir}b_{rs}\bigr)c_{sj}=[(AB)C]_{ij}\).
So a product of several matrices needs no parentheses. For \(A\in\R^{n\times n}\) we put
\(A^0=I_n\) and \(A^{k+1}=A^kA\) (\(k\geqslant0\)).

\begin{theorem}\label{r1-thm-views}
Let \(A\in\R^{m\times n}\) have columns \(a_1,\dots,a_n\) and \(B\in\R^{n\times p}\) columns
\(b_1,\dots,b_p\).
\begin{enumerate}
\item \(AB=[Ab_1\ \cdots\ Ab_p]\), that is, \(\col_j(AB)=Ab_j\).
\item For \(x\in\R^n\), \(Ax=x_1a_1+x_2a_2+\dots+x_na_n\). In particular \(Ae_j=a_j\).
\item \(\row_i(AB)=\row_i(A)B=a_{i1}\row_1(B)+\dots+a_{in}\row_n(B)\).
\item \((AB)_{ij}=\row_i(A)\,b_j\), and \(AB=a_1\row_1(B)+\dots+a_n\row_n(B)\), a sum of \(n\)
products of a column and a row (the \textbf{outer product expansion}).
\end{enumerate}
\end{theorem}
\begin{proof}
In each statement the entry in position \((i,j)\) (or \(i\), or \(j\)) of both sides is
\(\sum_r a_{ir}b_{rj}\); in (b) the \(i\)th entry of \(Ax\) is \(\sum_j a_{ij}x_j\), which is
the \(i\)th entry of \(\sum_jx_ja_j\).
\end{proof}

\begin{example}\label{r1-exm-products}
\begin{enumerate}
\item \(\begin{pmatrix}1&0\\-3&1\end{pmatrix}\begin{pmatrix}1&-1&2\\3&4&8\end{pmatrix}=\begin{pmatrix}1&-1&2\\0&7&2\end{pmatrix}\).
By Theorem~\ref{r1-thm-views}(c), its second row is \(-3\,(1,-1,2)+1\,(3,4,8)=(0,7,2)\).
\item \(\begin{pmatrix}2&1\\5&4\end{pmatrix}\begin{pmatrix}2\\-1\end{pmatrix}
=2\begin{pmatrix}2\\5\end{pmatrix}-1\begin{pmatrix}1\\4\end{pmatrix}=\begin{pmatrix}3\\6\end{pmatrix}\),
by Theorem~\ref{r1-thm-views}(b).
\item For \(N=\begin{pmatrix}0&1&0\\0&0&0\\0&0&0\end{pmatrix}\) and
\(M=\begin{pmatrix}1&-5&2\\2&3&4\\9&-1&3\end{pmatrix}\),
\[
NM=\begin{pmatrix}2&3&4\\0&0&0\\0&0&0\end{pmatrix},\qquad
MN=\begin{pmatrix}0&1&0\\0&2&0\\0&9&0\end{pmatrix}
\]
Row \(1\) of \(NM\) is row \(2\) of \(M\), and column \(2\) of \(MN\) is column \(1\) of \(M\).
So \(NM\ne MN\), although both are \(3\times3\).
\end{enumerate}
\end{example}

\begin{remark}\label{r1-rem-noncomm}
Matrix products differ from products of numbers. In general \(AB\ne BA\)
(Example~\ref{r1-exm-products}(c)); square matrices with \(AB=BA\) are said to
\textbf{commute}. A product can be \(0\) although no factor is \(0\): for
\[
X=\begin{pmatrix}1&1&3\\2&2&6\\-1&-1&-3\end{pmatrix},\qquad X^2=0,
\]
since \((X^2)_{ij}=\row_i(X)\col_j(X)\) (Theorem~\ref{r1-thm-views}(d)), rows \(1,2,3\) of \(X\) are
\(1,2,-1\) times \(\row_1(X)\), columns \(1,2,3\) are
\(1,1,3\) times \(\col_1(X)\), and \(\row_1(X)\col_1(X)=1+2-3=0\). Hence \(AB=0\) does not force
\(A=0\) or \(B=0\), and \(AB=AC\) does not force \(B=C\).
\end{remark}

\subsection*{Transpose and trace}

\begin{definition}\label{r1-def-transpose}
The \textbf{transpose} of \(A\in\R^{m\times n}\) is \(A'\in\R^{n\times m}\) with
\((A')_{ij}=a_{ji}\): the rows of \(A'\) are the transposed columns of \(A\). We write
column vectors as \(x=(x_1,\dots,x_n)'\). A square \(A\) is \textbf{symmetric} if
\(A'=A\).
\end{definition}

\begin{theorem}\label{r1-thm-transpose}
For matrices of sizes for which the operations are defined and \(c\in\R\),
\[
(A')'=A,\quad(A+B)'=A'+B',\quad(cA)'=cA',\quad
(AB)'=B'A',\quad(ABC)'=C'B'A'
\]
For \(x,y\in\R^n\), \(x'y=y'x=\sum_{i=1}^nx_iy_i\); in particular
\(x'x=\sum_ix_i^2\). For every \(B\in\R^{m\times n}\), \(B'B\) and \(BB'\) are
symmetric.
\end{theorem}
\begin{proof}
The first three rules hold entrywise. For \(A\in\R^{m\times n}\), \(B\in\R^{n\times p}\), the
\((k,i)\) entry of \(B'A'\) is \(\sum_j b_{jk}a_{ij}=(AB)_{ik}=((AB)')_{ki}\);
applying this twice gives \((ABC)'\). The \(1\times1\) matrix \(x'y=\sum_ix_iy_i\)
equals its transpose \(y'x\). Finally \((B'B)'=B'(B')'=B'B\), and
likewise for \(BB'\).
\end{proof}

\begin{definition}\label{r1-def-trace}
The \textbf{trace} of \(A\in\R^{n\times n}\) is \(\tr(A)=a_{11}+a_{22}+\dots+a_{nn}\).
\end{definition}

\begin{theorem}\label{r1-thm-trace}
\begin{enumerate}
\item \(\tr(A+B)=\tr(A)+\tr(B)\), \(\tr(cA)=c\tr(A)\) and \(\tr(A')=\tr(A)\) for
\(A,B\in\R^{n\times n}\).
\item \(\tr(AB)=\tr(BA)\) for \(A\in\R^{m\times n}\) and \(B\in\R^{n\times m}\).
\item \(\tr(ABC)=\tr(CAB)=\tr(BCA)\) for \(A\in\R^{m\times n}\), \(B\in\R^{n\times p}\),
\(C\in\R^{p\times m}\).
\item \(\tr(A'A)=\sum_{i=1}^m\sum_{j=1}^na_{ij}^2\) for \(A\in\R^{m\times n}\); hence
\(\tr(A'A)=0\) only if \(A=0\).
\end{enumerate}
\end{theorem}
\begin{proof}
(a) holds entrywise on the diagonal. (b) Both traces equal
\(\sum_{i=1}^m\sum_{j=1}^na_{ij}b_{ji}\). (c) Apply (b) to the pairs \(AB,C\) and \(A,BC\).
(d) The \(j\)th diagonal entry of \(A'A\) is \(\sum_ia_{ij}^2\); a sum of squares of real
numbers is \(0\) only if every term is \(0\).
\end{proof}

\subsection*{Special matrices; means and sums of squares}

\textbf{Diagonal and triangular matrices.} \(D=\diag(d_1,\dots,d_n)\) is the \(n\times n\)
matrix with diagonal entries \(d_1,\dots,d_n\) and all other entries \(0\). Left
multiplication by a diagonal matrix scales the rows, right multiplication scales the columns:
\[
\row_i(DA)=d_i\row_i(A),\qquad \col_j(AD)=d_j\col_j(A),
\]
since \((DA)_{ij}=d_ia_{ij}\) and \((AD)_{ij}=a_{ij}d_j\). In particular
\(\diag(d_1,\dots,d_n)\diag(f_1,\dots,f_n)=\diag(d_1f_1,\dots,d_nf_n)\). A square \(A\) is \textbf{upper
triangular} if \(a_{ij}=0\) for \(i>j\) and \textbf{lower triangular} if \(a_{ij}=0\) for
\(i<j\). A product of two upper (lower) triangular matrices is upper (lower) triangular, with
diagonal entries \(a_{ii}b_{ii}\): for upper triangular \(A,B\), each term \(a_{ik}b_{kj}\) of
\((AB)_{ij}\) has \(a_{ik}=0\) if \(k<i\) and \(b_{kj}=0\) if \(k>j\), so \((AB)_{ij}=0\) for
\(i>j\) and \((AB)_{ii}=a_{ii}b_{ii}\) (lower triangular: reverse the inequalities).

\begin{definition}\label{r1-def-ones}
The \textbf{summing vector} \(\one_n\in\R^n\) has all entries \(1\). Put
\[
J_n=\one_n\one_n'\in\R^{n\times n},\qquad \bar J_n=\tfrac1nJ_n,\qquad C_n=I_n-\bar J_n
\]
Every entry of \(J_n\) is \(1\), and \(C_n\) is the \textbf{centering matrix}. Subscripts are
dropped when the size is clear. A square matrix \(A\) with \(A^2=A\) is \textbf{idempotent}.
\end{definition}

\begin{proposition}\label{r1-prop-centering}
\(\one'\one=n\), \(J_n^2=nJ_n\), and \(\bar J_n\) and \(C=C_n\) are symmetric and
idempotent:
\[
\bar J_n'=\bar J_n=\bar J_n^2,\qquad C=C'=C^2,\qquad C\one=0
\]
\end{proposition}
\begin{proof}
\(\one'\one=1+\dots+1=n\), so \(J_n^2=\one(\one'\one)\one'=nJ_n\) and
\(\bar J_n^2=\frac1{n^2}nJ_n=\bar J_n\). Next \(J_n'=(\one\one')'=\one\one'=J_n\)
gives \(\bar J_n'=\bar J_n\) and \(C'=C\), and \(C^2=I-2\bar J_n+\bar J_n^2=I-\bar J_n=C\).
Finally \(C\one=\one-\frac1n\one(\one'\one)=\one-\one=0\).
\end{proof}

\begin{example}\label{r1-exm-centering}
For data \(x=(x_1,\dots,x_n)'\) the mean is \(\bar x=\frac1n\one'x\). Then
\[
Cx=x-\tfrac1n\one(\one'x)=x-\bar x\one=(x_1-\bar x,\ \dots,\ x_n-\bar x)',
\]
the vector of \textbf{deviations} from the mean. Moreover
\(x'Cx=x'x-\bar x(\one'x)=\sum_ix_i^2-n\bar x^2\), and, since
\(\sum_ix_i=n\bar x\),
\(\sum_i(x_i-\bar x)^2=\sum_ix_i^2-2\bar x\sum_ix_i+n\bar x^2=\sum_ix_i^2-n\bar x^2\). Hence the
sum of squares about the mean is
\[
\sum_{i=1}^n(x_i-\bar x)^2=x'x-n\bar x^2=x'Cx
\]
\end{example}

\begin{example}\label{r1-exm-data}
Four units (\(n=4\)) are measured on two variables (\(k=2\)); \(x_{ij}\) is the value of
variable \(j\) on unit \(i\):
\[
X_0=(x_{ij})=\begin{pmatrix}2&3\\0&3\\2&1\\0&1\end{pmatrix}\in\R^{4\times2},\qquad
\one'X_0=(4,\ 8),\qquad
\bar x=\tfrac14X_0'\one=\begin{pmatrix}1\\2\end{pmatrix}
\]
The vector \(\bar x=(\bar x_1,\bar x_2)'\) of column means \(\bar x_j=\frac1n\sum_ix_{ij}\) is
the \textbf{mean vector}. By
Example~\ref{r1-exm-centering}, applied to each column
(\(\col_j(CX_0)=C\col_j(X_0)\), Theorem~\ref{r1-thm-views}(a)), the centered data matrix and the
matrix of \textbf{corrected sums of squares and products} are
\[
CX_0=X_0-\one\bar x'=\begin{pmatrix}1&1\\-1&1\\1&-1\\-1&-1\end{pmatrix},\qquad
X_0'CX_0=(CX_0)'(CX_0)=\begin{pmatrix}4&0\\0&4\end{pmatrix},
\]
using \(C=C'C\) (Proposition~\ref{r1-prop-centering}). For instance \(4=1+1+1+1\) and
\(0=1-1-1+1\). The same matrix comes from the uncorrected sums: in general
\(X_0'CX_0=X_0'X_0-\frac1n(X_0'\one)(\one'X_0)=X_0'X_0-n\bar x\bar x'\),
and here
\[
X_0'X_0-4\bar x\bar x'=\begin{pmatrix}8&8\\8&20\end{pmatrix}-\begin{pmatrix}4&8\\8&16\end{pmatrix}
=\begin{pmatrix}4&0\\0&4\end{pmatrix}
\]
In regression one also uses the \textbf{design matrix} \(X=[\one_4\ X_0]\in\R^{4\times3}\), of
size \(n\times(k+1)\): its first column is \(\one_4\), and its \(i\)th row is
\(x_i'=(1,x_{i1},x_{i2})\). Write \(X_{\cdot1}\) and \(X_{\cdot2}\) for its other two columns,
the columns of \(X_0\). By Theorems~\ref{r1-thm-views}(d) and~\ref{r1-thm-transpose}, the
\((j,l)\) entry of \(X'X\) is \(\col_j(X)'\col_l(X)\), so
\[
X'X=\begin{pmatrix}\one'\one&\one'X_{\cdot1}&\one'X_{\cdot2}\\X_{\cdot1}'\one&X_{\cdot1}'X_{\cdot1}&X_{\cdot1}'X_{\cdot2}\\X_{\cdot2}'\one&X_{\cdot2}'X_{\cdot1}&X_{\cdot2}'X_{\cdot2}\end{pmatrix}
=\begin{pmatrix}4&4&8\\4&8&8\\8&8&20\end{pmatrix}:
\]
the entries are \(n\), the column sums \(\sum_ix_{i1}\), \(\sum_ix_{i2}\), and the uncorrected
sums of squares and products \(\sum_ix_{i1}^2\), \(\sum_ix_{i1}x_{i2}\), \(\sum_ix_{i2}^2\).
\end{example}

\subsection*{Exercises and solutions}

\begin{problem}\label{r1-ex-products}
Let \(A=\begin{pmatrix}1&2&1\\0&1&2\end{pmatrix}\) and
\(B=\begin{pmatrix}0&1&2\\1&2&1\end{pmatrix}\). Compute \(AB'\), \(BA'\), \(A'B\) and \(B'A\).
\end{problem}
\begin{solution}
\[
AB'=\begin{pmatrix}1&2&1\\0&1&2\end{pmatrix}\begin{pmatrix}0&1\\1&2\\2&1\end{pmatrix}
=\begin{pmatrix}4&6\\5&4\end{pmatrix},\qquad
BA'=\begin{pmatrix}0&1&2\\1&2&1\end{pmatrix}\begin{pmatrix}1&0\\2&1\\1&2\end{pmatrix}
=\begin{pmatrix}4&5\\6&4\end{pmatrix},
\]
\[
A'B=\begin{pmatrix}1&0\\2&1\\1&2\end{pmatrix}\begin{pmatrix}0&1&2\\1&2&1\end{pmatrix}
=\begin{pmatrix}0&1&2\\1&4&5\\2&5&4\end{pmatrix},\qquad
B'A=\begin{pmatrix}0&1\\1&2\\2&1\end{pmatrix}\begin{pmatrix}1&2&1\\0&1&2\end{pmatrix}
=\begin{pmatrix}0&1&2\\1&4&5\\2&5&4\end{pmatrix}
\]
For instance, \((AB')_{12}=1\cdot1+2\cdot2+1\cdot1=6\). Here \(A'B\) happens to be
symmetric while \(AB'\) is not. The four products have the same trace \(8\); this is no
coincidence: \(\tr(AB')=\tr(B'A)\) and \(\tr(BA')=\tr(A'B)\) by
Theorem~\ref{r1-thm-trace}(b), and \(\tr(BA')=\tr\bigl((AB')'\bigr)=\tr(AB')\) by
Theorems~\ref{r1-thm-transpose} and~\ref{r1-thm-trace}(a).
\end{solution}

\begin{problem}\label{r1-ex-products-drill}
Let \(A=\begin{pmatrix}1&0&2\\-1&3&1\end{pmatrix}\) and
\(B=\begin{pmatrix}2&1\\0&-1\\1&1\end{pmatrix}\).
\begin{enumerate}
\item Compute \(AB\) and \(BA\).
\item Write \(AB\) as the sum of the three products \(\col_j(A)\row_j(B)\) of
Theorem~\ref{r1-thm-views}(d).
\item Compute \(B'A'\), \(\tr(AB)\), \(\tr(BA)\) and \(\tr(A'A)\).
\end{enumerate}
\end{problem}
\begin{solution}
(a) By Theorem~\ref{r1-thm-views}(c), \(\row_i(AB)=\row_i(A)B\) is a combination of the rows of
\(B\):
\[
\begin{aligned}
\row_1(AB)&=1\,(2,1)+0\,(0,-1)+2\,(1,1)=(4,\ 3),\\
\row_2(AB)&=-1\,(2,1)+3\,(0,-1)+1\,(1,1)=(-1,\ -3)
\end{aligned}
\]
Likewise each row of \(BA\in\R^{3\times3}\) is a combination of the rows \((1,0,2)\) and
\((-1,3,1)\) of \(A\):
\[
\begin{aligned}
\row_1(BA)&=2\,(1,0,2)+1\,(-1,3,1)=(1,\ 3,\ 5),\\
\row_2(BA)&=0\,(1,0,2)-1\,(-1,3,1)=(1,\ -3,\ -1),\\
\row_3(BA)&=1\,(1,0,2)+1\,(-1,3,1)=(0,\ 3,\ 3)
\end{aligned}
\]
Hence
\[
AB=\begin{pmatrix}4&3\\-1&-3\end{pmatrix},\qquad
BA=\begin{pmatrix}1&3&5\\1&-3&-1\\0&3&3\end{pmatrix}
\]
The two products do not even have the same size.

(b)
\[
\begin{pmatrix}1\\-1\end{pmatrix}(2,\ 1)+\begin{pmatrix}0\\3\end{pmatrix}(0,\ -1)
+\begin{pmatrix}2\\1\end{pmatrix}(1,\ 1)
=\begin{pmatrix}2&1\\-2&-1\end{pmatrix}+\begin{pmatrix}0&0\\0&-3\end{pmatrix}
+\begin{pmatrix}2&2\\1&1\end{pmatrix}=\begin{pmatrix}4&3\\-1&-3\end{pmatrix}=AB
\]

(c) Here \(B'=\left(\begin{smallmatrix}2&0&1\\1&-1&1\end{smallmatrix}\right)\) and
\(A'=\left(\begin{smallmatrix}1&-1\\0&3\\2&1\end{smallmatrix}\right)\), so
\[
B'A'=\begin{pmatrix}2\cdot1+0\cdot0+1\cdot2&2\cdot(-1)+0\cdot3+1\cdot1\\
1\cdot1+(-1)\cdot0+1\cdot2&1\cdot(-1)+(-1)\cdot3+1\cdot1\end{pmatrix}
=\begin{pmatrix}4&-1\\3&-3\end{pmatrix}=(AB)',
\]
as Theorem~\ref{r1-thm-transpose} says. Next \(\tr(AB)=4-3=1\) and \(\tr(BA)=1-3+3=1\), equal by
Theorem~\ref{r1-thm-trace}(b); and by Theorem~\ref{r1-thm-trace}(d),
\(\tr(A'A)=1^2+0^2+2^2+(-1)^2+3^2+1^2=16\).
\end{solution}

\begin{problem}\label{r1-ex-unlike-scalars}
Give examples of real \(2\times2\) matrices such that:
(a) \(A^2=-I\);
(b) \(B^2=0\) (\(B\ne0\));
(c) \(CD=-DC\) (\(CD\ne0\));
(d) \(EF=0\) (no entries of \(E\) or \(F\) are zero).
\end{problem}
\begin{solution}
(a) No real number \(a\) has \(a^2=-1\), but
\(A=\left(\begin{smallmatrix}0&1\\-1&0\end{smallmatrix}\right)\) gives
\(A^2=\left(\begin{smallmatrix}0\cdot0+1\cdot(-1)&0\cdot1+1\cdot0\\(-1)\cdot0+0\cdot(-1)&(-1)\cdot1+0\cdot0\end{smallmatrix}\right)=-I\).

(b) \(B=\left(\begin{smallmatrix}0&1\\0&0\end{smallmatrix}\right)\ne0\) satisfies \(B^2=0\).

(c) For \(C=\left(\begin{smallmatrix}1&1\\1&-1\end{smallmatrix}\right)\) and
\(D=\left(\begin{smallmatrix}1&-1\\-1&-1\end{smallmatrix}\right)\),
\(CD=\left(\begin{smallmatrix}0&-2\\2&0\end{smallmatrix}\right)\) and
\(DC=\left(\begin{smallmatrix}0&2\\-2&0\end{smallmatrix}\right)=-CD\ne0\).

(d) For \(E=\left(\begin{smallmatrix}1&1\\1&1\end{smallmatrix}\right)\) and
\(F=\left(\begin{smallmatrix}1&1\\-1&-1\end{smallmatrix}\right)\), every entry of \(EF\) is
\(1\cdot1+1\cdot(-1)=0\), so \(EF=0\).
\end{solution}

\begin{problem}\label{r1-ex-difference-squares}
Let \(A,B\in\R^{n\times n}\). Show that \((A+B)(A-B)=A^2-B^2\) if and only if \(A\) and \(B\)
commute.
\end{problem}
\begin{solution}
By Theorem~\ref{r1-thm-rules}, \((A+B)(A-B)=A(A-B)+B(A-B)=A^2-AB+BA-B^2\). Thus
\((A+B)(A-B)=A^2-B^2\) if and only if \(-AB+BA=0\), that is, if and only if \(AB=BA\).
\end{solution}

\begin{problem}\label{r1-ex-commuting}
Let \(A=\begin{pmatrix}1&1\\3&4\end{pmatrix}\). Show that the matrices \(B\) with \(AB=BA\) are
exactly
\(B=\alpha\begin{pmatrix}1&0\\0&1\end{pmatrix}+\beta\begin{pmatrix}0&1\\3&3\end{pmatrix}\)
(\(\alpha,\beta\in\R\)).
\end{problem}
\begin{solution}
Let \(B=\left(\begin{smallmatrix}a&b\\c&d\end{smallmatrix}\right)\). Then \(AB=BA\) reads
\[
\begin{pmatrix}a+c&b+d\\3a+4c&3b+4d\end{pmatrix}=\begin{pmatrix}a+3b&a+4b\\c+3d&c+4d\end{pmatrix}
\]
Comparing entries, the \((1,1)\) and \((2,2)\) entries give the same equation, and the
conditions are
\[
3b-c=0,\qquad a+3b-d=0,\qquad a+c-d=0
\]
(the \((2,1)\) entries give \(3a+3c-3d=0\)). The first gives \(c=3b\), and then the third gives
\(d=a+3b\); with these values the second holds (\(a+3b-a-3b=0\)). So \(a\) and \(b\) are free,
and the solutions are
\[
B=\begin{pmatrix}a&b\\3b&a+3b\end{pmatrix}=a\begin{pmatrix}1&0\\0&1\end{pmatrix}+b\begin{pmatrix}0&1\\3&3\end{pmatrix},
\]
the stated form with \(\alpha=a\), \(\beta=b\). Conversely, every \(\alpha,\beta\) arise (take
\(a=\alpha\), \(b=\beta\)), so all these matrices commute with \(A\).
\end{solution}

\begin{problem}\label{r1-ex-symmetric-product}
\begin{enumerate}
\item Show that the product \(AB\) of two symmetric matrices \(A,B\in\R^{n\times n}\) is
symmetric if and only if \(A\) and \(B\) commute.
\item Give an example of two symmetric matrices of the same size whose product is not
symmetric.
\end{enumerate}
\end{problem}
\begin{solution}
(a) Since \(A\) and \(B\) are symmetric, \((AB)'=B'A'=BA\)
(Theorem~\ref{r1-thm-transpose}). So \(AB=(AB)'\) holds if and only if \(AB=BA\).

(b) Take \(A=\left(\begin{smallmatrix}1&0\\0&2\end{smallmatrix}\right)\) and
\(B=\left(\begin{smallmatrix}0&1\\1&0\end{smallmatrix}\right)\). Then
\(AB=\left(\begin{smallmatrix}0&1\\2&0\end{smallmatrix}\right)\ne\left(\begin{smallmatrix}0&2\\1&0\end{smallmatrix}\right)=BA\),
so by (a) \(AB\) is not symmetric.
\end{solution}

\begin{problem}\label{r1-ex-trace-transposes}
Show that for any \(A\in\R^{m\times n}\), \(B\in\R^{n\times p}\) and \(C\in\R^{p\times m}\),
\(\tr(ABC)=\tr(B'A'C')=\tr(A'C'B')\).
\end{problem}
\begin{solution}
By Theorem~\ref{r1-thm-trace}(c), then (a), then Theorem~\ref{r1-thm-transpose}, and finally
(c) again,
\[
\tr(ABC)=\tr(CAB)=\tr\bigl((CAB)'\bigr)=\tr(B'A'C')=\tr(A'C'B')
\]
In the last step the factors \(B'\in\R^{p\times n}\), \(A'\in\R^{n\times m}\),
\(C'\in\R^{m\times p}\) are moved cyclically.
\end{solution}

\begin{problem}\label{r1-ex-sum-entries}
Let \(A\in\R^{n\times n}\) and let \(D=\diag(a_{11},\dots,a_{nn})\). Show that the sum of all
entries of \(A\) is
\[
\one_n'A\one_n=\one_n'D\one_n+\tr\bigl((\one_n\one_n'-I_n)A\bigr)
\]
\end{problem}
\begin{solution}
By Theorem~\ref{r1-thm-trace}(a),(b),
\[
\tr\bigl((\one\one'-I_n)A\bigr)=\tr(\one\one'A)-\tr(A)=\tr(\one'A\one)-\tr(A)
=\one'A\one-\one'D\one
\]
Here \(\tr(\one\one'A)=\tr(\one'A\one)\) moves the \(n\times1\) factor \(\one\) past
the \(1\times n\) factor \(\one'A\); the trace of the \(1\times1\) matrix
\(\one'A\one\) is that number; and \(\one'D\one=\sum_ia_{ii}=\tr(A)\), because
\(D\one=(a_{11},\dots,a_{nn})'\). Rearranging gives the identity. Also
\(\one'A\one=\sum_i\sum_ja_{ij}\) by Theorem~\ref{r1-thm-views}(b),(c).
\end{solution}

\begin{problem}\label{r1-ex-idempotent}
\begin{enumerate}
\item Show that the only symmetric idempotent \(2\times2\) matrices are \(I_2\), \(0\), and
\(uu'\) with \(u\in\R^2\), \(u'u=1\).
\item Give an example of an \(n\times n\) idempotent matrix (\(n\geqslant2\)) that is not
symmetric.
\end{enumerate}
\end{problem}
\begin{solution}
(a) The symmetric matrix \(A=\left(\begin{smallmatrix}a&b\\b&d\end{smallmatrix}\right)\) is
idempotent if and only if
\(\left(\begin{smallmatrix}a^2+b^2&b(a+d)\\b(a+d)&b^2+d^2\end{smallmatrix}\right)=\left(\begin{smallmatrix}a&b\\b&d\end{smallmatrix}\right)\),
that is, \(a^2+b^2=a\), \(b(a+d)=b\), \(d^2+b^2=d\).
\begin{itemize}
\item If \(a+d\ne1\), the second equation gives \(b=0\); then \(a^2=a\) and \(d^2=d\), so
\(a,d\in\{0,1\}\), and \(a+d\ne1\) leaves \(a=d=0\) or \(a=d=1\): \(A=0\) or \(A=I_2\).
\item If \(a+d=1\), then \(d=1-a\) and the first equation gives \(b^2=a(1-a)\), so
\(0\leqslant a\leqslant1\) and \(b=\pm\sqrt{a(1-a)}\); the third equation then holds, since
\((1-a)^2+a(1-a)=1-a\). Hence
\[
A=\begin{pmatrix}a&\pm\sqrt{a(1-a)}\\\pm\sqrt{a(1-a)}&1-a\end{pmatrix}
=\begin{pmatrix}\sqrt a\\\pm\sqrt{1-a}\end{pmatrix}\begin{pmatrix}\sqrt a&\pm\sqrt{1-a}\end{pmatrix}=uu',
\]
with the same sign throughout, and \(u'u=a+(1-a)=1\).
\end{itemize}
Conversely, \(I_2\) and \(0\) are symmetric idempotent, and if \(u'u=1\), then \(uu'\)
is symmetric and \((uu')(uu')=u(u'u)u'=uu'\).

(b) Let \(a,b\in\R^n\) with \(b'a=1\), and \(A=ab'\). Then
\(A^2=a(b'a)b'=ab'=A\). This \(A\) is symmetric only when \(b\) is a multiple of
\(a\): if \(ab'=ba'\), multiplying on the right by \(a\) gives \(a(b'a)=b(a'a)\),
that is, \(a=(a'a)b\); here \(a\ne0\) because \(b'a=1\), so \(a'a>0\) and
\(b=a/(a'a)\). Conversely, if \(b=ca\), then \(ab'=c\,aa'\) is symmetric. For a
concrete example take \(a=e_1\) and \(b=e_1+e_2\): \(b'a=1\), and \(A=e_1e_1'+e_1e_2'\)
has \((1,1)\) and \((1,2)\) entries \(1\) and all other entries \(0\); for \(n=2\),
\(A=\left(\begin{smallmatrix}1&1\\0&0\end{smallmatrix}\right)\), with \(A^2=A\ne A'\).
\end{solution}

\begin{problem}\label{r1-ex-data-drill}
Four units are measured on two variables (\(n=4\), \(k=2\)), with data matrix
\[
X_0=\begin{pmatrix}1&0\\3&2\\2&0\\2&2\end{pmatrix}\in\R^{4\times2}
\]
Find
\begin{enumerate}
\item the mean vector \(\bar x\) and the centered data matrix \(C_4X_0\);
\item the matrix \(X_0'C_4X_0\) of corrected sums of squares and products, both as
\((C_4X_0)'(C_4X_0)\) and as \(X_0'X_0-4\bar x\bar x'\);
\item \(X'X\) for the design matrix \(X=[\one_4\ X_0]\).
\end{enumerate}
\end{problem}
\begin{solution}
(a) \(\one'X_0=(1+3+2+2,\ 0+2+0+2)=(8,\ 4)\), so \(\bar x=\frac14X_0'\one=(2,\ 1)'\). As in
Example~\ref{r1-exm-data}, \(C_4X_0=X_0-\one\bar x'\):
\[
C_4X_0=\begin{pmatrix}1-2&0-1\\3-2&2-1\\2-2&0-1\\2-2&2-1\end{pmatrix}
=\begin{pmatrix}-1&-1\\1&1\\0&-1\\0&1\end{pmatrix}
\]

(b) The \((j,l)\) entry of \((C_4X_0)'(C_4X_0)\) is \(\col_j(C_4X_0)'\col_l(C_4X_0)\)
(Theorems~\ref{r1-thm-views}(d) and~\ref{r1-thm-transpose}): \(1+1+0+0=2\) for \((1,1)\),
\(1+1+0+0=2\) for \((1,2)\) and \((2,1)\), and \(1+1+1+1=4\) for \((2,2)\). On the other hand
\[
X_0'X_0-4\bar x\bar x'=\begin{pmatrix}18&10\\10&8\end{pmatrix}-4\begin{pmatrix}4&2\\2&1\end{pmatrix}
=\begin{pmatrix}18&10\\10&8\end{pmatrix}-\begin{pmatrix}16&8\\8&4\end{pmatrix}
=\begin{pmatrix}2&2\\2&4\end{pmatrix}
\]
(for instance \(18=1+9+4+4\) and \(10=0+6+0+4\)). So
\(X_0'C_4X_0=\left(\begin{smallmatrix}2&2\\2&4\end{smallmatrix}\right)\): the sums of squares
about the means are \(\sum_i(x_{i1}-\bar x_1)^2=2\) and \(\sum_i(x_{i2}-\bar x_2)^2=4\), and the
corrected sum of products is \(\sum_i(x_{i1}-\bar x_1)(x_{i2}-\bar x_2)=2\).

(c) As in Example~\ref{r1-exm-data}, the entries of \(X'X\) are \(n=4\), the column sums
\(8\), \(4\) of \(X_0\) and the entries of \(X_0'X_0\):
\[
X'X=\begin{pmatrix}4&8&4\\8&18&10\\4&10&8\end{pmatrix}
\]
\end{solution}
\clearpage
\def\unitnumber{2}
\setcounter{theorem}{0}\setcounter{problem}{0}
\section*{Chapter 2. Gauss elimination, RREF, and inverses}

\subsection*{Linear systems and augmented matrices}

\begin{definition}\label{r2-def-system}
Let \(A=(a_{ij})\in\R^{m\times n}\) and \(b\in\R^m\). The system of \(m\) linear equations
in \(n\) unknowns
\[
\begin{aligned}
a_{11}x_1+a_{12}x_2+\dots+a_{1n}x_n&=b_1\\
&\ \,\vdots\\
a_{m1}x_1+a_{m2}x_2+\dots+a_{mn}x_n&=b_m
\end{aligned}
\]
is the matrix equation \(Ax=b\), since \((Ax)_i=\sum_ja_{ij}x_j\). \(A\) is the
\textbf{coefficient matrix} and \([A\mid b]\in\R^{m\times(n+1)}\), the matrix \(A\) with
the column \(b\) appended, is the \textbf{augmented matrix}. A \textbf{solution} is an
\(x\in\R^n\) with \(Ax=b\). The system is \textbf{consistent} if it has a solution and
\textbf{inconsistent} otherwise. It is \textbf{homogeneous} if \(b=0\); then \(x=0\) is a
solution, the \textbf{trivial solution}.
\end{definition}

\subsection*{Elementary row operations}

\begin{definition}\label{r2-def-rowops}
The \textbf{elementary row operations} on a matrix with \(m\) rows are, for row indices
\(r\ne s\):
\begin{enumerate}[label=(\arabic*)]
\item multiply row \(r\) by a number \(c\ne0\), written \(R_r\leftarrow cR_r\);
\item add \(c\) times row \(s\) to row \(r\) (\(c\in\R\)), written \(R_r\leftarrow R_r+cR_s\);
\item interchange rows \(r\) and \(s\), written \(R_r\leftrightarrow R_s\).
\end{enumerate}
When several operations label one arrow, they are performed one after another, in the order
written (upper label first). \(B\) is \textbf{row-equivalent} to \(A\) if \(B\) is obtained
from \(A\) by finitely many elementary row operations.
\end{definition}

\begin{theorem}\label{r2-thm-rowops}
\begin{enumerate}
\item Each elementary row operation is undone by an operation of the same type:
\(R_r\leftarrow cR_r\) by \(R_r\leftarrow c^{-1}R_r\), \(R_r\leftarrow R_r+cR_s\) by
\(R_r\leftarrow R_r-cR_s\), and \(R_r\leftrightarrow R_s\) by itself.
\item If row operations take \([A\mid b]\) to \([B\mid z]\), then the systems \(Ax=b\) and
\(Bx=z\) have exactly the same solutions.
\end{enumerate}
\end{theorem}
\begin{proof}
(a) is checked row by row. (b) It suffices to consider one operation. Each equation of the
new system is a multiple of an equation of the old one, or a sum of an old equation and a
multiple of another, so every solution of \(Ax=b\) solves \(Bx=z\). By (a), the old system
is obtained back from the new one by a row operation, so every solution of \(Bx=z\) solves
\(Ax=b\).
\end{proof}
So we solve a system by simplifying its augmented matrix with row operations; the next
definition says how far to simplify.

\subsection*{Echelon form and RREF}

\begin{definition}\label{r2-def-rref}
A matrix \(R\) is in \textbf{reduced row echelon form} (\textbf{RREF}) if
\begin{enumerate}
\item the first nonzero entry of each nonzero row is \(1\), its \textbf{pivot};
\item each column that contains a pivot has all its other entries \(0\);
\item every zero row lies below every nonzero row;
\item if rows \(1,\dots,r\) are the nonzero rows and the pivot of row \(i\) lies in column
\(k_i\), then \(k_1<k_2<\dots<k_r\).
\end{enumerate}
The columns \(k_1,\dots,k_r\) are the \textbf{pivot columns}; column \(k_i\) of \(R\) is the
unit vector \(e_i\). A matrix that satisfies (a), (c), (d), but not necessarily (b), is in
\textbf{echelon form}; then every entry below a pivot is \(0\).
\end{definition}
For example, \(I_n\) and \(0\) are in RREF, and so is
\[
R=\begin{pmatrix}0&1&-3&0&2\\0&0&0&1&2\\0&0&0&0&0\end{pmatrix}\qquad(r=2,\ k_1=2,\ k_2=4)
\]
The matrix \(\left(\begin{smallmatrix}1&0&0&0\\0&1&-1&0\\0&0&1&0\end{smallmatrix}\right)\) is
in echelon form but not in RREF (column \(3\) violates (b)), and
\(\left(\begin{smallmatrix}0&2&1\\1&0&-3\\0&0&0\end{smallmatrix}\right)\) violates (a) and (d).

\begin{definition}[Gauss--Jordan elimination]\label{r2-def-gauss-jordan}
Let \(M\) be any matrix (for instance an augmented matrix). Start with \(i=1\).
\begin{enumerate}[label=\textit{Step \arabic*.},leftmargin=*]
\item Find the first column that has a nonzero entry in rows \(i,i+1,\dots\); if there is
none, go to Step 5. This is the next pivot column.
\item If the entry of row \(i\) in this column is \(0\), interchange row \(i\) with a lower row
that has a nonzero entry there.
\item Multiply row \(i\) by the reciprocal of this entry, so that the pivot is \(1\).
\item Add suitable multiples of row \(i\) to the rows below it, so that every entry below the
pivot is \(0\). Increase \(i\) by \(1\) and return to Step 1.
\item (\textit{Backward elimination.}) For each pivot, from the last one to the first, add
suitable multiples of its row to the rows above it, so that every entry above the pivot is
\(0\).
\end{enumerate}
Steps 1--4 (\textbf{forward elimination}) produce an echelon form; Step 5 produces the RREF.
The scaling in Step 3 may also be postponed, to the end of forward elimination or even to the
very end (the pivots become \(1\) only then); this often avoids fractions. Every matrix is
thus row-equivalent to a matrix in RREF; this RREF is unique (proved after
Theorem~\ref{thm:r4-pivot-columns}), and is called the RREF of \(M\).
\end{definition}

\begin{example}\label{r2-exm-gauss-jordan}
Solve the homogeneous system
\[
x_1+x_2+x_3-2x_4=0,\qquad x_1+x_2+2x_3-3x_4=0,\qquad 2x_1+x_2+5x_3-4x_4=0
\]
\textit{Solution.} Row-reduce the coefficient matrix \(A\) (the column \(b=0\) never changes,
so it is not written).
\[
\begin{gathered}
\begin{pmatrix}1&1&1&-2\\1&1&2&-3\\2&1&5&-4\end{pmatrix}
\xrightarrow[R_3\leftarrow R_3-2R_1]{R_2\leftarrow R_2-R_1}
\begin{pmatrix}1&1&1&-2\\0&0&1&-1\\0&-1&3&0\end{pmatrix}
\xrightarrow{R_2\leftrightarrow R_3}
\begin{pmatrix}1&1&1&-2\\0&-1&3&0\\0&0&1&-1\end{pmatrix}
\\
\xrightarrow{R_2\leftarrow-R_2}
\begin{pmatrix}1&1&1&-2\\0&1&-3&0\\0&0&1&-1\end{pmatrix}
\xrightarrow[R_1\leftarrow R_1-R_3]{R_2\leftarrow R_2+3R_3}
\begin{pmatrix}1&1&0&-1\\0&1&0&-3\\0&0&1&-1\end{pmatrix}
\xrightarrow{R_1\leftarrow R_1-R_2}
\begin{pmatrix}1&0&0&2\\0&1&0&-3\\0&0&1&-1\end{pmatrix}
\end{gathered}
\]
After the first step the entry of row \(2\) in column \(2\) is \(0\), so Step 2 interchanges
rows \(2\) and \(3\), and Step 3 multiplies the new row \(2\) by \(-1\). The first matrix of the
second line is an echelon form (end of forward elimination); the last matrix is the RREF \(R\),
with pivot columns \(1,2,3\). The system \(Rx=0\) reads \(x_1+2x_4=0\), \(x_2-3x_4=0\),
\(x_3-x_4=0\), so \(x_4=t\) is free and all solutions are
\[
x=t\,(-2,\ 3,\ 1,\ 1)'\qquad(t\in\R)
\]
Check: \(-2+3+1-2=0\), \(-2+3+2-3=0\), \(-4+3+5-4=0\).
\end{example}

\subsection*{Reading off the solutions}

\begin{proposition}\label{r2-prop-solutions}
Let row operations take \([A\mid b]\) (\(A\in\R^{m\times n}\)) to \([R\mid z]\) with \(R\) in
RREF, nonzero rows \(1,\dots,r\) and pivot columns \(k_1<\dots<k_r\). Call \(x_{k_1},\dots,x_{k_r}\)
the \textbf{pivot unknowns} and the other \(n-r\) unknowns the \textbf{free unknowns}.
\begin{enumerate}
\item \(Ax=b\) is consistent if and only if \(z_i=0\) for every \(i>r\), that is, no row of
\([R\mid z]\) reads \((0\ \cdots\ 0\mid z_i)\) with \(z_i\ne0\).
\item If it is consistent, all solutions are obtained by choosing arbitrary values for the free
unknowns and computing each pivot unknown from its row:
\(x_{k_i}=z_i-\sum_{j\text{ free}}R_{ij}x_j\) (\(1\leqslant i\leqslant r\)).
\end{enumerate}
\end{proposition}
\begin{proof}
By Theorem~\ref{r2-thm-rowops}, \(Ax=b\) and \(Rx=z\) have the same solutions. (a) Row
\(i>r\) of \(Rx=z\) reads \(0=z_i\). So there is no solution if one of these \(z_i\) is not \(0\);
if all of them are \(0\), these rows hold for every \(x\), and the argument for (b) gives a
solution. (b) For \(i\leqslant r\), column \(k_l\) of \(R\) is
\(e_l\), so the pivot unknown \(x_{k_i}\) occurs only in equation \(i\), with coefficient \(1\);
each choice of the free unknowns therefore gives exactly one solution.
\end{proof}

\begin{example}\label{r2-exm-read}
For the RREF \(R\) displayed after Definition~\ref{r2-def-rref}, the system \(Rx=0\) is
\(x_2-3x_3+2x_5=0\), \(x_4+2x_5=0\); the pivot unknowns are \(x_2,x_4\) and the free
unknowns \(x_1,x_3,x_5\). With \(x_1=a\), \(x_3=b\), \(x_5=c\), all solutions are
\[
x=\begin{pmatrix}a\\3b-2c\\b\\-2c\\c\end{pmatrix}
=a\begin{pmatrix}1\\0\\0\\0\\0\end{pmatrix}+b\begin{pmatrix}0\\3\\1\\0\\0\end{pmatrix}+c\begin{pmatrix}0\\-2\\0\\-2\\1\end{pmatrix}
\qquad(a,b,c\in\R)
\]
\end{example}

\begin{example}\label{r2-exm-consistency}
For which \(y=(y_1,y_2,y_3)'\) is \(Ax=y\) consistent, where
\(A=\left(\begin{smallmatrix}1&-2&1\\2&-3&1\\0&1&-1\end{smallmatrix}\right)\)? Solve it for
\(y=(1,2,0)'\).

\textit{Solution.} Row-reduce the augmented matrix, keeping \(y_1,y_2,y_3\) as letters:
\[
\begin{gathered}
\left(\begin{array}{ccc|c}1&-2&1&y_1\\2&-3&1&y_2\\0&1&-1&y_3\end{array}\right)
\xrightarrow{R_2\leftarrow R_2-2R_1}
\left(\begin{array}{ccc|c}1&-2&1&y_1\\0&1&-1&y_2-2y_1\\0&1&-1&y_3\end{array}\right)
\\
\xrightarrow{R_3\leftarrow R_3-R_2}
\left(\begin{array}{ccc|c}1&-2&1&y_1\\0&1&-1&y_2-2y_1\\0&0&0&y_3-y_2+2y_1\end{array}\right)
\xrightarrow{R_1\leftarrow R_1+2R_2}
\left(\begin{array}{ccc|c}1&0&-1&2y_2-3y_1\\0&1&-1&y_2-2y_1\\0&0&0&y_3-y_2+2y_1\end{array}\right)
\end{gathered}
\]
Here \(r=2\) and \(z_3=y_3-y_2+2y_1\). By Proposition~\ref{r2-prop-solutions}, the system is
consistent if and only if \(2y_1-y_2+y_3=0\); then \(x_3=c\) is free and
\[
x_1=2y_2-3y_1+c,\qquad x_2=y_2-2y_1+c,\qquad x_3=c\qquad(c\in\R)
\]
For \(y=(1,2,0)'\), \(2-2+0=0\), and the solutions are \(x=(1,0,0)'+c\,(1,1,1)'\): the
solution \((1,0,0)'\) plus the multiples of \((1,1,1)'\), which solve \(Ax=0\). For
\(y=(1,0,0)'\), \(2y_1-y_2+y_3=2\ne0\), so \(Ax=y\) has no solution.
\end{example}

\subsection*{Homogeneous systems}

\begin{theorem}\label{r2-thm-homogeneous}
\begin{enumerate}
\item If \(A\in\R^{m\times n}\) with \(m<n\) (fewer equations than unknowns), then \(Ax=0\)
has a nontrivial solution.
\item For \(A\in\R^{n\times n}\), \(Ax=0\) has only the trivial solution if and only if the
RREF of \(A\) is \(I_n\).
\end{enumerate}
\end{theorem}
\begin{proof}
Let \(R\) be the RREF of \(A\), with \(r\) nonzero rows. (a) \(r\leqslant m<n\), so there is
a free unknown; giving it the value \(1\) and the other free unknowns the value \(0\) gives a
solution \(x\ne0\) (Proposition~\ref{r2-prop-solutions}(b)). (b) If \(R=I\), then \(Rx=x\),
so \(x=0\) is the only solution. If \(Ax=0\) has only the trivial solution, there is no free
unknown, so \(r=n\); then \(1\leqslant k_1<\dots<k_n\leqslant n\) forces \(k_i=i\), so column
\(i\) of \(R\) is \(e_i\) and \(R=I\).
\end{proof}

\subsection*{Inverses}

\begin{definition}\label{r2-def-inverse}
\(A\in\R^{n\times n}\) is \textbf{invertible} if \(AB=BA=I_n\) for some
\(B\in\R^{n\times n}\). Such a \(B\) is unique: if \(BA=I\) and \(AC=I\), then
\(B=B(AC)=(BA)C=C\). It is called the \textbf{inverse} of \(A\), written \(A^{-1}\). For
example, \(\diag(d_1,\dots,d_n)\) with all \(d_i\ne0\) has inverse
\(\diag(d_1^{-1},\dots,d_n^{-1})\), since the products in both orders are
\(\diag(1,\dots,1)=I\).
\end{definition}

\begin{theorem}\label{r2-thm-inverse-props}
Let \(A,B\in\R^{n\times n}\) be invertible.
\begin{enumerate}
\item \(A^{-1}\) is invertible and \((A^{-1})^{-1}=A\).
\item \(AB\) is invertible and \((AB)^{-1}=B^{-1}A^{-1}\); more generally
\((A_1\cdots A_k)^{-1}=A_k^{-1}\cdots A_1^{-1}\) for invertible \(A_1,\dots,A_k\).
\item \(A'\) is invertible and \((A')^{-1}=(A^{-1})'\).
\end{enumerate}
\end{theorem}
\begin{proof}
(a) \(AA^{-1}=A^{-1}A=I\) says that \(A\) is the inverse of \(A^{-1}\). (b)
\((AB)(B^{-1}A^{-1})=A(BB^{-1})A^{-1}=AA^{-1}=I\), and similarly
\((B^{-1}A^{-1})(AB)=I\); for \(k\) factors repeat this. (c)
\(A'(A^{-1})'=(A^{-1}A)'=I\) and \((A^{-1})'A'=(AA^{-1})'=I\) by
Theorem~\ref{r1-thm-transpose}.
\end{proof}

\subsection*{Elementary matrices}

\begin{definition}\label{r2-def-elementary}
An \textbf{elementary matrix} is a matrix \(E=\omega(I_m)\) obtained from \(I_m\) by one
elementary row operation \(\omega\). For \(r\ne s\) we write:
\begin{itemize}
\item \(D_r(c)\) (\(c\ne0\)): \(I_m\) with its \((r,r)\) entry replaced by \(c\)
(from \(R_r\leftarrow cR_r\));
\item \(T_{rs}(c)\): \(I_m\) with its \((r,s)\) entry replaced by \(c\)
(from \(R_r\leftarrow R_r+cR_s\));
\item \(P_{rs}\): \(I_m\) with rows \(r\) and \(s\) interchanged (from \(R_r\leftrightarrow R_s\)).
\end{itemize}
For example, for \(m=3\):
\[
P_{23}=\begin{pmatrix}1&0&0\\0&0&1\\0&1&0\end{pmatrix},\qquad
T_{31}(c)=\begin{pmatrix}1&0&0\\0&1&0\\c&0&1\end{pmatrix},\qquad
D_3(c)=\begin{pmatrix}1&0&0\\0&1&0\\0&0&c\end{pmatrix}
\]
\end{definition}

\begin{theorem}\label{r2-thm-elementary}
Let \(\omega\) be an elementary row operation on matrices with \(m\) rows and \(E=\omega(I_m)\).
\begin{enumerate}
\item \(\omega(A)=EA\) for every \(A\in\R^{m\times n}\): a row operation is a left
multiplication.
\item \(E\) is invertible, and \(E^{-1}\) is the elementary matrix of the inverse operation:
\(D_r(c)^{-1}=D_r(c^{-1})\), \(T_{rs}(c)^{-1}=T_{rs}(-c)\), \(P_{rs}^{-1}=P_{rs}\).
\item If row operations \(\omega_1,\dots,\omega_k\) (in this order) take \(A\) to \(B\), then
\(B=PA\) with \(P=E_k\cdots E_2E_1\), \(E_j=\omega_j(I_m)\); and \(P\) is invertible.
\end{enumerate}
\end{theorem}
\begin{proof}
(a) By Theorem~\ref{r1-thm-views}(c), \(\row_i(EA)=\row_i(E)A\). An unchanged row of \(E\) is
\(e_i'\), and \(e_i'A=\row_i(A)\). For type (2), \(\row_r(E)=e_r'+ce_s'\),
so \(\row_r(EA)=\row_r(A)+c\row_s(A)\); types (1) and (3) are the same computation. (b) Let
\(\tilde\omega\) be the inverse operation (Theorem~\ref{r2-thm-rowops}(a)) and \(\tilde E=\tilde\omega(I)\). By (a),
\(E\tilde E=\omega(\tilde\omega(I))=I\) and \(\tilde EE=\tilde\omega(\omega(I))=I\). (c) By (a), \(B=E_k(\cdots(E_1A))=(E_k\cdots E_1)A\),
and a product of invertible matrices is invertible (Theorem~\ref{r2-thm-inverse-props}(b)).
\end{proof}

\subsection*{Deciding invertibility and computing inverses}

\begin{theorem}\label{r2-thm-invertible}
For \(A\in\R^{n\times n}\) the following are equivalent:
\begin{enumerate}[label=(\roman*)]
\item \(A\) is invertible;
\item the RREF of \(A\) is \(I_n\);
\item \(Ax=0\) has only the trivial solution;
\item \(Ax=b\) has a solution for every \(b\in\R^n\);
\item \(A\) is a product of elementary matrices.
\end{enumerate}
Moreover, if \(BA=I\) or \(AB=I\) for some \(B\in\R^{n\times n}\), then \(A\) is invertible
and \(B=A^{-1}\).
\end{theorem}
\begin{proof}
Let \(R=PA\) be the RREF of \(A\), with \(P=E_k\cdots E_1\) (Theorem~\ref{r2-thm-elementary}(c)).
\begin{itemize}
\item If \(R=I\), then \(A=P^{-1}=E_1^{-1}\cdots E_k^{-1}\) is a product of elementary
matrices, hence invertible (Theorem~\ref{r2-thm-inverse-props}(b)); \(Ax=0\) has only the trivial
solution (Theorem~\ref{r2-thm-homogeneous}(b)); and \(x=A^{-1}b\) solves \(Ax=b\).
\item If \(R\ne I\), then \(R\) has a zero row, hence a zero last row (by (c) of
Definition~\ref{r2-def-rref}); otherwise \(r=n\) and \(R=I\) as in the proof of
Theorem~\ref{r2-thm-homogeneous}(b). Then \(Ax=0\) has a solution \(x\ne0\)
(Theorem~\ref{r2-thm-homogeneous}(b)), so \(A\) has no inverse (else \(x=A^{-1}Ax=0\)
\(\contradiction\)), and \(A\) is not a product of elementary matrices. For
\(b=P^{-1}e_n\), the system \(Ax=b\) has the same solutions as \(Rx=Pb=e_n\), whose last
equation reads \(0=1\); so (iv) fails.
\end{itemize}
Thus each of (i)--(v) holds exactly when \(R=I\). If \(BA=I\), then \(Ax=0\) gives
\(x=BAx=0\), so \(A\) is invertible by (iii), and \(B=BAA^{-1}=A^{-1}\). If \(AB=I\), then
\(B\) is invertible with \(B^{-1}=A\) by the case just shown, so \(A=B^{-1}\) is invertible
with \(A^{-1}=B\).
\end{proof}

\textbf{Computing \(A^{-1}\).} Row-reduce the \(n\times2n\) matrix \([A\mid I_n]\) by
Gauss--Jordan elimination.
\begin{itemize}
\item If the left half becomes \(I_n\), the right half is \(A^{-1}\):
\([A\mid I]\to[I\mid B]\) means \(E[A\mid I]=[EA\mid E]=[I\mid B]\) for the product \(E\) of
the elementary matrices used (Theorem~\ref{r2-thm-elementary}(c); \(E\) multiplies each column
separately, Theorem~\ref{r1-thm-views}(a)), so \(EA=I\) and \(B=E=A^{-1}\) by
Theorem~\ref{r2-thm-invertible}.
\item If a zero row appears in the left half, \(A\) is not invertible. The left half is then
\(EA\), with \(E\) the (invertible) product of the elementary matrices used so far
(Theorem~\ref{r2-thm-elementary}(c)). For every \(C\in\R^{n\times n}\), \(EAC\) has the same zero
row (Theorem~\ref{r1-thm-views}(c)), so \(EAC\ne I_n\) and \(EA\) is not invertible; if \(A\) were
invertible, so would be \(EA\) (Theorem~\ref{r2-thm-inverse-props}(b)).
\end{itemize}

\begin{example}\label{r2-exm-inverse-2x2}
Find the inverse of
\(A=\left(\begin{smallmatrix}1&-1&1\\-1&1&-2\\2&-1&2\end{smallmatrix}\right)\) and write \(A\)
and \(A^{-1}\) as products of elementary matrices.

\textit{Solution.}
\[
\begin{gathered}
\left(\begin{array}{ccc|ccc}1&-1&1&1&0&0\\-1&1&-2&0&1&0\\2&-1&2&0&0&1\end{array}\right)
\xrightarrow[R_3\leftarrow R_3-2R_1]{R_2\leftarrow R_2+R_1}
\left(\begin{array}{ccc|ccc}1&-1&1&1&0&0\\0&0&-1&1&1&0\\0&1&0&-2&0&1\end{array}\right)
\\
\xrightarrow{R_2\leftrightarrow R_3}
\left(\begin{array}{ccc|ccc}1&-1&1&1&0&0\\0&1&0&-2&0&1\\0&0&-1&1&1&0\end{array}\right)
\xrightarrow{R_3\leftarrow-R_3}
\left(\begin{array}{ccc|ccc}1&-1&1&1&0&0\\0&1&0&-2&0&1\\0&0&1&-1&-1&0\end{array}\right)
\\
\xrightarrow{R_1\leftarrow R_1-R_3}
\left(\begin{array}{ccc|ccc}1&-1&0&2&1&0\\0&1&0&-2&0&1\\0&0&1&-1&-1&0\end{array}\right)
\xrightarrow{R_1\leftarrow R_1+R_2}
\left(\begin{array}{ccc|ccc}1&0&0&0&1&1\\0&1&0&-2&0&1\\0&0&1&-1&-1&0\end{array}\right)
\end{gathered}
\]
After the first step the entry of row \(2\) in column \(2\) is \(0\), so Step 2 interchanges
rows \(2\) and \(3\). The left half has become \(I_3\), so
\(A^{-1}=\left(\begin{smallmatrix}0&1&1\\-2&0&1\\-1&-1&0\end{smallmatrix}\right)\). The six
operations have the elementary matrices \(E_1=T_{21}(1)\), \(E_2=T_{31}(-2)\), \(E_3=P_{23}\),
\(E_4=D_3(-1)\), \(E_5=T_{13}(-1)\), \(E_6=T_{12}(1)\), and by
Theorem~\ref{r2-thm-elementary}(c)
\[
A^{-1}=E_6E_5E_4E_3E_2E_1
=\begin{pmatrix}1&1&0\\0&1&0\\0&0&1\end{pmatrix}\begin{pmatrix}1&0&-1\\0&1&0\\0&0&1\end{pmatrix}
\begin{pmatrix}1&0&0\\0&1&0\\0&0&-1\end{pmatrix}\begin{pmatrix}1&0&0\\0&0&1\\0&1&0\end{pmatrix}
\begin{pmatrix}1&0&0\\0&1&0\\-2&0&1\end{pmatrix}\begin{pmatrix}1&0&0\\1&1&0\\0&0&1\end{pmatrix}
\]
By Theorem~\ref{r2-thm-inverse-props}(b) and Theorem~\ref{r2-thm-elementary}(b),
\[
A=E_1^{-1}E_2^{-1}E_3^{-1}E_4^{-1}E_5^{-1}E_6^{-1}=T_{21}(-1)\,T_{31}(2)\,P_{23}\,D_3(-1)\,T_{13}(1)\,T_{12}(-1)
\]
\end{example}

\begin{example}\label{r2-exm-xtx-inverse}
For the matrix \(X'X\) of Example~\ref{r1-exm-data}, find \((X'X)^{-1}\).

\textit{Solution.} We postpone the scaling to the end (Definition~\ref{r2-def-gauss-jordan}),
so that fractions appear only in the last step.
\[
\begin{gathered}
\left(\begin{array}{ccc|ccc}4&4&8&1&0&0\\4&8&8&0&1&0\\8&8&20&0&0&1\end{array}\right)
\xrightarrow[R_3\leftarrow R_3-2R_1]{R_2\leftarrow R_2-R_1}
\left(\begin{array}{ccc|ccc}4&4&8&1&0&0\\0&4&0&-1&1&0\\0&0&4&-2&0&1\end{array}\right)
\xrightarrow{R_1\leftarrow R_1-2R_3}
\left(\begin{array}{ccc|ccc}4&4&0&5&0&-2\\0&4&0&-1&1&0\\0&0&4&-2&0&1\end{array}\right)
\\
\xrightarrow{R_1\leftarrow R_1-R_2}
\left(\begin{array}{ccc|ccc}4&0&0&6&-1&-2\\0&4&0&-1&1&0\\0&0&4&-2&0&1\end{array}\right)
\xrightarrow[R_3\leftarrow\frac14R_3]{R_1\leftarrow\frac14R_1,\ R_2\leftarrow\frac14R_2}
\left(\begin{array}{ccc|ccc}1&0&0&\frac32&-\frac14&-\frac12\\0&1&0&-\frac14&\frac14&0\\0&0&1&-\frac12&0&\frac14\end{array}\right)
\end{gathered}
\]
So \((X'X)^{-1}=\frac14\left(\begin{smallmatrix}6&-1&-2\\-1&1&0\\-2&0&1\end{smallmatrix}\right)\).
The first step produces the corrected sums of Example~\ref{r1-exm-data}: row \(2\) of \(X'X\) is
\((n\bar x_1,\ X_{\cdot1}'X_{\cdot1},\ X_{\cdot1}'X_{\cdot2})=(4,8,8)\), and subtracting
\(\bar x_1=1\) times row \(1\), \((n,\ n\bar x_1,\ n\bar x_2)=(4,4,8)\), leaves
\((0,\ X_{\cdot1}'X_{\cdot1}-n\bar x_1^2,\ X_{\cdot1}'X_{\cdot2}-n\bar x_1\bar x_2)=(0,4,0)\);
likewise row \(3\) minus \(\bar x_2=2\) times row \(1\) is \((0,0,4)\). These are the rows of
\(X_0'CX_0\).
\end{example}

\begin{example}\label{r2-exm-singular}
The matrix \(A\) of Example~\ref{r2-exm-consistency} is not invertible: its RREF is
\(\left(\begin{smallmatrix}1&0&-1\\0&1&-1\\0&0&0\end{smallmatrix}\right)\ne I_3\) (the
left part of the last matrix there). Accordingly \(Ax=0\) has the nontrivial solution
\(x=(1,1,1)'\) (take \(c=1\)), and \(Ax=y\) has no solution for \(y=(1,0,0)'\).
\end{example}

\subsection*{Solving \(Ax=b\)}

In practice \(Ax=b\) is solved by forward elimination of \([A\mid b]\) alone: an echelon form
already shows whether there is a solution, and how many.

\begin{theorem}\label{thm:r2-consistency}
Let \(A\in\R^{m\times n}\) and \(b\in\R^m\), and let row operations take \([A\mid b]\) to a
matrix \([E\mid d]\) in echelon form (\(E\in\R^{m\times n}\), \(d\in\R^m\)). Then \(Ax=b\) is
consistent if and only if \([E\mid d]\) has no row \((0\ \cdots\ 0\mid c)\) with \(c\ne0\),
that is, if and only if the last column of \([E\mid d]\) is not a pivot column.
\end{theorem}
\begin{proof}
\begin{itemize}
\item By Theorem~\ref{r2-thm-rowops}, \(Ax=b\) and \(Ex=d\) have the same solutions. In an
echelon form, a row \((0\ \cdots\ 0\mid c)\) with \(c\ne0\) is exactly a row whose pivot lies
in the last column.
\item If \([E\mid d]\) has such a row, the corresponding equation reads \(0=c\), and no \(x\)
satisfies it. The same holds if such a row appears at any stage of the elimination, so one may
stop there.
\item Suppose there is no such row. Then every zero row of \(E\) is a zero row of \([E\mid d]\),
and its equation reads \(0=0\). Let rows \(1,\dots,r\) be the nonzero rows of \(E\), with pivots
in columns \(k_1<\dots<k_r\); call \(x_{k_1},\dots,x_{k_r}\) the pivot unknowns and the others
free, as in Proposition~\ref{r2-prop-solutions}.
\item Give the free unknowns any values. Row \(i\) of \(E=(e_{ij})\) is \(0\) before column
\(k_i\) and \(1\) there, so for \(i=r,r-1,\dots,1\) equation \(i\) reads
\(x_{k_i}+\sum_{j>k_i}e_{ij}x_j=d_i\). Each \(x_j\) with \(j>k_i\) is free or the pivot unknown
of a lower row, already found; so the equation determines \(x_{k_i}\) (\textbf{back
substitution}). The resulting \(x\) solves \(Ex=d\), hence \(Ax=b\).
\item Every solution arises in this way from the values of its free unknowns, since the
equations force the pivot unknowns.\qedhere
\end{itemize}
\end{proof}

\begin{theorem}\label{thm:r2-solution-structure}
Let \(A\in\R^{m\times n}\) and \(b\in\R^m\).
\begin{enumerate}
\item If \(x_p\) is one solution of \(Ax=b\) (a \textbf{particular solution}), then the
solutions of \(Ax=b\) are exactly the vectors \(x_p+z\) with \(Az=0\).
\item \(Ax=b\) has no solution, exactly one solution, or infinitely many solutions. If it is
consistent, it has exactly one solution if and only if there are no free unknowns, that is,
every column of \(E\) in Theorem~\ref{thm:r2-consistency} is a pivot column.
\item If \(A\) is square, \(Ax=b\) has exactly one solution if and only if \(A\) is invertible;
this solution is \(x=A^{-1}b\).
\end{enumerate}
\end{theorem}
\begin{proof}
\begin{enumerate}
\item If \(Az=0\), then \(A(x_p+z)=Ax_p+Az=b\). Conversely, if \(Ax=b\), then \(z=x-x_p\)
satisfies \(Az=b-b=0\), and \(x=x_p+z\).
\item Suppose \(Ax=b\) is consistent, and take \([E\mid d]\) as in
Theorem~\ref{thm:r2-consistency}.
\begin{itemize}
\item By the proof of that theorem, each choice of values of the free unknowns gives exactly one
solution, every solution arises so, and different choices give different solutions (they differ
in a free unknown, an entry of \(x\)).
\item Without free unknowns there is therefore exactly one solution; with a free unknown, each
of its values gives a different solution, so there are infinitely many.
\end{itemize}
\item
\begin{itemize}
\item If \(A\) is invertible, \(x=A^{-1}b\) is a solution, since \(A(A^{-1}b)=b\); and every
solution satisfies \(x=A^{-1}(Ax)=A^{-1}b\).
\item If \(Ax=b\) has exactly one solution \(x_p\), then \(Az=0\) only for \(z=0\), since
otherwise \(x_p+z\) would be a second solution by (a). So \(A\) is invertible by
Theorem~\ref{r2-thm-invertible} ((iii)\,\(\Rightarrow\)\,(i)).\qedhere
\end{itemize}
\end{enumerate}
\end{proof}
So, to solve \(Ax=b\): reduce \([A\mid b]\) by forward elimination to an echelon form; if a row
\((0\ \cdots\ 0\mid c)\) with \(c\ne0\) appears, there is no solution; otherwise choose the free
unknowns arbitrarily and find the pivot unknowns by back substitution (or continue to the RREF
and use Proposition~\ref{r2-prop-solutions}).

\begin{example}\label{r2-exm-unique}
Solve the system
\(x_1-x_2+x_3=1\), \(-x_1+x_2-2x_3=0\), \(2x_1-x_2+2x_3=3\).

\textit{Solution.} The coefficient matrix is the matrix \(A\) of
Example~\ref{r2-exm-inverse-2x2}, and \(b=(1,0,3)'\). Forward elimination of \([A\mid b]\)
uses the operations of the forward elimination there:
\[
\begin{gathered}
\left(\begin{array}{ccc|c}1&-1&1&1\\-1&1&-2&0\\2&-1&2&3\end{array}\right)
\xrightarrow[R_3\leftarrow R_3-2R_1]{R_2\leftarrow R_2+R_1}
\left(\begin{array}{ccc|c}1&-1&1&1\\0&0&-1&1\\0&1&0&1\end{array}\right)
\\
\xrightarrow{R_2\leftrightarrow R_3}
\left(\begin{array}{ccc|c}1&-1&1&1\\0&1&0&1\\0&0&-1&1\end{array}\right)
\xrightarrow{R_3\leftarrow-R_3}
\left(\begin{array}{ccc|c}1&-1&1&1\\0&1&0&1\\0&0&1&-1\end{array}\right)
\end{gathered}
\]
\begin{itemize}
\item The last matrix is an echelon form with no row \((0\ 0\ 0\mid c)\), \(c\ne0\), so the
system is consistent (Theorem~\ref{thm:r2-consistency}).
\item Every column of the left part is a pivot column, so there are no free unknowns and the
solution is unique (Theorem~\ref{thm:r2-solution-structure}(b)).
\item Back substitution: \(x_3=-1\); \(x_2=1\); \(x_1=1+x_2-x_3=1+1+1=3\). So
\(x=(3,1,-1)'\).
\item Check, with \(A^{-1}\) from Example~\ref{r2-exm-inverse-2x2}
(Theorem~\ref{thm:r2-solution-structure}(c)):
\(A^{-1}b=\left(\begin{smallmatrix}0&1&1\\-2&0&1\\-1&-1&0\end{smallmatrix}\right)
\left(\begin{smallmatrix}1\\0\\3\end{smallmatrix}\right)=\left(\begin{smallmatrix}3\\1\\-1\end{smallmatrix}\right)\).
\end{itemize}
\end{example}

\begin{example}\label{r2-exm-infinite}
Solve the system
\(x_1+x_2+x_3-2x_4=2\), \(x_1+x_2+2x_3-3x_4=3\), \(2x_1+x_2+5x_3-4x_4=7\).

\textit{Solution.} The coefficient matrix is the matrix \(A\) of
Example~\ref{r2-exm-gauss-jordan}, and \(b=(2,3,7)'\). Forward elimination of \([A\mid b]\)
uses the same operations as there:
\[
\begin{gathered}
\left(\begin{array}{cccc|c}1&1&1&-2&2\\1&1&2&-3&3\\2&1&5&-4&7\end{array}\right)
\xrightarrow[R_3\leftarrow R_3-2R_1]{R_2\leftarrow R_2-R_1}
\left(\begin{array}{cccc|c}1&1&1&-2&2\\0&0&1&-1&1\\0&-1&3&0&3\end{array}\right)
\\
\xrightarrow{R_2\leftrightarrow R_3}
\left(\begin{array}{cccc|c}1&1&1&-2&2\\0&-1&3&0&3\\0&0&1&-1&1\end{array}\right)
\xrightarrow{R_2\leftarrow-R_2}
\left(\begin{array}{cccc|c}1&1&1&-2&2\\0&1&-3&0&-3\\0&0&1&-1&1\end{array}\right)
\end{gathered}
\]
\begin{itemize}
\item The last matrix is an echelon form with no row \((0\ 0\ 0\ 0\mid c)\), \(c\ne0\): the
system is consistent (Theorem~\ref{thm:r2-consistency}).
\item The pivot columns are \(1,2,3\), so \(x_4\) is free and there are infinitely many
solutions (Theorem~\ref{thm:r2-solution-structure}(b)).
\item Back substitution with \(x_4=t\): \(x_3=1+t\); \(x_2=-3+3x_3=3t\);
\(x_1=2-x_2-x_3+2x_4=2-3t-(1+t)+2t=1-2t\). So
\[
x=\begin{pmatrix}1-2t\\3t\\1+t\\t\end{pmatrix}=\begin{pmatrix}1\\0\\1\\0\end{pmatrix}+t\begin{pmatrix}-2\\3\\1\\1\end{pmatrix}\qquad(t\in\R)
\]
\item This is the particular solution \(x_p=(1,0,1,0)'\) plus the solutions \(t(-2,3,1,1)'\)
of \(Ax=0\) found in Example~\ref{r2-exm-gauss-jordan}, as
Theorem~\ref{thm:r2-solution-structure}(a) says.
\end{itemize}
\end{example}

\begin{example}\label{r2-exm-inconsistent}
Show that the system
\(x_1-2x_2+x_3+2x_4=1\), \(x_1-x_2-x_3+x_4=2\), \(x_1+x_2-5x_3-x_4=3\) has no solution.

\textit{Solution.} Forward elimination of the augmented matrix \([A\mid b]\), \(b=(1,2,3)'\):
\[
\begin{gathered}
\left(\begin{array}{cccc|c}1&-2&1&2&1\\1&-1&-1&1&2\\1&1&-5&-1&3\end{array}\right)
\xrightarrow[R_3\leftarrow R_3-R_1]{R_2\leftarrow R_2-R_1}
\left(\begin{array}{cccc|c}1&-2&1&2&1\\0&1&-2&-1&1\\0&3&-6&-3&2\end{array}\right)
\\
\xrightarrow{R_3\leftarrow R_3-3R_2}
\left(\begin{array}{cccc|c}1&-2&1&2&1\\0&1&-2&-1&1\\0&0&0&0&-1\end{array}\right)
\xrightarrow{R_3\leftarrow-R_3}
\left(\begin{array}{cccc|c}1&-2&1&2&1\\0&1&-2&-1&1\\0&0&0&0&1\end{array}\right)
\end{gathered}
\]
The last matrix is an echelon form whose last column is a pivot column: row \(3\) reads
\(0=1\). So the system has no solution (Theorem~\ref{thm:r2-consistency}). Indeed, the
coefficient matrix \(A\) has \(\row_3(A)=3\row_2(A)-2\row_1(A)\), while
\(3b_2-2b_1=4\ne3=b_3\): the third equation contradicts the first two.
\end{example}

\subsection*{Exercises and solutions}

\begin{problem}\label{r2-ex-echelon}
Reduce the matrix
\(A=\begin{pmatrix}0&0&1&2&3\\2&4&-2&0&4\\2&4&-1&2&7\end{pmatrix}\)
to echelon form by elementary row operations, and record the elementary matrices used.
\end{problem}
\begin{solution}
Follow Definition~\ref{r2-def-gauss-jordan}, postponing the scaling (Step 3) to the end.
Column \(1\) is the first pivot column; since \(a_{11}=0\), interchange rows \(1\) and \(2\);
then clear column \(1\) below the pivot, and then column \(3\), the next pivot column:
\[
A\xrightarrow{R_1\leftrightarrow R_2}
\begin{pmatrix}2&4&-2&0&4\\0&0&1&2&3\\2&4&-1&2&7\end{pmatrix}
\xrightarrow{R_3\leftarrow R_3-R_1}
\begin{pmatrix}2&4&-2&0&4\\0&0&1&2&3\\0&0&1&2&3\end{pmatrix}
\xrightarrow{R_3\leftarrow R_3-R_2}
\begin{pmatrix}2&4&-2&0&4\\0&0&1&2&3\\0&0&0&0&0\end{pmatrix}
\]
Finally divide row \(1\) by \(2\):
\[
H=\begin{pmatrix}1&2&-1&0&2\\0&0&1&2&3\\0&0&0&0&0\end{pmatrix}
=D_1(\tfrac12)\,T_{32}(-1)\,T_{31}(-1)\,P_{12}\,A
\]
(Theorem~\ref{r2-thm-elementary}(c)). \(H\) is in echelon form with pivot columns \(1\) and
\(3\). One more operation, \(R_1\leftarrow R_1+R_2\), clears the entry above the second pivot
and gives the RREF
\(\left(\begin{smallmatrix}1&2&0&2&5\\0&0&1&2&3\\0&0&0&0&0\end{smallmatrix}\right)\).
\end{solution}

\begin{problem}\label{r2-ex-homogeneous-two}
Find the solutions of the equations
\(x_1+3x_2+2x_3=0\), \(2x_1+7x_2+3x_3=0\), \(x_1+4x_2+x_3=0\).
\end{problem}
\begin{solution}
Row-reduce the coefficient matrix:
\[
\begin{pmatrix}1&3&2\\2&7&3\\1&4&1\end{pmatrix}
\xrightarrow[R_3\leftarrow R_3-R_1]{R_2\leftarrow R_2-2R_1}
\begin{pmatrix}1&3&2\\0&1&-1\\0&1&-1\end{pmatrix}
\xrightarrow{R_3\leftarrow R_3-R_2}
\begin{pmatrix}1&3&2\\0&1&-1\\0&0&0\end{pmatrix}
\xrightarrow{R_1\leftarrow R_1-3R_2}
\begin{pmatrix}1&0&5\\0&1&-1\\0&0&0\end{pmatrix}
\]
Column \(3\) has no pivot, so \(x_3\) is a free unknown (Proposition~\ref{r2-prop-solutions}):
with \(x_3=\lambda\), the nonzero rows give \(x_1=-5\lambda\), \(x_2=\lambda\). All solutions are
\(x=\lambda(-5,1,1)'\), \(\lambda\in\R\); in particular the system has nontrivial solutions,
although it has as many equations as unknowns (the third equation is the second minus the first).
\end{solution}

\begin{problem}\label{r2-ex-homogeneous-four}
Solve the homogeneous system
\[
\begin{aligned}
x_1-x_2+2x_3+x_4&=0\\
x_1-x_3+3x_4&=0\\
3x_1-2x_2+3x_3+5x_4&=0\\
x_1+x_2-4x_3+5x_4&=0
\end{aligned}
\]
\end{problem}
\begin{solution}
Transform the coefficient matrix to echelon form:
\[
\begin{pmatrix}1&-1&2&1\\1&0&-1&3\\3&-2&3&5\\1&1&-4&5\end{pmatrix}
\xrightarrow[R_4\leftarrow R_4-R_1]{R_2\leftarrow R_2-R_1,\ R_3\leftarrow R_3-3R_1}
\begin{pmatrix}1&-1&2&1\\0&1&-3&2\\0&1&-3&2\\0&2&-6&4\end{pmatrix}
\xrightarrow[R_4\leftarrow R_4-2R_2]{R_3\leftarrow R_3-R_2}
\begin{pmatrix}1&-1&2&1\\0&1&-3&2\\0&0&0&0\\0&0&0&0\end{pmatrix}
\]
By Theorem~\ref{r2-thm-rowops}, the system has the same solutions as the two equations
\(x_1-x_2+2x_3+x_4=0\) and \(x_2-3x_3+2x_4=0\). For any values of \(x_3\) and \(x_4\), the
second equation determines \(x_2\) and then the first determines \(x_1\); so \(x_3\) and \(x_4\)
are free. Let \(x_3=\lambda\) and \(x_4=\mu\). Then \(x_2=3x_3-2x_4=3\lambda-2\mu\) and
\(x_1=x_2-2x_3-x_4=\lambda-3\mu\). Hence
\[
x=\lambda(1,3,1,0)'+\mu(-3,-2,0,1)'\qquad(\lambda,\mu\in\R)
\]
(One more step, \(R_1\leftarrow R_1+R_2\), gives the RREF, with rows \((1,0,-1,3)\) and
\((0,1,-3,2)\), from which the same solutions are read off.)
\end{solution}

\begin{problem}\label{r2-ex-two-systems}
Solve the equations
\[
\begin{pmatrix}1&-3&-2\\-1&2&1\\-1&2&2\end{pmatrix}\begin{pmatrix}x_1\\x_2\\x_3\end{pmatrix}=\begin{pmatrix}1\\-1\\-2\end{pmatrix}
\qquad\text{and}\qquad
\begin{pmatrix}1&-2&0\\1&-2&-1\\1&-1&1\end{pmatrix}\begin{pmatrix}x_1\\x_2\\x_3\end{pmatrix}=\begin{pmatrix}3\\1\\4\end{pmatrix}
\]
by Gaussian elimination.
\end{problem}
\begin{solution}
First system. Forward elimination:
\[
\begin{gathered}
\left(\begin{array}{ccc|c}1&-3&-2&1\\-1&2&1&-1\\-1&2&2&-2\end{array}\right)
\xrightarrow[R_3\leftarrow R_3+R_1]{R_2\leftarrow R_2+R_1}
\left(\begin{array}{ccc|c}1&-3&-2&1\\0&-1&-1&0\\0&-1&0&-1\end{array}\right)
\\
\xrightarrow{R_2\leftarrow-R_2}
\left(\begin{array}{ccc|c}1&-3&-2&1\\0&1&1&0\\0&-1&0&-1\end{array}\right)
\xrightarrow{R_3\leftarrow R_3+R_2}
\left(\begin{array}{ccc|c}1&-3&-2&1\\0&1&1&0\\0&0&1&-1\end{array}\right)
\end{gathered}
\]
Backward elimination:
\[
\xrightarrow[R_1\leftarrow R_1+2R_3]{R_2\leftarrow R_2-R_3}
\left(\begin{array}{ccc|c}1&-3&0&-1\\0&1&0&1\\0&0&1&-1\end{array}\right)
\xrightarrow{R_1\leftarrow R_1+3R_2}
\left(\begin{array}{ccc|c}1&0&0&2\\0&1&0&1\\0&0&1&-1\end{array}\right)
\]
By Proposition~\ref{r2-prop-solutions}, the unique solution is \((x_1,x_2,x_3)=(2,1,-1)\).
Check: \(2-3+2=1\), \(-2+2-1=-1\), \(-2+2-2=-2\).

Second system. Forward elimination; after the first step the entry of row \(2\) in column \(2\)
is \(0\), so rows \(2\) and \(3\) are interchanged (Step 2):
\[
\begin{gathered}
\left(\begin{array}{ccc|c}1&-2&0&3\\1&-2&-1&1\\1&-1&1&4\end{array}\right)
\xrightarrow[R_3\leftarrow R_3-R_1]{R_2\leftarrow R_2-R_1}
\left(\begin{array}{ccc|c}1&-2&0&3\\0&0&-1&-2\\0&1&1&1\end{array}\right)
\\
\xrightarrow{R_2\leftrightarrow R_3}
\left(\begin{array}{ccc|c}1&-2&0&3\\0&1&1&1\\0&0&-1&-2\end{array}\right)
\xrightarrow{R_3\leftarrow-R_3}
\left(\begin{array}{ccc|c}1&-2&0&3\\0&1&1&1\\0&0&1&2\end{array}\right)
\end{gathered}
\]
Backward elimination:
\[
\xrightarrow{R_2\leftarrow R_2-R_3}
\left(\begin{array}{ccc|c}1&-2&0&3\\0&1&0&-1\\0&0&1&2\end{array}\right)
\xrightarrow{R_1\leftarrow R_1+2R_2}
\left(\begin{array}{ccc|c}1&0&0&1\\0&1&0&-1\\0&0&1&2\end{array}\right)
\]
The unique solution is \((x_1,x_2,x_3)=(1,-1,2)\). Check: \(1+2+0=3\), \(1+2-2=1\),
\(1+1+2=4\).
\end{solution}

\begin{problem}\label{r2-ex-elementary}
Let \(A=\begin{pmatrix}1&2\\3&4\\5&6\end{pmatrix}\). Find an elementary matrix \(E\) with
\(B=EA\) when:
(a) \(B\) is obtained from \(A\) by interchanging the first and third rows;
(b) \(B\) is obtained from \(A\) by multiplying the second row by \(7\);
(c) \(B\) is obtained from \(A\) by adding five times the second row to the third row.
\end{problem}
\begin{solution}
Perform each operation on \(I_3\) (Theorem~\ref{r2-thm-elementary}(a)):
\[
P_{13}=\begin{pmatrix}0&0&1\\0&1&0\\1&0&0\end{pmatrix},\quad
D_2(7)=\begin{pmatrix}1&0&0\\0&7&0\\0&0&1\end{pmatrix},\quad
T_{32}(5)=\begin{pmatrix}1&0&0\\0&1&0\\0&5&1\end{pmatrix}
\]
and
\[
P_{13}A=\begin{pmatrix}5&6\\3&4\\1&2\end{pmatrix},\quad
D_2(7)A=\begin{pmatrix}1&2\\21&28\\5&6\end{pmatrix},\quad
T_{32}(5)A=\begin{pmatrix}1&2\\3&4\\20&26\end{pmatrix}
\]
\end{solution}

\begin{problem}\label{r2-ex-order}
Let \(A=\begin{pmatrix}1&2&3\\0&1&2\\1&0&1\end{pmatrix}\).
\begin{enumerate}
\item Find the matrix \(E\) such that \(EA\) is obtained from \(A\) by first interchanging rows
\(1\) and \(2\) and then multiplying row \(2\) by \(2\). Does the order of the two operations
matter?
\item Find the matrix \(E\) such that \(EA\) is obtained from \(A\) by first interchanging rows
\(1\) and \(2\) and then adding \(3\) times row \(1\) to row \(2\). Does the order matter?
\item Do elementary matrices commute?
\end{enumerate}
\end{problem}
\begin{solution}
By Theorem~\ref{r2-thm-elementary}, the operation done first is the factor next to \(A\).

(a) \(E=D_2(2)P_{12}=\begin{pmatrix}1&0&0\\0&2&0\\0&0&1\end{pmatrix}\begin{pmatrix}0&1&0\\1&0&0\\0&0&1\end{pmatrix}=\begin{pmatrix}0&1&0\\2&0&0\\0&0&1\end{pmatrix}\),
and
\[
D_2(2)P_{12}A=\begin{pmatrix}0&1&2\\2&4&6\\1&0&1\end{pmatrix}\ne\begin{pmatrix}0&2&4\\1&2&3\\1&0&1\end{pmatrix}=P_{12}D_2(2)A,
\]
since \(P_{12}D_2(2)=\left(\begin{smallmatrix}0&2&0\\1&0&0\\0&0&1\end{smallmatrix}\right)\). The order
matters.

(b) \(E=T_{21}(3)P_{12}=\left(\begin{smallmatrix}0&1&0\\1&3&0\\0&0&1\end{smallmatrix}\right)\), while
\(P_{12}T_{21}(3)=\left(\begin{smallmatrix}3&1&0\\1&0&0\\0&0&1\end{smallmatrix}\right)\). Accordingly
\(EA=\left(\begin{smallmatrix}0&1&2\\1&5&9\\1&0&1\end{smallmatrix}\right)\) but
\(P_{12}T_{21}(3)A=\left(\begin{smallmatrix}3&7&11\\1&2&3\\1&0&1\end{smallmatrix}\right)\): the order
matters.

(c) No: (a) and (b) give pairs of elementary matrices that do not commute.
\end{solution}

\begin{problem}\label{r2-ex-product-elementary}
Show that the product of two elementary matrices is, in general, not elementary.
\end{problem}
\begin{solution}
For example,
\[
P_{13}D_2(7)=\begin{pmatrix}0&0&1\\0&1&0\\1&0&0\end{pmatrix}\begin{pmatrix}1&0&0\\0&7&0\\0&0&1\end{pmatrix}
=\begin{pmatrix}0&0&1\\0&7&0\\1&0&0\end{pmatrix}
\]
This is not elementary: a matrix \(D_r(c)\) or \(T_{rs}(c)\) differs from \(I_3\) in one entry
only, while this product differs from \(I_3\) in five entries; and every \(P_{rs}\) has all
entries \(0\) or \(1\), while this product has the entry \(7\).
\end{solution}

\begin{problem}\label{r2-ex-inverse}
Use row reduction of \([A\mid I_3]\) to find the inverse of
\(A=\begin{pmatrix}0&1&1\\1&1&1\\3&4&3\end{pmatrix}\).
\end{problem}
\begin{solution}
\[
\begin{gathered}
\left(\begin{array}{ccc|ccc}0&1&1&1&0&0\\1&1&1&0&1&0\\3&4&3&0&0&1\end{array}\right)
\xrightarrow{R_1\leftrightarrow R_2}
\left(\begin{array}{ccc|ccc}1&1&1&0&1&0\\0&1&1&1&0&0\\3&4&3&0&0&1\end{array}\right)
\xrightarrow{R_3\leftarrow R_3-3R_1}
\left(\begin{array}{ccc|ccc}1&1&1&0&1&0\\0&1&1&1&0&0\\0&1&0&0&-3&1\end{array}\right)
\\
\xrightarrow[R_3\leftarrow -R_3]{R_3\leftarrow R_3-R_2}
\left(\begin{array}{ccc|ccc}1&1&1&0&1&0\\0&1&1&1&0&0\\0&0&1&1&3&-1\end{array}\right)
\xrightarrow[R_2\leftarrow R_2-R_3]{R_1\leftarrow R_1-R_3}
\left(\begin{array}{ccc|ccc}1&1&0&-1&-2&1\\0&1&0&0&-3&1\\0&0&1&1&3&-1\end{array}\right)
\\
\xrightarrow{R_1\leftarrow R_1-R_2}
\left(\begin{array}{ccc|ccc}1&0&0&-1&1&0\\0&1&0&0&-3&1\\0&0&1&1&3&-1\end{array}\right)
\end{gathered}
\]
Hence
\(A^{-1}=\left(\begin{smallmatrix}-1&1&0\\0&-3&1\\1&3&-1\end{smallmatrix}\right)\). Check:
\(\row_3(A)A^{-1}=(3,4,3)A^{-1}=(-3+3,\ 3-12+9,\ 4-3)=(0,0,1)\), and similarly for the other
rows.
\end{solution}

\begin{problem}\label{r2-ex-two-by-two}
Do \(A=\begin{pmatrix}1&2&1\\2&4&3\\3&6&4\end{pmatrix}\) and
\(B=\begin{pmatrix}1&3&1\\2&5&1\\1&2&1\end{pmatrix}\) have inverses? Find the inverse when it
exists.
\end{problem}
\begin{solution}
(\(A\)) Row-reduce \(A\):
\[
\begin{pmatrix}1&2&1\\2&4&3\\3&6&4\end{pmatrix}
\xrightarrow[R_3\leftarrow R_3-3R_1]{R_2\leftarrow R_2-2R_1}
\begin{pmatrix}1&2&1\\0&0&1\\0&0&1\end{pmatrix}
\xrightarrow{R_3\leftarrow R_3-R_2}
\begin{pmatrix}1&2&1\\0&0&1\\0&0&0\end{pmatrix}
\]
A zero row appears, so the RREF of \(A\) is not \(I_3\), and \(A\) has no inverse
(Theorem~\ref{r2-thm-invertible}). (Indeed row \(3\) of \(A\) is the sum of rows \(1\) and \(2\).)

(\(B\)) Row-reduce \([B\mid I_3]\):
\[
\left(\begin{array}{ccc|ccc}1&3&1&1&0&0\\2&5&1&0&1&0\\1&2&1&0&0&1\end{array}\right)
\xrightarrow[R_3\leftarrow R_3-R_1]{R_2\leftarrow R_2-2R_1}
\left(\begin{array}{ccc|ccc}1&3&1&1&0&0\\0&-1&-1&-2&1&0\\0&-1&0&-1&0&1\end{array}\right)
\]
\[
\xrightarrow{R_2\leftarrow -R_2}
\left(\begin{array}{ccc|ccc}1&3&1&1&0&0\\0&1&1&2&-1&0\\0&-1&0&-1&0&1\end{array}\right)
\xrightarrow{R_3\leftarrow R_3+R_2}
\left(\begin{array}{ccc|ccc}1&3&1&1&0&0\\0&1&1&2&-1&0\\0&0&1&1&-1&1\end{array}\right)
\]
\[
\xrightarrow[R_1\leftarrow R_1-R_3]{R_2\leftarrow R_2-R_3}
\left(\begin{array}{ccc|ccc}1&3&0&0&1&-1\\0&1&0&1&0&-1\\0&0&1&1&-1&1\end{array}\right)
\xrightarrow{R_1\leftarrow R_1-3R_2}
\left(\begin{array}{ccc|ccc}1&0&0&-3&1&2\\0&1&0&1&0&-1\\0&0&1&1&-1&1\end{array}\right)
\]
The left block is \(I_3\), so \(B\) is invertible and
\(B^{-1}=\left(\begin{smallmatrix}-3&1&2\\1&0&-1\\1&-1&1\end{smallmatrix}\right)\)
(Theorem~\ref{r2-thm-invertible}). Check: the first row of \(BB^{-1}\) is
\((-3+3+1,\ 1+0-1,\ 2-3+1)=(1,0,0)\), and similarly for the other rows.
\end{solution}

\begin{problem}\label{r2-ex-inverse-props}
Let \(A\in\R^{n\times n}\) be invertible. Show that:
\begin{enumerate}
\item \((\lambda A)^{-1}=\frac1\lambda A^{-1}\) (\(\lambda\ne0\));
\item if \(A\) and \(B\) commute and \(B\) is invertible, then \(A\) and \(B^{-1}\) commute.
\end{enumerate}
\end{problem}
\begin{solution}
(a) We check both products (Definition~\ref{r2-def-inverse}):
\((\lambda A)(\frac1\lambda A^{-1})=AA^{-1}=I\) and \((\frac1\lambda A^{-1})(\lambda A)=A^{-1}A=I\).
(b) Multiply \(AB=BA\) on the left and on the right by \(B^{-1}\):
\(B^{-1}ABB^{-1}=B^{-1}BAB^{-1}\), that is, \(B^{-1}A=AB^{-1}\).
\end{solution}

\begin{problem}\label{r2-ex-nilpotent}
Let \(A\in\R^{n\times n}\).
\begin{enumerate}
\item If \(A^2=0\), show that \(I-A\) is invertible.
\item If \(A^3=0\), show that \(I-A\) is invertible.
\item In general, if \(A^k=0\) for some positive integer \(k\), show that \(I-A\) is invertible.
\item Suppose that \(A^2+2A+I=0\). Show that \(A\) is invertible.
\item Suppose that \(A^3-A+I=0\). Show that \(A\) is invertible.
\end{enumerate}
\end{problem}
\begin{solution}
(a)--(c) For every positive integer \(k\), by Theorem~\ref{r1-thm-rules},
\[
(I-A)(I+A+A^2+\dots+A^{k-1})=(I+A+\dots+A^{k-1})-(A+A^2+\dots+A^k)=I-A^k,
\]
and in the same way \((I+A+\dots+A^{k-1})(I-A)=I-A^k\). If \(A^k=0\), both products equal
\(I\), so \(I-A\) is invertible with inverse \(I+A+\dots+A^{k-1}\). Parts (a) and (b) are the
cases \(k=2\) and \(k=3\): \((I-A)^{-1}=I+A\) and \(I+A+A^2\), respectively.

(d) \(A(-A-2I)=(-A-2I)A=-A^2-2A=I\), since \(A^2+2A+I=0\). So \(A^{-1}=-A-2I\).

(e) \(A(-A^2+I)=(-A^2+I)A=-A^3+A=I\), since \(A^3-A+I=0\). So \(A^{-1}=I-A^2\).
\end{solution}

\begin{problem}\label{r2-ex-inverses-b}
Find the inverses of \(\begin{pmatrix}a&b\\c&d\end{pmatrix}\) (\(ad-bc\ne0\)) and of
\(\begin{pmatrix}1&1&2\\3&2&3\\3&2&4\end{pmatrix}\).
\end{problem}
\begin{solution}
(\(2\times2\)) Put \(\Delta=ad-bc\). If \(a\ne0\),
\[
\begin{gathered}
\left(\begin{array}{cc|cc}a&b&1&0\\c&d&0&1\end{array}\right)
\xrightarrow{R_1\leftarrow\frac1aR_1}
\left(\begin{array}{cc|cc}1&\frac ba&\frac1a&0\\c&d&0&1\end{array}\right)
\xrightarrow{R_2\leftarrow R_2-cR_1}
\left(\begin{array}{cc|cc}1&\frac ba&\frac1a&0\\0&\frac\Delta a&-\frac ca&1\end{array}\right)
\\
\xrightarrow{R_2\leftarrow\frac a\Delta R_2}
\left(\begin{array}{cc|cc}1&\frac ba&\frac1a&0\\0&1&-\frac c\Delta&\frac a\Delta\end{array}\right)
\xrightarrow{R_1\leftarrow R_1-\frac baR_2}
\left(\begin{array}{cc|cc}1&0&\frac d\Delta&-\frac b\Delta\\0&1&-\frac c\Delta&\frac a\Delta\end{array}\right)
\end{gathered}
\]
using \(\frac1a+\frac{bc}{a\Delta}=\frac{\Delta+bc}{a\Delta}=\frac d\Delta\). The result
\[
\begin{pmatrix}a&b\\c&d\end{pmatrix}^{-1}=\frac1{ad-bc}\begin{pmatrix}d&-b\\-c&a\end{pmatrix}
\]
also holds when \(a=0\): in all cases
\(\left(\begin{smallmatrix}a&b\\c&d\end{smallmatrix}\right)\left(\begin{smallmatrix}d&-b\\-c&a\end{smallmatrix}\right)
=\left(\begin{smallmatrix}ad-bc&0\\0&ad-bc\end{smallmatrix}\right)
=\left(\begin{smallmatrix}d&-b\\-c&a\end{smallmatrix}\right)\left(\begin{smallmatrix}a&b\\c&d\end{smallmatrix}\right)\).

(\(3\times3\))
\[
\begin{gathered}
\left(\begin{array}{ccc|ccc}1&1&2&1&0&0\\3&2&3&0&1&0\\3&2&4&0&0&1\end{array}\right)
\xrightarrow[R_3\leftarrow R_3-3R_1]{R_2\leftarrow R_2-3R_1}
\left(\begin{array}{ccc|ccc}1&1&2&1&0&0\\0&-1&-3&-3&1&0\\0&-1&-2&-3&0&1\end{array}\right)
\\
\xrightarrow{R_2\leftarrow-R_2}
\left(\begin{array}{ccc|ccc}1&1&2&1&0&0\\0&1&3&3&-1&0\\0&-1&-2&-3&0&1\end{array}\right)
\xrightarrow{R_3\leftarrow R_3+R_2}
\left(\begin{array}{ccc|ccc}1&1&2&1&0&0\\0&1&3&3&-1&0\\0&0&1&0&-1&1\end{array}\right)
\\
\xrightarrow[R_1\leftarrow R_1-2R_3]{R_2\leftarrow R_2-3R_3}
\left(\begin{array}{ccc|ccc}1&1&0&1&2&-2\\0&1&0&3&2&-3\\0&0&1&0&-1&1\end{array}\right)
\xrightarrow{R_1\leftarrow R_1-R_2}
\left(\begin{array}{ccc|ccc}1&0&0&-2&0&1\\0&1&0&3&2&-3\\0&0&1&0&-1&1\end{array}\right)
\end{gathered}
\]
So the inverse is
\(\left(\begin{smallmatrix}-2&0&1\\3&2&-3\\0&-1&1\end{smallmatrix}\right)\). Check: the first
row \((1,1,2)\) of the matrix times the inverse is \((-2+3+0,\ 0+2-2,\ 1-3+2)=(1,0,0)\), and
similarly for the other rows.
\end{solution}

\begin{problem}\label{r2-ex-normal-equations}
For the four units of Example~\ref{r1-exm-data} a response \(y_i\) is also measured:
\(y=(9,3,5,3)'\). With \(X'X\) as in Example~\ref{r2-exm-xtx-inverse}, compute \(X'y\) and solve
the normal equations \(X'X\hat\beta=X'y\) for the least-squares estimate
\(\hat\beta=(\hat\beta_0,\hat\beta_1,\hat\beta_2)'\) by Gauss--Jordan elimination; check the
answer with \(\hat\beta=(X'X)^{-1}X'y\).
\end{problem}
\begin{solution}
The \(j\)th entry of \(X'y\) is \(\col_j(X)'y\) (Theorem~\ref{r1-thm-views}(d)):
\[
X'y=\begin{pmatrix}\sum_iy_i\\\sum_ix_{i1}y_i\\\sum_ix_{i2}y_i\end{pmatrix}
=\begin{pmatrix}9+3+5+3\\2\cdot9+0\cdot3+2\cdot5+0\cdot3\\3\cdot9+3\cdot3+1\cdot5+1\cdot3\end{pmatrix}
=\begin{pmatrix}20\\28\\44\end{pmatrix}
\]
As in Example~\ref{r2-exm-xtx-inverse}, postpone the scaling to the end:
\[
\begin{gathered}
\left(\begin{array}{ccc|c}4&4&8&20\\4&8&8&28\\8&8&20&44\end{array}\right)
\xrightarrow[R_3\leftarrow R_3-2R_1]{R_2\leftarrow R_2-R_1}
\left(\begin{array}{ccc|c}4&4&8&20\\0&4&0&8\\0&0&4&4\end{array}\right)
\xrightarrow{R_1\leftarrow R_1-2R_3}
\left(\begin{array}{ccc|c}4&4&0&12\\0&4&0&8\\0&0&4&4\end{array}\right)
\\
\xrightarrow{R_1\leftarrow R_1-R_2}
\left(\begin{array}{ccc|c}4&0&0&4\\0&4&0&8\\0&0&4&4\end{array}\right)
\xrightarrow[R_3\leftarrow\frac14R_3]{R_1\leftarrow\frac14R_1,\ R_2\leftarrow\frac14R_2}
\left(\begin{array}{ccc|c}1&0&0&1\\0&1&0&2\\0&0&1&1\end{array}\right)
\end{gathered}
\]
So \(\hat\beta=(1,2,1)'\) is the unique solution. Check:
\[
(X'X)^{-1}X'y=\frac14\begin{pmatrix}6&-1&-2\\-1&1&0\\-2&0&1\end{pmatrix}\begin{pmatrix}20\\28\\44\end{pmatrix}
=\frac14\begin{pmatrix}120-28-88\\-20+28\\-40+44\end{pmatrix}=\frac14\begin{pmatrix}4\\8\\4\end{pmatrix}=\begin{pmatrix}1\\2\\1\end{pmatrix}
\]
\end{solution}

\begin{problem}\label{r2-ex-which-b}
Let \(A\) be the coefficient matrix of Example~\ref{r2-exm-inconsistent}.
\begin{enumerate}
\item For which \(b=(b_1,b_2,b_3)'\) is \(Ax=b\) consistent?
\item Find all solutions of \(Ax=b\) for \(b=(1,2,4)'\).
\end{enumerate}
\end{problem}
\begin{solution}
(a) Row-reduce \([A\mid b]\), keeping \(b_1,b_2,b_3\) as letters:
\[
\begin{gathered}
\left(\begin{array}{cccc|c}1&-2&1&2&b_1\\1&-1&-1&1&b_2\\1&1&-5&-1&b_3\end{array}\right)
\xrightarrow[R_3\leftarrow R_3-R_1]{R_2\leftarrow R_2-R_1}
\left(\begin{array}{cccc|c}1&-2&1&2&b_1\\0&1&-2&-1&b_2-b_1\\0&3&-6&-3&b_3-b_1\end{array}\right)
\\
\xrightarrow{R_3\leftarrow R_3-3R_2}
\left(\begin{array}{cccc|c}1&-2&1&2&b_1\\0&1&-2&-1&b_2-b_1\\0&0&0&0&b_3-3b_2+2b_1\end{array}\right)
\xrightarrow{R_1\leftarrow R_1+2R_2}
\left(\begin{array}{cccc|c}1&0&-3&0&2b_2-b_1\\0&1&-2&-1&b_2-b_1\\0&0&0&0&b_3-3b_2+2b_1\end{array}\right)
\end{gathered}
\]
The left part of the last matrix is the RREF of \(A\), with \(r=2\). By
Proposition~\ref{r2-prop-solutions}(a), \(Ax=b\) is consistent if and only if
\(2b_1-3b_2+b_3=0\). (In Example~\ref{r2-exm-inconsistent}, \(2-6+3=-1\ne0\).)

(b) For \(b=(1,2,4)'\), \(2-6+4=0\), so the system is consistent, and the last matrix becomes
\[
\left(\begin{array}{cccc|c}1&0&-3&0&3\\0&1&-2&-1&1\\0&0&0&0&0\end{array}\right)
\]
The free unknowns are \(x_3=s\), \(x_4=t\), and by Proposition~\ref{r2-prop-solutions}(b)
\(x_1=3+3s\), \(x_2=1+2s+t\). So
\[
x=\begin{pmatrix}3\\1\\0\\0\end{pmatrix}+s\begin{pmatrix}3\\2\\1\\0\end{pmatrix}+t\begin{pmatrix}0\\1\\0\\1\end{pmatrix}\qquad(s,t\in\R),
\]
infinitely many solutions (Theorem~\ref{thm:r2-solution-structure}(b)): the particular solution
\((3,1,0,0)'\) plus the solutions of \(Ax=0\) (Theorem~\ref{thm:r2-solution-structure}(a)).
Check: \(A(3,1,0,0)'=(3-2,\ 3-1,\ 3+1)'=(1,2,4)'\).
\end{solution}

\begin{problem}\label{r2-ex-system-drill}
Let
\[
A=\begin{pmatrix}1&3&5&-1\\1&2&2&0\\0&1&3&-1\\0&0&1&-1\end{pmatrix},\qquad
b=\begin{pmatrix}6\\4\\2\\0\end{pmatrix},\qquad c=\begin{pmatrix}6\\4\\3\\0\end{pmatrix}
\]
\begin{enumerate}
\item Find all solutions of \(Ax=b\).
\item Show that \(Ax=c\) has no solution.
\end{enumerate}
\end{problem}
\begin{solution}
(a) Row-reduce \([A\mid b]\):
\[
\begin{gathered}
\left(\begin{array}{cccc|c}1&3&5&-1&6\\1&2&2&0&4\\0&1&3&-1&2\\0&0&1&-1&0\end{array}\right)
\xrightarrow{R_2\leftarrow R_2-R_1}
\left(\begin{array}{cccc|c}1&3&5&-1&6\\0&-1&-3&1&-2\\0&1&3&-1&2\\0&0&1&-1&0\end{array}\right)
\xrightarrow{R_2\leftarrow-R_2}
\left(\begin{array}{cccc|c}1&3&5&-1&6\\0&1&3&-1&2\\0&1&3&-1&2\\0&0&1&-1&0\end{array}\right)
\\
\xrightarrow{R_3\leftarrow R_3-R_2}
\left(\begin{array}{cccc|c}1&3&5&-1&6\\0&1&3&-1&2\\0&0&0&0&0\\0&0&1&-1&0\end{array}\right)
\xrightarrow{R_3\leftrightarrow R_4}
\left(\begin{array}{cccc|c}1&3&5&-1&6\\0&1&3&-1&2\\0&0&1&-1&0\\0&0&0&0&0\end{array}\right)
\\
\xrightarrow[R_1\leftarrow R_1-5R_3]{R_2\leftarrow R_2-3R_3}
\left(\begin{array}{cccc|c}1&3&0&4&6\\0&1&0&2&2\\0&0&1&-1&0\\0&0&0&0&0\end{array}\right)
\xrightarrow{R_1\leftarrow R_1-3R_2}
\left(\begin{array}{cccc|c}1&0&0&-2&0\\0&1&0&2&2\\0&0&1&-1&0\\0&0&0&0&0\end{array}\right)
\end{gathered}
\]
After the third step the entry of row \(3\) in column \(3\) is \(0\), so Step 2 interchanges rows
\(3\) and \(4\). The left part of the last matrix is the RREF \(R\) of \(A\), with \(r=3\) and pivot
columns \(1,2,3\); the unknown \(x_4\) is free. The last matrix is \([R\mid z]\) with
\(z=(0,2,0,0)'\). As \(z_4=0\), the system is consistent (Proposition~\ref{r2-prop-solutions}(a)),
and with \(x_4=t\), Proposition~\ref{r2-prop-solutions}(b) gives \(x_1=2t\), \(x_2=2-2t\), \(x_3=t\):
\[
x=\begin{pmatrix}0\\2\\0\\0\end{pmatrix}+t\begin{pmatrix}2\\-2\\1\\1\end{pmatrix}\qquad(t\in\R),
\]
infinitely many solutions (Theorem~\ref{thm:r2-solution-structure}(a),(b)). Check:
\(A(0,2,0,0)'=2\col_2(A)=(6,4,2,0)'=b\) and
\(A(2,-2,1,1)'=(2-6+5-1,\ 2-4+2,\ -2+3-1,\ 1-1)'=0\).

(b) Row-reduce \([A\mid c]\) by the same first three operations:
\[
\begin{gathered}
\left(\begin{array}{cccc|c}1&3&5&-1&6\\1&2&2&0&4\\0&1&3&-1&3\\0&0&1&-1&0\end{array}\right)
\xrightarrow{R_2\leftarrow R_2-R_1}
\left(\begin{array}{cccc|c}1&3&5&-1&6\\0&-1&-3&1&-2\\0&1&3&-1&3\\0&0&1&-1&0\end{array}\right)
\\
\xrightarrow{R_2\leftarrow-R_2}
\left(\begin{array}{cccc|c}1&3&5&-1&6\\0&1&3&-1&2\\0&1&3&-1&3\\0&0&1&-1&0\end{array}\right)
\xrightarrow{R_3\leftarrow R_3-R_2}
\left(\begin{array}{cccc|c}1&3&5&-1&6\\0&1&3&-1&2\\0&0&0&0&1\\0&0&1&-1&0\end{array}\right)
\end{gathered}
\]
The third row of the last matrix reads \((0\ 0\ 0\ 0\mid1)\), that is, \(0=1\); so \(Ax=c\) has no
solution (Theorem~\ref{thm:r2-consistency}). Indeed \(\row_3(A)=\row_1(A)-\row_2(A)\), while
\(c_3=3\ne2=c_1-c_2\).
\end{solution}

\begin{problem}\label{r2-ex-parameters}
For which values of \(k\) and \(h\) does the system
\[
x_1+x_2+x_3=1,\qquad x_1+2x_2+3x_3=3,\qquad x_1+3x_2+kx_3=h
\]
have no solution, exactly one solution, or infinitely many solutions? Give the solutions when
they exist.
\end{problem}
\begin{solution}
Forward elimination:
\[
\left(\begin{array}{ccc|c}1&1&1&1\\1&2&3&3\\1&3&k&h\end{array}\right)
\xrightarrow[R_3\leftarrow R_3-R_1]{R_2\leftarrow R_2-R_1}
\left(\begin{array}{ccc|c}1&1&1&1\\0&1&2&2\\0&2&k-1&h-1\end{array}\right)
\xrightarrow{R_3\leftarrow R_3-2R_2}
\left(\begin{array}{ccc|c}1&1&1&1\\0&1&2&2\\0&0&k-5&h-5\end{array}\right)
\]
\begin{itemize}
\item \(k\ne5\): \(R_3\leftarrow\frac1{k-5}R_3\) gives an echelon form with pivots in columns
\(1,2,3\) and none in the last column. So there is exactly one solution
(Theorems~\ref{thm:r2-consistency} and~\ref{thm:r2-solution-structure}(b)): back substitution
gives \(x_3=\frac{h-5}{k-5}\), \(x_2=2-2x_3\), \(x_1=1-x_2-x_3=-1+x_3\), that is,
\(x=(-1,2,0)'+\frac{h-5}{k-5}(1,-2,1)'\).
\item \(k=5\), \(h\ne5\): row \(3\) reads \(0=h-5\ne0\), so there is no solution.
\item \(k=5\), \(h=5\): row \(3\) is \(0\), and the first two rows form an echelon form with
pivots in columns \(1,2\); \(x_3=t\) is free, so there are infinitely many solutions:
\(x_2=2-2t\), \(x_1=1-x_2-x_3=-1+t\), that is, \(x=(-1,2,0)'+t(1,-2,1)'\) (\(t\in\R\)).
\end{itemize}
\end{solution}

\clearpage
\def\unitnumber{3}
\setcounter{theorem}{0}\setcounter{problem}{0}
\section*{Chapter 3. Vector spaces, subspaces, span, and bases}

\subsection*{Vector spaces}

\begin{definition}\label{r3-def-vector-space}
A (real) \textbf{vector space} is a set \(V\), whose elements are called \textbf{vectors}, with
an \textbf{addition} that gives \(u+v\in V\) for \(u,v\in V\) and a \textbf{scalar
multiplication} that gives \(cv\in V\) for \(c\in\R\) and \(v\in V\), such that for all
\(u,v,w\in V\) and \(c,d\in\R\):
\begin{enumerate}[label=(V\arabic*)]
\item \(u+v=v+u\) and (V2) \(u+(v+w)=(u+v)+w\);
\setcounter{enumi}{2}
\item there is a \textbf{zero vector} \(0\in V\) with \(v+0=v\) for every \(v\in V\);
\item for each \(v\in V\) there is \(-v\in V\) with \(v+(-v)=0\);
\item \(1v=v\) and (V6) \((cd)v=c(dv)\);
\setcounter{enumi}{6}
\item \(c(u+v)=cu+cv\) and (V8) \((c+d)v=cv+dv\).
\end{enumerate}
A vector \(c_1v_1+\dots+c_kv_k\) (\(c_i\in\R\)) is a \textbf{linear combination} of
\(v_1,\dots,v_k\).
\end{definition}
From these rules, \(0v=0\) and \((-1)v=-v\): indeed \(0v=(0+0)v=0v+0v\), and adding \(-(0v)\)
to both sides gives \(0=0v\); then \(v+(-1)v=(1+(-1))v=0v=0\), and \(-v\) is the only vector \(w\)
with \(v+w=0\), since such a \(w\) equals \(w+(v+(-v))=(w+v)+(-v)=0+(-v)=-v\). Likewise \(c0=0\), and
\(cv=0\) only if \(c=0\) or \(v=0\) (if \(c\ne0\), then \(v=c^{-1}(cv)=c^{-1}0=0\)).

\begin{example}\label{r3-exm-spaces}
\(\R^n\), \(\R^{1\times n}\) and \(\R^{m\times n}\), with the entrywise sum and scalar multiple
of Definition~\ref{r1-def-matrix}, are vector spaces: (V1)--(V8) hold entry by entry, the zero vector is the zero
matrix, and \(-A=(-1)A\). In this course every vector space is one of these or a subspace of
one of them.
\end{example}

\subsection*{Subspaces}

\begin{definition}\label{r3-def-subspace}
A \textbf{subspace} of a vector space \(V\) is a subset \(W\subseteq V\) that is itself a
vector space with the addition and scalar multiplication of \(V\).
\end{definition}

\begin{theorem}\label{r3-thm-subspace-test}
For a subset \(W\) of a vector space \(V\) the following are equivalent:
\begin{enumerate}
\item \(W\) is a subspace of \(V\);
\item \(0\in W\), and \(u+v\in W\) and \(cu\in W\) for all \(u,v\in W\) and \(c\in\R\)
(\(W\) is \textbf{closed under addition and scalar multiplication});
\item \(0\in W\), and \(cu+v\in W\) for all \(u,v\in W\) and \(c\in\R\).
\end{enumerate}
\end{theorem}
\begin{proof}
(a)\,\(\Rightarrow\)\,(b): the operations of the vector space \(W\) are those of \(V\) and give
vectors of \(W\); \(W\) has a zero vector, so it contains some \(w\), and then
\(0=0w\in W\). (b)\,\(\Rightarrow\)\,(c): \(cu\in W\), so \(cu+v\in W\).
(c)\,\(\Rightarrow\)\,(a): taking \(v=0\) gives \(cu\in W\), and \(c=1\) gives \(u+v\in W\);
so the operations of \(V\) give operations on \(W\). The rules (V1), (V2), (V5)--(V8) hold in
\(W\) because they hold in \(V\); (V3) holds since \(0\in W\), and (V4) since
\(-w=(-1)w\in W\).
\end{proof}
So, to show that \(W\) is a subspace, check that it contains \(0\) and is closed under the two
operations; to show that it is not, exhibit one failure.

\begin{example}\label{r3-exm-subspaces}
\begin{enumerate}
\item \(\{0\}\) and \(V\) itself are subspaces of \(V\).
\item In \(\R^n\), \(W=\{x:x_1=0\}\) is a subspace: \(0\in W\), and
\((cx+y)_1=cx_1+y_1=0\) for \(x,y\in W\). For \(n\geqslant2\), \(\{x:x_1=1+x_2\}\) is not,
since it does not contain \(0\).
\item In \(\R^2\), \(\{(x,y)':2x-5y=0\}\) is a subspace, since
\(2(cx_1+x_2)-5(cy_1+y_2)=c(2x_1-5y_1)+(2x_2-5y_2)\) for \((x_1,y_1)'\) and \((x_2,y_2)'\) in
it; but for \(k\ne0\) the set
\(\{(x,y)':ax+by=k\}\) is not, since it does not contain \(0\).
\item The symmetric \(n\times n\) matrices form a subspace of \(\R^{n\times n}\): \(0'=0\),
and \((cA+B)'=cA'+B'=cA+B\) for symmetric \(A,B\). So do the \textbf{skew-symmetric}
matrices, those with \(A'=-A\): \(0'=-0\), and \((cA+B)'=cA'+B'=-(cA+B)\) if \(A'=-A\) and
\(B'=-B\).
\item For \(A\in\R^{m\times n}\), the solution set \(\{x\in\R^n:Ax=0\}\) of a homogeneous
system is a subspace of \(\R^n\): \(A0=0\), and \(A(cu+v)=cAu+Av=0\) if \(Au=Av=0\). For
\(b\ne0\), the solutions of \(Ax=b\) do not form a subspace (\(A0=0\ne b\)).
\item In statistics, the \textbf{centered} vectors \(W_0=\{x\in\R^n:\one'x=0\}\), whose
entries add up to \(0\), form a subspace: this is (e) with the \(1\times n\) matrix
\(A=\one'\). For every \(x\), the deviation vector \(Cx\) of Example~\ref{r1-exm-centering} lies in \(W_0\),
since \(\one'Cx=(C\one)'x=0\).
\item The set \(\{(x,y,z)'\in\R^3:xy=0\}\) is not a subspace: \((1,0,0)'\) and
\((0,1,0)'\) lie in it, but their sum \((1,1,0)'\) does not. It is the union of the
subspaces \(\{x=0\}\) and \(\{y=0\}\); a union of subspaces need not be a subspace.
\end{enumerate}
\end{example}

\subsection*{Span}

\begin{definition}\label{r3-def-span}
Let \(v_1,\dots,v_k\) be vectors of a vector space \(V\). Their \textbf{span} is the set of all
their linear combinations,
\[
\Span\{v_1,\dots,v_k\}=\{c_1v_1+\dots+c_kv_k:\ c_1,\dots,c_k\in\R\},
\]
and \(\Span\varnothing=\{0\}\). If \(\Span\{v_1,\dots,v_k\}=W\), we say that \(v_1,\dots,v_k\)
\textbf{span} \(W\).
\end{definition}

\begin{theorem}\label{r3-thm-span}
\(\Span\{v_1,\dots,v_k\}\) is a subspace of \(V\) that contains \(v_1,\dots,v_k\), and it is
contained in every subspace that contains \(v_1,\dots,v_k\): it is the smallest subspace
containing them.
\end{theorem}
\begin{proof}
\(0=0v_1+\dots+0v_k\) lies in the span, and
\(c\sum_ia_iv_i+\sum_ib_iv_i=\sum_i(ca_i+b_i)v_i\); so the span is a subspace
(Theorem~\ref{r3-thm-subspace-test}(c)). It contains \(v_j=1v_j+\sum_{i\ne j}0v_i\). A subspace
that contains \(v_1,\dots,v_k\) contains every \(c_iv_i\) and every sum of these, by closure.
\end{proof}

For vectors \(v_1,\dots,v_k\in\R^m\), put \(V=[v_1\ \cdots\ v_k]\in\R^{m\times k}\). Since
\(Vc=c_1v_1+\dots+c_kv_k\) (Theorem~\ref{r1-thm-views}(b)),
\[
\Span\{v_1,\dots,v_k\}=\{Vc:\ c\in\R^k\},\qquad
b\in\Span\{v_1,\dots,v_k\}\iff\text{the system }Vc=b\text{ is consistent}
\]
So membership in a span is decided by row-reducing \([V\mid b]\)
(Proposition~\ref{r2-prop-solutions}, Theorem~\ref{thm:r2-consistency}). For instance, by
Example~\ref{r2-exm-consistency}, a vector \(y\in\R^3\) lies in the span of the columns of the
matrix \(A\) there if and only if \(2y_1-y_2+y_3=0\).

\begin{example}\label{r3-exm-span}
Let \(W=\Span\{v_1,v_2,v_3\}\subseteq\R^5\), where \(v_1=(1,2,0,3,0)'\),
\(v_2=(0,0,1,4,0)'\) and \(v_3=(0,0,0,0,1)'\). Which \(w=(w_1,\dots,w_5)'\) lie in \(W\)?

\textit{Solution.} Row-reduce \([v_1\ v_2\ v_3\mid w]\), keeping the \(w_i\) as letters:
\[
\left(\begin{array}{ccc|c}1&0&0&w_1\\2&0&0&w_2\\0&1&0&w_3\\3&4&0&w_4\\0&0&1&w_5\end{array}\right)
\xrightarrow[R_4\leftarrow R_4-3R_1]{R_2\leftarrow R_2-2R_1}
\left(\begin{array}{ccc|c}1&0&0&w_1\\0&0&0&w_2-2w_1\\0&1&0&w_3\\0&4&0&w_4-3w_1\\0&0&1&w_5\end{array}\right)
\xrightarrow{R_4\leftarrow R_4-4R_3}
\left(\begin{array}{ccc|c}1&0&0&w_1\\0&0&0&w_2-2w_1\\0&1&0&w_3\\0&0&0&w_4-3w_1-4w_3\\0&0&1&w_5\end{array}\right)
\]
\[
\xrightarrow[R_3\leftrightarrow R_5]{R_2\leftrightarrow R_3}
\left(\begin{array}{ccc|c}1&0&0&w_1\\0&1&0&w_3\\0&0&1&w_5\\0&0&0&w_4-3w_1-4w_3\\0&0&0&w_2-2w_1\end{array}\right)
\]
By Proposition~\ref{r2-prop-solutions}, \(w\in W\) if and only if
\[
w_2=2w_1,\qquad w_4=3w_1+4w_3,
\]
and then \(w=w_1v_1+w_3v_2+w_5v_3\). For example, \((-3,-6,1,-5,2)'\in W\), since
\(-6=2(-3)\) and \(-5=3(-3)+4\cdot1\); indeed it equals \(-3v_1+v_2+2v_3\). But
\((2,4,6,7,8)'\notin W\), since \(3\cdot2+4\cdot6=30\ne7\).
\end{example}

\begin{example}\label{r3-exm-constant}
\(\Span\{\one_n\}=\{c\one_n:c\in\R\}\) is the subspace of constant vectors. Every
\(x\in\R^n\) is the sum of a constant vector and a centered vector:
\[
x=\bar x\one+Cx,\qquad \bar x\one\in\Span\{\one\},\quad Cx\in W_0,
\]
since \(Cx=x-\bar x\one\) (Example~\ref{r1-exm-centering}) and \(Cx\in W_0\)
(Example~\ref{r3-exm-subspaces}(f)). For \(x=(1,3,5,7,9)'\):
\(x=5\one+(-4,-2,0,2,4)'\).
\end{example}

\subsection*{Linear independence}

\begin{definition}\label{r3-def-independence}
Vectors \(v_1,\dots,v_k\) of a vector space are \textbf{linearly dependent} if
\(c_1v_1+\dots+c_kv_k=0\) for some scalars \(c_1,\dots,c_k\), not all \(0\) (a \textbf{linear
relation}). Otherwise they are \textbf{linearly independent}: \(c_1v_1+\dots+c_kv_k=0\) only
for \(c_1=\dots=c_k=0\). The empty list is independent.
\end{definition}

\begin{proposition}\label{r3-prop-independence-test}
Let \(v_1,\dots,v_k\in\R^m\) and \(V=[v_1\ \cdots\ v_k]\in\R^{m\times k}\). Then
\(v_1,\dots,v_k\) are linearly independent if and only if \(Vc=0\) has only the solution
\(c=0\), that is, if and only if every column of the RREF of \(V\) is a pivot column (no free
unknowns). In particular, more than \(m\) vectors of \(\R^m\) are linearly dependent.
\end{proposition}
\begin{proof}
\(Vc=c_1v_1+\dots+c_kv_k\), so the linear relations among the \(v_i\) are the solutions of
\(Vc=0\). This system is consistent (\(c=0\)), so by
Theorem~\ref{thm:r2-solution-structure}(b) \(c=0\) is its only solution exactly when there is no
free unknown. If \(k>m\), there is a free unknown, by Theorem~\ref{r2-thm-homogeneous}(a).
\end{proof}

\begin{proposition}\label{r3-prop-dependence}
Let \(v_1,\dots,v_k\) be vectors of a vector space.
\begin{enumerate}
\item If some \(v_i=0\), or \(v_i=v_j\) for some \(i\ne j\), then \(v_1,\dots,v_k\) are
dependent. A single vector \(v\) is independent if and only if \(v\ne0\).
\item A sublist of an independent list is independent.
\item For \(k\geqslant2\), \(v_1,\dots,v_k\) are dependent if and only if some \(v_j\) is a
linear combination of the others. In particular, two vectors are dependent if and only if one
of them is a multiple of the other.
\end{enumerate}
\end{proposition}
\begin{proof}
(a) \(1\cdot0=0\) and \(v_i-v_j=0\) are linear relations with a coefficient \(1\); and if
\(v\ne0\), then \(cv=0\) forces \(c=0\). (b) A relation among some of the vectors is a
relation among all of them, with coefficient \(0\) for the others. (c) If
\(v_j=\sum_{i\ne j}x_iv_i\), then \(v_j-\sum_{i\ne j}x_iv_i=0\) is a relation with coefficient
\(1\) on \(v_j\). Conversely, if \(\sum_ix_iv_i=0\) with \(x_j\ne0\), then
\(v_j=\sum_{i\ne j}(-x_i/x_j)v_i\).
\end{proof}

\begin{example}\label{r3-exm-dependent}
Let \(a_1,a_2,a_3,a_4\) be the columns of the matrix \(A\) of Example~\ref{r2-exm-gauss-jordan}:
\[
a_1=\begin{pmatrix}1\\1\\2\end{pmatrix},\quad a_2=\begin{pmatrix}1\\1\\1\end{pmatrix},\quad
a_3=\begin{pmatrix}1\\2\\5\end{pmatrix},\quad a_4=\begin{pmatrix}-2\\-3\\-4\end{pmatrix}
\]
Four vectors of \(\R^3\) are dependent (Proposition~\ref{r3-prop-independence-test}). A relation
comes from the solutions \(x=t(-2,3,1,1)'\) of \(Ax=0\) found there:
\[
-2a_1+3a_2+a_3+a_4=0,\qquad\text{so}\qquad a_4=2a_1-3a_2-a_3
\]
On the other hand \(a_1,a_2,a_3\) are independent: the same row operations take
\([a_1\ a_2\ a_3]\) to the first three columns of the RREF found there, which form \(I_3\), so
every column is a pivot column.
\end{example}

\subsection*{Bases}

\begin{definition}\label{r3-def-basis}
Vectors \(v_1,\dots,v_k\) of a vector space \(V\) form a \textbf{basis} of a subspace \(W\) if
they lie in \(W\), are linearly independent, and span \(W\).
\begin{itemize}
\item The unit vectors \(e_1,\dots,e_n\) form a basis of \(\R^n\), the \textbf{standard basis}:
\(x=x_1e_1+\dots+x_ne_n\) for every \(x\in\R^n\), and this is \(0\) only for \(x=0\).
\item Let \(E_{ij}\in\R^{m\times n}\) have \((i,j)\) entry \(1\) and all other entries \(0\)
(two indices; not the elementary matrices \(E_1,E_2,\dots\) of Theorem~\ref{r2-thm-elementary}(c)). The \(mn\) matrices
\(E_{ij}\) form a basis of \(\R^{m\times n}\): \(A=\sum_i\sum_ja_{ij}E_{ij}\), and
\(\sum_i\sum_jc_{ij}E_{ij}=(c_{ij})\) is \(0\) only if every \(c_{ij}=0\).
\item The empty list is a basis of \(\{0\}\).
\end{itemize}
\end{definition}

\begin{theorem}\label{r3-thm-unique-coef}
If \(v_1,\dots,v_k\) form a basis of \(W\), every \(w\in W\) can be written as
\(w=c_1v_1+\dots+c_kv_k\) in exactly one way; the numbers \(c_1,\dots,c_k\) are the
\textbf{coordinates} of \(w\) in this basis.
\end{theorem}
\begin{proof}
Such an expression exists because the \(v_i\) span \(W\). If
\(\sum_ix_iv_i=\sum_iy_iv_i\), then \(\sum_i(x_i-y_i)v_i=0\), and independence gives
\(x_i=y_i\) for every \(i\).
\end{proof}

\begin{example}\label{r3-exm-invertible-basis}
If \(P\in\R^{n\times n}\) is invertible, its columns \(p_1,\dots,p_n\) form a basis of \(\R^n\),
and the coordinates of \(y\in\R^n\) are \(x=P^{-1}y\). Indeed \(Px=\sum_jx_jp_j\); \(Px=0\) only
for \(x=0\) (Theorem~\ref{r2-thm-invertible}), so the columns are independent; and
\(y=P(P^{-1}y)\), so they span. For example, the columns \((1,-1,2)'\), \((-1,1,-1)'\),
\((1,-2,2)'\) of the matrix \(A\) of Example~\ref{r2-exm-inverse-2x2} form a basis of \(\R^3\).
The coordinates of \(y=(1,0,3)'\) are \(A^{-1}y=(3,1,-1)'\), the solution of \(Ax=y\) found in
Example~\ref{r2-exm-unique}:
\(y=3(1,-1,2)'+(-1,1,-1)'-(1,-2,2)'\).
\end{example}

\begin{example}\label{r3-exm-two-of-three}
Let \(u_1=(1,2,0)'\), \(u_2=(1,3,0)'\), \(u_3=(2,5,0)'\), and
\(W=\{(a,b,0)':a,b\in\R\}\), a subspace of \(\R^3\). Any two of \(u_1,u_2,u_3\) form a basis
of \(W\). For \(w=(a,b,0)'\in W\), the equation \(cu_i+du_j=w\) has third entries \(0=0\),
and its first two entries give
\begin{itemize}
\item \(u_1,u_2\): \(c+d=a\), \(2c+3d=b\), so \(d=b-2a\), \(c=3a-b\);
\item \(u_1,u_3\): \(c+2d=a\), \(2c+5d=b\), so \(d=b-2a\), \(c=5a-2b\);
\item \(u_2,u_3\): \(c+2d=a\), \(3c+5d=b\), so \(d=3a-b\), \(c=2b-5a\).
\end{itemize}
In each case a solution exists for every \(w\), so the pair spans \(W\), and it is unique, so
for \(w=0\) only \(c=d=0\): the pair is independent. All three vectors together also span
\(W\), but they are dependent (\(u_1+u_2-u_3=0\)), so they do not form a basis. A subspace has
many bases.
\end{example}

\begin{example}\label{r3-exm-symmetric-basis}
The symmetric \(2\times2\) matrices (for instance \(2\times2\) covariance matrices) form a
subspace of \(\R^{2\times2}\) (Example~\ref{r3-exm-subspaces}(d)). The matrices
\(E_{11}\), \(E_{12}+E_{21}\), \(E_{22}\) form a basis of it: they are symmetric, every
symmetric matrix is
\[
\begin{pmatrix}a&b\\b&d\end{pmatrix}=a\begin{pmatrix}1&0\\0&0\end{pmatrix}+b\begin{pmatrix}0&1\\1&0\end{pmatrix}+d\begin{pmatrix}0&0\\0&1\end{pmatrix},
\]
and this combination is \(0\) only if \(a=b=d=0\).
\end{example}

\begin{example}\label{r3-exm-centered-basis}
In \(W_0=\{x\in\R^3:x_1+x_2+x_3=0\}\), the vectors \(d_1=(1,-1,0)'\) and
\(d_2=(0,1,-1)'\) form a basis. They lie in \(W_0\). They span \(W_0\): for \(x\in W_0\),
\(x_3=-x_1-x_2\), so
\[
x=x_1d_1+(x_1+x_2)d_2=\bigl(x_1,\ -x_1+x_1+x_2,\ -x_1-x_2\bigr)'
\]
They are independent: \(c_1d_1+c_2d_2=(c_1,\ c_2-c_1,\ -c_2)'=0\) forces \(c_1=c_2=0\).
\end{example}

\subsection*{Sums and direct sums of subspaces}

\begin{definition}\label{r3-def-sum}
Let \(U\) and \(W\) be subspaces of a vector space \(V\). Their \textbf{intersection}
\(U\cap W\) is the set of vectors that lie in both \(U\) and \(W\), and their \textbf{sum} is
\[
U+W=\{u+w:\ u\in U,\ w\in W\}
\]
The sum is \textbf{direct}, written \(U\oplus W\), if \(U\cap W=\{0\}\), that is, if \(0\) is the
only vector that \(U\) and \(W\) have in common. Thus \(V=U\oplus W\) means that \(V=U+W\) and
\(U\cap W=\{0\}\).
\end{definition}

\begin{theorem}\label{r3-thm-sum}
Let \(U\) and \(W\) be subspaces of a vector space \(V\).
\begin{enumerate}
\item \(U\cap W\) and \(U+W\) are subspaces of \(V\).
\item \(U+W\) contains \(U\) and \(W\), and it is contained in every subspace of \(V\) that
contains \(U\cup W\): it is the smallest subspace of \(V\) containing \(U\cup W\). Moreover
\(U+W=\Span(U\cup W)\), the set of all linear combinations \(c_1x_1+\dots+c_kx_k\) of finitely
many vectors \(x_1,\dots,x_k\in U\cup W\).
\item If \(U=\Span\{u_1,\dots,u_p\}\) and \(W=\Span\{w_1,\dots,w_q\}\), then
\(U+W=\Span\{u_1,\dots,u_p,w_1,\dots,w_q\}\).
\end{enumerate}
\end{theorem}
\begin{proof}
(a)
\begin{itemize}
\item Both sets contain \(0\): \(0\in U\), \(0\in W\), and \(0=0+0\).
\item Let \(x,y\in U\cap W\) and \(c\in\R\). Then \(cx+y\in U\) and \(cx+y\in W\), since \(U\)
and \(W\) are subspaces; so \(cx+y\in U\cap W\).
\item Let \(x=u_1+w_1\) and \(y=u_2+w_2\) with \(u_1,u_2\in U\), \(w_1,w_2\in W\), and let
\(c\in\R\). Then \(cx+y=(cu_1+u_2)+(cw_1+w_2)\) with \(cu_1+u_2\in U\) and \(cw_1+w_2\in W\);
so \(cx+y\in U+W\).
\item By Theorem~\ref{r3-thm-subspace-test}(c), \(U\cap W\) and \(U+W\) are subspaces.
\end{itemize}
(b)
\begin{itemize}
\item Every \(u\in U\) is \(u=u+0\), and every \(w\in W\) is \(w=0+w\); so \(U\) and \(W\) lie
in \(U+W\).
\item If a subspace \(X\) of \(V\) contains \(U\cup W\), then it contains \(u\) and \(w\) for all
\(u\in U\), \(w\in W\), hence also \(u+w\); so \(U+W\subseteq X\).
\item In a combination \(c_1x_1+\dots+c_kx_k\) with all \(x_i\in U\cup W\), the terms with
\(x_i\in U\) add up to a vector \(u\in U\), and the other terms, whose \(x_i\) lie in \(W\), add
up to a vector \(w\in W\) (an empty sum is \(0\)). So the combination is \(u+w\in U+W\).
Conversely, \(u+w=1u+1w\) is such a combination.
\end{itemize}
(c)
\begin{itemize}
\item If \(u=\sum_ia_iu_i\in U\) and \(w=\sum_jb_jw_j\in W\), then
\(u+w=\sum_ia_iu_i+\sum_jb_jw_j\) is a combination of \(u_1,\dots,u_p,w_1,\dots,w_q\).
\item Conversely, \(\sum_ia_iu_i+\sum_jb_jw_j=u+w\) with \(u=\sum_ia_iu_i\in U\) and
\(w=\sum_jb_jw_j\in W\).\qedhere
\end{itemize}
\end{proof}

For matrices: let \(A\in\R^{m\times p}\) have columns \(a_1,\dots,a_p\), let \(B\in\R^{m\times q}\)
have columns \(b_1,\dots,b_q\), and let \([A\ B]\in\R^{m\times(p+q)}\) be the matrix with columns
\(a_1,\dots,a_p,b_1,\dots,b_q\). By (c), the columns of \([A\ B]\) span the sum of the subspaces
spanned by the columns of \(A\) and by the columns of \(B\),
\[
\Span\{a_1,\dots,a_p,b_1,\dots,b_q\}=\Span\{a_1,\dots,a_p\}+\Span\{b_1,\dots,b_q\}
\]
and, by Theorem~\ref{r1-thm-views}(b), the vectors of this subspace are the vectors \(Ax+By\)
(\(x\in\R^p\), \(y\in\R^q\)).

\begin{theorem}\label{r3-thm-direct-sum}
Let \(U\) and \(W\) be subspaces of a vector space \(V\). Then \(U\cap W=\{0\}\) if and only if
every \(x\in U+W\) can be written in exactly one way as \(x=u+w\) with \(u\in U\) and \(w\in W\),
that is: if \(u+w=u_1+w_1\) with \(u,u_1\in U\) and \(w,w_1\in W\), then \(u=u_1\) and \(w=w_1\).
\end{theorem}
\begin{proof}
\begin{itemize}
\item Let \(U\cap W=\{0\}\), and let \(u+w=u_1+w_1\) as above. Then \(u-u_1=w_1-w\). The left
side lies in \(U\) and the right side in \(W\), so this vector lies in \(U\cap W=\{0\}\). Hence
\(u=u_1\) and \(w=w_1\).
\item Let \(U\cap W\ne\{0\}\), and take \(x\ne0\) in \(U\cap W\). Then \(x=x+0\) (with \(x\in U\),
\(0\in W\)) and \(x=0+x\) (with \(0\in U\), \(x\in W\)) are two different ways of writing
\(x\in U+W\).\qedhere
\end{itemize}
\end{proof}
So, if \(V=U\oplus W\), each \(v\in V\) is \(v=u+w\) for exactly one \(u\in U\) and one
\(w\in W\) (the parts of \(v\) in \(U\) and in \(W\)).

\begin{theorem}\label{r3-thm-direct-sum-basis}
Let \(U\) and \(W\) be subspaces of a vector space \(V\), let \(u_1,\dots,u_p\) form a basis of
\(U\), and let \(w_1,\dots,w_q\) form a basis of \(W\). If \(U\cap W=\{0\}\), then
\(u_1,\dots,u_p,w_1,\dots,w_q\) form a basis of \(U\oplus W\). If \(U\cap W\ne\{0\}\), then
\(u_1,\dots,u_p,w_1,\dots,w_q\) are linearly dependent. Hence the sum \(U+W\) is direct if and
only if \(u_1,\dots,u_p,w_1,\dots,w_q\) are linearly independent.
\end{theorem}
\begin{proof}
\begin{itemize}
\item The vectors \(u_1,\dots,u_p,w_1,\dots,w_q\) lie in \(U+W\) and span it, by
Theorem~\ref{r3-thm-sum}(b),(c).
\item Let \(U\cap W=\{0\}\) and \(\sum_ia_iu_i+\sum_jb_jw_j=0\). Then
\(x=\sum_ia_iu_i=\sum_j(-b_j)w_j\) lies in \(U\) and in \(W\), so \(x=0\). Hence
\(\sum_ia_iu_i=0\) and \(\sum_jb_jw_j=0\), and since each list is independent, every \(a_i\) and
every \(b_j\) is \(0\). So the vectors form a basis of \(U+W=U\oplus W\).
\item Let \(U\cap W\ne\{0\}\), take \(x\ne0\) in \(U\cap W\), and write
\(x=\sum_ia_iu_i=\sum_jb_jw_j\). Then \(\sum_ia_iu_i+\sum_j(-b_j)w_j=0\), and not every \(a_i\)
is \(0\), since \(x\ne0\): a linear relation among \(u_1,\dots,u_p,w_1,\dots,w_q\).\qedhere
\end{itemize}
\end{proof}

\begin{example}\label{r3-exm-sum-r3}
In \(\R^3\) let \(U=\{x:x_3=0\}=\Span\{e_1,e_2\}\) (the subspace \(W\) of
Example~\ref{r3-exm-two-of-three}) and \(W_0=\{x:x_1+x_2+x_3=0\}=\Span\{d_1,d_2\}\), with
\(d_1=(1,-1,0)'\) and \(d_2=(0,1,-1)'\) (Example~\ref{r3-exm-centered-basis}); both are planes
through \(0\).
\begin{enumerate}
\item \(U+W_0=\R^3\), but the sum is not direct.
\begin{itemize}
\item Every \(x\in\R^3\) is \(x=(x_1,\ x_2+x_3,\ 0)'+(0,\ -x_3,\ x_3)'\), with the first
vector in \(U\) and the second, \(-x_3d_2\), in \(W_0\). So \(U+W_0=\R^3\).
\item \(x\in U\cap W_0\) means \(x_3=0\) and \(x_1+x_2+x_3=0\), that is, \(x=(t,-t,0)'=td_1\)
for some \(t\in\R\). So \(U\cap W_0=\Span\{d_1\}\), a line, and the sum is not direct.
\item Accordingly (Theorem~\ref{r3-thm-direct-sum}) the parts are not unique: for example
\((1,2,3)'=(1,5,0)'+(0,-3,3)'=(4,2,0)'+(-3,0,3)'\), and the two parts in \(U\) differ by
\(3d_1\in U\cap W_0\). Likewise \(e_1,e_2,d_1,d_2\) are dependent, since \(d_1=e_1-e_2\)
(Theorem~\ref{r3-thm-direct-sum-basis}).
\end{itemize}
\item \(\R^3=W_0\oplus\Span\{e_3\}\).
\begin{itemize}
\item \(ce_3=(0,0,c)'\) lies in \(W_0\) only for \(c=0\); so \(W_0\cap\Span\{e_3\}=\{0\}\).
\item For \(x\in\R^3\) and \(c\in\R\), \(x-ce_3=(x_1,\ x_2,\ x_3-c)'\) lies in \(W_0\) exactly
when \(c=x_1+x_2+x_3\). So every \(x\) is
\(x=(x_1,\ x_2,\ -x_1-x_2)'+(x_1+x_2+x_3)\,e_3\) with the first vector in \(W_0\), and in no
other way. For example, \((1,2,3)'=(1,2,-3)'+6e_3\).
\item By Theorem~\ref{r3-thm-direct-sum-basis}, \(d_1,d_2,e_3\) form a basis of \(\R^3\).
\end{itemize}
\end{enumerate}
\end{example}

\begin{example}\label{r3-exm-constant-centered}
The constant vectors and the centered vectors \(W_0=\{x\in\R^n:\one'x=0\}\)
(Example~\ref{r3-exm-subspaces}(f)) give a direct sum:
\[
\R^n=\Span\{\one_n\}\oplus W_0,\qquad x=\bar x\one+Cx
\]
\begin{itemize}
\item Sum: every \(x\in\R^n\) is \(x=\bar x\one+Cx\) with \(\bar x\one\in\Span\{\one\}\) and
\(Cx\in W_0\) (Example~\ref{r3-exm-constant}).
\item Intersection: if \(c\one\in W_0\), then \(0=\one'(c\one)=cn\)
(Proposition~\ref{r1-prop-centering}), so \(c=0\).
\item Hence \(\bar x\one\) is the only constant vector whose difference from \(x\) is centered
(Theorem~\ref{r3-thm-direct-sum}). In matrix form, \(x=\bar J_nx+C_nx\) with
\(\bar J_nx=\frac1n\one(\one'x)=\bar x\one\), and \(\bar J_n+C_n=I_n\)
(Definition~\ref{r1-def-ones}). For \(x=(1,3,5,7,9)'\) the two parts are \(\bar x\one=5\one\)
and \(Cx=(-4,-2,0,2,4)'\).
\item For \(n=3\), \(\one_3\) together with the basis \(d_1,d_2\) of \(W_0\)
(Example~\ref{r3-exm-centered-basis}) is a basis of \(\R^3\)
(Theorem~\ref{r3-thm-direct-sum-basis}).
\item The second subspace matters: also \(\R^3=W_0\oplus\Span\{e_3\}\)
(Example~\ref{r3-exm-sum-r3}(b)), and the parts of \(x=(1,2,3)'\) in \(W_0\) differ: here
\(\bar x=2\) and \(Cx=(-1,0,1)'\), there \((1,2,-3)'\).
\end{itemize}
\end{example}

\subsection*{Exercises and solutions}

\begin{problem}\label{r3-ex-subsets-r2}
Is the set of vectors \((x,y)'\in\R^2\) with \(x\geqslant0\) and \(y\geqslant0\) a subspace of
\(\R^2\)?
\end{problem}
\begin{solution}
No: \((1,1)'\) lies in the set, but \((-1)(1,1)'=(-1,-1)'\) does not; so the set is not closed
under scalar multiplication (Theorem~\ref{r3-thm-subspace-test}(b)).
\end{solution}

\begin{problem}\label{r3-ex-matrix-subspaces}
In the vector space \(\R^{3\times3}\), show that the set of lower triangular matrices is a
subspace.
\end{problem}
\begin{solution}
Let \(A,B\in\R^{3\times3}\) be lower triangular, that is, \(a_{ij}=b_{ij}=0\) for \(i<j\), and
let \(c\in\R\). For \(i<j\), \((A+B)_{ij}=a_{ij}+b_{ij}=0\) and \((cA)_{ij}=ca_{ij}=0\); so
\(A+B\) and \(cA\) are lower triangular, and so is \(0\). By
Theorem~\ref{r3-thm-subspace-test}(b), the lower triangular matrices form a subspace.
\end{solution}

\begin{problem}\label{r3-ex-same-span}
Consider \(e_1=(1,0,0)'\), \(e_2=(0,1,0)'\) and \(v=(0,-2,0)'\) in \(\R^3\). Show
that \(e_1,e_2,v\) span \(W=\{(x,y,0)':x,y\in\R\}\), that \(e_1\) and \(e_2\) alone span the
same subspace \(W\), and that \(e_2\) and \(v\) span only \(\{(0,y,0)':y\in\R\}\).
\end{problem}
\begin{solution}
For scalars \(a_1,a_2,a_3\) put \(\lambda=a_2-2a_3\). Since \(v=-2e_2\),
\[
a_1e_1+a_2e_2+a_3v=a_1e_1+\lambda e_2=(a_1,\lambda,0)',\qquad
a_2e_2+a_3v=\lambda e_2=(0,\lambda,0)'
\]
So \(\Span\{e_1,e_2,v\}\) consists of the vectors \((a_1,\lambda,0)'\). Each lies in \(W\),
and every \((x,y,0)'\in W\) is obtained with \(a_1=x\), \(a_2=y\), \(a_3=0\); so
\(\Span\{e_1,e_2,v\}=W\). The same computation with \(a_3=0\) gives \(\Span\{e_1,e_2\}=W\), and
\(\Span\{e_2,v\}=\{\lambda e_2:\lambda\in\R\}=\{(0,y,0)':y\in\R\}\).
\end{solution}

\begin{problem}\label{r3-ex-span-drill}
Let \(W=\Span\{v_1,v_2,v_3\}\subseteq\R^4\), where \(v_1=(1,2,1,0)'\), \(v_2=(0,1,1,1)'\) and
\(v_3=(1,3,3,2)'\).
\begin{enumerate}
\item Which \(w=(w_1,w_2,w_3,w_4)'\) lie in \(W\)?
\item Show that \((3,6,4,1)'\in W\), and write it as a combination of \(v_1,v_2,v_3\).
\item Does \((1,1,1,1)'\) lie in \(W\)?
\end{enumerate}
\end{problem}
\begin{solution}
(a) With \(V=[v_1\ v_2\ v_3]\), \(w\in W\) if and only if \(Vc=w\) is consistent (the paragraph
after Theorem~\ref{r3-thm-span}). Row-reduce \([V\mid w]\), keeping the \(w_i\) as letters:
\[
\begin{gathered}
\left(\begin{array}{ccc|c}1&0&1&w_1\\2&1&3&w_2\\1&1&3&w_3\\0&1&2&w_4\end{array}\right)
\xrightarrow[R_3\leftarrow R_3-R_1]{R_2\leftarrow R_2-2R_1}
\left(\begin{array}{ccc|c}1&0&1&w_1\\0&1&1&w_2-2w_1\\0&1&2&w_3-w_1\\0&1&2&w_4\end{array}\right)
\\
\xrightarrow[R_4\leftarrow R_4-R_2]{R_3\leftarrow R_3-R_2}
\left(\begin{array}{ccc|c}1&0&1&w_1\\0&1&1&w_2-2w_1\\0&0&1&w_1-w_2+w_3\\0&0&1&2w_1-w_2+w_4\end{array}\right)
\xrightarrow{R_4\leftarrow R_4-R_3}
\left(\begin{array}{ccc|c}1&0&1&w_1\\0&1&1&w_2-2w_1\\0&0&1&w_1-w_2+w_3\\0&0&0&w_1-w_3+w_4\end{array}\right)
\end{gathered}
\]
The last matrix is an echelon form. By Theorem~\ref{thm:r2-consistency}, \(Vc=w\) is consistent
if and only if its last row is not \((0\ 0\ 0\mid d)\) with \(d\ne0\):
\[
w\in W\iff w_1-w_3+w_4=0
\]
Then every column of the left part is a pivot column, so the coefficients are unique
(Theorem~\ref{thm:r2-solution-structure}(b)); back substitution gives \(c_3=w_1-w_2+w_3\),
\(c_2=(w_2-2w_1)-c_3=-3w_1+2w_2-w_3\) and \(c_1=w_1-c_3=w_2-w_3\).

(b) For \(w=(3,6,4,1)'\), \(3-4+1=0\), so \(w\in W\), with \(c_3=3-6+4=1\), \(c_2=-9+12-4=-1\) and
\(c_1=6-4=2\): \(w=2v_1-v_2+v_3\). Check: \((2,4,2,0)'-(0,1,1,1)'+(1,3,3,2)'=(3,6,4,1)'\).

(c) No: \(w_1-w_3+w_4=1-1+1=1\ne0\).
\end{solution}

\begin{problem}\label{r3-ex-independent-examples}
Show that the following vectors are linearly independent.
(a) \((1,2)'\) and \((1,3)'\).
(b) \((1,1,0)'\), \((1,1,1)'\) and \((0,1,-1)'\).
(c) \((0,1,1)'\), \((0,2,1)'\) and \((1,5,3)'\).
\end{problem}
\begin{solution}
(a) If \(a(1,2)'+b(1,3)'=0\), then \(a+b=0\) and \(2a+3b=0\). The second equation minus
twice the first gives \(b=0\), and then \(a=0\).

(b) If \(a(1,1,0)'+b(1,1,1)'+c(0,1,-1)'=0\), then \(a+b=0\), \(a+b+c=0\),
\(b-c=0\). The second equation minus the first gives \(c=0\); then the third gives \(b=0\), and
the first \(a=0\).

(c) If \(a(0,1,1)'+b(0,2,1)'+c(1,5,3)'=0\), then \(c=0\), \(a+2b+5c=0\),
\(a+b+3c=0\). Subtracting the third equation from the second gives \(b+2c=0\), so \(b=0\) (as
\(c=0\)), and then \(a=0\).
\end{solution}

\begin{problem}\label{r3-ex-independence-drill}
Let \(u_1=(1,2,3,1)'\), \(u_2=(0,1,1,1)'\), \(u_3=(1,2,3,0)'\) and \(u_4=(0,2,2,1)'\).
\begin{enumerate}
\item Are \(u_1,u_2,u_3\) linearly independent?
\item Are \(u_1,u_2,u_3,u_4\) linearly independent? If not, find a linear relation among them, and
write \(u_4\) as a combination of \(u_1,u_2,u_3\).
\end{enumerate}
\end{problem}
\begin{solution}
Row-reduce \(V=[u_1\ u_2\ u_3\ u_4]\):
\[
\begin{gathered}
\begin{pmatrix}1&0&1&0\\2&1&2&2\\3&1&3&2\\1&1&0&1\end{pmatrix}
\xrightarrow[R_4\leftarrow R_4-R_1]{R_2\leftarrow R_2-2R_1,\ R_3\leftarrow R_3-3R_1}
\begin{pmatrix}1&0&1&0\\0&1&0&2\\0&1&0&2\\0&1&-1&1\end{pmatrix}
\xrightarrow[R_4\leftarrow R_4-R_2]{R_3\leftarrow R_3-R_2}
\begin{pmatrix}1&0&1&0\\0&1&0&2\\0&0&0&0\\0&0&-1&-1\end{pmatrix}
\\
\xrightarrow{R_3\leftrightarrow R_4}
\begin{pmatrix}1&0&1&0\\0&1&0&2\\0&0&-1&-1\\0&0&0&0\end{pmatrix}
\xrightarrow{R_3\leftarrow-R_3}
\begin{pmatrix}1&0&1&0\\0&1&0&2\\0&0&1&1\\0&0&0&0\end{pmatrix}
\xrightarrow{R_1\leftarrow R_1-R_3}
\begin{pmatrix}1&0&0&-1\\0&1&0&2\\0&0&1&1\\0&0&0&0\end{pmatrix}
\end{gathered}
\]
After the second step the entry of row \(3\) in column \(3\) is \(0\), so Step 2 interchanges rows
\(3\) and \(4\). The last matrix is the RREF of \(V\), with pivot columns \(1,2,3\).

(a) Yes. Row operations act on each column separately (Theorem~\ref{r2-thm-elementary}(c) and
Theorem~\ref{r1-thm-views}(a)), so the same operations take \([u_1\ u_2\ u_3]\) to the first three
columns of the last matrix, whose columns are the unit vectors \(e_1,e_2,e_3\) of \(\R^4\). This is
the RREF of \([u_1\ u_2\ u_3]\), and every column is a pivot column; so \(u_1,u_2,u_3\) are
independent (Proposition~\ref{r3-prop-independence-test}).

(b) No. Column \(4\) is not a pivot column, so \(Vc=0\) has a free unknown and a solution
\(c\ne0\) (Proposition~\ref{r3-prop-independence-test}). The RREF gives \(c_1-c_4=0\),
\(c_2+2c_4=0\), \(c_3+c_4=0\); with \(c_4=t\), \(c=t\,(1,-2,-1,1)'\). For \(t=1\),
\[
u_1-2u_2-u_3+u_4=0,\qquad\text{so}\qquad u_4=-u_1+2u_2+u_3
\]
Check: \(-(1,2,3,1)'+2(0,1,1,1)'+(1,2,3,0)'=(0,2,2,1)'=u_4\).
\end{solution}

\begin{problem}\label{r3-ex-dependent-k}
For what values of the scalar \(k\) are the three row vectors \((k,1,0)\), \((1,k,1)\) and
\((0,1,k)\) linearly dependent, and for what values are they linearly independent?
\end{problem}
\begin{solution}
Let \(x_1,x_2,x_3\) be scalars with \(x_1(k,1,0)+x_2(1,k,1)+x_3(0,1,k)=0\), that is,
\[
kx_1+x_2=0,\qquad x_1+kx_2+x_3=0,\qquad x_2+kx_3=0,
\]
or equivalently (the first and third equations give (1), the second gives (2))
\[
x_2=-kx_3=-kx_1\quad(1),\qquad kx_2=-x_1-x_3\quad(2)
\]
\begin{itemize}
\item \(k=0\): (1) and (2) become \(x_2=0\) and \(x_3=-x_1\); for instance
\((x_1,x_2,x_3)=(1,0,-1)\) is a solution other than \(0\).
\item \(k\ne0\): (1) gives \(kx_3=kx_1\), so \(x_3=x_1\); then (1) and (2) give
\(-k^2x_1=kx_2=-2x_1\), so \((2-k^2)x_1=0\). Hence \(k^2=2\), or else \(x_1=0\) and then
\(x_3=x_2=0\).
\item If \(k=\sqrt2\), conditions (1), (2) are equivalent to \(x_3=x_1\) and
\(x_2=-\sqrt2\,x_1\); if \(k=-\sqrt2\), to \(x_3=x_1\) and \(x_2=\sqrt2\,x_1\). In both cases
\(x_1=1\) gives a solution other than \(0\).
\end{itemize}
Thus there are solutions other than \(x_1=x_2=x_3=0\) if and only if \(k=0\) or
\(k=\pm\sqrt2\). The three vectors are linearly dependent if \(k\in\{0,\sqrt2,-\sqrt2\}\) and
linearly independent otherwise.
\end{solution}

\begin{problem}\label{r3-ex-column-row-relations}
Show that the columns of the matrix
\(A=\begin{pmatrix}1&3&2&0\\1&-5&-3&1\\5&-1&0&2\end{pmatrix}\) are linearly dependent. How
about the rows?
\end{problem}
\begin{solution}
The columns satisfy, for example, \(3\col_1(A)-\col_2(A)-8\col_4(A)=0\):
\[
3\begin{pmatrix}1\\1\\5\end{pmatrix}-\begin{pmatrix}3\\-5\\-1\end{pmatrix}-8\begin{pmatrix}0\\1\\2\end{pmatrix}
=\begin{pmatrix}3-3-0\\3+5-8\\15+1-16\end{pmatrix}=0
\]
The coefficients \(3,-1,0,-8\) are not all \(0\), so the columns are dependent; equivalently,
\(Ax=0\) for \(x=(3,-1,0,-8)'\ne0\) (Theorem~\ref{r1-thm-views}(b)). The rows satisfy
\(3\row_1(A)+2\row_2(A)-\row_3(A)=0\):
\((3,9,6,0)+(2,-10,-6,2)-(5,-1,0,2)=(0,0,0,0)\), that is, \(y'A=0\) for
\(y=(3,2,-1)'\) (Theorem~\ref{r1-thm-views}(c)). So the rows are dependent too.
\end{solution}

\begin{problem}\label{r3-ex-nonunique}
Suppose that \(v_1,\dots,v_k\) span a vector space \(V\) but do not form a basis of \(V\). Show
that, for any \(v\in V\), the representation of \(v\) as a linear combination of
\(v_1,\dots,v_k\) is not unique.
\end{problem}
\begin{solution}
Let \(v\in V\). Since \(v_1,\dots,v_k\) span \(V\), there are scalars \(x_1,\dots,x_k\) with
\(v=\sum_ix_iv_i\). Since \(v_1,\dots,v_k\) span \(V\) but do not form a basis, they are
linearly dependent (Definition~\ref{r3-def-basis}): there are scalars \(z_1,\dots,z_k\), not
all \(0\), with \(\sum_iz_iv_i=0\). Let \(y_i=x_i+z_i\). Then
\(\sum_iy_iv_i=\sum_ix_iv_i+\sum_iz_iv_i=v+0=v\), and
\((y_1,\dots,y_k)\ne(x_1,\dots,x_k)\) because some \(z_i\ne0\). So \(v\) has two different
representations. (Compare Theorem~\ref{r3-thm-unique-coef}.)
\end{solution}

\begin{problem}\label{r3-ex-sum-intersection}
Let \(U\) and \(W\) be the subspaces of \(\R^3\) spanned by the columns of
\[
A=\begin{pmatrix}1&0\\0&1\\0&0\end{pmatrix}\qquad\text{and}\qquad
B=\begin{pmatrix}0&2\\1&1\\2&3\end{pmatrix},
\]
respectively. Find (a) a basis of \(U+W\); (b) a basis of \(U\cap W\) (is the sum \(U+W\)
direct?); (c) a vector of \(U+W\) that is not in \(U\cup W\).
\end{problem}
\begin{solution}
Let \(b_1=(0,1,2)'\) and \(b_2=(2,1,3)'\) be the columns of \(B\); the columns of \(A\) are
\(e_1,e_2\).

(a)
\begin{itemize}
\item By Theorem~\ref{r3-thm-sum}(c), \(U+W=\Span\{e_1,e_2,b_1,b_2\}\).
\item \(e_1,e_2,b_1\) are independent: \(c_1e_1+c_2e_2+c_3b_1=(c_1,\ c_2+c_3,\ 2c_3)'=0\) gives
\(c_3=0\), then \(c_2=0\) and \(c_1=0\).
\item So \(M=[e_1\ e_2\ b_1]\in\R^{3\times3}\) has \(Mc=0\) only for \(c=0\). By
Theorem~\ref{r2-thm-invertible}, \(M\) is invertible, and by
Example~\ref{r3-exm-invertible-basis} its columns form a basis of \(\R^3\).
\item These columns lie in the subspace \(U+W\), so every vector of \(\R^3\), a combination of
them, lies in \(U+W\). Hence \(U+W=\R^3\), and \(e_1,e_2,b_1\) form a basis of \(U+W\).
\end{itemize}
(b)
\begin{itemize}
\item \(U\) consists of the vectors \((x_1,x_2,0)'\), and \(W\) of the vectors
\(y_1b_1+y_2b_2=(2y_2,\ y_1+y_2,\ 2y_1+3y_2)'\) (\(x_1,x_2,y_1,y_2\in\R\)).
\item So \(U\cap W\) consists of the vectors \(y_1b_1+y_2b_2\) with \(2y_1+3y_2=0\). This
equation holds exactly when \(y_1=3t\) and \(y_2=-2t\) for some \(t\in\R\): these values give
\(6t-6t=0\); conversely, if \(2y_1+3y_2=0\), then \(t=y_1+y_2\) has \(3t=3y_1-2y_1=y_1\) and
\(-2t=3y_2-2y_2=y_2\).
\item For \(y_1=3t\), \(y_2=-2t\): \(y_1b_1+y_2b_2=(-4t,\ t,\ 0)'=t\,(-4,1,0)'\). Hence
\(U\cap W=\Span\{(-4,1,0)'\}\), and the single nonzero vector \((-4,1,0)'\) is a basis of it
(Proposition~\ref{r3-prop-dependence}(a)).
\item Since \(U\cap W\ne\{0\}\), the sum \(U+W\) is not direct (Definition~\ref{r3-def-sum}).
\end{itemize}
(c) Take \(x=(0,2,2)'\); it lies in \(U+W=\R^3\).
\begin{itemize}
\item \(x\notin U\), since its third entry is not \(0\).
\item If \(x=(2y_2,\ y_1+y_2,\ 2y_1+3y_2)'\), the first two entries give \(y_2=0\) and \(y_1=2\),
and then the third entry would be \(4\ne2\). So \(x\notin W\).
\end{itemize}
\end{solution}

\begin{problem}\label{r3-ex-direct-sum-r3}
In \(\R^3\) let \(W=\Span\{e_1\}\) and \(U=\Span\{u_1,u_2\}\), where \(u_1=(1,1,0)'\) and
\(u_2=(0,1,1)'\). Show that \(\R^3=W\oplus U\), and write \(x=(1,2,3)'\) as \(x=w+u\) with
\(w\in W\) and \(u\in U\).
\end{problem}
\begin{solution}
The vectors \(e_1,u_1,u_2\) are linearly independent: if \(ae_1+bu_1+cu_2=0\), then
\[
a+b=0,\qquad b+c=0,\qquad c=0,
\]
so \(c=0\), then \(b=0\) and \(a=0\).
\begin{itemize}
\item \(W\cap U=\{0\}\): \(e_1\) is a basis of \(W\), and \(u_1,u_2\) (independent by
Proposition~\ref{r3-prop-dependence}(b)) form a basis of \(U\). Their combined list
\(e_1,u_1,u_2\) is independent, so the sum \(W+U\) is direct
(Theorem~\ref{r3-thm-direct-sum-basis}).
\item \(W+U=\R^3\): the matrix \(M=[e_1\ u_1\ u_2]\in\R^{3\times3}\) has \(Mc=0\) only for
\(c=0\), so it is invertible (Theorem~\ref{r2-thm-invertible}), and its columns span \(\R^3\)
(Example~\ref{r3-exm-invertible-basis}). By Theorem~\ref{r3-thm-sum}(c),
\(W+U=\Span\{e_1,u_1,u_2\}=\R^3\).
\end{itemize}
Hence \(\R^3=W\oplus U\). To write \(x=(1,2,3)'\) as \(ae_1+bu_1+cu_2\), solve the system whose
augmented matrix is already in echelon form,
\[
\left(\begin{array}{ccc|c}1&1&0&1\\0&1&1&2\\0&0&1&3\end{array}\right)
\]
Back substitution (Theorem~\ref{thm:r2-consistency}) gives \(c=3\), \(b=2-3=-1\),
\(a=1-(-1)=2\). So \(w=2e_1=(2,0,0)'\) and \(u=-u_1+3u_2=(-1,2,3)'\); indeed
\(w+u=(1,2,3)'\). By Theorem~\ref{r3-thm-direct-sum}, this is the only such decomposition.
\end{solution}

\begin{problem}\label{r3-ex-sym-skew}
Show that every \(A\in\R^{n\times n}\) can be written in exactly one way as \(A=B+C\) with \(B\)
symmetric and \(C\) skew-symmetric, and conclude that \(\R^{n\times n}\) is the direct sum of the
subspace of symmetric matrices and the subspace of skew-symmetric matrices.
\end{problem}
\begin{solution}
Both sets are subspaces of \(\R^{n\times n}\) (Example~\ref{r3-exm-subspaces}(d)).
\begin{itemize}
\item Existence: let \(B=\frac12(A+A')\) and \(C=\frac12(A-A')\). Then \(B+C=A\), and, since
\((A')'=A\) (Theorem~\ref{r1-thm-transpose}), \(B'=\frac12(A'+A)=B\) and
\(C'=\frac12(A'-A)=-C\).
\item Uniqueness: let also \(A=B_1+C_1\) with \(B_1\) symmetric and \(C_1\) skew-symmetric. Then
\(B_1-B=C-C_1\), where \(B_1-B\) is symmetric and \(C-C_1\) is skew-symmetric. So
\(B_1-B=(B_1-B)'=(C-C_1)'=-(C-C_1)=-(B_1-B)\), hence \(2(B_1-B)=0\). Thus \(B_1=B\), and then
\(C_1=A-B_1=C\).
\item So the sum of the two subspaces is all of \(\R^{n\times n}\), and by
Theorem~\ref{r3-thm-direct-sum} their intersection is \(\{0\}\): \(\R^{n\times n}\) is their
direct sum.
\end{itemize}
For example, \(B=\frac12(A+A')\) and \(C=\frac12(A-A')\) give
\[
A=\begin{pmatrix}1&3&0\\1&2&4\\2&-2&5\end{pmatrix}
=\begin{pmatrix}1&2&1\\2&2&1\\1&1&5\end{pmatrix}
+\begin{pmatrix}0&1&-1\\-1&0&3\\1&-3&0\end{pmatrix}=B+C
\]
\end{solution}
\clearpage
\def\unitnumber{4}
\setcounter{theorem}{0}\setcounter{problem}{0}
\section*{Chapter 4. Bases, dimension, and rank}

\subsection*{Bases and dimension of subspaces of \(\R^n\)}

\begin{theorem}\label{thm:r4-pivot-test}
Let \(R\) be the RREF of \(A\in\R^{m\times n}\), with \(r\) nonzero rows. The following are
equivalent:
\begin{enumerate}[label=(\roman*)]
\item the columns \(a_1,\dots,a_n\) of \(A\) are linearly independent;
\item the system \(Ax=0\) has only the solution \(x=0\);
\item \(r=n\), that is, every column of \(R\) is a pivot column.
\end{enumerate}
In particular, if \(n>m\), then any \(n\) vectors of \(\R^m\) are linearly dependent.
\end{theorem}
We call this the \textbf{pivot test}.
\begin{proof}
\begin{itemize}
\item (i)\,\(\Leftrightarrow\)\,(ii): since \(Ax=x_1a_1+\dots+x_na_n\) (Theorem~\ref{r1-thm-views}(b)), the statement
``\(x_1a_1+\dots+x_na_n=0\) only for \(x=0\)'' says exactly that \(Ax=0\) only for \(x=0\).
\item (ii)\,\(\Leftrightarrow\)\,(iii): \(Ax=0\) and \(Rx=0\) have the same solutions (Theorem~\ref{r2-thm-rowops}(b)). If \(r<n\),
there is a free unknown; giving it the value \(1\) and the other free unknowns the value \(0\)
yields a solution \(x\ne0\) (Proposition~\ref{r2-prop-solutions}(b)). If \(r=n\), every unknown is a pivot unknown, there are no free
unknowns, and the equations of \(Rx=0\) read \(x_{k_i}=0\) (\(1\leqslant i\leqslant n\)); so
\(x=0\).
\item If \(n>m\), put the \(n\) vectors into a matrix \(A\in\R^{m\times n}\). Its RREF has
\(r\leqslant m<n\) nonzero rows, so (iii) fails and the vectors are dependent.\qedhere
\end{itemize}
\end{proof}

\begin{remark}\label{rem:r4-echelon}
For the pivot test one may stop at an echelon form \(E\) of \(A\) (Definitions~\ref{r2-def-rref}
and~\ref{r2-def-gauss-jordan}).
Backward elimination adds multiples of a row \(i\), whose leading \(1\) is in column \(k_i\), to rows
\(l<i\), whose leading \(1\) is in a column \(k_l<k_i\). Row \(i\) is \(0\) before column \(k_i\), so
row \(l\) changes only in columns \(\geqslant k_i\) and keeps its leading \(1\) in column \(k_l\).
Hence the RREF of \(A\) reached from \(E\) has the same leading columns as \(E\): the pivot columns
of \(A\), and their number \(r\), can be read off from any echelon form of \(A\). In particular,
(iii) of Theorem~\ref{thm:r4-pivot-test} says that every column of \(E\) contains a leading \(1\).
The same holds before the leading entries are scaled to \(1\): scaling a row does not move its
leading entry.
\end{remark}

\begin{theorem}\label{thm:r4-pivot-columns}
Let \(R\) be the RREF of \(A\in\R^{m\times n}\), with \(r\) nonzero rows, pivot columns
\(k_1<\dots<k_r\) and entries \(R_{ij}\).
\begin{enumerate}
\item The columns of \(A\) and those of \(R\) satisfy the same linear relations: for
\(x\in\R^n\), \(x_1a_1+\dots+x_na_n=0\) if and only if \(x_1\col_1(R)+\dots+x_n\col_n(R)=0\).
\item If \(r\geqslant1\), the pivot columns \(a_{k_1},\dots,a_{k_r}\) of \(A\) form a basis of
\(\Span\{a_1,\dots,a_n\}\), and for every \(j=1,\dots,n\)
\[
a_j=R_{1j}a_{k_1}+R_{2j}a_{k_2}+\dots+R_{rj}a_{k_r}
\]
If \(r=0\), then \(A=0\) and \(\Span\{a_1,\dots,a_n\}=\{0\}\).
\end{enumerate}
In particular, every finite list that spans a subspace \(W\) contains a basis of \(W\).
\end{theorem}
\begin{proof}
(a) The two sums are \(Ax\) and \(Rx\), and \(Ax=0\) if and only if \(Rx=0\) (Theorem~\ref{r2-thm-rowops}(b)).

(b) If \(r=0\), then \(R=0\), and \(A=P^{-1}R=0\), where \(R=PA\) with \(P\in\R^{m\times m}\)
invertible (Theorem~\ref{r2-thm-elementary}(c)). Let \(r\geqslant1\).
\begin{itemize}
\item Columns of \(R\): \(\col_{k_i}(R)=e_i\), and rows \(r+1,\dots,m\) of \(R\) are zero, so
\(\col_j(R)=\sum_{i=1}^rR_{ij}e_i=\sum_{i=1}^rR_{ij}\col_{k_i}(R)\) for every \(j\).
\item Relations: for a column \(j\notin\{k_1,\dots,k_r\}\) let \(x\in\R^n\) have \(x_j=1\),
\(x_{k_i}=-R_{ij}\) (\(1\leqslant i\leqslant r\)) and all other entries \(0\). By the first item,
\(Rx=\col_j(R)-\sum_iR_{ij}\col_{k_i}(R)=0\); by (a), \(Ax=a_j-\sum_iR_{ij}a_{k_i}=0\). For
\(j=k_l\) the formula reads \(a_{k_l}=a_{k_l}\), since \(R_{ik_l}\) is \(1\) for \(i=l\) and \(0\)
otherwise.
\item Spanning: each \(a_j\) lies in the subspace \(U=\Span\{a_{k_1},\dots,a_{k_r}\}\), so
\(\Span\{a_1,\dots,a_n\}\subseteq U\) (Theorem~\ref{r3-thm-span}); the reverse inclusion is clear.
\item Independence: if \(c_1a_{k_1}+\dots+c_ra_{k_r}=0\), then by (a)
\(c_1\col_{k_1}(R)+\dots+c_r\col_{k_r}(R)=c_1e_1+\dots+c_re_r=0\), so every \(c_i=0\).
\end{itemize}
For the last sentence, write the spanning vectors as the columns of a matrix and apply (b).
\end{proof}

By Theorem~\ref{thm:r4-pivot-columns}(a), \(a_j\) is a linear combination of \(a_1,\dots,a_{j-1}\) if
and only if \(\col_j(R)\) is a linear combination of \(\col_1(R),\dots,\col_{j-1}(R)\) (take
\(x_j=-1\) and \(x_l=0\) for \(l>j\)). In \(R\) this happens exactly when \(j\) is not a pivot column:
\(\col_{k_i}(R)=e_i\), while the columns before \(k_i\) are \(0\) in rows \(i,\dots,m\); and a
non-pivot column \(j\) equals \(\sum_{k_i<j}R_{ij}\col_{k_i}(R)\). So the pivot columns of \(A\) are
the columns that are not combinations of the earlier ones; they do not depend on the row operations
used. Hence the RREF of \(A\) is unique, as stated in Definition~\ref{r2-def-gauss-jordan}: two RREFs
\(R\), \(S\) of \(A\) have the same pivot columns \(k_1<\dots<k_r\), hence the same zero rows
\(r+1,\dots,m\); and for each \(j\), both \(R_{1j},\dots,R_{rj}\) and \(S_{1j},\dots,S_{rj}\) are
coefficients of \(a_j\) in the basis \(a_{k_1},\dots,a_{k_r}\)
(Theorem~\ref{thm:r4-pivot-columns}(b)), so they are equal (Theorem~\ref{r3-thm-unique-coef}). Thus
\(R=S\).

\begin{theorem}\label{thm:r4-dimension}
Every subspace \(W\) of \(\R^n\) has a basis, and any two bases of \(W\) have the same number of
elements. This number is the \textbf{dimension} of \(W\), written \(\dim W\). For example,
\(\dim\R^n=n\) (standard basis) and \(\dim\{0\}=0\) (the empty list is a basis of \(\{0\}\)).
\end{theorem}
\begin{proof}
\begin{itemize}
\item If \(W\) is spanned by \(w_1,\dots,w_m\), then any \(k>m\) vectors \(v_1,\dots,v_k\) of \(W\)
are linearly dependent. Indeed, write \(v_j=a_{1j}w_1+\dots+a_{mj}w_m\) (\(1\leqslant j\leqslant k\))
and \(A=(a_{ij})\in\R^{m\times k}\). As \(k>m\), the columns of \(A\) are dependent
(Theorem~\ref{thm:r4-pivot-test}): \(Ax=0\) for some \(x\ne0\). Then
\[
x_1v_1+\dots+x_kv_k=\sum_{j=1}^kx_j\sum_{i=1}^ma_{ij}w_i=\sum_{i=1}^m(Ax)_iw_i=0
\]
\item Same size: if \(v_1,\dots,v_k\) and \(w_1,\dots,w_m\) are bases of \(W\), each list spans
\(W\) and is independent, so the first item gives \(k\leqslant m\) and \(m\leqslant k\).
\item Existence: start with the empty list and, as long as the vectors chosen do not span \(W\),
adjoin a vector \(v\in W\) outside their span. The list stays independent: in a relation
\(c_1v_1+\dots+c_kv_k+cv=0\), \(c\ne0\) would put \(v\) in the span of \(v_1,\dots,v_k\), so \(c=0\),
and then every \(c_i=0\). More than \(n\) vectors of \(\R^n\) are dependent
(Theorem~\ref{thm:r4-pivot-test}), so the process stops after at most \(n\) steps, at a basis of
\(W\).\qedhere
\end{itemize}
\end{proof}

By Theorem~\ref{thm:r4-pivot-columns}, \(\dim\Span\{a_1,\dots,a_n\}=r\), the number of pivots of
the RREF of \([a_1\ \cdots\ a_n]\).

\begin{theorem}\label{thm:r4-counting}
Let \(W\) be a subspace of \(\R^n\) with \(\dim W=d\).
\begin{enumerate}
\item If \(U\) is a subspace of \(W\), then \(\dim U\leqslant d\), and \(\dim U=d\) only for
\(U=W\). In particular \(d\leqslant n\), and \(d=n\) only for \(W=\R^n\).
\item Every linearly independent list of vectors of \(W\) can be extended to a basis of \(W\).
\item More than \(d\) vectors of \(W\) are linearly dependent; fewer than \(d\) vectors do not span
\(W\).
\item If \(d\geqslant1\), then \(d\) vectors of \(W\) are linearly independent \(\Leftrightarrow\)
they span \(W\) \(\Leftrightarrow\) they form a basis of \(W\).
\end{enumerate}
\end{theorem}
\begin{proof}
If \(d=0\), then \(W=\{0\}\), and (a)--(c) are immediate: every subspace \(U\) of \(W\) is \(\{0\}=W\);
every vector of \(W\) is \(0\), so a nonempty list of vectors of \(W\) is dependent
(Proposition~\ref{r3-prop-dependence}(a)); and the empty list, the only independent one, is a basis of
\(W\) (Definition~\ref{r3-def-basis}). Let \(d\geqslant1\), and let \(w_1,\dots,w_d\) be a basis of \(W\)
(Theorem~\ref{thm:r4-dimension}). We prove (c) first.

\needspace{3\baselineskip}
(c)
\begin{itemize}
\item The vectors \(w_1,\dots,w_d\) span \(W\), so any \(k>d\) vectors of \(W\) are linearly dependent, by
the first step of the proof of Theorem~\ref{thm:r4-dimension} (with \(m=d\)).
\item If \(k<d\) vectors spanned \(W\), the same step, applied to them, would make the \(d>k\) vectors
\(w_1,\dots,w_d\) of \(W\) dependent \(\contradiction\)
\end{itemize}
(b)
\begin{itemize}
\item Let \(v_1,\dots,v_k\) be linearly independent vectors of \(W\). If they span \(W\), they form a basis
of \(W\).
\item Otherwise take \(v\in W\) outside \(\Span\{v_1,\dots,v_k\}\). Then \(v_1,\dots,v_k,v\) are independent:
in a relation \(c_1v_1+\dots+c_kv_k+cv=0\), \(c\ne0\) would put \(v\) in the span, so \(c=0\), and then
every \(c_i=0\).
\item Repeat this. Each list obtained consists of independent vectors of \(W\), so it has at most \(d\)
vectors by (c); hence the process stops, at a basis of \(W\) that contains \(v_1,\dots,v_k\).
\end{itemize}
(a)
\begin{itemize}
\item A basis of \(U\) (Theorem~\ref{thm:r4-dimension}) consists of \(\dim U\) independent vectors of \(W\), so
\(\dim U\leqslant d\) by (c).
\item If \(U\ne W\), take \(v\in W\) outside \(U\). Adjoining \(v\) to a basis of \(U\) gives, as in (b),
\(\dim U+1\) independent vectors of \(W\), so \(\dim U+1\leqslant d\) by (c). Hence \(\dim U=d\) only for
\(U=W\).
\item For the last sentence, apply this to \(W\) as a subspace of \(\R^n\), which has dimension \(n\)
(Theorem~\ref{thm:r4-dimension}).
\end{itemize}
(d) Let \(v_1,\dots,v_d\in W\). A basis of \(W\) is independent and spans \(W\)
(Definition~\ref{r3-def-basis}).
\begin{itemize}
\item Let \(v_1,\dots,v_d\) be independent, and let \(w\in W\). By (c), the \(d+1\) vectors
\(w,v_1,\dots,v_d\) are dependent: \(x_0w+x_1v_1+\dots+x_dv_d=0\) with not all \(x_i=0\). Here
\(x_0\ne0\), since \(v_1,\dots,v_d\) are independent; so \(w=-(x_1/x_0)v_1-\dots-(x_d/x_0)v_d\). Thus
\(v_1,\dots,v_d\) span \(W\), and they form a basis.
\item Let \(v_1,\dots,v_d\) span \(W\). By Theorem~\ref{thm:r4-pivot-columns}, some of them form a basis of
\(W\), which has \(d\) vectors (Theorem~\ref{thm:r4-dimension}); so it consists of all of
\(v_1,\dots,v_d\), and they form a basis.\qedhere
\end{itemize}
\end{proof}

Transposition carries a subspace of \(\R^{1\times n}\) onto a subspace of \(\R^n\) and preserves linear
combinations, hence independence, spanning and bases; so Theorems~\ref{thm:r4-dimension}
and~\ref{thm:r4-counting} also hold for subspaces of \(\R^{1\times n}\), such as the row space
\(\mathcal R(A)\) below.

\begin{example}\label{exm:r4-span}
Let \(W=\Span\{u_1,u_2,u_3,u_4\}\subseteq\R^3\), where
\[
u_1=\begin{pmatrix}1\\2\\-1\end{pmatrix},\quad
u_2=\begin{pmatrix}-2\\-4\\2\end{pmatrix},\quad
u_3=\begin{pmatrix}1\\3\\0\end{pmatrix},\quad
u_4=\begin{pmatrix}3\\7\\-2\end{pmatrix}
\]
(a) Are \(u_1,\dots,u_4\) linearly independent? (b) Find a basis of \(W\) among
\(u_1,\dots,u_4\), and \(\dim W\). (c) Express the remaining vectors in this basis.

\textit{Solution.} Row-reduce \(A=[u_1\ u_2\ u_3\ u_4]\):
\[
\begin{gathered}
\begin{pmatrix}1&-2&1&3\\2&-4&3&7\\-1&2&0&-2\end{pmatrix}
\xrightarrow[R_3\leftarrow R_3+R_1]{R_2\leftarrow R_2-2R_1}
\begin{pmatrix}1&-2&1&3\\0&0&1&1\\0&0&1&1\end{pmatrix}
\xrightarrow{R_3\leftarrow R_3-R_2}
\begin{pmatrix}1&-2&1&3\\0&0&1&1\\0&0&0&0\end{pmatrix}
\\
\xrightarrow{R_1\leftarrow R_1-R_2}
\begin{pmatrix}1&-2&0&2\\0&0&1&1\\0&0&0&0\end{pmatrix}=R
\end{gathered}
\]
So \(r=2\), with pivot columns \(k_1=1\), \(k_2=3\).
\begin{enumerate}
\item No: \(r=2<4\) (Theorem~\ref{thm:r4-pivot-test}); indeed any four vectors of \(\R^3\) are
dependent.
\item By Theorem~\ref{thm:r4-pivot-columns}, \(u_1,u_3\) form a basis of \(W\), and
\(\dim W=2\). As \(2<3\), \(W\neq\R^3\) (Theorem~\ref{thm:r4-counting}(a)).
\item Columns \(2\) and \(4\) of \(R\) give \(u_2=-2u_1\) and \(u_4=2u_1+u_3\)
(Theorem~\ref{thm:r4-pivot-columns}(b)). Check:
\(2u_1+u_3=(2,4,-2)'+(1,3,0)'=(3,7,-2)'=u_4\). Equivalently,
\(2u_1+u_2=0\) and \(2u_1+u_3-u_4=0\).
\end{enumerate}
\end{example}

\subsection*{The four subspaces of a matrix}

\begin{definition}\label{def:r4-spaces}
For \(A\in\R^{m\times n}\),
\[
\begin{aligned}
\mathcal C(A)&=\{Ax:x\in\R^n\}\subseteq\R^m, &\qquad
\mathcal R(A)&=\{y'A:y\in\R^m\}\subseteq\R^{1\times n},\\
\mathcal N(A)&=\{x\in\R^n:Ax=0\}\subseteq\R^n, &
\mathcal N(A')&=\{y\in\R^m:A'y=0\}\subseteq\R^m
\end{aligned}
\]
are the \textbf{column space}, \textbf{row space}, \textbf{null space} and \textbf{left null space}
of \(A\). The columns of \(A\) span \(\mathcal C(A)\), the rows of \(A\) span \(\mathcal R(A)\) (a
space of row vectors), \(\mathcal N(A)\) is the set of solutions of \(Ax=0\), and
\(y\in\mathcal N(A')\) if and only if \(y'A=0\) (transpose both sides). All four are subspaces
(Theorem~\ref{r3-thm-span} for \(\mathcal C(A)\) and \(\mathcal R(A)\), Example~\ref{r3-exm-subspaces}(e) for
\(\mathcal N(A)\) and \(\mathcal N(A')\)).
\end{definition}

Let \(R\) be the RREF of \(A\in\R^{m\times n}\), with \(r\) nonzero rows, pivot columns
\(k_1<\dots<k_r\) and other columns \(f_1<\dots<f_{n-r}\), so that \(x_{f_1},\dots,x_{f_{n-r}}\) are
the free unknowns of \(Ax=0\) (Proposition~\ref{r2-prop-solutions}(b)).
For \(j=1,\dots,n-r\), the \textbf{special solution} \(s_j\in\R^n\) is the solution of \(Ax=0\)
with free unknowns \(x_{f_j}=1\) and \(x_{f_l}=0\) (\(l\ne j\)). By
Proposition~\ref{r2-prop-solutions}(b), its entries are
\[
(s_j)_{f_l}=\delta_{jl}\quad(1\leqslant l\leqslant n-r),\qquad
(s_j)_{k_i}=-R_{if_j}\quad(1\leqslant i\leqslant r)
\]

\begin{theorem}\label{thm:r4-four-bases}
Let \(A\in\R^{m\times n}\), and let row operations take \([A\mid I_m]\) to \([R\mid P]\), where
\(R\) is the RREF of \(A\), with nonzero rows \(\rho_1,\dots,\rho_r\), pivot columns
\(k_1<\dots<k_r\) and free columns \(f_1<\dots<f_{n-r}\). Then \(P\in\R^{m\times m}\) is invertible,
\(PA=R\), and:
\begin{enumerate}
\item if \(r\geqslant1\), the pivot columns \(a_{k_1},\dots,a_{k_r}\) of \(A\) form a basis of
\(\mathcal C(A)\);
\item if \(r\geqslant1\), the nonzero rows \(\rho_1,\dots,\rho_r\) of \(R\) form a basis of
\(\mathcal R(A)\);
\item the special solutions \(s_1,\dots,s_{n-r}\) form a basis of \(\mathcal N(A)\) (if \(r=n\),
then \(\mathcal N(A)=\{0\}\), with the empty basis);
\item the transposed rows \(\row_{r+1}(P)',\dots,\row_m(P)'\) form a basis of \(\mathcal N(A')\) (if
\(r=m\), then \(\mathcal N(A')=\{0\}\), with the empty basis).
\end{enumerate}
Hence \(\dim\mathcal C(A)=\dim\mathcal R(A)=r\), \(\dim\mathcal N(A)=n-r\) and
\(\dim\mathcal N(A')=m-r\) (if \(r=0\), then \(A=0\) and \(\mathcal C(A)=\mathcal R(A)=\{0\}\)).
\end{theorem}
\begin{proof}
The row operations multiply \([A\mid I_m]\) on the left by an invertible matrix \(E\)
(Theorem~\ref{r2-thm-elementary}(c)), and \(E[A\mid I_m]=[EA\mid E]\), column by column
(Theorem~\ref{r1-thm-views}(a)). So \(P=E\) is invertible and
\(PA=R\).

(a) \(\mathcal C(A)=\Span\{a_1,\dots,a_n\}\); apply Theorem~\ref{thm:r4-pivot-columns}.

(b)
\begin{itemize}
\item Same row space: for \(y\in\R^m\), \(y'R=(y'P)A\) and
\(y'A=(y'P^{-1})R\), where \(y'P=(P'y)'\) and
\(y'P^{-1}=((P^{-1})'y)'\). So \(\mathcal R(R)\subseteq\mathcal R(A)\) and
\(\mathcal R(A)\subseteq\mathcal R(R)\), that is, \(\mathcal R(R)=\mathcal R(A)\).
\item Spanning: the rows of \(R\) span \(\mathcal R(R)\), and rows \(r+1,\dots,m\) are zero, so
\(\rho_1,\dots,\rho_r\) span \(\mathcal R(R)=\mathcal R(A)\).
\item Independent: \(\col_{k_i}(R)=e_i\), so entry \(k_i\) of \(c_1\rho_1+\dots+c_r\rho_r\) is
\(c_i\); if this combination is \(0\), every \(c_i=0\).
\end{itemize}
(c) The solutions of \(Ax=0\) are those of \(Rx=0\) (Theorem~\ref{r2-thm-rowops}(b)).
\begin{itemize}
\item Independent: entry \(f_l\) of \(\sum_jc_js_j\) is \(c_l\), so \(\sum_jc_js_j=0\) forces every
\(c_l=0\).
\item Spanning: let \(t\in\mathcal N(A)\) and \(\tilde t=\sum_jt_{f_j}s_j\). Then \(\tilde t\in\mathcal N(A)\),
and \(\tilde t\) has the same free unknowns as \(t\). A solution is determined by its free unknowns
(Proposition~\ref{r2-prop-solutions}(b)), so
\(t=\tilde t\in\Span\{s_1,\dots,s_{n-r}\}\).
\item If \(r=n\), there are no free unknowns and \(x=0\) is the only solution.
\end{itemize}
(d)
\begin{itemize}
\item In \(\mathcal N(A')\): for \(i>r\), \(\row_i(P)A=\row_i(PA)=\row_i(R)=0\) (Theorem~\ref{r1-thm-views}(c)), so
\(A'\row_i(P)'=(\row_i(P)A)'=0\).
\item Independent: if \(c'P=c_1\row_1(P)+\dots+c_m\row_m(P)=0\), then \(c'=c'PP^{-1}=0\); so the
rows of \(P\) are independent, and so are rows \(r+1,\dots,m\)
(Proposition~\ref{r3-prop-dependence}(b)).
\item Spanning: let \(A'y=0\), so \(y'A=0\), and put \(z'=y'P^{-1}\), so that
\(y'=z'P=\sum_{i=1}^mz_i\row_i(P)\). Then \(0=y'A=z'PA=z'R=z_1\rho_1+\dots+z_r\rho_r\), and (b)
gives \(z_1=\dots=z_r=0\). So \(y'=\sum_{i>r}z_i\row_i(P)\) (and \(y=0\) if \(r=m\)).\qedhere
\end{itemize}
\end{proof}

\begin{remark}\label{rem:r4-warning}
Row operations do not change \(\mathcal R(A)\) or \(\mathcal N(A)\), so \(R\) may be used for these
two spaces. They do change \(\mathcal C(A)\) in general: a basis of \(\mathcal C(A)\) consists of the
pivot columns of \(A\) itself, not of \(R\) (Example~\ref{exm:r4-four-bases}).
\end{remark}

\begin{example}\label{exm:r4-four-bases}
Find bases of the four subspaces of
\[
A=\begin{pmatrix}1&2&1&3\\1&3&0&5\\2&5&1&8\end{pmatrix}
\]
\textit{Solution.} Row-reduce \([A\mid I_3]\):
\[
\begin{gathered}
\left(\begin{array}{cccc|ccc}1&2&1&3&1&0&0\\1&3&0&5&0&1&0\\2&5&1&8&0&0&1\end{array}\right)
\xrightarrow[R_3\leftarrow R_3-2R_1]{R_2\leftarrow R_2-R_1}
\left(\begin{array}{cccc|ccc}1&2&1&3&1&0&0\\0&1&-1&2&-1&1&0\\0&1&-1&2&-2&0&1\end{array}\right)
\\
\xrightarrow{R_3\leftarrow R_3-R_2}
\left(\begin{array}{cccc|ccc}1&2&1&3&1&0&0\\0&1&-1&2&-1&1&0\\0&0&0&0&-1&-1&1\end{array}\right)
\xrightarrow{R_1\leftarrow R_1-2R_2}
\left(\begin{array}{cccc|ccc}1&0&3&-1&3&-2&0\\0&1&-1&2&-1&1&0\\0&0&0&0&-1&-1&1\end{array}\right)
=[R\mid P]
\end{gathered}
\]
So \(r=2\), with pivot columns \(1,2\) and free unknowns \(x_3,x_4\). (The pivots are already
visible in the echelon form after the second arrow; Remark~\ref{rem:r4-echelon}.)
\begin{itemize}
\item \(\mathcal C(A)\): \(a_1=(1,1,2)'\) and \(a_2=(2,3,5)'\) form a basis;
\(\dim\mathcal C(A)=2\). Columns \(3\), \(4\) of \(R\) give
\[
a_3=3a_1-a_2,\qquad a_4=-a_1+2a_2
\]
Check: \(3(1,1,2)'-(2,3,5)'=(1,0,1)'\) and \(-(1,1,2)'+2(2,3,5)'=(3,5,8)'\).
\item \(\mathcal R(A)\): \(\rho_1=(1,0,3,-1)\) and \(\rho_2=(0,1,-1,2)\) form a basis;
\(\dim\mathcal R(A)=2\).
\item \(\mathcal N(A)\): \(Rx=0\) reads \(x_1+3x_3-x_4=0\), \(x_2-x_3+2x_4=0\). The special
solutions
\[
s_1=\begin{pmatrix}-3\\1\\1\\0\end{pmatrix}\ (x_3=1,\ x_4=0),\qquad
s_2=\begin{pmatrix}1\\-2\\0\\1\end{pmatrix}\ (x_3=0,\ x_4=1)
\]
form a basis of \(\mathcal N(A)\); \(\dim\mathcal N(A)=2=4-2\). Note \(As_1=-3a_1+a_2+a_3=0\) is the
relation \(a_3=3a_1-a_2\): each special solution is a column relation.
\item \(\mathcal N(A')\): row \(3\) of \(P\) gives the basis \(y=(-1,-1,1)'\);
\(\dim\mathcal N(A')=1=3-2\). Indeed
\(y'A=-\row_1(A)-\row_2(A)+\row_3(A)=0\), that is, \(\row_3(A)=\row_1(A)+\row_2(A)\): each
vector of \(\mathcal N(A')\) is a row relation.
\item Warning: every column of \(R\) has third entry \(0\), so every vector of \(\mathcal C(R)\)
does; but \(a_1\) has third entry \(2\). Thus \(\mathcal C(R)\ne\mathcal C(A)\), and the pivot
columns \(e_1,e_2\) of \(R\) do not span \(\mathcal C(A)\).
\end{itemize}
\end{example}

\subsection*{Rank}

\begin{definition}\label{def:r4-rank}
By Theorem~\ref{thm:r4-four-bases}, \(\dim\mathcal C(A)=\dim\mathcal R(A)=r\), the number of
pivots of the RREF of \(A\) (also when \(r=0\)). This common number is the \textbf{rank} of \(A\),
written \(\rank A\); by Remark~\ref{rem:r4-echelon} it is also the
number of nonzero rows of any echelon form of \(A\). \(A\) has \textbf{full column rank} if
\(\rank A=n\), and \textbf{full row rank} if \(\rank A=m\). The number \(\dim\mathcal N(A)\) is the
\textbf{nullity} of \(A\).
\end{definition}

\begin{theorem}\label{thm:r4-rank-nullity}
For every \(A\in\R^{m\times n}\),
\[
\rank A+\dim\mathcal N(A)=n,\qquad\rank A+\dim\mathcal N(A')=m
\]
\end{theorem}
\begin{proof}
Each of the \(n\) columns of \(R\) is either a pivot column (\(r\) of them) or a free column
(\(n-r\) of them), and by Theorem~\ref{thm:r4-four-bases}(c), \(\dim\mathcal N(A)=n-r\), the number
of free unknowns. Likewise \(\dim\mathcal N(A')=m-r\), the number of zero rows of \(R\)
(Theorem~\ref{thm:r4-four-bases}(d)). Since \(\rank A=r\), both equalities follow.
\end{proof}

\begin{theorem}\label{thm:r4-rank-rules}
Let \(A\in\R^{m\times n}\).
\begin{enumerate}
\item \(\rank A'=\rank A\), and \(\rank A\leqslant\min\{m,n\}\).
\item \(A\) has full column rank \(\Leftrightarrow\) its columns are linearly independent
\(\Leftrightarrow\) \(\mathcal N(A)=\{0\}\).
\item \(A\) has full row rank \(\Leftrightarrow\) its rows are linearly independent
\(\Leftrightarrow\) \(\mathcal N(A')=\{0\}\) \(\Leftrightarrow\) \(\mathcal C(A)=\R^m\).
\item If \(m=n\), then \(A\) is invertible \(\Leftrightarrow\) \(\rank A=n\).
\end{enumerate}
\end{theorem}
\begin{proof}
(a) For \(y\in\R^m\), \((y'A)'=A'y\), so transposing maps \(\mathcal R(A)\) onto
\(\mathcal C(A')\). It preserves linear combinations and is undone by transposing again, so it
carries a basis of \(\mathcal R(A)\) to a basis of \(\mathcal C(A')\), and
\(\rank A'=\dim\mathcal C(A')=\dim\mathcal R(A)=\rank A\). Also \(\rank A=r\leqslant m\), as
\(R\) has \(m\) rows, and \(r\leqslant n\), as \(k_1<\dots<k_r\) lie in \(\{1,\dots,n\}\).

(b) Since \(\rank A=r\), this is Theorem~\ref{thm:r4-pivot-test}: (iii)\,\(\Leftrightarrow\)\,(i)
\(\Leftrightarrow\)\,(ii).

(c)
\begin{itemize}
\item \(\rank A=m\) if and only if \(\dim\mathcal N(A')=m-\rank A=0\)
(Theorem~\ref{thm:r4-rank-nullity}), that is, \(\mathcal N(A')=\{0\}\).
\item \(y'A=y_1\row_1(A)+\dots+y_m\row_m(A)\) (Theorem~\ref{r1-thm-views}(c)), and \(y'A=0\) if and only if
\(A'y=0\). So \(\mathcal N(A')=\{0\}\) says exactly that the rows of \(A\) are independent.
\item \(\rank A=\dim\mathcal C(A)\), and a subspace of \(\R^m\) has dimension \(m\) only if it is
\(\R^m\) (Theorem~\ref{thm:r4-counting}(a)).
\end{itemize}
(d) A square matrix \(A\) is invertible if and only if \(Ax=0\) has only the solution \(x=0\)
(Theorem~\ref{r2-thm-invertible}), that is, \(\mathcal N(A)=\{0\}\); by (b), this means \(\rank A=n\).
\end{proof}

\begin{corollary}\label{cor:r4-ax-b}
Let \(A\in\R^{m\times n}\), \(b\in\R^m\), and let \([A\mid b]\in\R^{m\times(n+1)}\) be the augmented
matrix.
\begin{enumerate}
\item \(Ax=b\) is consistent \(\Leftrightarrow\) \(b\in\mathcal C(A)\) \(\Leftrightarrow\)
\(\rank[A\mid b]=\rank A\).
\item If \(Ax=b\) is consistent, its solution is unique \(\Leftrightarrow\) \(\rank A=n\).
\end{enumerate}
\end{corollary}
\begin{proof}
(a)
\begin{itemize}
\item \(Ax=x_1a_1+\dots+x_na_n\) (Theorem~\ref{r1-thm-views}(b)), so \(Ax=b\) has a solution if and only if \(b\) is
a linear combination of the columns of \(A\), that is, \(b\in\mathcal C(A)\).
\item \(\mathcal C(A)\subseteq\mathcal C([A\mid b])\): a combination of \(a_1,\dots,a_n\) is one of
\(a_1,\dots,a_n,b\), with coefficient \(0\) on \(b\).
\item If \(b=Ac\), then \(x_1a_1+\dots+x_na_n+tb=A(x+tc)\in\mathcal C(A)\); so
\(\mathcal C([A\mid b])=\mathcal C(A)\), and the ranks are equal.
\item If the ranks are equal, the subspace \(\mathcal C(A)\) of \(\mathcal C([A\mid b])\) has the
same dimension, so the two are equal (Theorem~\ref{thm:r4-counting}(a)), and
\(b\in\mathcal C([A\mid b])=\mathcal C(A)\).
\end{itemize}
(b) By Theorem~\ref{thm:r2-solution-structure}(a), the solutions are \(x_p+z\), where \(x_p\) is
one solution and \(z\) runs through the solutions of \(Ax=0\), that is, through \(\mathcal N(A)\). So
the solution is unique if and only if \(\mathcal N(A)=\{0\}\), that is, \(\rank A=n\)
(Theorem~\ref{thm:r4-rank-rules}(b)).
\end{proof}
In practice, consistency is read off from an echelon form of \([A\mid b]\)
(Theorem~\ref{thm:r2-consistency}).

\subsection*{Rank equalities and inequalities}

\begin{theorem}\label{thm:r4-rank-product}
Let \(A\in\R^{m\times n}\) and \(B\in\R^{n\times p}\).
\begin{enumerate}
\item \(\mathcal C(AB)\subseteq\mathcal C(A)\) and \(\mathcal N(B)\subseteq\mathcal N(AB)\); and
\(\rank(AB)\leqslant\min\{\rank A,\rank B\}\).
\item If \(A\) has full column rank, then \(\mathcal N(AB)=\mathcal N(B)\) and \(\rank(AB)=\rank B\).
\item If \(B\) has full row rank, then \(\mathcal C(AB)=\mathcal C(A)\) and \(\rank(AB)=\rank A\).
\item \(\rank(PAQ)=\rank A\) for all invertible \(P\in\R^{m\times m}\) and \(Q\in\R^{n\times n}\).
\end{enumerate}
\end{theorem}
\begin{proof}
(a)
\begin{itemize}
\item \(ABx=A(Bx)\in\mathcal C(A)\), and \(Bx=0\) gives \(ABx=A0=0\).
\item So \(\rank(AB)=\dim\mathcal C(AB)\leqslant\dim\mathcal C(A)=\rank A\)
(Theorem~\ref{thm:r4-counting}(a)).
\item By this, \((AB)'=B'A'\) (Theorem~\ref{r1-thm-transpose}) and Theorem~\ref{thm:r4-rank-rules}(a),
\(\rank(AB)=\rank(B'A')\leqslant\rank B'=\rank B\).
\end{itemize}
(b)
\begin{itemize}
\item If \(ABx=0\), then \(Bx\in\mathcal N(A)=\{0\}\) (Theorem~\ref{thm:r4-rank-rules}(b)), so
\(x\in\mathcal N(B)\). With (a), \(\mathcal N(AB)=\mathcal N(B)\).
\item \(AB\) and \(B\) both have \(p\) columns, so by Theorem~\ref{thm:r4-rank-nullity},
\(\rank(AB)=p-\dim\mathcal N(AB)=p-\dim\mathcal N(B)=\rank B\).
\end{itemize}
(c)
\begin{itemize}
\item \(\mathcal C(B)=\R^n\) (Theorem~\ref{thm:r4-rank-rules}(c)): every \(x\in\R^n\) is \(Bw\) for
some \(w\in\R^p\), and then \(Ax=(AB)w\in\mathcal C(AB)\).
\item So \(\mathcal C(A)\subseteq\mathcal C(AB)\); with (a), \(\mathcal C(AB)=\mathcal C(A)\), and the
ranks are equal.
\end{itemize}
(d) \(Px=0\) gives \(x=P^{-1}Px=0\), so \(P\) has full column rank
(Theorem~\ref{thm:r4-rank-rules}(b)); and \(\rank Q=n\) (Theorem~\ref{thm:r4-rank-rules}(d)), so
\(Q\) has full row rank. By (b) and (c), \(\rank(P(AQ))=\rank(AQ)=\rank A\).
\end{proof}

\begin{theorem}\label{thm:r4-partitioned}
\begin{enumerate}
\item For \(A\in\R^{m\times n}\) and \(B\in\R^{m\times p}\), let \([A\ B]\in\R^{m\times(n+p)}\) have
the columns of \(A\) followed by those of \(B\). Then
\[
\max\{\rank A,\rank B\}\leqslant\rank[A\ B]\leqslant\rank A+\rank B
\]
\item For \(A\in\R^{m\times n}\) and \(C\in\R^{q\times n}\), let
\(\left(\begin{smallmatrix}A\\C\end{smallmatrix}\right)\in\R^{(m+q)\times n}\) have the rows of \(A\)
followed by those of \(C\). Then
\[
\max\{\rank A,\rank C\}\leqslant\rank\begin{pmatrix}A\\C\end{pmatrix}\leqslant\rank A+\rank C
\]
\item For \(A\in\R^{m\times n}\) and \(B\in\R^{q\times p}\), let
\(\diag(A,B)=\left(\begin{smallmatrix}A&0\\0&B\end{smallmatrix}\right)\in\R^{(m+q)\times(n+p)}\), with
zero blocks of sizes \(m\times p\) and \(q\times n\). Then \(\rank\diag(A,B)=\rank A+\rank B\).
\end{enumerate}
\end{theorem}
\begin{proof}
(a)
\begin{itemize}
\item A combination of the columns of \(A\), or of those of \(B\), is a combination of the columns
of \([A\ B]\) (the others with coefficient \(0\)). So \(\mathcal C(A)\) and \(\mathcal C(B)\) are
subspaces of \(\mathcal C([A\ B])\), and Theorem~\ref{thm:r4-counting}(a) gives the lower bound.
\item Let \(r=\rank A\), \(s=\rank B\), and take bases \(v_1,\dots,v_r\) of \(\mathcal C(A)\) and
\(w_1,\dots,w_s\) of \(\mathcal C(B)\). Every column of \([A\ B]\) is a combination of these
\(r+s\) vectors of \(\mathcal C([A\ B])\), so they span \(\mathcal C([A\ B])\) (Theorem~\ref{r3-thm-span}), and
\(\rank[A\ B]\leqslant r+s\) (Theorem~\ref{thm:r4-counting}(c)).
\end{itemize}
(b) Transposing turns the rows of \(\left(\begin{smallmatrix}A\\C\end{smallmatrix}\right)\) into the
columns of \([A'\ C']\), so \(\left(\begin{smallmatrix}A\\C\end{smallmatrix}\right)'=[A'\ C']\).
Apply (a) to \(A'\) and \(C'\), and use Theorem~\ref{thm:r4-rank-rules}(a).

(c) Let \(M=\diag(A,B)\), \(r=\rank A\) and \(s=\rank B\).
\begin{itemize}
\item Upper bound: \(M=[M_1\ M_2]\) with \(M_1=\left(\begin{smallmatrix}A\\0\end{smallmatrix}\right)\)
and \(M_2=\left(\begin{smallmatrix}0\\B\end{smallmatrix}\right)\). By (b), and as a zero matrix has
rank \(0\), \(r\leqslant\rank M_1\leqslant r+0\); so \(\rank M_1=r\), and likewise
\(\rank M_2=s\). By (a), \(\rank M\leqslant r+s\).
\item Lower bound: let \(a_{k_1},\dots,a_{k_r}\) and \(b_{l_1},\dots,b_{l_s}\) be the pivot columns of
\(A\) and of \(B\), bases of \(\mathcal C(A)\) and \(\mathcal C(B)\)
(Theorem~\ref{thm:r4-four-bases}(a)). The corresponding columns
\(\left(\begin{smallmatrix}a_{k_i}\\0\end{smallmatrix}\right)\) and
\(\left(\begin{smallmatrix}0\\b_{l_j}\end{smallmatrix}\right)\) of \(M\) are independent: a combination
of them is \(\left(\begin{smallmatrix}\sum_i\alpha_ia_{k_i}\\\sum_j\beta_jb_{l_j}\end{smallmatrix}\right)\),
which is \(0\) only if all \(\alpha_i\) and \(\beta_j\) are \(0\). These \(r+s\) vectors lie in
\(\mathcal C(M)\), so \(\rank M\geqslant r+s\) (Theorem~\ref{thm:r4-counting}(c)).\qedhere
\end{itemize}
\end{proof}

\begin{corollary}\label{cor:r4-rank-sum}
For \(A,B\in\R^{m\times n}\), \(\rank(A+B)\leqslant\rank A+\rank B\).
\end{corollary}
\begin{proof}
For \(x\in\R^n\), \((A+B)x=Ax+Bx\) is a combination of the columns of \(A\) and of \(B\), so it
lies in \(\mathcal C([A\ B])\). Hence \(\mathcal C(A+B)\subseteq\mathcal C([A\ B])\), so
\(\rank(A+B)\leqslant\rank[A\ B]\) (Theorem~\ref{thm:r4-counting}(a)), and
\(\rank[A\ B]\leqslant\rank A+\rank B\) (Theorem~\ref{thm:r4-partitioned}(a)).
\end{proof}

\begin{theorem}\label{thm:r4-rank-ata}
Let \(A\in\R^{m\times n}\).
\begin{enumerate}
\item \(\mathcal N(A'A)=\mathcal N(A)\).
\item \(\rank(A'A)=\rank A=\rank(AA')\).
\item \(\mathcal C(A'A)=\mathcal C(A')\) and \(\mathcal C(AA')=\mathcal C(A)\).
\end{enumerate}
\end{theorem}
\begin{proof}
(a)
\begin{itemize}
\item If \(Ax=0\), then \(A'Ax=A'0=0\).
\item If \(A'Ax=0\), then \((Ax)'(Ax)=x'A'Ax=0\). But \((Ax)'(Ax)\) is the sum of the squares of
the entries of \(Ax\) (Theorem~\ref{r1-thm-transpose}), so \(Ax=0\).
\end{itemize}
(b) \(A'A\) and \(A\) both have \(n\) columns, so by (a) and Theorem~\ref{thm:r4-rank-nullity},
\(\rank(A'A)=n-\dim\mathcal N(A'A)=n-\dim\mathcal N(A)=\rank A\). Applied to \(A'\) in place of
\(A\), this gives \(\rank(AA')=\rank A'=\rank A\) (Theorem~\ref{thm:r4-rank-rules}(a)).

(c) \(\mathcal C(A'A)\subseteq\mathcal C(A')\) (Theorem~\ref{thm:r4-rank-product}(a)), and both have
dimension \(\rank A\), by (b) and Theorem~\ref{thm:r4-rank-rules}(a); so they are equal
(Theorem~\ref{thm:r4-counting}(a)). Likewise \(\mathcal C(AA')\subseteq\mathcal C(A)\), both of
dimension \(\rank A\).
\end{proof}

\begin{theorem}[Sylvester's inequality]\label{thm:r4-sylvester}
For \(A\in\R^{m\times n}\) and \(B\in\R^{n\times p}\),
\[
\dim\mathcal N(AB)\leqslant\dim\mathcal N(A)+\dim\mathcal N(B),\qquad\text{equivalently}\qquad
\rank(AB)\geqslant\rank A+\rank B-n
\]
In particular, if \(AB=0\), then \(\rank A+\rank B\leqslant n\).
\end{theorem}
\begin{proof}
\begin{itemize}
\item Let \(z_1,\dots,z_s\) be a basis of \(\mathcal N(B)\). Since
\(\mathcal N(B)\subseteq\mathcal N(AB)\) (Theorem~\ref{thm:r4-rank-product}(a)), it extends to a
basis \(z_1,\dots,z_s,w_1,\dots,w_t\) of \(\mathcal N(AB)\) (Theorem~\ref{thm:r4-counting}(b));
so \(\dim\mathcal N(AB)=s+t\).
\item The vectors \(Bw_1,\dots,Bw_t\) lie in \(\mathcal N(A)\), since \(A(Bw_j)=(AB)w_j=0\).
\item They are independent: if \(c_1Bw_1+\dots+c_tBw_t=0\), then \(B(c_1w_1+\dots+c_tw_t)=0\), so
\(c_1w_1+\dots+c_tw_t\in\mathcal N(B)\) equals some \(d_1z_1+\dots+d_sz_s\); as
\(z_1,\dots,z_s,w_1,\dots,w_t\) are independent, every \(c_j=0\).
\item Hence \(t\leqslant\dim\mathcal N(A)\) (Theorem~\ref{thm:r4-counting}(c)), and
\(\dim\mathcal N(AB)=s+t\leqslant\dim\mathcal N(B)+\dim\mathcal N(A)\).
\item By Theorem~\ref{thm:r4-rank-nullity}, \(\dim\mathcal N(AB)=p-\rank(AB)\),
\(\dim\mathcal N(B)=p-\rank B\) and \(\dim\mathcal N(A)=n-\rank A\). Substituting gives
\(p-\rank(AB)\leqslant(p-\rank B)+(n-\rank A)\), which is the rank form. If \(AB=0\), then
\(\rank(AB)=0\).\qedhere
\end{itemize}
\end{proof}

\begin{example}\label{exm:r4-bounds}
Let \(A=\diag(1,1,0)\) and \(B=\diag(0,1,1)\) in \(\R^{3\times3}\). Their columns are \(e_1,e_2,0\)
and \(0,e_2,e_3\), so \(\mathcal C(A)=\Span\{e_1,e_2\}\), \(\mathcal C(B)=\Span\{e_2,e_3\}\) and
\(\rank A=\rank B=2\).
\begin{itemize}
\item \(AB=\diag(0,1,0)\) (row \(i\) of \(AB\) is \(a_{ii}\row_i(B)\), Theorem~\ref{r1-thm-views}(c)) has rank
\(1\): less than \(\min\{2,2\}\)
(Theorem~\ref{thm:r4-rank-product}(a)), and equal to \(\rank A+\rank B-3=1\), so Sylvester's bound
is attained (Theorem~\ref{thm:r4-sylvester}).
\item \(A+B=\diag(1,2,1)\) has rank \(3<2+2\) (Corollary~\ref{cor:r4-rank-sum}); and
\(A+(-A)=0\) has rank \(0\).
\item \([A\ B]\in\R^{3\times6}\) has the columns \(e_1,e_2,0,0,e_2,e_3\), so
\(\mathcal C([A\ B])=\R^3\) and \(\rank[A\ B]=3\), strictly between \(\max\{2,2\}\) and \(2+2\)
(Theorem~\ref{thm:r4-partitioned}(a)).
\item \(\diag(A,B)=\diag(1,1,0,0,1,1)\in\R^{6\times6}\) has rank \(4=2+2\)
(Theorem~\ref{thm:r4-partitioned}(c)).
\item \(A'A=A^2=A\) has rank \(2=\rank A\) (Theorem~\ref{thm:r4-rank-ata}).
\end{itemize}
\end{example}

\begin{theorem}\label{thm:r4-full-rank-factorization}
Let \(A\in\R^{m\times n}\) have rank \(r\). If \(r\geqslant1\), a product \(A=BC\) with
\(B\in\R^{m\times r}\) of full column rank and \(C\in\R^{r\times n}\) of full row rank is a
\textbf{full-rank factorization} of \(A\).
\begin{enumerate}
\item If \(r\geqslant1\), let \(R\) be the RREF of \(A\), with pivot columns \(k_1<\dots<k_r\); let
\(B=[a_{k_1}\ \cdots\ a_{k_r}]\in\R^{m\times r}\) consist of the pivot columns of \(A\), and
\(C\in\R^{r\times n}\) of the \(r\) nonzero rows of \(R\). Then \(A=BC\) is a full-rank
factorization. If \(r=0\), then \(A=0\), and there is nothing to factor.
\item For every full-rank factorization \(A=BC\), \(\mathcal C(A)=\mathcal C(B)\) and
\(\mathcal N(A)=\mathcal N(C)\).
\item For every full-rank factorization \(A=BC\), the \(r\times r\) matrices \(B'B\) and \(CC'\) are
invertible, and
\[
L=(B'B)^{-1}B'\in\R^{r\times m},\qquad K=C'(CC')^{-1}\in\R^{n\times r}
\]
satisfy \(LB=I_r\) and \(CK=I_r\): \(B\) has a left inverse and \(C\) has a right inverse.
\end{enumerate}
\end{theorem}
\begin{proof}
(a)
\begin{itemize}
\item Column \(j\) of \(BC\) is \(B\col_j(C)=R_{1j}a_{k_1}+\dots+R_{rj}a_{k_r}\)
(Theorem~\ref{r1-thm-views}(a),(b)), which
is \(a_j\) by Theorem~\ref{thm:r4-pivot-columns}(b). So \(BC=A\).
\item The columns of \(B\) form a basis of \(\mathcal C(A)\), so they are independent: \(B\) has full
column rank (Theorem~\ref{thm:r4-rank-rules}(b)).
\item The rows of \(C\) form a basis of \(\mathcal R(A)\) (Theorem~\ref{thm:r4-four-bases}(b)), so
they are independent: \(C\) has full row rank (Theorem~\ref{thm:r4-rank-rules}(c)).
\item If \(r=0\), then \(A=0\) (Theorem~\ref{thm:r4-four-bases}).
\end{itemize}
(b) Theorem~\ref{thm:r4-rank-product}(b) gives \(\mathcal N(BC)=\mathcal N(C)\), as \(B\) has full
column rank, and Theorem~\ref{thm:r4-rank-product}(c) gives \(\mathcal C(BC)=\mathcal C(B)\), as
\(C\) has full row rank.

(c)
\begin{itemize}
\item \(\rank(B'B)=\rank B=r\) and \(\rank(CC')=\rank C=r\) (Theorem~\ref{thm:r4-rank-ata}(b)), so
the \(r\times r\) matrices \(B'B\) and \(CC'\) are invertible (Theorem~\ref{thm:r4-rank-rules}(d)).
\item \(LB=(B'B)^{-1}(B'B)=I_r\) and \(CK=(CC')(CC')^{-1}=I_r\).\qedhere
\end{itemize}
\end{proof}

\begin{example}\label{exm:r4-factorization}
For \(A\) of Example~\ref{exm:r4-four-bases}, Theorem~\ref{thm:r4-full-rank-factorization}(a) gives
\[
\begin{pmatrix}1&2&1&3\\1&3&0&5\\2&5&1&8\end{pmatrix}
=\begin{pmatrix}1&2\\1&3\\2&5\end{pmatrix}\begin{pmatrix}1&0&3&-1\\0&1&-1&2\end{pmatrix}
\]
\begin{itemize}
\item Columns \(1\), \(2\) of \(BC\) are \(a_1\), \(a_2\), and columns \(3\), \(4\) are
\(3a_1-a_2=a_3\) and \(-a_1+2a_2=a_4\) (Example~\ref{exm:r4-four-bases}).
\item By Theorem~\ref{thm:r4-full-rank-factorization}(b),
\(\mathcal C(A)=\mathcal C(B)=\Span\{a_1,a_2\}\) and
\(\mathcal N(A)=\mathcal N(C)=\Span\{s_1,s_2\}\): the system \(Cx=0\) is \(Rx=0\) without its zero
row.
\item With \(N=[s_1\ s_2]\in\R^{4\times2}\), \(AN=0\) and \(\rank N=2\), and Sylvester's inequality
\(\rank A+\rank N\leqslant4\) holds with equality: \(2+2=4\).
\end{itemize}
\end{example}

\subsection*{Exercises and solutions}

\begin{problem}\label{ex:r4-which-bases}
Let \(a_1,\dots,a_9\) be the columns of
\[
A=\begin{pmatrix}1&1&1&1&1&1&2&2&5\\0&1&1&2&2&3&-1&3&3\\1&1&2&3&5&5&1&0&4\end{pmatrix}
\]
Which of the following form a basis of \(\R^3\)? (a) \(a_1,a_2\); (b) \(a_1,a_4,a_6,a_8\);
(c) \(a_2,a_4,a_7\); (d) \(a_3,a_5,a_9\).
\end{problem}
\begin{solution}
(a), (b) No: \(\dim\R^3=3\), so two vectors cannot span \(\R^3\), and four vectors of \(\R^3\) are
dependent (Theorem~\ref{thm:r4-counting}(c)).

(c) Yes. Three vectors of \(\R^3\) form a basis if and only if they are independent
(Theorem~\ref{thm:r4-counting}(d)). Row-reduce \([a_2\ a_4\ a_7]\):
\[
\begin{gathered}
\begin{pmatrix}1&1&2\\1&2&-1\\1&3&1\end{pmatrix}
\xrightarrow[R_3\leftarrow R_3-R_1]{R_2\leftarrow R_2-R_1}
\begin{pmatrix}1&1&2\\0&1&-3\\0&2&-1\end{pmatrix}
\xrightarrow{R_3\leftarrow R_3-2R_2}
\begin{pmatrix}1&1&2\\0&1&-3\\0&0&5\end{pmatrix}
\\
\xrightarrow[R_2\leftarrow R_2+3R_3]{R_3\leftarrow\frac15R_3,\ R_1\leftarrow R_1-2R_3}
\begin{pmatrix}1&1&0\\0&1&0\\0&0&1\end{pmatrix}
\xrightarrow{R_1\leftarrow R_1-R_2}
\begin{pmatrix}1&0&0\\0&1&0\\0&0&1\end{pmatrix}
\end{gathered}
\]
All three columns are pivot columns, so \(a_2,a_4,a_7\) are independent
(Theorem~\ref{thm:r4-pivot-test}).

(d) No: the three vectors are dependent, since
\[
7a_3-2a_5-a_9=7\begin{pmatrix}1\\1\\2\end{pmatrix}-2\begin{pmatrix}1\\2\\5\end{pmatrix}-\begin{pmatrix}5\\3\\4\end{pmatrix}
=\begin{pmatrix}7-2-5\\7-4-3\\14-10-4\end{pmatrix}=0
\]
\end{solution}

\begin{problem}\label{ex:r4-column-basis}
Let \(W=\mathcal C(A)\subseteq\R^4\) for
\[
A=\begin{pmatrix}1&2&3\\-2&3&8\\5&1&-3\\-3&-4&-5\end{pmatrix}
\]
(a) Find a basis of \(W\). (b) What is \(\dim W\)?
\end{problem}
\begin{solution}
Let \(a_1,a_2,a_3\) be the columns of \(A\); they span \(W\).

(a) \(a_1,a_2\) are independent: otherwise, as \(a_1\ne0\), \(a_2=ca_1\) for some \(c\)
(Proposition~\ref{r3-prop-dependence}(c)); the first entries give \(c=2\), and then the second entries give \(3=2(-2)\)
\(\contradiction\). But
\[
a_1-2a_2+a_3=\begin{pmatrix}1-4+3\\-2-6+8\\5-2-3\\-3+8-5\end{pmatrix}=0,
\]
so \(a_3=2a_2-a_1\), and every combination \(x_1a_1+x_2a_2+x_3a_3=(x_1-x_3)a_1+(x_2+2x_3)a_2\) is a
combination of \(a_1,a_2\). Thus \(a_1,a_2\) span \(W\) and form a basis of \(W\). (Check by
Theorem~\ref{thm:r4-pivot-columns}: the RREF of \(A\) has rows \((1,0,-1)\), \((0,1,2)\) and two
zero rows; its pivot columns are \(1,2\), and column \(3\) gives \(a_3=-a_1+2a_2\).)

(b) \(\dim W=2\).
\end{solution}

\begin{problem}\label{ex:r4-extend-drill}
Let \(u_1=(1,-2,1,-1)'\) and \(u_2=(0,1,-1,1)'\) in \(\R^4\), and let \(e_1,\dots,e_4\) be the
unit vectors.
\begin{enumerate}
\item Row-reduce \([u_1\ u_2\ e_1\ e_2\ e_3\ e_4]\in\R^{4\times6}\), and find a basis of \(\R^4\) that
consists of \(u_1\), \(u_2\) and two of the unit vectors.
\item Do \(u_1,u_2,e_1,e_2\) form a basis of \(\R^4\)?
\end{enumerate}
\end{problem}
\begin{solution}
(a) Call the matrix \(A=[u_1\ u_2\ e_1\ e_2\ e_3\ e_4]\); its columns span \(\R^4\), since they
include \(e_1,\dots,e_4\).
\[
\begin{gathered}
\left(\begin{array}{cc|cccc}1&0&1&0&0&0\\-2&1&0&1&0&0\\1&-1&0&0&1&0\\-1&1&0&0&0&1\end{array}\right)
\xrightarrow[R_4\leftarrow R_4+R_1]{R_2\leftarrow R_2+2R_1,\ R_3\leftarrow R_3-R_1}
\left(\begin{array}{cc|cccc}1&0&1&0&0&0\\0&1&2&1&0&0\\0&-1&-1&0&1&0\\0&1&1&0&0&1\end{array}\right)
\\
\xrightarrow[R_4\leftarrow R_4-R_2]{R_3\leftarrow R_3+R_2}
\left(\begin{array}{cc|cccc}1&0&1&0&0&0\\0&1&2&1&0&0\\0&0&1&1&1&0\\0&0&-1&-1&0&1\end{array}\right)
\xrightarrow{R_4\leftarrow R_4+R_3}
\left(\begin{array}{cc|cccc}1&0&1&0&0&0\\0&1&2&1&0&0\\0&0&1&1&1&0\\0&0&0&0&1&1\end{array}\right)
\end{gathered}
\]
The last matrix is an echelon form, with leading \(1\)s in columns \(1,2,3,5\); these are the pivot
columns of \(A\) (Remark~\ref{rem:r4-echelon}). By Theorem~\ref{thm:r4-pivot-columns}(b),
\(u_1,u_2,e_1,e_3\) form a basis of the span of the columns, which is \(\R^4\).

(b) No. Backward elimination finishes the RREF:
\[
\xrightarrow{R_3\leftarrow R_3-R_4}
\left(\begin{array}{cc|cccc}1&0&1&0&0&0\\0&1&2&1&0&0\\0&0&1&1&0&-1\\0&0&0&0&1&1\end{array}\right)
\xrightarrow[R_2\leftarrow R_2-2R_3]{R_1\leftarrow R_1-R_3}
\left(\begin{array}{cc|cccc}1&0&0&-1&0&1\\0&1&0&-1&0&2\\0&0&1&1&0&-1\\0&0&0&0&1&1\end{array}\right)
\]
Column \(4\) of the RREF gives \(e_2=-u_1-u_2+e_1\) (Theorem~\ref{thm:r4-pivot-columns}(b)). So
\(u_1,u_2,e_1,e_2\) are linearly dependent and do not form a basis. Check:
\(-(1,-2,1,-1)'-(0,1,-1,1)'+(1,0,0,0)'=(0,1,0,0)'=e_2\).
\end{solution}

\begin{problem}\label{ex:r4-solution-spaces}
Find the dimension of the space of solutions \((x,y,z)'\in\R^3\) of each system, and a basis of
it: (a) \(x-y+z=0\); (b) \(x+y+z=0\), \(x-y=0\), \(y+z=0\).
\end{problem}
\begin{solution}
The space of solutions is \(\mathcal N(A)\) for the coefficient matrix \(A\), of dimension
\(3-\rank A\) (Theorem~\ref{thm:r4-rank-nullity}).

(a) \(A=(1,-1,1)\) is a nonzero RREF with one row, so \(\rank A=1\) and the dimension is \(2\). The
vectors \((1,1,0)'\) and \((0,1,1)'\) are solutions (\(1-1+0=0\), \(0-1+1=0\)) and are
independent (\(c_1(1,1,0)'+c_2(0,1,1)'=(c_1,c_1+c_2,c_2)'=0\) forces \(c_1=c_2=0\)); so
they form a basis (Theorem~\ref{thm:r4-counting}(d)).

(b) Row-reduce the coefficient matrix:
\[
\begin{pmatrix}1&1&1\\1&-1&0\\0&1&1\end{pmatrix}
\xrightarrow{R_2\leftarrow R_2-R_1}
\begin{pmatrix}1&1&1\\0&-2&-1\\0&1&1\end{pmatrix}
\xrightarrow{R_2\leftrightarrow R_3}
\begin{pmatrix}1&1&1\\0&1&1\\0&-2&-1\end{pmatrix}
\xrightarrow{R_3\leftarrow R_3+2R_2}
\begin{pmatrix}1&1&1\\0&1&1\\0&0&1\end{pmatrix}
\]
This echelon form has three nonzero rows, so \(\rank A=3\) (Definition~\ref{def:r4-rank}): the
space of solutions is \(\{0\}\), of dimension \(0\), with the empty basis.
\end{solution}

\begin{problem}\label{ex:r4-four-subspaces-drill}
Find \(\rank A\) and bases of \(\mathcal C(A)\), \(\mathcal R(A)\), \(\mathcal N(A)\) and
\(\mathcal N(A')\) for
\[
A=\begin{pmatrix}1&2&1&0\\-1&-2&0&1\\1&2&2&1\end{pmatrix},
\]
and write the non-pivot columns of \(A\) as combinations of the pivot columns.
\end{problem}
\begin{solution}
Row-reduce \([A\mid I_3]\) (Theorem~\ref{thm:r4-four-bases}):
\[
\begin{gathered}
\left(\begin{array}{cccc|ccc}1&2&1&0&1&0&0\\-1&-2&0&1&0&1&0\\1&2&2&1&0&0&1\end{array}\right)
\xrightarrow[R_3\leftarrow R_3-R_1]{R_2\leftarrow R_2+R_1}
\left(\begin{array}{cccc|ccc}1&2&1&0&1&0&0\\0&0&1&1&1&1&0\\0&0&1&1&-1&0&1\end{array}\right)
\\
\xrightarrow{R_3\leftarrow R_3-R_2}
\left(\begin{array}{cccc|ccc}1&2&1&0&1&0&0\\0&0&1&1&1&1&0\\0&0&0&0&-2&-1&1\end{array}\right)
\xrightarrow{R_1\leftarrow R_1-R_2}
\left(\begin{array}{cccc|ccc}1&2&0&-1&0&-1&0\\0&0&1&1&1&1&0\\0&0&0&0&-2&-1&1\end{array}\right)
=[R\mid P]
\end{gathered}
\]
So \(\rank A=2\), with pivot columns \(1,3\) and free unknowns \(x_2,x_4\).
\begin{itemize}
\item \(\mathcal C(A)\): \(a_1=(1,-1,1)'\) and \(a_3=(1,0,2)'\). Columns \(2\) and \(4\) of \(R\) give
\(a_2=2a_1\) and \(a_4=-a_1+a_3\) (Theorem~\ref{thm:r4-pivot-columns}(b)). Check:
\(2(1,-1,1)'=(2,-2,2)'\) and \(-(1,-1,1)'+(1,0,2)'=(0,1,1)'\).
\item \(\mathcal R(A)\): \(\rho_1=(1,2,0,-1)\) and \(\rho_2=(0,0,1,1)\).
\item \(\mathcal N(A)\): \(Rx=0\) reads \(x_1+2x_2-x_4=0\), \(x_3+x_4=0\). The special solutions
\(s_1=(-2,1,0,0)'\) (\(x_2=1\), \(x_4=0\)) and \(s_2=(1,0,-1,1)'\) (\(x_2=0\), \(x_4=1\)) form a
basis; \(\dim\mathcal N(A)=2=4-2\).
\item \(\mathcal N(A')\): row \(3\) of \(P\) gives the basis \(y=(-2,-1,1)'\); \(\dim\mathcal N(A')=1=3-2\).
Indeed \(y'A=-2\row_1(A)-\row_2(A)+\row_3(A)=0\), that is,
\(\row_3(A)=2\row_1(A)+\row_2(A)\).
\end{itemize}
\end{solution}

\begin{problem}\label{ex:r4-rank-drill}
Find \(\rank A\), bases of \(\mathcal C(A)\), \(\mathcal R(A)\) and \(\mathcal N(A)\), and
\(\dim\mathcal N(A')\) for
\[
A=\begin{pmatrix}1&2&1&1\\0&1&1&-1\\1&1&0&2\\0&-2&-2&3\end{pmatrix},
\]
and write the non-pivot column of \(A\) as a combination of the pivot columns.
\end{problem}
\begin{solution}
Row-reduce \(A\):
\[
\begin{gathered}
A\xrightarrow{R_3\leftarrow R_3-R_1}
\begin{pmatrix}1&2&1&1\\0&1&1&-1\\0&-1&-1&1\\0&-2&-2&3\end{pmatrix}
\xrightarrow[R_4\leftarrow R_4+2R_2]{R_3\leftarrow R_3+R_2}
\begin{pmatrix}1&2&1&1\\0&1&1&-1\\0&0&0&0\\0&0&0&1\end{pmatrix}
\xrightarrow{R_3\leftrightarrow R_4}
\begin{pmatrix}1&2&1&1\\0&1&1&-1\\0&0&0&1\\0&0&0&0\end{pmatrix}
\\
\xrightarrow[R_1\leftarrow R_1-R_3]{R_2\leftarrow R_2+R_3}
\begin{pmatrix}1&2&1&0\\0&1&1&0\\0&0&0&1\\0&0&0&0\end{pmatrix}
\xrightarrow{R_1\leftarrow R_1-2R_2}
\begin{pmatrix}1&0&-1&0\\0&1&1&0\\0&0&0&1\\0&0&0&0\end{pmatrix}=R
\end{gathered}
\]
After the second step column \(3\) has no nonzero entry in rows \(3,4\), and the next pivot column,
\(4\), has its nonzero entry in row \(4\); so rows \(3\) and \(4\) are interchanged. The echelon form
after the third arrow has three nonzero rows: \(\rank A=3\), with pivot columns \(1,2,4\).
\begin{itemize}
\item \(\mathcal C(A)\): \(a_1=(1,0,1,0)'\), \(a_2=(2,1,1,-2)'\), \(a_4=(1,-1,2,3)'\). Column \(3\) of \(R\)
gives \(a_3=-a_1+a_2\); check: \(-(1,0,1,0)'+(2,1,1,-2)'=(1,1,0,-2)'=a_3\).
\item \(\mathcal R(A)\): \((1,0,-1,0)\), \((0,1,1,0)\), \((0,0,0,1)\).
\item \(\mathcal N(A)\): \(Rx=0\) reads \(x_1-x_3=0\), \(x_2+x_3=0\), \(x_4=0\); the special solution
\(s=(1,-1,1,0)'\) (\(x_3=1\)) is a basis, and \(\dim\mathcal N(A)=1=4-3\).
\item \(\dim\mathcal N(A')=4-\rank A=1\) (Theorem~\ref{thm:r4-rank-nullity}).
\end{itemize}
\end{solution}

\begin{problem}\label{ex:r4-sum-rank-one}
Let \(A\in\R^{m\times n}\) have rank \(r\). Show that \(A\) is the sum of \(r\) matrices of
rank \(1\).
\end{problem}
\begin{solution}
If \(r=0\), then \(A=0\), the empty sum. Let \(r\geqslant1\), and let \(A=BC\) be the full-rank
factorization of Theorem~\ref{thm:r4-full-rank-factorization}(a), where \(B\in\R^{m\times r}\) has columns
\(b_1,\dots,b_r\) and \(C\in\R^{r\times n}\). By the outer product expansion
(Theorem~\ref{r1-thm-views}(d)),
\[
A=\sum_{j=1}^rA_j,\qquad A_j=b_j\row_j(C)\in\R^{m\times n}
\]
\begin{itemize}
\item The columns of \(B\) are independent, and so are the rows of \(C\)
(Theorem~\ref{thm:r4-rank-rules}(b),(c)); hence every \(b_j\) and every \(\row_j(C)\) is nonzero,
since a list containing \(0\) is dependent (Proposition~\ref{r3-prop-dependence}(a)).
\item If the \(l\)th entry \(c_{jl}\) of \(\row_j(C)\) is nonzero, then
\(\col_l(A_j)=c_{jl}b_j\ne0\); so \(A_j\ne0\), and \(\rank A_j\geqslant1\).
\item \(\rank A_j\leqslant\rank b_j\leqslant1\) (Theorems~\ref{thm:r4-rank-product}(a)
and~\ref{thm:r4-rank-rules}(a)), as \(b_j\) has one column.
\end{itemize}
So each \(A_j\) has rank \(1\).
\end{solution}

\begin{problem}\label{ex:r4-rank-atab}
Show that \(\rank(A'AB)=\rank(AB)=\rank(ABB')\) for all \(A\in\R^{m\times n}\) and
\(B\in\R^{n\times p}\).
\end{problem}
\begin{solution}
Apply Theorem~\ref{thm:r4-rank-ata}(b) to \(AB\), with \((AB)'=B'A'\) (Theorem~\ref{r1-thm-transpose}),
and Theorem~\ref{thm:r4-rank-product}(a) twice:
\[
\begin{aligned}
\rank(AB)&=\rank\bigl((AB)'AB\bigr)=\rank\bigl(B'(A'AB)\bigr)\leqslant\rank(A'AB)
=\rank\bigl(A'(AB)\bigr)\leqslant\rank(AB),\\
\rank(AB)&=\rank\bigl(AB(AB)'\bigr)=\rank\bigl((ABB')A'\bigr)\leqslant\rank(ABB')
=\rank\bigl((AB)B'\bigr)\leqslant\rank(AB)
\end{aligned}
\]
So equality holds throughout.
\end{solution}

\begin{problem}\label{ex:r4-no-left-inverse}
Let \(A\in\R^{m\times n}\) with \(m<n\). Show that there is no \(B\in\R^{m\times n}\) with
\(B'A=I_n\).
\end{problem}
\begin{solution}
If \(B'A=I_n\), then by Theorems~\ref{thm:r4-rank-rules}(a),(d) and~\ref{thm:r4-rank-product}(a),
\[
n=\rank I_n=\rank(B'A)\leqslant\rank A\leqslant m<n\ \contradiction
\]
\end{solution}

\begin{problem}\label{ex:r4-involutory-rank}
Let \(A\in\R^{n\times n}\) satisfy \(A^2=I\). Show that \(\rank(I-A)+\rank(I+A)=n\).
\end{problem}
\begin{solution}
\begin{itemize}
\item \(2I=(I-A)+(I+A)\), and \(2I=\diag(2,\dots,2)\) is invertible (Definition~\ref{r2-def-inverse}). By
Theorem~\ref{thm:r4-rank-rules}(d) and Corollary~\ref{cor:r4-rank-sum},
\(n=\rank(2I)\leqslant\rank(I-A)+\rank(I+A)\).
\item \((I-A)(I+A)=I+A-A-A^2=I-A^2=0\) (Theorem~\ref{r1-thm-rules}), so
\(\rank(I-A)+\rank(I+A)\leqslant n\) by Theorem~\ref{thm:r4-sylvester}.
\end{itemize}
Together, \(\rank(I-A)+\rank(I+A)=n\).
\end{solution}

\clearpage
\def\unitnumber{5}
\setcounter{theorem}{0}\setcounter{problem}{0}
\section*{Chapter 5. Block matrices and determinants}
\subsection*{Block matrices}

\begin{definition}\label{r5b-def-partition}
A \textbf{partition} of \(A\in\R^{m\times n}\) splits the rows into consecutive groups of
\(m_1,\dots,m_r\) rows (\(m_1+\dots+m_r=m\)) and the columns into consecutive groups of
\(n_1,\dots,n_c\) columns (\(n_1+\dots+n_c=n\)). The \textbf{block}
\(A_{st}\in\R^{m_s\times n_t}\) consists of the entries in row group \(s\) and column group
\(t\), and we write
\[
A=\begin{pmatrix}A_{11}&\cdots&A_{1c}\\\vdots&&\vdots\\A_{r1}&\cdots&A_{rc}\end{pmatrix}
\]
Lines between rows and between columns may mark the partition. Two matrices are equal if and
only if their corresponding blocks are equal. Partitions of \(A,B\in\R^{m\times n}\) are
\textbf{conformal for addition} if they use the same row and column groups; partitions of
\(A\in\R^{m\times n}\) and \(B\in\R^{n\times p}\) are \textbf{conformal for multiplication}
if the column groups of \(A\) are the row groups of \(B\).
\end{definition}

\begin{theorem}\label{r5b-thm-block}
Let the partitions below be conformal.
\begin{enumerate}
\item \(A+B=(A_{st}+B_{st})\) and \(cA=(cA_{st})\) (\(c\in\R\)).
\item \(A'\) has \(c\) block rows and \(r\) block columns, and its \((t,s)\) block is
\(A_{st}'\): transpose the array of blocks and also each block.
\item If \(A\in\R^{m\times n}\) has blocks \(A_{st}\) (\(s\leqslant r\), \(t\leqslant c\)) and
\(B\in\R^{n\times p}\) has blocks \(B_{tu}\) (\(t\leqslant c\), \(u\leqslant v\)), then \(AB\)
has blocks
\[
(AB)_{su}=\sum_{t=1}^cA_{st}B_{tu}=A_{s1}B_{1u}+A_{s2}B_{2u}+\dots+A_{sc}B_{cu}
\]
\end{enumerate}
So conformal block matrices are added and multiplied as if the blocks were numbers, but the
order of the factors in each \(A_{st}B_{tu}\) must be kept. In particular, for \(2\times2\)
blocks,
\[
\begin{pmatrix}A_{11}&A_{12}\\A_{21}&A_{22}\end{pmatrix}
\begin{pmatrix}B_{11}&B_{12}\\B_{21}&B_{22}\end{pmatrix}
=\begin{pmatrix}A_{11}B_{11}+A_{12}B_{21}&A_{11}B_{12}+A_{12}B_{22}\\
A_{21}B_{11}+A_{22}B_{21}&A_{21}B_{12}+A_{22}B_{22}\end{pmatrix},
\]
and \([A_1\ A_2]\begin{pmatrix}B_1\\B_2\end{pmatrix}=A_1B_1+A_2B_2\),
\(\begin{pmatrix}A_1\\A_2\end{pmatrix}[B_1\ B_2]=\begin{pmatrix}A_1B_1&A_1B_2\\A_2B_1&A_2B_2\end{pmatrix}\),
\(A[B_1\ B_2]=[AB_1\ \ AB_2]\).
\end{theorem}
\begin{proof}
(a) and (b) hold entry by entry. (c) Let \(i\) lie in row group \(s\) of \(A\) and \(j\) in
column group \(u\) of \(B\). In \((AB)_{ij}=\sum_{h=1}^na_{ih}b_{hj}\), group the terms
according to the column group \(t\) of \(A\) that contains \(h\):
\((AB)_{ij}=\sum_{t=1}^c\sum_{h\text{ in group }t}a_{ih}b_{hj}\). The inner sum runs over a
row of \(A_{st}\) and a column of \(B_{tu}\), so it is the corresponding entry of
\(A_{st}B_{tu}\).
\end{proof}

\begin{example}\label{r5b-exm-block}
Partition
\[
A=\left(\begin{array}{cc|c}1&0&3\\2&-1&0\\\hline0&4&1\end{array}\right),\qquad
B=\left(\begin{array}{cc|c}3&3&1\\-4&2&-1\\\hline-2&0&0\end{array}\right)
\]
into blocks of sizes \(2\times2\), \(2\times1\), \(1\times2\), \(1\times1\). By
Theorem~\ref{r5b-thm-block}(c),
\[
\begin{aligned}
A_{11}B_{11}+A_{12}B_{21}&=\begin{pmatrix}3&3\\10&4\end{pmatrix}+\begin{pmatrix}-6&0\\0&0\end{pmatrix}=\begin{pmatrix}-3&3\\10&4\end{pmatrix},&
A_{11}B_{12}+A_{12}B_{22}&=\begin{pmatrix}1\\3\end{pmatrix}+\begin{pmatrix}0\\0\end{pmatrix},\\
A_{21}B_{11}+A_{22}B_{21}&=(-16,\ 8)+(-2,\ 0)=(-18,\ 8),&
A_{21}B_{12}+A_{22}B_{22}&=-4+0
\end{aligned}
\]
so \(AB=\left(\begin{array}{cc|c}-3&3&1\\10&4&3\\\hline-18&8&-4\end{array}\right)\), the same
matrix as the ordinary product. For instance
\(A_{12}B_{21}=\left(\begin{smallmatrix}3\\0\end{smallmatrix}\right)(-2,\ 0)=\left(\begin{smallmatrix}-6&0\\0&0\end{smallmatrix}\right)\).
By Theorem~\ref{r5b-thm-block}(b),
\(B'=\left(\begin{smallmatrix}B_{11}'&B_{21}'\\B_{12}'&B_{22}'\end{smallmatrix}\right)\).
\end{example}

\begin{definition}\label{r5b-def-block-tri}
Let \(A\in\R^{n\times n}\) be partitioned with the same groups \(n_1,\dots,n_r\) for rows and
columns, so that the \textbf{diagonal blocks} \(A_{ss}\) are square. \(A\) is \textbf{block
diagonal}, written \(A=\diag(A_{11},\dots,A_{rr})\), if \(A_{st}=0\) for \(s\ne t\);
\textbf{upper block triangular} if \(A_{st}=0\) for \(s>t\); and \textbf{lower block
triangular} if \(A_{st}=0\) for \(s<t\).
\end{definition}

\begin{proposition}\label{r5b-prop-block-tri}
Let \(A,B\in\R^{n\times n}\) be partitioned as in Definition~\ref{r5b-def-block-tri}, with the
same groups.
\begin{enumerate}
\item If \(A\) and \(B\) are upper (lower) block triangular, so is \(AB\), and its diagonal
blocks are \(A_{ss}B_{ss}\).
\item \(\diag(A_{11},\dots,A_{rr})\diag(B_{11},\dots,B_{rr})=\diag(A_{11}B_{11},\dots,A_{rr}B_{rr})\).
If every \(A_{ss}\) is invertible, then
\(\diag(A_{11},\dots,A_{rr})^{-1}=\diag(A_{11}^{-1},\dots,A_{rr}^{-1})\).
\item If \(A_{11}\in\R^{p\times p}\) and \(A_{22}\in\R^{q\times q}\) are invertible and
\(A_{12}\in\R^{p\times q}\), then
\[
\begin{pmatrix}A_{11}&A_{12}\\0&A_{22}\end{pmatrix}^{-1}
=\begin{pmatrix}A_{11}^{-1}&-A_{11}^{-1}A_{12}A_{22}^{-1}\\0&A_{22}^{-1}\end{pmatrix}
\]
\end{enumerate}
\end{proposition}
\begin{proof}
(a) The \((s,u)\) block of \(AB\) is \(\sum_tA_{st}B_{tu}\). For upper block triangular
\(A,B\), the term \(t\) can be nonzero only if \(s\leqslant t\leqslant u\): no term survives
when \(s>u\), and only \(t=s\) survives when \(s=u\). The lower case is the same with the
inequalities reversed.
(b) Block diagonal means upper and lower block triangular, so (a) gives the product. Since
\(I_n=\diag(I_{n_1},\dots,I_{n_r})\), the products of \(\diag(A_{ss})\) and
\(\diag(A_{ss}^{-1})\) in either order are \(I_n\) (Definition~\ref{r2-def-inverse}).
(c) Call the two matrices \(M\) and \(X\). By the \(2\times2\) block formula,
\[
MX=\begin{pmatrix}I_p&-A_{12}A_{22}^{-1}+A_{12}A_{22}^{-1}\\0&I_q\end{pmatrix}=I,\qquad
XM=\begin{pmatrix}I_p&A_{11}^{-1}A_{12}-A_{11}^{-1}A_{12}A_{22}^{-1}A_{22}\\0&I_q\end{pmatrix}=I\qedhere
\]
\end{proof}

\begin{example}\label{r5b-exm-block-inverse}
Find the inverse of
\[
M=\left(\begin{array}{cc|cc}2&1&1&3\\1&1&2&1\\\hline0&0&1&2\\0&0&0&1\end{array}\right)
\]

\textit{Solution.} Here \(A_{11}=\left(\begin{smallmatrix}2&1\\1&1\end{smallmatrix}\right)\),
\(A_{12}=\left(\begin{smallmatrix}1&3\\2&1\end{smallmatrix}\right)\),
\(A_{22}=\left(\begin{smallmatrix}1&2\\0&1\end{smallmatrix}\right)\), and by the \(2\times2\)
formula of Exercise~\ref{r2-ex-inverses-b},
\(A_{11}^{-1}=\left(\begin{smallmatrix}1&-1\\-1&2\end{smallmatrix}\right)\) and
\(A_{22}^{-1}=\left(\begin{smallmatrix}1&-2\\0&1\end{smallmatrix}\right)\). Then
\[
A_{11}^{-1}A_{12}=\begin{pmatrix}-1&2\\3&-1\end{pmatrix},\qquad
-A_{11}^{-1}A_{12}A_{22}^{-1}=-\begin{pmatrix}-1&2\\3&-1\end{pmatrix}\begin{pmatrix}1&-2\\0&1\end{pmatrix}
=\begin{pmatrix}1&-4\\-3&7\end{pmatrix},
\]
and by Proposition~\ref{r5b-prop-block-tri}(c)
\[
M^{-1}=\left(\begin{array}{cc|cc}1&-1&1&-4\\-1&2&-3&7\\\hline0&0&1&-2\\0&0&0&1\end{array}\right)
\]
\end{example}

\subsection*{Block elimination and the Schur complement}

\begin{proposition}\label{r5b-prop-block-elementary}
Let \(M=\begin{pmatrix}A&B\\C&D\end{pmatrix}\) with \(A\in\R^{m\times p}\),
\(B\in\R^{m\times q}\), \(C\in\R^{n\times p}\), \(D\in\R^{n\times q}\), and let
\(V\in\R^{n\times m}\), \(F\in\R^{m\times n}\), \(H\in\R^{p\times q}\). Then
\[
\begin{pmatrix}I_m&0\\V&I_n\end{pmatrix}M=\begin{pmatrix}A&B\\C+VA&D+VB\end{pmatrix},\qquad
\begin{pmatrix}I_m&F\\0&I_n\end{pmatrix}M=\begin{pmatrix}A+FC&B+FD\\C&D\end{pmatrix},
\]
\[
M\begin{pmatrix}I_p&H\\0&I_q\end{pmatrix}=\begin{pmatrix}A&B+AH\\C&D+CH\end{pmatrix}
\]
These \textbf{block elementary matrices} are invertible:
\(\left(\begin{smallmatrix}I&0\\V&I\end{smallmatrix}\right)^{-1}=\left(\begin{smallmatrix}I&0\\-V&I\end{smallmatrix}\right)\)
and
\(\left(\begin{smallmatrix}I&F\\0&I\end{smallmatrix}\right)^{-1}=\left(\begin{smallmatrix}I&-F\\0&I\end{smallmatrix}\right)\).
\end{proposition}
Each identity is the \(2\times2\) block formula; for example, the second block row of the
first product is \((VA+I_nC,\ VB+I_nD)\). Thus left multiplication adds \(V\) times the first
block row to the second, the block analogue of the row operation \(R_2\leftarrow R_2+cR_1\)
of Definition~\ref{r2-def-rowops}. The inverses follow from the products
\(\left(\begin{smallmatrix}I&0\\-V&I\end{smallmatrix}\right)\left(\begin{smallmatrix}I&0\\V&I\end{smallmatrix}\right)=I\)
and \(\left(\begin{smallmatrix}I&0\\V&I\end{smallmatrix}\right)\left(\begin{smallmatrix}I&0\\-V&I\end{smallmatrix}\right)=I\),
and likewise for \(F\).

\begin{theorem}\label{r5b-thm-schur}
Let \(M=\begin{pmatrix}A&B\\C&D\end{pmatrix}\) with \(A\in\R^{p\times p}\) invertible,
\(B\in\R^{p\times q}\), \(C\in\R^{n\times p}\), \(D\in\R^{n\times q}\). Then
\[
\begin{pmatrix}I_p&0\\-CA^{-1}&I_n\end{pmatrix}\begin{pmatrix}A&B\\C&D\end{pmatrix}
=\begin{pmatrix}A&B\\0&D-CA^{-1}B\end{pmatrix},\qquad
\begin{pmatrix}A&B\\C&D\end{pmatrix}\begin{pmatrix}I_p&-A^{-1}B\\0&I_q\end{pmatrix}
=\begin{pmatrix}A&0\\C&D-CA^{-1}B\end{pmatrix},
\]
and hence
\(\left(\begin{smallmatrix}I_p&0\\-CA^{-1}&I_n\end{smallmatrix}\right)M\left(\begin{smallmatrix}I_p&-A^{-1}B\\0&I_q\end{smallmatrix}\right)
=\left(\begin{smallmatrix}A&0\\0&D-CA^{-1}B\end{smallmatrix}\right)\). The matrix
\(D-CA^{-1}B\) is the \textbf{Schur complement} of \(A\) in \(M\).
\end{theorem}
\begin{proof}
In Proposition~\ref{r5b-prop-block-elementary} choose \(V\) so that the \((2,1)\) block
\(C+VA\) is \(0\): \(V=-CA^{-1}\); then \(D+VB=D-CA^{-1}B\). Likewise \(H=-A^{-1}B\) makes
\(B+AH=0\) and \(D+CH=D-CA^{-1}B\). Applying the column step to the result of the row step,
\(\left(\begin{smallmatrix}A&B\\0&S\end{smallmatrix}\right)\left(\begin{smallmatrix}I&-A^{-1}B\\0&I\end{smallmatrix}\right)
=\left(\begin{smallmatrix}A&-AA^{-1}B+B\\0&S\end{smallmatrix}\right)=\left(\begin{smallmatrix}A&0\\0&S\end{smallmatrix}\right)\)
with \(S=D-CA^{-1}B\).
\end{proof}
This is Gaussian elimination done one block at a time. When \(D\) is invertible, the matrix
\(A-BD^{-1}C\) is likewise called the \textbf{Schur complement} of \(D\) in \(M\).

\begin{example}\label{r5b-exm-schur}
With the data matrix \(X_0\in\R^{4\times2}\) of Example~\ref{r1-exm-data} (\(n=4\), \(k=2\)), let
\(X=[\one_4\ X_0]\) be the design matrix. By Theorem~\ref{r5b-thm-block},
\[
X'X=\begin{pmatrix}\one'\\X_0'\end{pmatrix}[\one\ \ X_0]
=\begin{pmatrix}\one'\one&\one'X_0\\X_0'\one&X_0'X_0\end{pmatrix}
=\left(\begin{array}{c|cc}4&4&8\\\hline4&8&8\\8&8&20\end{array}\right)
\]
Here \(A=(4)\), \(C=X_0'\one=(4,8)'\) and \(-CA^{-1}=-(1,2)'=-\bar x\), so
Theorem~\ref{r5b-thm-schur} gives
\[
\left(\begin{array}{c|cc}1&0&0\\\hline-1&1&0\\-2&0&1\end{array}\right)
\left(\begin{array}{c|cc}4&4&8\\\hline4&8&8\\8&8&20\end{array}\right)
=\left(\begin{array}{c|cc}4&4&8\\\hline0&4&0\\0&0&4\end{array}\right)
\]
The Schur complement of \(n=4\) is
\(D-CA^{-1}B=X_0'X_0-\frac14(X_0'\one)(\one'X_0)=\left(\begin{smallmatrix}4&0\\0&4\end{smallmatrix}\right)\),
the matrix \(X_0'C_4X_0\) of corrected sums of squares and products of
Example~\ref{r1-exm-data} (compare the first step of Example~\ref{r2-exm-xtx-inverse}). The
same holds for any data matrix \(X_0\in\R^{n\times k}\): for the design matrix
\(X=[\one_n\ X_0]\in\R^{n\times(k+1)}\), the Schur complement of \(\one'\one=n\) in \(X'X\) is
\(X_0'X_0-\frac1nX_0'\one\one'X_0=X_0'C_nX_0\) (the centering matrix keeps its subscript
here, since \(C\) is a block).
\end{example}

\subsection*{Determinants: definition and computation}

\begin{definition}\label{r5d-def-det}
Let \(A=(a_{ij})\in\R^{n\times n}\). For \(n\geqslant2\) and \(1\leqslant i,j\leqslant n\), let
\(A(i|j)\in\R^{(n-1)\times(n-1)}\) be the matrix obtained from \(A\) by deleting row \(i\) and
column \(j\). The \textbf{determinant} \(\det A\) is the number defined by \(\det(a_{11})=a_{11}\) if
\(n=1\) and, recursively, for \(n\geqslant2\) by
\[
\det A=\sum_{j=1}^n(-1)^{1+j}a_{1j}\det A(1|j)=a_{11}\det A(1|1)-a_{12}\det A(1|2)+\dots+(-1)^{1+n}a_{1n}\det A(1|n)
\]
The number \(\det A(i|j)\) is the \textbf{minor} of \(a_{ij}\), and \(c_{ij}=(-1)^{i+j}\det A(i|j)\)
is its \textbf{cofactor}. Thus \(\det A=a_{11}c_{11}+a_{12}c_{12}+\dots+a_{1n}c_{1n}\), the
\textbf{expansion along the first row}. Many books write \(\lvert A\rvert\) for \(\det A\).
\end{definition}
For \(n=2\), \(A(1|1)=(a_{22})\) and \(A(1|2)=(a_{21})\), so
\[
\det\begin{pmatrix}a_{11}&a_{12}\\a_{21}&a_{22}\end{pmatrix}=a_{11}a_{22}-a_{12}a_{21},
\]
the number \(ad-bc\) of the \(2\times2\) inverse formula of Exercise~\ref{r2-ex-inverses-b}. The signs
\((-1)^{i+j}\) form a checkerboard with \(+\) on the diagonal, for \(n=3\)
\(\left(\begin{smallmatrix}+&-&+\\-&+&-\\+&-&+\end{smallmatrix}\right)\).

\begin{theorem}\label{r5d-thm-expansion}
Let \(A\in\R^{n\times n}\) with \(n\geqslant2\), and let \(c_{ij}\) be the cofactors of \(A\). For every
row \(i\) and every column \(j\),
\[
\det A=a_{i1}c_{i1}+a_{i2}c_{i2}+\dots+a_{in}c_{in}=a_{1j}c_{1j}+a_{2j}c_{2j}+\dots+a_{nj}c_{nj}
\]
(\textbf{expansion along row \(i\)} and \textbf{along column \(j\)}).
\end{theorem}
We omit the proof: it is an induction on \(n\) in which the minors are expanded once more, and its only
difficulty is the bookkeeping of the signs \((-1)^{i+j}\). In practice one expands along a row or
column with many zeros.

\begin{example}\label{r5d-exm-expansion}
Compute \(\det A\) for \(A=\begin{pmatrix}1&2&3\\4&5&6\\2&3&5\end{pmatrix}\) by expansion along the
first row, and check the value by expansion along the second row.

\textit{Solution.} Along the first row,
\[
\det A=1\det\begin{pmatrix}5&6\\3&5\end{pmatrix}-2\det\begin{pmatrix}4&6\\2&5\end{pmatrix}
+3\det\begin{pmatrix}4&5\\2&3\end{pmatrix}=1(25-18)-2(20-12)+3(12-10)=-3
\]
so the cofactors of row \(1\) are \(c_{11}=7\), \(c_{12}=-8\), \(c_{13}=2\). Along the second row
(Theorem~\ref{r5d-thm-expansion}),
\[
\det A=-4\det\begin{pmatrix}2&3\\3&5\end{pmatrix}+5\det\begin{pmatrix}1&3\\2&5\end{pmatrix}
-6\det\begin{pmatrix}1&2\\2&3\end{pmatrix}=-4(10-9)+5(5-6)-6(3-4)=-3
\]
\end{example}

\begin{theorem}\label{r5d-thm-triangular}
Let \(A\in\R^{n\times n}\).
\begin{enumerate}
\item If \(A\) is upper or lower triangular, then \(\det A=a_{11}a_{22}\cdots a_{nn}\). In particular
\(\det\diag(d_1,\dots,d_n)=d_1d_2\cdots d_n\) and \(\det I_n=1\).
\item \(\det A'=\det A\).
\end{enumerate}
\end{theorem}
\begin{proof}
Both parts are proved by induction on \(n\); for \(n=1\) there is nothing to prove. Let \(n\geqslant2\).

(a)
\begin{itemize}
\item If \(A\) is lower (upper) triangular, the entries of its first row (column) other than
\(a_{11}\) are \(0\), so the definition (the expansion along column \(1\),
Theorem~\ref{r5d-thm-expansion}) gives \(\det A=a_{11}\det A(1|1)\).
\item \(A(1|1)\) is triangular of the same kind with diagonal entries \(a_{22},\dots,a_{nn}\), so
\(\det A(1|1)=a_{22}\cdots a_{nn}\) by induction.
\end{itemize}

(b)
\begin{itemize}
\item Row \(1\) of \(A'\) is \((a_{11},a_{21},\dots,a_{n1})\), and \(A'(1|j)=A(j|1)'\): both consist of the
entries \(a_{lk}\) with \(l\ne j\) and \(k\ne1\), arranged as transposes of each other.
\item By the definition and the induction hypothesis,
\(\det A'=\sum_{j=1}^n(-1)^{1+j}a_{j1}\det A(j|1)'=\sum_{j=1}^n(-1)^{j+1}a_{j1}\det A(j|1)\), which is
the expansion of \(\det A\) along column \(1\) (Theorem~\ref{r5d-thm-expansion}).\qedhere
\end{itemize}
\end{proof}

\begin{theorem}\label{r5d-thm-rowops}
Let \(A\in\R^{n\times n}\) with \(n\geqslant2\), let \(r\ne s\) be row indices, and let \(c\in\R\).
\begin{enumerate}
\item Multiplying row \(r\) of \(A\) by \(c\) multiplies \(\det A\) by \(c\). In particular, if \(A\) has
a zero row, then \(\det A=0\).
\item Interchanging rows \(r\) and \(s\) changes the sign of \(\det A\).
\item If two rows of \(A\) are equal, then \(\det A=0\).
\item Adding \(c\) times row \(s\) to row \(r\) does not change \(\det A\).
\item For the elementary matrices of Definition~\ref{r2-def-elementary}, \(\det D_r(c)=c\),
\(\det T_{rs}(c)=1\) and \(\det P_{rs}=-1\); and \(\det(EA)=\det E\cdot\det A\) for every elementary
matrix \(E\in\R^{n\times n}\).
\item Statements (a)--(d) hold with columns in place of rows.
\end{enumerate}
\end{theorem}
\begin{proof}
(a) Let \(B\) be \(A\) with row \(r\) multiplied by \(c\). Deleting row \(r\) gives \(B(r|j)=A(r|j)\), so
the expansion along row \(r\) (Theorem~\ref{r5d-thm-expansion}) gives
\(\det B=\sum_j(-1)^{r+j}ca_{rj}\det A(r|j)=c\det A\). A zero row is \(0\) times itself, so then
\(\det A=0\cdot\det A=0\).

(b) Let \(B\) be \(A\) with rows \(r\) and \(s\) interchanged; induction on \(n\).
\begin{itemize}
\item For \(n=2\),
\(\det\left(\begin{smallmatrix}a_{21}&a_{22}\\a_{11}&a_{12}\end{smallmatrix}\right)=a_{21}a_{12}-a_{22}a_{11}=-\det A\).
\item For \(n\geqslant3\), expand \(\det A\) and \(\det B\) along a row \(i\notin\{r,s\}\), which is the
same in both. Each \(B(i|j)\) is \(A(i|j)\) with two rows interchanged, so
\(\det B(i|j)=-\det A(i|j)\) by induction; hence \(\det B=-\det A\).
\end{itemize}

(c) Interchanging the two equal rows does not change \(A\), but by (b) it changes the sign of
\(\det A\). So \(\det A=-\det A\), that is, \(\det A=0\).

(d) Let \(B\) be \(A\) with row \(r\) replaced by \(\row_r(A)+c\row_s(A)\), and \(\tilde A\) be \(A\)
with row \(r\) replaced by \(\row_s(A)\). Since \(B(r|j)=\tilde A(r|j)=A(r|j)\), expanding along row
\(r\),
\[
\det B=\sum_{j=1}^n(-1)^{r+j}(a_{rj}+ca_{sj})\det A(r|j)=\det A+c\det\tilde A
\]
Rows \(r\) and \(s\) of \(\tilde A\) are equal, so \(\det\tilde A=0\) by (c).

(e) \(D_r(c)\) is diagonal and \(T_{rs}(c)\) is triangular with diagonal entries \(1\), so
\(\det D_r(c)=c\) and \(\det T_{rs}(c)=1\) (Theorem~\ref{r5d-thm-triangular}(a)); \(P_{rs}\) is \(I_n\)
with rows \(r\), \(s\) interchanged, so \(\det P_{rs}=-1\) by (b). \(EA\) is obtained from \(A\) by the
row operation of \(E\) (Theorem~\ref{r2-thm-elementary}(a)), so by (a), (d), (b) \(\det(EA)\) is
\(c\det A\), \(\det A\) or \(-\det A\), that is, \(\det E\cdot\det A\).

(f) The columns of \(A\) are the transposed rows of \(A'\), and \(\det A=\det A'\)
(Theorem~\ref{r5d-thm-triangular}(b)). An operation on the columns of \(A\) is the same operation on
the rows of \(A'\); so (a)--(d) for the rows of \(A'\) give (a)--(d) for the columns of \(A\).
\end{proof}

\textbf{Computing \(\det A\) by row reduction.} Forward elimination (Steps 1, 2 and 4 of
Definition~\ref{r2-def-gauss-jordan}, without making the pivots \(1\)) turns \(A\in\R^{n\times n}\) into
an upper triangular matrix \(U\) (the first nonzero entry of row \(i\) lies in column \(i\) or
later), whose determinant is the product of its diagonal entries
(Theorem~\ref{r5d-thm-triangular}(a)). By Theorem~\ref{r5d-thm-rowops}, each interchange of two rows
changes the sign and adding a multiple of a row to another row changes nothing, so
\(\det A=(-1)^tu_{11}u_{22}\cdots u_{nn}\), where \(t\) is the number of interchanges; if a row is
multiplied by \(c\ne0\) on the way, divide by \(c\). Column operations may be used in the same way, and
a row or column with a single nonzero entry can be expanded at once.

\begin{example}\label{r5d-exm-rowred3}
Compute \(\det A\) for \(A=\begin{pmatrix}0&2&-1\\1&3&2\\2&4&1\end{pmatrix}\).

\textit{Solution.} Since \(a_{11}=0\), begin with an interchange:
\[
A\xrightarrow{R_1\leftrightarrow R_2}
\begin{pmatrix}1&3&2\\0&2&-1\\2&4&1\end{pmatrix}
\xrightarrow{R_3\leftarrow R_3-2R_1}
\begin{pmatrix}1&3&2\\0&2&-1\\0&-2&-3\end{pmatrix}
\xrightarrow{R_3\leftarrow R_3+R_2}
\begin{pmatrix}1&3&2\\0&2&-1\\0&0&-4\end{pmatrix}=U
\]
The interchange changes the sign and the two additions change nothing, so
\(\det A=-\det U=-\bigl(1\cdot2\cdot(-4)\bigr)=8\). Check by expansion along column \(1\), whose entries are
\(0,1,2\):
\(\det A=-1\det\left(\begin{smallmatrix}2&-1\\4&1\end{smallmatrix}\right)
+2\det\left(\begin{smallmatrix}2&-1\\3&2\end{smallmatrix}\right)=-6+14=8\).
\end{example}

\begin{example}\label{r5d-exm-rowred4}
Compute \(\det A\) for
\(A=\begin{pmatrix}1&-1&2&3\\2&2&0&2\\4&1&-1&-1\\1&2&3&0\end{pmatrix}\).

\textit{Solution.}
\[
\begin{gathered}
A\xrightarrow[R_4\leftarrow R_4-R_1]{R_2\leftarrow R_2-2R_1,\ R_3\leftarrow R_3-4R_1}
\begin{pmatrix}1&-1&2&3\\0&4&-4&-4\\0&5&-9&-13\\0&3&1&-3\end{pmatrix}
\xrightarrow{R_2\leftarrow\frac14R_2}
\begin{pmatrix}1&-1&2&3\\0&1&-1&-1\\0&5&-9&-13\\0&3&1&-3\end{pmatrix}
\\
\xrightarrow[R_4\leftarrow R_4-3R_2]{R_3\leftarrow R_3-5R_2}
\begin{pmatrix}1&-1&2&3\\0&1&-1&-1\\0&0&-4&-8\\0&0&4&0\end{pmatrix}
\xrightarrow{R_4\leftarrow R_4+R_3}
\begin{pmatrix}1&-1&2&3\\0&1&-1&-1\\0&0&-4&-8\\0&0&0&-8\end{pmatrix}=U
\end{gathered}
\]
Only the second step changes the determinant: it multiplies it by \(\frac14\)
(Theorem~\ref{r5d-thm-rowops}(a)). Hence
\(\det A=4\det U=4\cdot1\cdot1\cdot(-4)\cdot(-8)=128\).
\end{example}

\subsection*{Main properties}

\begin{theorem}\label{r5d-thm-product}
Let \(A,B\in\R^{n\times n}\).
\begin{enumerate}
\item \(A\) is invertible \(\Leftrightarrow\) \(\det A\ne0\) \(\Leftrightarrow\) \(\rank A=n\).
\item \(\det(AB)=\det A\cdot\det B\).
\item If \(A\) is invertible, then \(\det A^{-1}=1/\det A\). Moreover \(\det(cA)=c^n\det A\) for
\(c\in\R\), and \(\det A^k=(\det A)^k\) for \(k=1,2,\dots\)
\end{enumerate}
\end{theorem}
\begin{proof}
For \(n=1\), \(\det(a)=a\) and all statements are clear; let \(n\geqslant2\).

(a)
\begin{itemize}
\item Let \(R=PA\) be the RREF of \(A\), where \(P=E_k\cdots E_1\) is a product of elementary matrices
(Theorem~\ref{r2-thm-elementary}(c)). By Theorem~\ref{r5d-thm-rowops}(e), applied \(k\) times,
\(\det R=\det E_k\cdots\det E_1\cdot\det A\), and each \(\det E_i\) is \(c\ne0\), \(1\) or \(-1\).
\item If \(A\) is invertible, then \(R=I_n\) (Theorem~\ref{r2-thm-invertible}), so \(\det A\ne0\). If
not, then \(\rank A<n\) (Theorem~\ref{thm:r4-rank-rules}(d)), so \(R\) has fewer than \(n\) nonzero
rows (Definition~\ref{def:r4-rank}), that is, a zero row; then \(\det R=0\)
(Theorem~\ref{r5d-thm-rowops}(a)) and \(\det A=0\). Finally, \(A\) is invertible if and only if
\(\rank A=n\) (Theorem~\ref{thm:r4-rank-rules}(d)).
\end{itemize}

(b)
\begin{itemize}
\item If \(A\) is not invertible, neither is \(AB\): an inverse \(C\) of \(AB\) would give
\(A(BC)=I_n\), and \(A\) would be invertible (Theorem~\ref{r2-thm-invertible}). By (a),
\(\det(AB)=0=\det A\cdot\det B\).
\item If \(A\) is invertible, then \(A=F_1\cdots F_k\) with elementary matrices \(F_i\)
(Theorem~\ref{r2-thm-invertible}). By Theorem~\ref{r5d-thm-rowops}(e), applied \(k\) times,
\(\det(AB)=\det F_1\cdots\det F_k\cdot\det B\); for \(B=I_n\) this reads
\(\det A=\det F_1\cdots\det F_k\), so \(\det(AB)=\det A\cdot\det B\).
\end{itemize}

(c) By (b), \(\det A\cdot\det A^{-1}=\det(AA^{-1})=\det I_n=1\). Next, \(cA=(cI_n)A\) and \(cI_n\) is
diagonal, so \(\det(cA)=c^n\det A\) by (b) and Theorem~\ref{r5d-thm-triangular}(a). Finally
\(\det A^{k+1}=\det(A^kA)=\det A^k\cdot\det A\) by (b), which gives the powers by induction on \(k\).
\end{proof}

\begin{theorem}\label{r5d-thm-adjugate}
Let \(A\in\R^{n\times n}\) with \(n\geqslant2\), and let \(c_{ij}\) be the cofactors of \(A\). The
\textbf{adjugate} of \(A\) is the transpose of the matrix of cofactors: \(\operatorname{adj}A=(c_{ij})'\),
that is, \((\operatorname{adj}A)_{ij}=c_{ji}\). Then
\[
A\operatorname{adj}A=(\det A)\,I_n
\]
Hence, if \(\det A\ne0\), then \(A^{-1}=\frac1{\det A}\operatorname{adj}A\).
\end{theorem}
\begin{proof}
\begin{itemize}
\item The \((i,k)\) entry of \(A\operatorname{adj}A\) is
\(\sum_{j=1}^na_{ij}(\operatorname{adj}A)_{jk}=\sum_{j=1}^na_{ij}c_{kj}\).
\item For \(k=i\), this is the expansion of \(\det A\) along row \(i\) (Theorem~\ref{r5d-thm-expansion}).
\item For \(k\ne i\), let \(\tilde A\) be \(A\) with row \(k\) replaced by row \(i\). Since
\(\tilde A(k|j)=A(k|j)\), the cofactors of row \(k\) of \(\tilde A\) are \(c_{k1},\dots,c_{kn}\), and the
sum is the expansion of \(\det\tilde A\) along row \(k\). As \(\tilde A\) has two equal rows, the sum is
\(0\) (Theorem~\ref{r5d-thm-rowops}(c)).
\item If \(\det A\ne0\), then \(A\bigl(\frac1{\det A}\operatorname{adj}A\bigr)=I_n\), so
\(\frac1{\det A}\operatorname{adj}A\) is the inverse of \(A\) (Theorem~\ref{r2-thm-invertible}).\qedhere
\end{itemize}
\end{proof}
For \(n=2\), \(\operatorname{adj}\left(\begin{smallmatrix}a&b\\c&d\end{smallmatrix}\right)
=\left(\begin{smallmatrix}d&-b\\-c&a\end{smallmatrix}\right)\), and the theorem gives again the formula of
Exercise~\ref{r2-ex-inverses-b}.

\begin{example}\label{r5d-exm-adjugate}
Find \(A^{-1}\) for \(A=\begin{pmatrix}0&1&1\\1&1&1\\3&4&3\end{pmatrix}\) by
Theorem~\ref{r5d-thm-adjugate}.

\textit{Solution.} The cofactors are
\[
\begin{aligned}
c_{11}&=\det\begin{pmatrix}1&1\\4&3\end{pmatrix}=-1,&
c_{12}&=-\det\begin{pmatrix}1&1\\3&3\end{pmatrix}=0,&
c_{13}&=\det\begin{pmatrix}1&1\\3&4\end{pmatrix}=1,\\
c_{21}&=-\det\begin{pmatrix}1&1\\4&3\end{pmatrix}=1,&
c_{22}&=\det\begin{pmatrix}0&1\\3&3\end{pmatrix}=-3,&
c_{23}&=-\det\begin{pmatrix}0&1\\3&4\end{pmatrix}=3,\\
c_{31}&=\det\begin{pmatrix}1&1\\1&1\end{pmatrix}=0,&
c_{32}&=-\det\begin{pmatrix}0&1\\1&1\end{pmatrix}=1,&
c_{33}&=\det\begin{pmatrix}0&1\\1&1\end{pmatrix}=-1
\end{aligned}
\]
Along row \(1\), \(\det A=0\cdot(-1)+1\cdot0+1\cdot1=1\). Hence
\[
A^{-1}=\frac11\operatorname{adj}A=\begin{pmatrix}c_{11}&c_{21}&c_{31}\\c_{12}&c_{22}&c_{32}\\c_{13}&c_{23}&c_{33}\end{pmatrix}
=\begin{pmatrix}-1&1&0\\0&-3&1\\1&3&-1\end{pmatrix},
\]
the inverse found by row reduction in Exercise~\ref{r2-ex-inverse}. Check:
\(\row_3(A)\operatorname{adj}A=(3,4,3)\operatorname{adj}A=(-3+0+3,\ 3-12+9,\ 0+4-3)=(0,0,1)\).
\end{example}

\subsection*{Determinants of block matrices and the Schur complement}

\begin{theorem}\label{r5d-thm-block-tri}
Let \(A\in\R^{p\times p}\), \(B\in\R^{p\times q}\), \(C\in\R^{q\times p}\) and \(D\in\R^{q\times q}\). Then
\[
\det\begin{pmatrix}A&B\\0&D\end{pmatrix}=\det A\cdot\det D=\det\begin{pmatrix}A&0\\C&D\end{pmatrix}
\]
In particular \(\det\diag(A,D)=\det A\cdot\det D\).
\end{theorem}
\begin{proof}
\begin{itemize}
\item The first row of \(\left(\begin{smallmatrix}I_p&0\\0&D\end{smallmatrix}\right)\) is
\((1,0,\dots,0)\), and deleting row \(1\) and column \(1\) leaves
\(\left(\begin{smallmatrix}I_{p-1}&0\\0&D\end{smallmatrix}\right)\); expanding along the first row \(p\)
times gives \(\det\left(\begin{smallmatrix}I_p&0\\0&D\end{smallmatrix}\right)=\det D\). Likewise,
expanding \(\left(\begin{smallmatrix}A&0\\0&I_q\end{smallmatrix}\right)\) \(q\) times along its last row
\((0,\dots,0,1)\), whose entry \(1\) is a diagonal entry (sign \(+1\)), gives \(\det A\).
\item \(\left(\begin{smallmatrix}I_p&B\\0&I_q\end{smallmatrix}\right)\) is upper triangular with
diagonal entries \(1\), so its determinant is \(1\) (Theorem~\ref{r5d-thm-triangular}(a)).
\item By the block product formula (Theorem~\ref{r5b-thm-block}),
\[
\begin{pmatrix}I_p&0\\0&D\end{pmatrix}\begin{pmatrix}I_p&B\\0&I_q\end{pmatrix}\begin{pmatrix}A&0\\0&I_q\end{pmatrix}
=\begin{pmatrix}I_p&0\\0&D\end{pmatrix}\begin{pmatrix}A&B\\0&I_q\end{pmatrix}=\begin{pmatrix}A&B\\0&D\end{pmatrix}
\]
and by Theorem~\ref{r5d-thm-product}(b) its determinant is \(\det D\cdot1\cdot\det A\).
\item The transpose of \(\left(\begin{smallmatrix}A&0\\C&D\end{smallmatrix}\right)\) is
\(\left(\begin{smallmatrix}A'&C'\\0&D'\end{smallmatrix}\right)\) (Theorem~\ref{r5b-thm-block}(b)); by
Theorem~\ref{r5d-thm-triangular}(b) and the case just proved, its determinant is
\(\det A'\cdot\det D'=\det A\cdot\det D\). Finally, \(\diag(A,D)\) is the case \(B=0\).\qedhere
\end{itemize}
\end{proof}

\begin{theorem}\label{r5d-thm-schur-det}
Let \(M=\begin{pmatrix}A&B\\C&D\end{pmatrix}\) with \(A\in\R^{p\times p}\), \(B\in\R^{p\times q}\),
\(C\in\R^{q\times p}\) and \(D\in\R^{q\times q}\).
\begin{enumerate}
\item If \(A\) is invertible, then \(\det M=\det A\cdot\det(D-CA^{-1}B)\).
\item If \(D\) is invertible, then \(\det M=\det D\cdot\det(A-BD^{-1}C)\).
\end{enumerate}
In particular, if \(A\) is invertible, then \(M\) is invertible if and only if its Schur complement
\(D-CA^{-1}B\) is invertible.
\end{theorem}
\begin{proof}
(a) By Theorem~\ref{r5b-thm-schur},
\(\left(\begin{smallmatrix}I_p&0\\-CA^{-1}&I_q\end{smallmatrix}\right)M
=\left(\begin{smallmatrix}A&B\\0&D-CA^{-1}B\end{smallmatrix}\right)\). The left factor is lower
triangular with diagonal entries \(1\), so its determinant is \(1\). By
Theorems~\ref{r5d-thm-product}(b) and~\ref{r5d-thm-block-tri},
\(\det M=\det A\cdot\det(D-CA^{-1}B)\).

(b) By the block product formula,
\(\left(\begin{smallmatrix}I_p&-BD^{-1}\\0&I_q\end{smallmatrix}\right)M
=\left(\begin{smallmatrix}A-BD^{-1}C&B-BD^{-1}D\\C&D\end{smallmatrix}\right)
=\left(\begin{smallmatrix}A-BD^{-1}C&0\\C&D\end{smallmatrix}\right)\). The left factor is upper
triangular with diagonal entries \(1\); as in (a), \(\det M=\det(A-BD^{-1}C)\cdot\det D\).

If \(A\) is invertible, then \(\det A\ne0\), so by (a) \(\det M\ne0\) if and only if
\(\det(D-CA^{-1}B)\ne0\); apply Theorem~\ref{r5d-thm-product}(a) to \(M\) and to \(D-CA^{-1}B\).
\end{proof}
For \(p=q=1\) and \(a\ne0\), (a) reads
\(\det\left(\begin{smallmatrix}a&b\\c&d\end{smallmatrix}\right)=a\bigl(d-\frac{cb}a\bigr)=ad-bc\).

\begin{theorem}\label{r5d-thm-det-lemma}
Let \(A\in\R^{n\times n}\) be invertible and \(u,v\in\R^n\). Then
\[
\det(A+uv')=\det A\cdot(1+v'A^{-1}u)
\]
In particular \(\det(I_n+uv')=1+v'u\). (This is the \textbf{matrix determinant lemma}.)
\end{theorem}
\begin{proof}
Apply Theorem~\ref{r5d-thm-schur-det} to
\(M=\left(\begin{smallmatrix}A&-u\\v'&1\end{smallmatrix}\right)\in\R^{(n+1)\times(n+1)}\) (\(p=n\), \(q=1\)).
\begin{itemize}
\item By (a), with the invertible block \(A\), whose Schur complement is the \(1\times1\) matrix
\(1-v'A^{-1}(-u)\): \(\det M=\det A\cdot(1+v'A^{-1}u)\).
\item By (b), with the invertible block \(D=(1)\): \(\det M=\det(1)\cdot\det\bigl(A-(-u)(1)^{-1}v'\bigr)=\det(A+uv')\).
\item Comparing gives the formula; for \(A=I_n\), \(A^{-1}=I_n\).\qedhere
\end{itemize}
\end{proof}

\begin{example}\label{r5d-exm-equicorrelation}
Let \(n\geqslant2\) and \(\rho\in\R\). The \textbf{equicorrelation matrix}
\(R=(1-\rho)I_n+\rho J_n\in\R^{n\times n}\), with diagonal entries \(1\) and all other entries
\(\rho\), is the correlation matrix of \(n\) variables any two of which have correlation \(\rho\).
Find \(\det R\), and the values of \(\rho\) for which \(R\) is invertible.

\textit{Solution.}
\begin{itemize}
\item If \(\rho=1\), then \(R=J_n\) has equal rows (\(n\geqslant2\)), so \(\det R=0\)
(Theorem~\ref{r5d-thm-rowops}(c)).
\item Let \(\rho\ne1\). As \(J_n=\one_n\one_n'\), \(R=A+uv'\) with \(A=(1-\rho)I_n\), \(u=\rho\one_n\) and
\(v=\one_n\). Here \(\det A=(1-\rho)^n\) (Theorem~\ref{r5d-thm-triangular}(a)),
\(A^{-1}=(1-\rho)^{-1}I_n\) and \(v'A^{-1}u=\frac{\rho}{1-\rho}\one'\one=\frac{n\rho}{1-\rho}\)
(Proposition~\ref{r1-prop-centering}). By Theorem~\ref{r5d-thm-det-lemma},
\[
\det R=(1-\rho)^n\Bigl(1+\frac{n\rho}{1-\rho}\Bigr)=(1-\rho)^{n-1}(1-\rho+n\rho)
=(1-\rho)^{n-1}\bigl(1+(n-1)\rho\bigr)
\]
This formula also gives \(0\) for \(\rho=1\).
\item By Theorem~\ref{r5d-thm-product}(a), \(R\) is invertible if and only if \(\det R\ne0\), that is,
\(\rho\ne1\) and \(\rho\ne-1/(n-1)\). For instance, for \(n=3\) and \(\rho=-\frac12\) every row of
\(R\) has sum \(0\), so \(R\one_3=0\).
\end{itemize}
\end{example}

\subsection*{Rank and the Schur complement}

\begin{theorem}\label{r5d-thm-schur-rank}
Let \(M=\begin{pmatrix}A&B\\C&D\end{pmatrix}\).
\begin{enumerate}
\item If \(A\in\R^{p\times p}\) is invertible, and \(B\in\R^{p\times q}\), \(C\in\R^{n\times p}\),
\(D\in\R^{n\times q}\), then \(\rank M=p+\rank(D-CA^{-1}B)\).
\item If \(D\in\R^{q\times q}\) is invertible, and \(A\in\R^{m\times p}\), \(B\in\R^{m\times q}\),
\(C\in\R^{q\times p}\), then \(\rank M=q+\rank(A-BD^{-1}C)\).
\end{enumerate}
\end{theorem}
\begin{proof}
(a)
\begin{itemize}
\item By Theorem~\ref{r5b-thm-schur},
\(\left(\begin{smallmatrix}I_p&0\\-CA^{-1}&I_n\end{smallmatrix}\right)M
\left(\begin{smallmatrix}I_p&-A^{-1}B\\0&I_q\end{smallmatrix}\right)
=\left(\begin{smallmatrix}A&0\\0&D-CA^{-1}B\end{smallmatrix}\right)\), and the two outer factors are
invertible (Proposition~\ref{r5b-prop-block-elementary}).
\item Invertible factors do not change the rank (Theorem~\ref{thm:r4-rank-product}(d)), and the rank
of a block-diagonal matrix is the sum of the ranks of its blocks
(Theorem~\ref{thm:r4-partitioned}(c)); so \(\rank M=\rank A+\rank(D-CA^{-1}B)\), where \(\rank A=p\)
(Theorem~\ref{thm:r4-rank-rules}(d)).
\end{itemize}

(b)
\begin{itemize}
\item By the block product formula,
\(\left(\begin{smallmatrix}I_m&-BD^{-1}\\0&I_q\end{smallmatrix}\right)M
=\left(\begin{smallmatrix}A-BD^{-1}C&0\\C&D\end{smallmatrix}\right)\) (as in the proof of
Theorem~\ref{r5d-thm-schur-det}(b)), and multiplying on the right by
\(\left(\begin{smallmatrix}I_p&0\\-D^{-1}C&I_q\end{smallmatrix}\right)\) gives
\(\left(\begin{smallmatrix}A-BD^{-1}C&0\\C-DD^{-1}C&D\end{smallmatrix}\right)
=\left(\begin{smallmatrix}A-BD^{-1}C&0\\0&D\end{smallmatrix}\right)\).
\item Both outer factors are invertible (Proposition~\ref{r5b-prop-block-elementary}); as in (a),
\(\rank M=\rank(A-BD^{-1}C)+\rank D=\rank(A-BD^{-1}C)+q\).\qedhere
\end{itemize}
\end{proof}

\begin{example}\label{r5d-exm-oneway}
Let \(d_1=(1,1,1,0,0,0)'\), \(d_2=(0,0,0,1,1,0)'\) and \(d_3=(0,0,0,0,0,1)'\) be the indicators of three
groups of \(3\), \(2\) and \(1\) observations, and let \(X=[\one_6\ d_1\ d_2\ d_3]\in\R^{6\times4}\). Find
\(\rank(X'X)\) and \(\det(X'X)\) with a Schur complement, and \(\rank X\).

\textit{Solution.} The \((j,l)\) entry of \(X'X\) is \(\col_j(X)'\col_l(X)\) (Example~\ref{r1-exm-data}):
\(\one'\one=6\), \(\one'd_l=d_l'd_l=n_l\), the size of group \(l\) (\(n_1=3\), \(n_2=2\), \(n_3=1\)), and
\(d_l'd_m=0\) for \(l\ne m\), as no observation lies in two groups. Partition after the first row and
column:
\[
X'X=\left(\begin{array}{c|ccc}6&3&2&1\\\hline3&3&0&0\\2&0&2&0\\1&0&0&1\end{array}\right)
=\begin{pmatrix}A&B\\C&D\end{pmatrix},\qquad A=(6),\quad B=(3,2,1),\quad C=B',\quad D=\diag(3,2,1)
\]
\(D\) is invertible with \(D^{-1}=\diag(\frac13,\frac12,1)\), so \(BD^{-1}=(1,1,1)\), and the Schur
complement of \(D\) is \(A-BD^{-1}C=6-(1,1,1)(3,2,1)'=6-6=0\). Hence
\(\rank(X'X)=3+\rank(0)=3\) (Theorem~\ref{r5d-thm-schur-rank}(b)) and
\(\det(X'X)=\det D\cdot0=0\) (Theorem~\ref{r5d-thm-schur-det}(b)): \(X'X\) is not invertible. Also
\(\rank X=\rank(X'X)=3\) (Theorem~\ref{thm:r4-rank-ata}(b)), so the four columns of \(X\) are linearly
dependent (Theorem~\ref{thm:r4-rank-rules}(b)); indeed \(d_1+d_2+d_3=\one_6\), that is,
\(X(-1,1,1,1)'=0\).
\end{example}

\subsection*{Exercises and solutions}

\begin{problem}\label{r5b-ex-partitioned}
Let \(A\) and \(B\) be the \(3\times5\) matrices and \(C\) the \(5\times4\) matrix partitioned
as
\[
A=\left(\begin{array}{ccc|cc}1&3&-2&1&2\\2&1&0&-1&3\\\hline0&0&1&2&1\end{array}\right),\quad
B=\left(\begin{array}{ccc|cc}1&-3&-2&2&1\\1&2&3&2&0\\\hline1&0&2&0&1\end{array}\right),\quad
C=\left(\begin{array}{cc|cc}1&0&2&1\\0&2&0&0\\-1&0&3&1\\\hline3&1&0&2\\2&-1&1&1\end{array}\right)
\]
With blocks \(A=(A_{st})\), \(B=(B_{st})\), \(C=(C_{st})\) (\(s,t=1,2\)), show that
\[
A+B=(A_{st}+B_{st}),\qquad
A'=\begin{pmatrix}A_{11}'&A_{21}'\\A_{12}'&A_{22}'\end{pmatrix},\qquad
AC=\begin{pmatrix}A_{11}C_{11}+A_{12}C_{21}&A_{11}C_{12}+A_{12}C_{22}\\A_{21}C_{11}+A_{22}C_{21}&A_{21}C_{12}+A_{22}C_{22}\end{pmatrix}
\]
\end{problem}
\begin{solution}
Adding entrywise,
\(A+B=\left(\begin{array}{ccc|cc}2&0&-4&3&3\\3&3&3&1&3\\\hline1&0&3&2&2\end{array}\right)\),
and each block is the sum of the corresponding blocks of \(A\) and \(B\). Transposing,
\[
A'=\left(\begin{array}{cc|c}1&2&0\\3&1&0\\-2&0&1\\\hline1&-1&2\\2&3&1\end{array}\right)
=\begin{pmatrix}A_{11}'&A_{21}'\\A_{12}'&A_{22}'\end{pmatrix}
\]
Multiplying in the ordinary way,
\(AC=\left(\begin{array}{cc|cc}10&5&-2&3\\5&-2&7&3\\\hline7&1&4&6\end{array}\right)\).
The blockwise products give the same four blocks:
\[
\begin{aligned}
A_{11}C_{11}+A_{12}C_{21}&=\begin{pmatrix}3&6\\2&2\end{pmatrix}+\begin{pmatrix}7&-1\\3&-4\end{pmatrix}=\begin{pmatrix}10&5\\5&-2\end{pmatrix},\\
A_{11}C_{12}+A_{12}C_{22}&=\begin{pmatrix}-4&-1\\4&2\end{pmatrix}+\begin{pmatrix}2&4\\3&1\end{pmatrix}=\begin{pmatrix}-2&3\\7&3\end{pmatrix},\\
A_{21}C_{11}+A_{22}C_{21}&=(-1,\ 0)+(8,\ 1)=(7,\ 1),\\
A_{21}C_{12}+A_{22}C_{22}&=(3,\ 1)+(1,\ 5)=(4,\ 6)
\end{aligned}
\]
For instance,
\(A_{12}C_{21}=\left(\begin{smallmatrix}1&2\\-1&3\end{smallmatrix}\right)\left(\begin{smallmatrix}3&1\\2&-1\end{smallmatrix}\right)=\left(\begin{smallmatrix}7&-1\\3&-4\end{smallmatrix}\right)\).
This agrees with Theorem~\ref{r5b-thm-block}.
\end{solution}

\begin{problem}\label{r5b-ex-inverse-blocks}
Let \(A\in\R^{n\times n}\) be invertible and partitioned by columns as \(A=[A_1\ A_2]\), where
\(A_1\) has \(k\) columns, and let \(A^{-1}\) be partitioned by rows into \(B_1\) (\(k\) rows)
and \(B_2\). Show that
\[
B_1A_1=I,\quad B_1A_2=0,\quad B_2A_1=0,\quad B_2A_2=I,\qquad
A_1B_1=I-A_2B_2,\quad A_2B_2=I-A_1B_1
\]
\end{problem}
\begin{solution}
With the partition \((k,n-k)\) of \(I_n\), Theorem~\ref{r5b-thm-block}(c) gives
\[
\begin{pmatrix}B_1A_1&B_1A_2\\B_2A_1&B_2A_2\end{pmatrix}
=\begin{pmatrix}B_1\\B_2\end{pmatrix}[A_1\ A_2]=A^{-1}A=I_n=\begin{pmatrix}I_k&0\\0&I_{n-k}\end{pmatrix},
\]
and comparing blocks gives the first four identities. Also
\(A_1B_1+A_2B_2=[A_1\ A_2]\left(\begin{smallmatrix}B_1\\B_2\end{smallmatrix}\right)=AA^{-1}=I_n\),
which gives the last two.
\end{solution}

\begin{problem}\label{r5b-ex-unipotent}
Show that the matrix
\(Z=\begin{pmatrix}-I_m&B\\C&I_n\end{pmatrix}\) (\(B\in\R^{m\times n}\), \(C\in\R^{n\times m}\))
satisfies \(Z^2=I_{m+n}\) if \(B=0\) or \(C=0\). Is this condition necessary?
\end{problem}
\begin{solution}
By the \(2\times2\) block formula (Theorem~\ref{r5b-thm-block}),
\[
Z^2=\begin{pmatrix}(-I)(-I)+BC&-B+B\\-C+C&CB+I\end{pmatrix}=\begin{pmatrix}I_m+BC&0\\0&I_n+CB\end{pmatrix},
\]
so \(Z^2=I_{m+n}\) if and only if \(BC=0\) and \(CB=0\). This holds if \(B=0\) or \(C=0\).
The condition is not necessary when \(m,n\geqslant2\): for \(m=n=2\) and
\(B=C=\left(\begin{smallmatrix}0&1\\0&0\end{smallmatrix}\right)\ne0\) we have \(BC=CB=0\)
(Exercise~\ref{r1-ex-unlike-scalars}(b)); for larger \(m,n\) put this \(2\times2\) matrix in
the top left corner of zero matrices. If \(m=1\), then \(CB\in\R^{n\times n}\) has entries
\(c_ib_j\), so \(CB=0\) forces \(B=0\) or \(C=0\); the case \(n=1\) is the same with \(BC\).
\end{solution}

\begin{problem}\label{r5b-ex-commuting}
Show that
\(Z_1=\begin{pmatrix}I_m&B_1\\0&I_n\end{pmatrix}\) and
\(Z_2=\begin{pmatrix}I_m&B_2\\0&I_n\end{pmatrix}\) (\(B_1,B_2\in\R^{m\times n}\)) commute.
\end{problem}
\begin{solution}
Write \(Z_1=I_{m+n}+X_1\) and \(Z_2=I_{m+n}+X_2\) with
\(X_i=\left(\begin{smallmatrix}0&B_i\\0&0\end{smallmatrix}\right)\). By the block formula,
\(X_1X_2=\left(\begin{smallmatrix}0\cdot0+B_1\cdot0&0\cdot B_2+B_1\cdot0\\0&0\end{smallmatrix}\right)=0\),
and likewise \(X_2X_1=0\). Hence
\[
Z_1Z_2=I+X_1+X_2+X_1X_2=I+X_1+X_2=I+X_2+X_1+X_2X_1=Z_2Z_1
\]
\end{solution}

\begin{problem}\label{r5b-ex-block-tri-inverse}
If \(A\in\R^{m\times m}\) and \(D\in\R^{n\times n}\) are invertible and \(C\in\R^{n\times m}\), show that
\[
\begin{pmatrix}A&0\\C&D\end{pmatrix}^{-1}=\begin{pmatrix}A^{-1}&0\\-D^{-1}CA^{-1}&D^{-1}\end{pmatrix}
\]
\end{problem}
\begin{solution}
Look for the inverse in the form
\(\left(\begin{smallmatrix}A^{-1}&0\\R&D^{-1}\end{smallmatrix}\right)\). By the block formula,
\[
\begin{pmatrix}A&0\\C&D\end{pmatrix}\begin{pmatrix}A^{-1}&0\\R&D^{-1}\end{pmatrix}
=\begin{pmatrix}I_m&0\\CA^{-1}+DR&I_n\end{pmatrix},
\]
which is \(I\) if and only if \(CA^{-1}+DR=0\), that is, \(R=-D^{-1}CA^{-1}\). With this \(R\),
the product in the other order is
\(\left(\begin{smallmatrix}A^{-1}A&0\\RA+D^{-1}C&D^{-1}D\end{smallmatrix}\right)
=\left(\begin{smallmatrix}I&0\\-D^{-1}C+D^{-1}C&I\end{smallmatrix}\right)=I\).
\end{solution}

\begin{problem}\label{r5b-ex-block-numeric}
Use Exercise~\ref{r5b-ex-block-tri-inverse} to find the inverse of
\[
N=\left(\begin{array}{cc|cc}1&0&0&0\\2&1&0&0\\\hline1&0&3&1\\4&1&2&1\end{array}\right),
\]
and check the answer by multiplication.
\end{problem}
\begin{solution}
Here \(A=\left(\begin{smallmatrix}1&0\\2&1\end{smallmatrix}\right)\),
\(C=\left(\begin{smallmatrix}1&0\\4&1\end{smallmatrix}\right)\),
\(D=\left(\begin{smallmatrix}3&1\\2&1\end{smallmatrix}\right)\). By the \(2\times2\) formula of
Exercise~\ref{r2-ex-inverses-b}, \(A^{-1}=\left(\begin{smallmatrix}1&0\\-2&1\end{smallmatrix}\right)\)
and \(D^{-1}=\left(\begin{smallmatrix}1&-1\\-2&3\end{smallmatrix}\right)\). Then
\[
CA^{-1}=\begin{pmatrix}1&0\\4&1\end{pmatrix}\begin{pmatrix}1&0\\-2&1\end{pmatrix}=\begin{pmatrix}1&0\\2&1\end{pmatrix},\qquad
-D^{-1}CA^{-1}=-\begin{pmatrix}1&-1\\-2&3\end{pmatrix}\begin{pmatrix}1&0\\2&1\end{pmatrix}=\begin{pmatrix}1&1\\-4&-3\end{pmatrix},
\]
so
\[
N^{-1}=\left(\begin{array}{cc|cc}1&0&0&0\\-2&1&0&0\\\hline1&1&1&-1\\-4&-3&-2&3\end{array}\right)
\]
Check: the \((2,1)\) block of \(NN^{-1}\) is
\(CA^{-1}+D\left(\begin{smallmatrix}1&1\\-4&-3\end{smallmatrix}\right)
=\left(\begin{smallmatrix}1&0\\2&1\end{smallmatrix}\right)+\left(\begin{smallmatrix}-1&0\\-2&-1\end{smallmatrix}\right)=0\),
the \((1,2)\) block is \(A\cdot0+0\cdot D^{-1}=0\), and the diagonal blocks are \(AA^{-1}=I\)
and \(DD^{-1}=I\); so \(NN^{-1}=I_4\), and \(N^{-1}\) is the inverse by
Theorem~\ref{r2-thm-invertible}.
\end{solution}

\begin{problem}\label{r5b-ex-partitioned-inverse}
Let \(M=\begin{pmatrix}A&B\\C&D\end{pmatrix}\) with \(A\in\R^{p\times p}\), \(B\in\R^{p\times q}\),
\(C\in\R^{q\times p}\) and \(D\in\R^{q\times q}\).
\begin{enumerate}
\item If \(A\) and its Schur complement \(E=D-CA^{-1}B\) are invertible, show that \(M\) is invertible
and
\[
M^{-1}=\begin{pmatrix}A^{-1}+A^{-1}BE^{-1}CA^{-1}&-A^{-1}BE^{-1}\\-E^{-1}CA^{-1}&E^{-1}\end{pmatrix}
\]
\item If \(D\) and \(F=A-BD^{-1}C\) are invertible, show that \(M\) is invertible and
\[
M^{-1}=\begin{pmatrix}F^{-1}&-F^{-1}BD^{-1}\\-D^{-1}CF^{-1}&D^{-1}+D^{-1}CF^{-1}BD^{-1}\end{pmatrix}
\]
\item Let \(X=[\one_n\ X_0]\in\R^{n\times(k+1)}\) be a design matrix with \(X'X\) invertible, let
\(\bar x=\frac1nX_0'\one\) be the mean vector, and let \(S=X_0'C_nX_0\). Show that \(S\) is invertible
and
\[
(X'X)^{-1}=\begin{pmatrix}\frac1n+\bar x'S^{-1}\bar x&-\bar x'S^{-1}\\-S^{-1}\bar x&S^{-1}\end{pmatrix},
\]
and check this for the data of Example~\ref{r1-exm-data} against Example~\ref{r2-exm-xtx-inverse}.
\end{enumerate}
\end{problem}
\begin{solution}
(a)
\begin{itemize}
\item By Theorem~\ref{r5b-thm-schur}, \(LM=U\) with
\(L=\left(\begin{smallmatrix}I_p&0\\-CA^{-1}&I_q\end{smallmatrix}\right)\) and
\(U=\left(\begin{smallmatrix}A&B\\0&E\end{smallmatrix}\right)\).
\item \(L\) is invertible with inverse \(\left(\begin{smallmatrix}I_p&0\\CA^{-1}&I_q\end{smallmatrix}\right)\)
(Proposition~\ref{r5b-prop-block-elementary}), so \(M=L^{-1}U\); and \(U\) is invertible with
\(U^{-1}=\left(\begin{smallmatrix}A^{-1}&-A^{-1}BE^{-1}\\0&E^{-1}\end{smallmatrix}\right)\)
(Proposition~\ref{r5b-prop-block-tri}(c)).
\item Hence \(M\) is invertible and \(M^{-1}=U^{-1}(L^{-1})^{-1}=U^{-1}L\)
(Theorem~\ref{r2-thm-inverse-props}(a),(b)). By the block product formula
(Theorem~\ref{r5b-thm-block}(c)),
\[
U^{-1}L=\begin{pmatrix}A^{-1}&-A^{-1}BE^{-1}\\0&E^{-1}\end{pmatrix}
\begin{pmatrix}I_p&0\\-CA^{-1}&I_q\end{pmatrix}
=\begin{pmatrix}A^{-1}+A^{-1}BE^{-1}CA^{-1}&-A^{-1}BE^{-1}\\-E^{-1}CA^{-1}&E^{-1}\end{pmatrix}
\]
\end{itemize}

(b) Look for \(X=\left(\begin{smallmatrix}P&Q\\R&T\end{smallmatrix}\right)\) with
\(MX=I_{p+q}\); by the block product formula this means
\[
AP+BR=I_p,\qquad CP+DR=0,\qquad AQ+BT=0,\qquad CQ+DT=I_q
\]
\begin{itemize}
\item The second equation gives \(R=-D^{-1}CP\); then the first becomes
\((A-BD^{-1}C)P=FP=I_p\), so \(P=F^{-1}\) and \(R=-D^{-1}CF^{-1}\).
\item The fourth equation gives \(T=D^{-1}-D^{-1}CQ\); then the third becomes
\(AQ+BD^{-1}-BD^{-1}CQ=FQ+BD^{-1}=0\), so \(Q=-F^{-1}BD^{-1}\) and
\(T=D^{-1}+D^{-1}CF^{-1}BD^{-1}\).
\item Conversely, these blocks satisfy the four equations:
\(AP+BR=(A-BD^{-1}C)F^{-1}=I_p\), \(CP+DR=CF^{-1}-CF^{-1}=0\),
\(AQ+BT=-(A-BD^{-1}C)F^{-1}BD^{-1}+BD^{-1}=0\) and
\(CQ+DT=-CF^{-1}BD^{-1}+I_q+CF^{-1}BD^{-1}=I_q\).
\item So \(MX=I_{p+q}\), and since \(M\) is square, \(M\) is invertible with \(M^{-1}=X\)
(Theorem~\ref{r2-thm-invertible}).
\end{itemize}

(c)
\begin{itemize}
\item By Example~\ref{r5b-exm-schur}, \(X'X\) has the blocks \(A=\one'\one=(n)\), \(B=\one'X_0\),
\(C=X_0'\one\), \(D=X_0'X_0\), and the Schur complement of \(A\) is \(X_0'C_nX_0=S\). As \(A\) and
\(X'X\) are invertible, so is \(S\) (Theorem~\ref{r5d-thm-schur-det}).
\item Here \(A^{-1}B=\frac1n\one'X_0=\bar x'\) and \(CA^{-1}=\frac1nX_0'\one=\bar x\), so (a) gives the
displayed formula: \(A^{-1}+A^{-1}BS^{-1}CA^{-1}=\frac1n+\bar x'S^{-1}\bar x\),
\(-A^{-1}BS^{-1}=-\bar x'S^{-1}\) and \(-S^{-1}CA^{-1}=-S^{-1}\bar x\). For \(k=1\) and \(X_0=x\),
\(S=x'C_nx=\sum_i(x_i-\bar x)^2\) (Example~\ref{r1-exm-centering}), so the last diagonal entry of
\((X'X)^{-1}\) is \(1/\sum_i(x_i-\bar x)^2\).
\item In Example~\ref{r1-exm-data}, \(n=4\), \(\bar x=(1,2)'\) and \(S=4I_2\), so
\(S^{-1}=\frac14I_2\), \(\frac14+\bar x'S^{-1}\bar x=\frac14+\frac{1+4}4=\frac32\) and
\(-\bar x'S^{-1}=(-\frac14,\ -\frac12)\): these are the entries of
\((X'X)^{-1}=\frac14\left(\begin{smallmatrix}6&-1&-2\\-1&1&0\\-2&0&1\end{smallmatrix}\right)\) found in
Example~\ref{r2-exm-xtx-inverse}.
\end{itemize}
\end{solution}

\begin{problem}\label{r5b-ex-partitioned-inverse-drill}
Let
\[
M=\left(\begin{array}{cc|cc}2&1&1&1\\1&1&1&0\\\hline1&1&2&-1\\2&1&1&2\end{array}\right)
=\begin{pmatrix}A&B\\C&D\end{pmatrix}
\]
Find the Schur complement \(E\) of \(A\) and \(\det M\), and use
Exercise~\ref{r5b-ex-partitioned-inverse}(a) to find \(M^{-1}\).
\end{problem}
\begin{solution}
Here \(A=\left(\begin{smallmatrix}2&1\\1&1\end{smallmatrix}\right)\),
\(B=\left(\begin{smallmatrix}1&1\\1&0\end{smallmatrix}\right)\),
\(C=\left(\begin{smallmatrix}1&1\\2&1\end{smallmatrix}\right)\),
\(D=\left(\begin{smallmatrix}2&-1\\1&2\end{smallmatrix}\right)\); \(\det A=1\), and
\(A^{-1}=\left(\begin{smallmatrix}1&-1\\-1&2\end{smallmatrix}\right)\) by the \(2\times2\) formula of
Exercise~\ref{r2-ex-inverses-b}. Then
\[
A^{-1}B=\begin{pmatrix}1&-1\\-1&2\end{pmatrix}\begin{pmatrix}1&1\\1&0\end{pmatrix}=\begin{pmatrix}0&1\\1&-1\end{pmatrix},\qquad
CA^{-1}=\begin{pmatrix}1&1\\2&1\end{pmatrix}\begin{pmatrix}1&-1\\-1&2\end{pmatrix}=\begin{pmatrix}0&1\\1&0\end{pmatrix},
\]
\[
E=D-CA^{-1}B=\begin{pmatrix}2&-1\\1&2\end{pmatrix}-\begin{pmatrix}0&1\\1&0\end{pmatrix}\begin{pmatrix}1&1\\1&0\end{pmatrix}
=\begin{pmatrix}2&-1\\1&2\end{pmatrix}-\begin{pmatrix}1&0\\1&1\end{pmatrix}=\begin{pmatrix}1&-1\\0&1\end{pmatrix}
\]
So \(\det M=\det A\cdot\det E=1\cdot1=1\) (Theorem~\ref{r5d-thm-schur-det}(a)), and
\(E^{-1}=\left(\begin{smallmatrix}1&1\\0&1\end{smallmatrix}\right)\). Next
\[
A^{-1}BE^{-1}=\begin{pmatrix}0&1\\1&-1\end{pmatrix}\begin{pmatrix}1&1\\0&1\end{pmatrix}=\begin{pmatrix}0&1\\1&0\end{pmatrix},\qquad
E^{-1}CA^{-1}=\begin{pmatrix}1&1\\0&1\end{pmatrix}\begin{pmatrix}0&1\\1&0\end{pmatrix}=\begin{pmatrix}1&1\\1&0\end{pmatrix},
\]
and \(A^{-1}BE^{-1}CA^{-1}=\left(\begin{smallmatrix}0&1\\1&0\end{smallmatrix}\right)\left(\begin{smallmatrix}0&1\\1&0\end{smallmatrix}\right)=I_2\).
By Exercise~\ref{r5b-ex-partitioned-inverse}(a),
\[
M^{-1}=\begin{pmatrix}A^{-1}+I_2&-A^{-1}BE^{-1}\\-E^{-1}CA^{-1}&E^{-1}\end{pmatrix}
=\left(\begin{array}{cc|cc}2&-1&0&-1\\-1&3&-1&0\\\hline-1&-1&1&1\\-1&0&0&1\end{array}\right)
\]
Check: \(\row_1(M)M^{-1}=2(2,-1,0,-1)+(-1,3,-1,0)+(-1,-1,1,1)+(-1,0,0,1)=(1,0,0,0)\) and
\(\row_4(M)M^{-1}=2(2,-1,0,-1)+(-1,3,-1,0)+(-1,-1,1,1)+2(-1,0,0,1)=(0,0,0,1)\); rows \(2\) and \(3\)
likewise give \((0,1,0,0)\) and \((0,0,1,0)\).
\end{solution}

\begin{problem}\label{r5d-ex-three-by-three}
Show that
\[
\det\begin{pmatrix}1&5&-5\\3&2&-5\\6&-2&-5\end{pmatrix}=-5\qquad\text{and}\qquad
\det\begin{pmatrix}2&6&5\\-2&-3&-5\\2&3&9\end{pmatrix}=24,
\]
in each case (i) by expansion along the first row and (ii) by row and column operations.
\end{problem}
\begin{solution}
Call the two matrices \(A\) and \(B\).

(\(A\))
\begin{itemize}
\item (i)
\(\det A=1\det\left(\begin{smallmatrix}2&-5\\-2&-5\end{smallmatrix}\right)
-5\det\left(\begin{smallmatrix}3&-5\\6&-5\end{smallmatrix}\right)
-5\det\left(\begin{smallmatrix}3&2\\6&-2\end{smallmatrix}\right)
=(-10-10)-5(-15+30)-5(-6-12)=-20-75+90=-5\).
\item (ii) Column \(3\) is \(-5\) times \((1,1,1)'\). Taking the factor \(-5\) out of column \(3\)
(Theorem~\ref{r5d-thm-rowops}(a),(f)) and then subtracting row \(1\) from rows \(2\) and \(3\),
\[
\det A=-5\det\begin{pmatrix}1&5&1\\3&2&1\\6&-2&1\end{pmatrix},\qquad
\begin{pmatrix}1&5&1\\3&2&1\\6&-2&1\end{pmatrix}
\xrightarrow[R_3\leftarrow R_3-R_1]{R_2\leftarrow R_2-R_1}
\begin{pmatrix}1&5&1\\2&-3&0\\5&-7&0\end{pmatrix}
\]
The last matrix has the single nonzero entry \(1\) in column \(3\), at position \((1,3)\); expanding
along column \(3\), its determinant is \((-1)^{1+3}\det\left(\begin{smallmatrix}2&-3\\5&-7\end{smallmatrix}\right)=-14+15=1\).
Hence \(\det A=-5\cdot1=-5\).
\end{itemize}

(\(B\))
\begin{itemize}
\item (i)
\(\det B=2\det\left(\begin{smallmatrix}-3&-5\\3&9\end{smallmatrix}\right)
-6\det\left(\begin{smallmatrix}-2&-5\\2&9\end{smallmatrix}\right)
+5\det\left(\begin{smallmatrix}-2&-3\\2&3\end{smallmatrix}\right)
=2(-27+15)-6(-18+10)+5(-6+6)=-24+48+0=24\).
\item (ii)
\[
B\xrightarrow[R_3\leftarrow R_3-R_1]{R_2\leftarrow R_2+R_1}
\begin{pmatrix}2&6&5\\0&3&0\\0&-3&4\end{pmatrix}
\xrightarrow{R_3\leftarrow R_3+R_2}
\begin{pmatrix}2&6&5\\0&3&0\\0&0&4\end{pmatrix}
\]
No step changes the determinant, so \(\det B=2\cdot3\cdot4=24\).
\end{itemize}
\end{solution}

\begin{problem}\label{r5d-ex-four-by-four}
Compute the determinant of
\[
A=\begin{pmatrix}0&4&0&5\\1&0&-1&2\\0&3&0&-2\\0&0&-6&0\end{pmatrix}
\]
(a) by repeated cofactor expansion: expand \(\det A\) in terms of determinants of \(3\times3\)
matrices, these in terms of determinants of \(2\times2\) matrices, and finally these in terms of
determinants of \(1\times1\) matrices; (b) by row reduction to triangular form.
\end{problem}
\begin{solution}
(a)
\begin{itemize}
\item Column \(1\) has the single nonzero entry \(a_{21}=1\); expanding along column \(1\),
\(\det A=1\cdot(-1)^{2+1}\det\left(\begin{smallmatrix}4&0&5\\3&0&-2\\0&-6&0\end{smallmatrix}\right)\).
\item Row \(3\) of this \(3\times3\) matrix has the single nonzero entry \(-6\), in column \(2\):
\(\det\left(\begin{smallmatrix}4&0&5\\3&0&-2\\0&-6&0\end{smallmatrix}\right)
=(-6)(-1)^{3+2}\det\left(\begin{smallmatrix}4&5\\3&-2\end{smallmatrix}\right)\).
\item Expanding the \(2\times2\) determinant along its first row,
\(4(-1)^{1+1}\det(-2)+5(-1)^{1+2}\det(3)=-8-15=-23\).
\item Hence \(\det A=(-1)\cdot6\cdot(-23)=138\).
\end{itemize}
(b)
\[
\begin{gathered}
A\xrightarrow{R_1\leftrightarrow R_2}
\begin{pmatrix}1&0&-1&2\\0&4&0&5\\0&3&0&-2\\0&0&-6&0\end{pmatrix}
\xrightarrow{R_2\leftarrow R_2-R_3}
\begin{pmatrix}1&0&-1&2\\0&1&0&7\\0&3&0&-2\\0&0&-6&0\end{pmatrix}
\\
\xrightarrow{R_3\leftarrow R_3-3R_2}
\begin{pmatrix}1&0&-1&2\\0&1&0&7\\0&0&0&-23\\0&0&-6&0\end{pmatrix}
\xrightarrow{R_3\leftrightarrow R_4}
\begin{pmatrix}1&0&-1&2\\0&1&0&7\\0&0&-6&0\\0&0&0&-23\end{pmatrix}
\end{gathered}
\]
There are two interchanges, so \(\det A=(-1)^2\cdot1\cdot1\cdot(-6)\cdot(-23)=138\).
\end{solution}

\begin{problem}\label{r5d-ex-alpha}
Let \(\alpha\in\R\) and \(A=\begin{pmatrix}1&\alpha&0\\4&5&3\\1&0&2\end{pmatrix}\). For which \(\alpha\) is
\(\det A=0\)? Show that the columns of \(A\) are linearly dependent in that case.
\end{problem}
\begin{solution}
Expanding along the first row,
\(\det A=1\det\left(\begin{smallmatrix}5&3\\0&2\end{smallmatrix}\right)
-\alpha\det\left(\begin{smallmatrix}4&3\\1&2\end{smallmatrix}\right)=10-5\alpha\), so \(\det A=0\) exactly
when \(\alpha=2\) (and \(A\) is invertible exactly when \(\alpha\ne2\),
Theorem~\ref{r5d-thm-product}(a)). For \(\alpha=2\) and \(x=(2,-1,-1)'\),
\(Ax=(2-2+0,\ 8-5-3,\ 2-0-2)'=0\); that is, \(2a_1-a_2-a_3=0\), a nontrivial relation among the columns.
\end{solution}

\begin{problem}\label{r5d-ex-k-power}
Let \(K\in\R^{12\times12}\) satisfy \(K^5=3K\). What is the numerical value of \(\det K\)? Is there more
than one possibility? If so, give all possible values.
\end{problem}
\begin{solution}
\begin{itemize}
\item By Theorem~\ref{r5d-thm-product}(c), \(\det K^5=(\det K)^5\) and \(\det(3K)=3^{12}\det K\). So
\((\det K)^5=3^{12}\det K\), that is, \(\det K\cdot\bigl((\det K)^4-3^{12}\bigr)=0\).
\item Hence \(\det K=0\) or \((\det K)^4=3^{12}\), that is,
\(\bigl((\det K)^2-3^6\bigr)\bigl((\det K)^2+3^6\bigr)=0\). The second factor is positive, so
\((\det K)^2=3^6\) and \(\det K=\pm3^3=\pm27\).
\item All three values occur. Let \(t\) be the positive number with \(t^4=3\), so \(t^5=3t\). For
\(K=0\), \(\det K=0\). For \(K=tI_{12}\), \(K^5=t^5I_{12}=3K\) and \(\det K=t^{12}=3^3=27\). For
\(K=\diag(-t,t,\dots,t)\), \(K^5=\diag(-t^5,t^5,\dots,t^5)=3K\) and \(\det K=-t^{12}=-27\).
\end{itemize}
So \(\det K\) is \(0\), \(27\) or \(-27\).
\end{solution}

\begin{problem}\label{r5d-ex-integer}
Let \(A\in\R^{n\times n}\) be invertible. Show that if all entries of \(A\) and of \(A^{-1}\) are
integers, then \(\det A=\pm1\).
\end{problem}
\begin{solution}
\begin{itemize}
\item A matrix with integer entries has an integer determinant: for \(n=1\) it is the entry, and for
\(n\geqslant2\) Definition~\ref{r5d-def-det} gives \(\det A\) as a sum of products of entries and of
determinants of \((n-1)\times(n-1)\) matrices with integer entries; use induction on \(n\). So
\(d=\det A\) and \(\det A^{-1}\) are integers.
\item By Theorem~\ref{r5d-thm-product}(a),(c), \(d\ne0\) and \(\det A^{-1}=1/d\). So \(d\) and \(1/d\)
are nonzero integers.
\item A nonzero integer has absolute value at least \(1\): \(\lvert d\rvert\geqslant1\) and
\(\lvert1/d\rvert\geqslant1\), that is, \(\lvert d\rvert\leqslant1\). Hence \(\lvert d\rvert=1\), and
\(\det A=\pm1\).\qedhere
\end{itemize}
\end{solution}

\begin{problem}\label{r5d-ex-cofactors}
Let \(A=\begin{pmatrix}1&0&-1\\-1&3&-3\\0&-4&5\end{pmatrix}\).
\begin{enumerate}
\item Compute the cofactor \(c_{ij}\) of each entry of \(A\).
\item Compute \(\det A\) by expansion along the second row, and check the answer by expansion along the
second column.
\item Use Theorem~\ref{r5d-thm-adjugate} to compute \(A^{-1}\).
\end{enumerate}
\end{problem}
\begin{solution}
(a)
\[
\begin{aligned}
c_{11}&=\det\begin{pmatrix}3&-3\\-4&5\end{pmatrix}=3,&
c_{12}&=-\det\begin{pmatrix}-1&-3\\0&5\end{pmatrix}=5,&
c_{13}&=\det\begin{pmatrix}-1&3\\0&-4\end{pmatrix}=4,\\
c_{21}&=-\det\begin{pmatrix}0&-1\\-4&5\end{pmatrix}=4,&
c_{22}&=\det\begin{pmatrix}1&-1\\0&5\end{pmatrix}=5,&
c_{23}&=-\det\begin{pmatrix}1&0\\0&-4\end{pmatrix}=4,\\
c_{31}&=\det\begin{pmatrix}0&-1\\3&-3\end{pmatrix}=3,&
c_{32}&=-\det\begin{pmatrix}1&-1\\-1&-3\end{pmatrix}=4,&
c_{33}&=\det\begin{pmatrix}1&0\\-1&3\end{pmatrix}=3
\end{aligned}
\]

(b) Along row \(2\): \(\det A=(-1)\cdot4+3\cdot5+(-3)\cdot4=-1\). Along column \(2\):
\(\det A=0\cdot5+3\cdot5+(-4)\cdot4=-1\) (Theorem~\ref{r5d-thm-expansion}).

(c) Since \(\det A=-1\ne0\),
\[
A^{-1}=\frac1{-1}\operatorname{adj}A=-\begin{pmatrix}3&4&3\\5&5&4\\4&4&3\end{pmatrix}
=\begin{pmatrix}-3&-4&-3\\-5&-5&-4\\-4&-4&-3\end{pmatrix}
\]
Check: \(\row_1(A)A^{-1}=(1,0,-1)A^{-1}=(-3+4,\ -4+4,\ -3+3)=(1,0,0)\); likewise
\(\row_2(A)A^{-1}=(3-15+12,\ 4-15+12,\ 3-12+9)=(0,1,0)\) and
\(\row_3(A)A^{-1}=(20-20,\ 20-20,\ 16-15)=(0,0,1)\).
\end{solution}

\begin{problem}\label{r5d-ex-cramer}
\begin{enumerate}
\item Let \(A\in\R^{n\times n}\) (\(n\geqslant2\)) with \(\det A\ne0\), and let \(b\in\R^n\). Show that
the solution \(x\) of \(Ax=b\) has the entries
\[
x_j=\frac{\det A_j}{\det A}\qquad(1\leqslant j\leqslant n),
\]
where \(A_j\) is \(A\) with column \(j\) replaced by \(b\) (\textbf{Cramer's rule}).
\item Use (a) to solve the first system of Exercise~\ref{r2-ex-two-systems},
\(\left(\begin{smallmatrix}1&-3&-2\\-1&2&1\\-1&2&2\end{smallmatrix}\right)x=(1,-1,-2)'\).
\end{enumerate}
\end{problem}
\begin{solution}
(a)
\begin{itemize}
\item By Theorem~\ref{r5d-thm-adjugate}, \(A\) is invertible, so \(x=A^{-1}b=\frac1{\det A}(\operatorname{adj}A)b\)
is the only solution.
\item Entry \(j\) of \((\operatorname{adj}A)b\) is \(\sum_{i=1}^n(\operatorname{adj}A)_{ji}b_i=\sum_{i=1}^nc_{ij}b_i\).
\item \(A_j\) differs from \(A\) only in column \(j\), so \(A_j(i|j)=A(i|j)\), and the \((i,j)\) entry
of \(A_j\) is \(b_i\). Expanding \(\det A_j\) along column \(j\) (Theorem~\ref{r5d-thm-expansion}) gives
\(\sum_{i=1}^nb_i(-1)^{i+j}\det A(i|j)=\sum_{i=1}^nc_{ij}b_i\). Hence \(x_j=\det A_j/\det A\).
\end{itemize}
(b) Expanding along the first row,
\[
\begin{aligned}
\det A&=1(4-2)+3(-2+1)-2(-2+2)=2-3+0=-1,\\
\det A_1&=\det\begin{pmatrix}1&-3&-2\\-1&2&1\\-2&2&2\end{pmatrix}=1(4-2)+3(-2+2)-2(-2+4)=2+0-4=-2,\\
\det A_2&=\det\begin{pmatrix}1&1&-2\\-1&-1&1\\-1&-2&2\end{pmatrix}=1(-2+2)-1(-2+1)-2(2-1)=0+1-2=-1,\\
\det A_3&=\det\begin{pmatrix}1&-3&1\\-1&2&-1\\-1&2&-2\end{pmatrix}=1(-4+2)+3(2-1)+1(-2+2)=-2+3+0=1
\end{aligned}
\]
So \(x=\bigl(\frac{-2}{-1},\ \frac{-1}{-1},\ \frac1{-1}\bigr)'=(2,1,-1)'\), the solution found by
elimination in Exercise~\ref{r2-ex-two-systems}.
\end{solution}

\begin{problem}\label{r5d-ex-sylvester}
Let \(A\in\R^{m\times n}\) and \(B\in\R^{n\times m}\). Show that \(\det(I_m-AB)=\det(I_n-BA)\).
\end{problem}
\begin{solution}
By the block product formula (Theorem~\ref{r5b-thm-block}),
\[
\begin{gathered}
\begin{pmatrix}I_m-AB&A\\0&I_n\end{pmatrix}\begin{pmatrix}I_m&0\\B&I_n\end{pmatrix}
=\begin{pmatrix}I_m-AB+AB&A\\B&I_n\end{pmatrix}=\begin{pmatrix}I_m&A\\B&I_n\end{pmatrix},\\
\begin{pmatrix}I_m&0\\B&I_n\end{pmatrix}\begin{pmatrix}I_m&A\\0&I_n-BA\end{pmatrix}
=\begin{pmatrix}I_m&A\\B&BA+I_n-BA\end{pmatrix}=\begin{pmatrix}I_m&A\\B&I_n\end{pmatrix}
\end{gathered}
\]
So the two products are equal. By Theorems~\ref{r5d-thm-product}(b) and~\ref{r5d-thm-block-tri}, the
first has determinant \(\det(I_m-AB)\det I_n\cdot\det I_m\det I_n=\det(I_m-AB)\), and the second
\(\det I_m\det I_n\cdot\det I_m\det(I_n-BA)=\det(I_n-BA)\).
\end{solution}

\begin{problem}\label{r5d-ex-schur-numeric}
For \(t\in\R\) let
\[
M=\left(\begin{array}{cc|cc}2&1&1&0\\1&1&1&1\\\hline1&1&t&2\\0&1&2&t+1\end{array}\right)
=\begin{pmatrix}A&B\\C&D\end{pmatrix}
\]
(a) Find \(\det M\) with the Schur complement of \(A\). (b) For which \(t\) is \(M\) invertible?
(c) Find \(\rank M\) for these exceptional values of \(t\).
\end{problem}
\begin{solution}
(a) Here \(A=\left(\begin{smallmatrix}2&1\\1&1\end{smallmatrix}\right)\),
\(B=\left(\begin{smallmatrix}1&0\\1&1\end{smallmatrix}\right)\), \(C=B'\) and
\(D=\left(\begin{smallmatrix}t&2\\2&t+1\end{smallmatrix}\right)\). Here \(\det A=2\cdot1-1\cdot1=1\),
and \(A^{-1}=\left(\begin{smallmatrix}1&-1\\-1&2\end{smallmatrix}\right)\) by the \(2\times2\) formula of
Exercise~\ref{r2-ex-inverses-b}. Then
\[
CA^{-1}=\begin{pmatrix}1&1\\0&1\end{pmatrix}\begin{pmatrix}1&-1\\-1&2\end{pmatrix}=\begin{pmatrix}0&1\\-1&2\end{pmatrix},\qquad
CA^{-1}B=\begin{pmatrix}0&1\\-1&2\end{pmatrix}\begin{pmatrix}1&0\\1&1\end{pmatrix}=\begin{pmatrix}1&1\\1&2\end{pmatrix},
\]
so the Schur complement is
\(S=D-CA^{-1}B=\left(\begin{smallmatrix}t-1&1\\1&t-1\end{smallmatrix}\right)\). By
Theorem~\ref{r5d-thm-schur-det}(a),
\(\det M=\det A\cdot\det S=(t-1)^2-1=t(t-2)\).

(b) By Theorem~\ref{r5d-thm-product}(a), \(M\) is invertible if and only if \(t\ne0\) and \(t\ne2\).

(c) For \(t=0\), \(S=\left(\begin{smallmatrix}-1&1\\1&-1\end{smallmatrix}\right)\); for \(t=2\),
\(S=\left(\begin{smallmatrix}1&1\\1&1\end{smallmatrix}\right)\). In both cases the second row of \(S\) is
a multiple of the first and \(S\ne0\), so \(\rank S=1\). By Theorem~\ref{r5d-thm-schur-rank}(a),
\(\rank M=2+1=3\).
\end{solution}

\begin{problem}\label{r5d-ex-equicorrelation-inverse}
Let \(R=(1-\rho)I_n+\rho J_n\) be the equicorrelation matrix of Example~\ref{r5d-exm-equicorrelation},
with \(\rho\ne1\) and \(\rho\ne-1/(n-1)\). Show that
\[
R^{-1}=\frac1{1-\rho}\Bigl(I_n-\frac{\rho}{1+(n-1)\rho}J_n\Bigr)
\]
\end{problem}
\begin{solution}
Let \(\tau=\rho/(1+(n-1)\rho)\), so that \(\tau(1-\rho+n\rho)=\rho\). Using \(J_n^2=nJ_n\)
(Proposition~\ref{r1-prop-centering}),
\[
\begin{aligned}
\bigl((1-\rho)I_n+\rho J_n\bigr)(I_n-\tau J_n)&=(1-\rho)I_n+\bigl(\rho-(1-\rho)\tau-n\rho\tau\bigr)J_n\\
&=(1-\rho)I_n+\bigl(\rho-\tau(1-\rho+n\rho)\bigr)J_n=(1-\rho)I_n
\end{aligned}
\]
Dividing by \(1-\rho\ne0\) gives \(R\cdot\frac1{1-\rho}(I_n-\tau J_n)=I_n\), so
\(\frac1{1-\rho}(I_n-\tau J_n)=R^{-1}\) (Theorem~\ref{r2-thm-invertible}).
\end{solution}

\clearpage
\def\unitnumber{6}
\setcounter{theorem}{0}\setcounter{problem}{0}
\section*{Chapter 6. Orthogonality, projections, and least squares}

\subsection*{Inner products, lengths, and orthogonality}

\begin{definition}\label{r6-def-inner}
For \(x,y\in\R^n\), the \textbf{inner product} of \(x\) and \(y\) is the number
\(x'y=x_1y_1+\dots+x_ny_n\); it equals \(y'x\) (Theorem~\ref{r1-thm-transpose}). The \textbf{length}
(or norm) of \(x\) is \(\lVert x\rVert=\sqrt{x'x}=(x_1^2+\dots+x_n^2)^{1/2}\), and \(\lVert x-y\rVert\) is the
\textbf{distance} between \(x\) and \(y\). The vectors \(x\) and \(y\) are \textbf{orthogonal}, written
\(x\perp y\), if \(x'y=0\).
\end{definition}
Since \(x'x\) is a sum of squares, \(\lVert x\rVert\geqslant0\), and \(\lVert x\rVert=0\) only for \(x=0\); also
\(\lVert cx\rVert=\lvert c\rvert\,\lVert x\rVert\) for \(c\in\R\). A vector of length \(1\) is a \textbf{unit vector};
for \(x\ne0\), \(x/\lVert x\rVert\) is one. The zero vector is orthogonal to every vector, and it is the only
vector orthogonal to itself.

\begin{proposition}\label{r6-prop-pythagoras}
For \(x,y\in\R^n\), \(\lVert x+y\rVert^2=\lVert x\rVert^2+2x'y+\lVert y\rVert^2\). In particular, if \(x\perp y\), then
\(\lVert x+y\rVert^2=\lVert x\rVert^2+\lVert y\rVert^2\) (\textbf{Pythagoras}).
\end{proposition}
\begin{proof}
\((x+y)'(x+y)=x'x+x'y+y'x+y'y\), and \(y'x=x'y\).
\end{proof}

\begin{theorem}\label{r6-thm-cauchy-schwarz}
Let \(x,y\in\R^n\).
\begin{enumerate}
\item (\textbf{Cauchy--Schwarz inequality}) \(\lvert x'y\rvert\leqslant\lVert x\rVert\,\lVert y\rVert\), with equality if and
only if one of \(x,y\) is a multiple of the other.
\item (\textbf{Triangle inequality}) \(\lVert x+y\rVert\leqslant\lVert x\rVert+\lVert y\rVert\).
\end{enumerate}
\end{theorem}
\begin{proof}
(a) If \(x=0\), both sides are \(0\), and \(x=0y\). Let \(x\ne0\).
\begin{itemize}
\item Put \(c=x'y/\lVert x\rVert^2\) and \(w=y-cx\). Then \(x'w=x'y-c\,x'x=0\), so \(cx\perp w\).
\item Pythagoras for \(y=cx+w\) gives
\(\lVert y\rVert^2=c^2\lVert x\rVert^2+\lVert w\rVert^2\geqslant c^2\lVert x\rVert^2=(x'y)^2/\lVert x\rVert^2\). Multiply by
\(\lVert x\rVert^2\) and take square roots.
\item Equality holds if and only if \(w=0\), that is, \(y=cx\). Conversely, if \(y=ax\), then \(c=a\) and
\(w=0\); if \(x=ay\) with \(x\ne0\), then \(a\ne0\) and \(y=a^{-1}x\).
\end{itemize}
(b) By Proposition~\ref{r6-prop-pythagoras} and (a),
\(\lVert x+y\rVert^2\leqslant\lVert x\rVert^2+2\lVert x\rVert\,\lVert y\rVert+\lVert y\rVert^2=(\lVert x\rVert+\lVert y\rVert)^2\), and both
\(\lVert x+y\rVert\) and \(\lVert x\rVert+\lVert y\rVert\) are \(\geqslant0\).
\end{proof}

\begin{definition}\label{r6-def-complement}
Let \(W\) be a subspace of \(\R^n\). A vector \(y\in\R^n\) is \textbf{orthogonal to} \(W\), written
\(y\perp W\), if \(y'w=0\) for every \(w\in W\). The \textbf{orthogonal complement} of \(W\) is
\(W^\perp=\{y\in\R^n:y\perp W\}\).
\end{definition}

\begin{theorem}\label{r6-thm-complement}
Let \(W\) be a subspace of \(\R^n\).
\begin{enumerate}
\item If \(w_1,\dots,w_k\) span \(W\) and \(A=[w_1\ \cdots\ w_k]\in\R^{n\times k}\), then \(y\perp W\) if and
only if \(y'w_1=\dots=y'w_k=0\), that is, \(A'y=0\). So \(W^\perp=\mathcal N(A')\), a subspace of \(\R^n\).
\item \(W\cap W^\perp=\{0\}\).
\item \(\dim W+\dim W^\perp=n\) and \(\R^n=W\oplus W^\perp\): every \(y\in\R^n\) is \(y=w+v\) with \(w\in W\)
and \(v\in W^\perp\) in exactly one way.
\end{enumerate}
\end{theorem}
\begin{proof}
(a) If \(y\perp W\), then \(y'w_j=0\), as \(w_j\in W\). Conversely, let all \(y'w_j=0\). Every \(w\in W\)
is \(w=Ac\) for some \(c\in\R^k\) (Definition~\ref{r3-def-span} and Theorem~\ref{r1-thm-views}(b)), and
\(y'w=(A'y)'c=0\). The \(j\)th entry of \(A'y\) is \(w_j'y\) (Theorem~\ref{r1-thm-views}(d)).

(b) If \(y\in W\cap W^\perp\), then \(y\perp y\), so \(y=0\).

(c) Let \(d=\dim W\). If \(d=0\), then \(W=\{0\}\) and \(W^\perp=\R^n\). Let \(d\geqslant1\).
\begin{itemize}
\item Let \(A\in\R^{n\times d}\) have a basis of \(W\) as its columns, so \(\rank A=d\). By (a) and
Theorem~\ref{thm:r4-rank-nullity}, \(\dim W^\perp=\dim\mathcal N(A')=n-\rank A=n-d\).
\item By (b), the sum \(W+W^\perp\) is direct (Definition~\ref{r3-def-sum}), and a basis of \(W\)
together with a basis of \(W^\perp\), \(d+(n-d)=n\) vectors, is a basis of \(W\oplus W^\perp\)
(Theorem~\ref{r3-thm-direct-sum-basis}). So \(W\oplus W^\perp\) is a subspace of \(\R^n\)
(Theorem~\ref{r3-thm-sum}(a)) of dimension \(n\), hence equal to \(\R^n\)
(Theorem~\ref{thm:r4-counting}(a)). The decomposition \(y=w+v\) is unique by
Theorem~\ref{r3-thm-direct-sum}.\qedhere
\end{itemize}
\end{proof}

\begin{corollary}\label{r6-cor-four-subspaces}
For \(A\in\R^{m\times n}\),
\[
\mathcal N(A)=\mathcal C(A')^\perp,\quad \R^n=\mathcal C(A')\oplus\mathcal N(A);\qquad
\mathcal N(A')=\mathcal C(A)^\perp,\quad \R^m=\mathcal C(A)\oplus\mathcal N(A')
\]
So \(Ax=0\) if and only if \(x\) is orthogonal to every (transposed) row of \(A\), and the four
subspaces of \(A\) (Definition~\ref{def:r4-spaces}) are orthogonal in pairs: the transposed row space \(\mathcal C(A')\) and
the null space in \(\R^n\), the column space and the left null space in \(\R^m\).
\end{corollary}
\begin{proof}
The columns of \(A'\) span \(\mathcal C(A')\), so Theorem~\ref{r6-thm-complement}(a) gives
\(\mathcal C(A')^\perp=\mathcal N((A')')=\mathcal N(A)\); applied to the columns of \(A\), it gives
\(\mathcal C(A)^\perp=\mathcal N(A')\). The direct sums follow from
Theorem~\ref{r6-thm-complement}(c).
\end{proof}

\begin{example}\label{r6-exm-four-subspaces}
For \(A=\left(\begin{smallmatrix}1&2&1&3\\1&3&0&5\\2&5&1&8\end{smallmatrix}\right)\) of
Example~\ref{exm:r4-four-bases}, the transposed nonzero rows \(\rho_1'=(1,0,3,-1)'\),
\(\rho_2'=(0,1,-1,2)'\) of its RREF form a basis of \(\mathcal C(A')\), and \(s_1=(-3,1,1,0)'\),
\(s_2=(1,-2,0,1)'\) a basis of \(\mathcal N(A)\). Indeed
\[
\rho_1s_1=-3+0+3+0=0,\quad \rho_1s_2=1+0+0-1=0,\quad \rho_2s_1=0+1-1+0=0,\quad \rho_2s_2=0-2+0+2=0,
\]
and \(2+2=4\). In \(\R^3\), \(y=(-1,-1,1)'\) is a basis of \(\mathcal N(A')\), orthogonal to the basis
\(a_1=(1,1,2)'\), \(a_2=(2,3,5)'\) of \(\mathcal C(A)\): \(-1-1+2=0\) and \(-2-3+5=0\); and \(2+1=3\).
\end{example}

\subsection*{Orthonormal bases, Gram--Schmidt, and QR}

\begin{proposition}\label{r6-prop-orthogonal-list}
Let \(u_1,\dots,u_k\in\R^n\) be nonzero and pairwise orthogonal (\(u_i'u_j=0\) for \(i\ne j\)).
\begin{enumerate}
\item If \(y=c_1u_1+\dots+c_ku_k\), then \(c_j=u_j'y/\lVert u_j\rVert^2\) for \(j=1,\dots,k\).
\item \(u_1,\dots,u_k\) are linearly independent.
\end{enumerate}
\end{proposition}
\begin{proof}
(a) \(u_j'y=\sum_ic_iu_j'u_i=c_j\lVert u_j\rVert^2\), and \(\lVert u_j\rVert^2>0\). (b) If \(\sum_ic_iu_i=0\), then (a)
with \(y=0\) gives every \(c_j=0\).
\end{proof}

\begin{definition}\label{r6-def-orthonormal}
Vectors \(q_1,\dots,q_p\in\R^n\) are \textbf{orthonormal} if \(q_i'q_j=0\) for \(i\ne j\) and \(\lVert q_j\rVert=1\)
for every \(j\). A matrix \(Q=[q_1\ \cdots\ q_p]\in\R^{n\times p}\) has \textbf{orthonormal columns} if
\(Q'Q=I_p\), since \((Q'Q)_{ij}=q_i'q_j\) (Theorem~\ref{r1-thm-views}(d)). A square matrix \(Q\in\R^{n\times n}\) with
\(Q'Q=I_n\) is an \textbf{orthogonal matrix}. A basis of a subspace consisting of orthonormal vectors is
an \textbf{orthonormal basis}.
\end{definition}

\begin{proposition}\label{r6-prop-orthonormal}
Let \(Q\in\R^{n\times p}\) satisfy \(Q'Q=I_p\).
\begin{enumerate}
\item \((Qx)'(Qy)=x'y\) and \(\lVert Qx\rVert=\lVert x\rVert\) for all \(x,y\in\R^p\): \(Q\) preserves inner products and
lengths.
\item If \(p=n\), then \(Q\) is invertible, \(Q^{-1}=Q'\) and \(QQ'=I_n\): the rows of \(Q\) are orthonormal
too.
\end{enumerate}
\end{proposition}
\begin{proof}
(a) \((Qx)'(Qy)=x'Q'Qy=x'y\); take \(y=x\). (b) \(Q'\) is a left inverse of the square matrix \(Q\), so
\(Q^{-1}=Q'\) (Theorem~\ref{r2-thm-invertible}), and \(QQ'=QQ^{-1}=I_n\).
\end{proof}

\begin{theorem}[Gram--Schmidt]\label{r6-thm-gram-schmidt}
Let \(a_1,\dots,a_p\in\R^n\) be linearly independent. Put \(b_1=a_1\) and, for \(j=2,\dots,p\),
\[
b_j=a_j-t_{1j}b_1-t_{2j}b_2-\dots-t_{j-1,j}b_{j-1},\qquad t_{ij}=\frac{a_j'b_i}{\lVert b_i\rVert^2}\quad(i<j)
\]
Then \(b_1,\dots,b_p\) are nonzero and pairwise orthogonal, and
\(\Span\{b_1,\dots,b_j\}=\Span\{a_1,\dots,a_j\}\) for every \(j\). Hence the vectors
\(q_j=b_j/\lVert b_j\rVert\) (\(1\leqslant j\leqslant p\)) form an orthonormal basis of
\(\Span\{a_1,\dots,a_p\}\).
\end{theorem}
\begin{proof}
Induction on \(j\): \(b_1=a_1\ne0\), since an independent list does not contain \(0\). Let
\(j\geqslant2\), and let \(b_1,\dots,b_{j-1}\) be nonzero, pairwise orthogonal, with span
\(W_{j-1}=\Span\{a_1,\dots,a_{j-1}\}\); so \(b_j\) is defined.
\begin{itemize}
\item Orthogonal: for \(l<j\), \(b_l'b_j=b_l'a_j-\sum_{i<j}t_{ij}b_l'b_i=b_l'a_j-t_{lj}\lVert b_l\rVert^2=0\).
\item Nonzero: otherwise \(a_j=\sum_{i<j}t_{ij}b_i\in W_{j-1}\), and \(a_1,\dots,a_j\) would be dependent
(Proposition~\ref{r3-prop-dependence}(c)) \(\contradiction\)
\item Span: \(b_j\) is a combination of \(a_j\) and \(b_1,\dots,b_{j-1}\in W_{j-1}\), and
\(a_j=b_j+\sum_{i<j}t_{ij}b_i\); so each of the lists \(a_1,\dots,a_j\) and \(b_1,\dots,b_j\) lies in the
span of the other, and the two spans are equal (Theorem~\ref{r3-thm-span}).
\end{itemize}
Finally, the \(q_j\) are pairwise orthogonal unit vectors, hence independent
(Proposition~\ref{r6-prop-orthogonal-list}), and they span \(\Span\{b_1,\dots,b_p\}=\Span\{a_1,\dots,a_p\}\).
\end{proof}

\begin{theorem}[QR decomposition]\label{r6-thm-qr}
Let \(A=[a_1\ \cdots\ a_p]\in\R^{n\times p}\) have full column rank \(p\). Then \(A=QR\), where
\(Q\in\R^{n\times p}\) has orthonormal columns (\(Q'Q=I_p\)) and \(R\in\R^{p\times p}\) is upper triangular with
positive diagonal entries; moreover \(R=Q'A\).
\end{theorem}
\begin{proof}
\begin{itemize}
\item The columns of \(A\) are independent (Theorem~\ref{thm:r4-rank-rules}(b)). Apply
Theorem~\ref{r6-thm-gram-schmidt} to them, and let \(Q=[q_1\ \cdots\ q_p]\).
\item The construction gives \(a_j=b_j+\sum_{i<j}t_{ij}b_i=\sum_{i\leqslant j}r_{ij}q_i\), where
\(r_{jj}=\lVert b_j\rVert>0\) and \(r_{ij}=t_{ij}\lVert b_i\rVert\) for \(i<j\). Let \(R=(r_{ij})\), with \(r_{ij}=0\) for
\(i>j\). Then \(\col_j(QR)=Q\col_j(R)=\sum_{i\leqslant j}r_{ij}q_i=a_j\) (Theorem~\ref{r1-thm-views}(a),(b)), so
\(A=QR\).
\item \(Q'Q=I_p\), as the \(q_j\) are orthonormal; hence \(Q'A=Q'QR=R\).\qedhere
\end{itemize}
\end{proof}

\begin{example}\label{r6-exm-qr}
Find the QR decomposition of
\(A=\begin{pmatrix}1&3&2\\2&1&3\\0&2&4\\2&2&5\end{pmatrix}\).

\textit{Solution.} Apply Gram--Schmidt (Theorem~\ref{r6-thm-gram-schmidt}) to the columns
\(a_1,a_2,a_3\).
\begin{itemize}
\item \(b_1=a_1=(1,2,0,2)'\), \(\lVert b_1\rVert^2=1+4+0+4=9\).
\item \(a_2'b_1=3+2+0+4=9\), so \(t_{12}=1\) and \(b_2=a_2-b_1=(3,1,2,2)'-(1,2,0,2)'=(2,-1,2,0)'\), with
\(\lVert b_2\rVert^2=4+1+4+0=9\).
\item \(a_3'b_1=2+6+0+10=18\) and \(a_3'b_2=4-3+8+0=9\), so \(t_{13}=2\), \(t_{23}=1\) and
\[
b_3=a_3-2b_1-b_2=\begin{pmatrix}2\\3\\4\\5\end{pmatrix}-\begin{pmatrix}2\\4\\0\\4\end{pmatrix}-\begin{pmatrix}2\\-1\\2\\0\end{pmatrix}
=\begin{pmatrix}-2\\0\\2\\1\end{pmatrix},\qquad\lVert b_3\rVert^2=4+0+4+1=9
\]
\item Check: \(b_1'b_2=2-2+0+0=0\), \(b_1'b_3=-2+0+0+2=0\), \(b_2'b_3=-4+0+4+0=0\).
\item All three lengths are \(3\), so \(q_j=\frac13b_j\), \(r_{jj}=3\), \(r_{12}=3t_{12}=3\), \(r_{13}=3t_{13}=6\)
and \(r_{23}=3t_{23}=3\):
\[
Q=\frac13\begin{pmatrix}1&2&-2\\2&-1&0\\0&2&2\\2&0&1\end{pmatrix},\qquad
R=\begin{pmatrix}3&3&6\\0&3&3\\0&0&3\end{pmatrix}
\]
\end{itemize}
As a check, \(R=Q'A\); for instance \(r_{23}=q_2'a_3=\frac13(4-3+8+0)=3\).
\end{example}

\subsection*{Orthogonal projections}

\begin{definition}\label{r6-def-projection}
Let \(W\) be a subspace of \(\R^n\) and \(y\in\R^n\). By Theorem~\ref{r6-thm-complement}(c),
\(y=\hat w+(y-\hat w)\) with \(\hat w\in W\) and \(y-\hat w\in W^\perp\) for exactly one vector \(\hat w\), the
\textbf{orthogonal projection} of \(y\) onto \(W\). Thus \(\hat w\) is the only vector of \(W\) with
\(y-\hat w\perp W\).
\end{definition}

\begin{theorem}\label{r6-thm-best-approximation}
Let \(\hat w\) be the orthogonal projection of \(y\in\R^n\) onto a subspace \(W\) of \(\R^n\). Then
\[
\lVert y-w\rVert^2=\lVert y-\hat w\rVert^2+\lVert \hat w-w\rVert^2\qquad(w\in W)
\]
Hence \(\hat w\) is the vector of \(W\) closest to \(y\): \(\lVert y-w\rVert>\lVert y-\hat w\rVert\) for every \(w\in W\)
with \(w\ne\hat w\).
\end{theorem}
\begin{proof}
\(y-w=(y-\hat w)+(\hat w-w)\), where \(\hat w-w\in W\) and \(y-\hat w\perp W\); apply Pythagoras
(Proposition~\ref{r6-prop-pythagoras}). If \(w\ne\hat w\), then \(\lVert \hat w-w\rVert^2>0\).
\end{proof}

\begin{theorem}\label{r6-thm-hat}
Let \(X\in\R^{n\times p}\) have full column rank \(p\). Then \(X'X\) is invertible. Put
\[
H=X(X'X)^{-1}X'\in\R^{n\times n}\ \ \text{(the \textbf{hat matrix})},\qquad
M=I_n-H\ \ \text{(the \textbf{residual maker})}
\]
\begin{enumerate}
\item For every \(y\in\R^n\), \(Hy\) is the orthogonal projection of \(y\) onto \(\mathcal C(X)\), and
\(y=Hy+My\) is the decomposition of \(y\) along \(\R^n=\mathcal C(X)\oplus\mathcal C(X)^\perp\).
\item \(H'=H=H^2\) and \(M'=M=M^2\): both are symmetric and idempotent.
\item \(HX=X\), \(MX=0\) and \(HM=MH=0\).
\end{enumerate}
\end{theorem}
\begin{proof}
\(\rank(X'X)=\rank X=p\) (Theorem~\ref{thm:r4-rank-ata}(b)), so the \(p\times p\) matrix \(X'X\) is
invertible (Theorem~\ref{thm:r4-rank-rules}(d)).

(a) \(Hy=X\bigl((X'X)^{-1}X'y\bigr)\in\mathcal C(X)\), and \(X'(y-Hy)=X'y-(X'X)(X'X)^{-1}X'y=0\), so
\(My=y-Hy\in\mathcal C(X)^\perp\) (Theorem~\ref{r6-thm-complement}(a)). By
Definition~\ref{r6-def-projection}, \(Hy\) is the projection.

(b) \(X'X\) is symmetric, hence so is its inverse:
\(((X'X)^{-1})'=((X'X)')^{-1}=(X'X)^{-1}\) (Theorem~\ref{r2-thm-inverse-props}(c)). So
\(H'=X((X'X)^{-1})'X'=H\) and \(H^2=X(X'X)^{-1}(X'X)(X'X)^{-1}X'=H\). Then \(M'=I-H'=M\) and
\(M^2=I-2H+H^2=M\).

(c) \(HX=X(X'X)^{-1}X'X=X\), \(MX=X-HX=0\), and \(HM=H-H^2=0=MH\).
\end{proof}

\begin{example}\label{r6-exm-centering}
For \(X=\one_n\) (\(p=1\)), \(X'X=\one'\one=n\), so \(H=\frac1n\one\one'=\bar J_n\) and
\(M=I_n-\bar J_n=C_n\), the centering matrix of Definition~\ref{r1-def-ones}: \(Hy=\bar y\one\) and \(My=y-\bar y\one\). The centered vectors form
\(W_0=\{x:\one'x=0\}=\Span\{\one\}^\perp\) (Theorem~\ref{r6-thm-complement}(a)), so the decomposition
\(y=\bar y\one+C_ny\) along \(\R^n=\Span\{\one\}\oplus W_0\) of Example~\ref{r3-exm-constant-centered} is
the orthogonal one, and \(\bar y\one\) is the constant vector closest to \(y\)
(Theorem~\ref{r6-thm-best-approximation}). Likewise, for any nonzero \(a\in\R^n\) the projection of \(y\)
onto \(\Span\{a\}\) is \(a(a'a)^{-1}a'y=\frac{a'y}{a'a}\,a\).
\end{example}

\begin{example}\label{r6-exm-hat}
Let \(X=[\one_4\ X_0]\) be the design matrix of Example~\ref{r1-exm-data} (\(n=4\), \(p=3\)), and let
\(y=(9,3,5,3)'\) be a response measured on the four units. With \((X'X)^{-1}\) of
Example~\ref{r2-exm-xtx-inverse},
\[
(X'X)^{-1}X'=\frac14\begin{pmatrix}6&-1&-2\\-1&1&0\\-2&0&1\end{pmatrix}\begin{pmatrix}1&1&1&1\\2&0&2&0\\3&3&1&1\end{pmatrix}
=\frac14\begin{pmatrix}-2&0&2&4\\1&-1&1&-1\\1&1&-1&-1\end{pmatrix},
\]
\[
H=X(X'X)^{-1}X'=\begin{pmatrix}1&2&3\\1&0&3\\1&2&1\\1&0&1\end{pmatrix}\frac14\begin{pmatrix}-2&0&2&4\\1&-1&1&-1\\1&1&-1&-1\end{pmatrix}
=\frac14\begin{pmatrix}3&1&1&-1\\1&3&-1&1\\1&-1&3&1\\-1&1&1&3\end{pmatrix},
\]
and \(M=I_4-H=\frac14\left(\begin{smallmatrix}1&-1&-1&1\\-1&1&1&-1\\-1&1&1&-1\\1&-1&-1&1\end{smallmatrix}\right)\). Hence
\[
Hy=\frac14\begin{pmatrix}27+3+5-3\\9+9-5+3\\9-3+15+3\\-9+3+5+9\end{pmatrix}=\begin{pmatrix}8\\4\\6\\2\end{pmatrix},\qquad
My=y-Hy=\begin{pmatrix}1\\-1\\-1\\1\end{pmatrix}
\]
Check: \(X'(My)=(1-1-1+1,\ 2-2,\ 3-3-1+1)'=0\), and \((Hy)'(My)=8-4-6+2=0\).
\end{example}

\subsection*{Idempotent matrices: rank equals trace}

\begin{theorem}\label{r6-thm-idempotent}
Let \(A\in\R^{n\times n}\) be idempotent, \(A^2=A\).
\begin{enumerate}
\item \(Av=v\) for every \(v\in\mathcal C(A)\), and \(\R^n=\mathcal N(A)\oplus\mathcal C(A)\): every \(v\in\R^n\) is
\(v=(v-Av)+Av\) with \(v-Av\in\mathcal N(A)\) and \(Av\in\mathcal C(A)\).
\item \(I_n-A\) is idempotent.
\item \(\rank A=\tr A\).
\item If moreover \(A'=A\), then \(Av\) is the orthogonal projection of \(v\) onto \(\mathcal C(A)\), for every
\(v\in\R^n\).
\end{enumerate}
\end{theorem}
\begin{proof}
(a)
\begin{itemize}
\item If \(v=Aw\), then \(Av=A^2w=Aw=v\). For every \(v\), \(A(v-Av)=Av-A^2v=0\).
\item If \(v\in\mathcal N(A)\cap\mathcal C(A)\), then \(v=Av=0\). So the sum \(\mathcal N(A)+\mathcal C(A)\) is direct
(Definition~\ref{r3-def-sum}), and it is \(\R^n\) by the displayed decomposition.
\end{itemize}
(b) \((I-A)^2=I-2A+A^2=I-A\).

(c) If \(A=0\), then \(\rank A=0=\tr A\). Let \(r=\rank A\geqslant1\).
\begin{itemize}
\item Take a full-rank factorization \(A=BC\), with \(B\in\R^{n\times r}\), \(C\in\R^{r\times n}\), and
\(L\in\R^{r\times n}\), \(K\in\R^{n\times r}\) such that \(LB=I_r\) and \(CK=I_r\)
(Theorem~\ref{thm:r4-full-rank-factorization}(a),(c)).
\item \(A^2=A\) reads \(BCBC=BC\). Multiplying on the left by \(L\) and on the right by \(K\),
\(L(BCBC)K=(LB)(CB)(CK)=CB\) and \(L(BC)K=(LB)(CK)=I_r\); so \(CB=I_r\).
\item Hence \(\tr A=\tr(BC)=\tr(CB)=\tr I_r=r\) (Theorem~\ref{r1-thm-trace}(b)).
\end{itemize}
(d) By Corollary~\ref{r6-cor-four-subspaces}, \(\mathcal N(A)=\mathcal C(A')^\perp=\mathcal C(A)^\perp\). So in (a),
\(Av\in\mathcal C(A)\) and \(v-Av\in\mathcal C(A)^\perp\): by Definition~\ref{r6-def-projection}, \(Av\) is the
orthogonal projection of \(v\) onto \(\mathcal C(A)\).
\end{proof}

\begin{corollary}\label{r6-cor-trace-hat}
Let \(X\in\R^{n\times p}\) have full column rank, with hat matrix \(H\) and residual maker \(M\). Then
\(\rank H=\tr H=p=\rank X\) and \(\rank M=\tr M=n-p\). In particular \(\rank C_n=\tr C_n=n-1\).
\end{corollary}
\begin{proof}
By Theorem~\ref{r1-thm-trace}(a),(b),
\(\tr H=\tr\bigl(X(X'X)^{-1}X'\bigr)=\tr\bigl((X'X)^{-1}X'X\bigr)=\tr I_p=p\) and \(\tr M=\tr I_n-\tr H=n-p\).
Both matrices are idempotent (Theorem~\ref{r6-thm-hat}(b)), so their ranks equal their traces
(Theorem~\ref{r6-thm-idempotent}(c)). For \(X=\one_n\), \(M=C_n\) (Example~\ref{r6-exm-centering}).
\end{proof}
In Example~\ref{r6-exm-hat}, \(\tr H=\frac14(3+3+3+3)=3\) and \(\tr M=\frac14(1+1+1+1)=1\); indeed every
column of \(M\) is a multiple of \((1,-1,-1,1)'\), so \(\rank M=1\).

\subsection*{Least squares}

In the \textbf{linear regression model}, the response vector \(y=(y_1,\dots,y_n)'\) holds \(n\)
observations, and \(x_{ij}\) is the value of the \(j\)th explanatory variable on observation \(i\)
(\(1\leqslant j\leqslant k\)):
\[
y_i=\beta_0+\beta_1x_{i1}+\dots+\beta_kx_{ik}+\varepsilon_i\quad(1\leqslant i\leqslant n),\qquad\text{that is,}\qquad
y=X\beta+\varepsilon
\]
The \textbf{design matrix} \(X\in\R^{n\times(k+1)}\) has \(i\)th row \(x_i'=(1,x_{i1},\dots,x_{ik})\), so its
first column is \(\one_n\), the column of the \textbf{intercept} \(\beta_0\); \(\beta=(\beta_0,\dots,\beta_k)'\), and
\(\varepsilon=(\varepsilon_1,\dots,\varepsilon_n)'\) is the vector of errors. A \textbf{least-squares
estimate} of \(\beta\) is a vector \(\hat\beta\in\R^{k+1}\) that minimizes
\[
S(\beta)=\lVert y-X\beta\rVert^2=\sum_{i=1}^n(y_i-x_i'\beta)^2\qquad(\beta\in\R^{k+1}),
\]
and the \(k+1\) equations \(X'X\beta=X'y\) are the \textbf{normal equations}. The results of this section
and the next hold for every \(X\in\R^{n\times(k+1)}\), with or without the column \(\one_n\), unless an
intercept is assumed.

\begin{theorem}\label{r6-thm-normal-equations}
Let \(X\in\R^{n\times(k+1)}\) and \(y\in\R^n\).
\begin{enumerate}
\item The normal equations \(X'X\beta=X'y\) (unknown \(\beta\in\R^{k+1}\)) are consistent.
\item \(X'X\) is invertible if and only if \(X\) has full column rank, \(\rank X=k+1\). Then the normal
equations have exactly one solution, \(\hat\beta=(X'X)^{-1}X'y\).
\item If \(\rank X<k+1\) and \(\hat\beta\) is one solution, the solutions are the vectors \(\hat\beta+z\) with
\(z\in\mathcal N(X)\); there are infinitely many, and all of them give the same vector \(X\hat\beta\).
\end{enumerate}
\end{theorem}
\begin{proof}
(a) \(X'y\in\mathcal C(X')=\mathcal C(X'X)\) (Theorem~\ref{thm:r4-rank-ata}(c)), so the system is consistent
(Corollary~\ref{cor:r4-ax-b}(a)).

(b) \(X'X\in\R^{(k+1)\times(k+1)}\) is invertible if and only if \(\rank(X'X)=k+1\)
(Theorem~\ref{thm:r4-rank-rules}(d)), that is, \(\rank X=k+1\) (Theorem~\ref{thm:r4-rank-ata}(b)). Then
multiplying \(X'X\beta=X'y\) by \((X'X)^{-1}\) gives \(\beta=(X'X)^{-1}X'y\), and this vector is a solution.

(c) By Theorem~\ref{thm:r2-solution-structure}(a), the solutions are \(\hat\beta+z\) with
\(z\in\mathcal N(X'X)=\mathcal N(X)\) (Theorem~\ref{thm:r4-rank-ata}(a)), and
\(X(\hat\beta+z)=X\hat\beta+Xz=X\hat\beta\). As \(\rank X<k+1\), some \(z\ne0\) lies in \(\mathcal N(X)\)
(Theorem~\ref{thm:r4-rank-rules}(b)), and the solutions \(\hat\beta+tz\) (\(t\in\R\)) are pairwise distinct,
since \(t_1z=t_2z\) forces \(t_1=t_2\).
\end{proof}

\begin{theorem}\label{r6-thm-least-squares}
Let \(X\in\R^{n\times(k+1)}\) and \(y\in\R^n\).
\begin{enumerate}
\item \(\hat\beta\) is a least-squares estimate if and only if \(X'X\hat\beta=X'y\). For every solution
\(\hat\beta\) of the normal equations,
\[
S(\beta)=S(\hat\beta)+\lVert X(\hat\beta-\beta)\rVert^2\qquad(\beta\in\R^{k+1})
\]
\item All least-squares estimates give the same vector of \textbf{fitted values} \(\hat y=X\hat\beta\), the
orthogonal projection of \(y\) onto \(\mathcal C(X)\); the vector of \textbf{residuals} \(e=y-\hat y\) lies in
\(\mathcal C(X)^\perp\).
\item If \(\rank X=k+1\), the least-squares estimate is unique, \(\hat\beta=(X'X)^{-1}X'y\), and
\(\hat y=Hy\), \(e=My\).
\end{enumerate}
\end{theorem}
\begin{proof}
(a)
\begin{itemize}
\item Let \(X'X\hat\beta=X'y\). Then \(X'(y-X\hat\beta)=0\), so \(y-X\hat\beta\perp\mathcal C(X)\)
(Theorem~\ref{r6-thm-complement}(a)); as \(X\hat\beta\in\mathcal C(X)\), it is the orthogonal projection of \(y\)
onto \(\mathcal C(X)\) (Definition~\ref{r6-def-projection}).
\item For \(\beta\in\R^{k+1}\), \(X\beta\in\mathcal C(X)\), and Theorem~\ref{r6-thm-best-approximation} gives
\(\lVert y-X\beta\rVert^2=\lVert y-X\hat\beta\rVert^2+\lVert X\hat\beta-X\beta\rVert^2\), the identity. Hence
\(S(\beta)\geqslant S(\hat\beta)\): \(\hat\beta\) is a least-squares estimate.
\item Solutions exist (Theorem~\ref{r6-thm-normal-equations}(a)). If \(\beta^*\) is a least-squares
estimate and \(\hat\beta\) a solution, then \(S(\beta^*)\leqslant S(\hat\beta)\), so the identity gives
\(\lVert X(\hat\beta-\beta^*)\rVert^2=0\), that is, \(X\beta^*=X\hat\beta\); then \(X'X\beta^*=X'X\hat\beta=X'y\).
\end{itemize}
(b) By the last step, every least-squares estimate \(\beta^*\) gives \(X\beta^*=X\hat\beta\), the projection
found in the first step; and \(e=y-X\hat\beta\perp\mathcal C(X)\).

(c) By Theorem~\ref{r6-thm-normal-equations}(b) and (a), \((X'X)^{-1}X'y\) is the only least-squares
estimate, and then \(\hat y=X(X'X)^{-1}X'y=Hy\) and \(e=y-Hy=My\) (Theorem~\ref{r6-thm-hat}).
\end{proof}
So least squares needs no calculus: \(\hat y\) is the projection of \(y\) onto \(\mathcal C(X)\). The formula
\(\hat\beta=(X'X)^{-1}X'y\) needs a design matrix of full column rank, \(\rank X=k+1\); since
\(\rank X\leqslant n\) (Theorem~\ref{thm:r4-rank-rules}(a)), this requires at least \(k+1\) observations.
Then Corollary~\ref{r6-cor-trace-hat} with \(p=k+1\) gives \(\tr H=k+1=\rank X\) and \(\tr M=n-k-1\).

\begin{proposition}\label{r6-prop-residuals}
Let \(\hat\beta\) be a least-squares estimate, \(\hat y=X\hat\beta\) and \(e=y-\hat y\).
\begin{enumerate}
\item \(X'e=0\): the residuals are orthogonal to every column of \(X\).
\item If \(\one_n\in\mathcal C(X)\) (for instance, if the model has an intercept), then
\(\one'e=e_1+\dots+e_n=0\); hence \(\hat y\) and \(y\) have the same mean \(\bar y\).
\item \(\hat y'e=0\).
\end{enumerate}
\end{proposition}
\begin{proof}
(a) This is \(X'(y-X\hat\beta)=0\), the normal equations (Theorem~\ref{r6-thm-least-squares}(a)). (b) If
\(\one=Xc\), then \(\one'e=c'X'e=0\) by (a), and \(\one'\hat y=\one'y-\one'e=\one'y\). (c)
\(\hat y'e=\hat\beta'X'e=0\).
\end{proof}

\begin{example}\label{r6-exm-least-squares}
For \(X\) and \(y\) of Example~\ref{r6-exm-hat}, \(X\) has full column rank, as \(X'X\) is invertible
(Example~\ref{r2-exm-xtx-inverse}, Theorem~\ref{r6-thm-normal-equations}(b)); so the least-squares
estimate is
\[
\hat\beta=(X'X)^{-1}X'y=\frac14\begin{pmatrix}-2&0&2&4\\1&-1&1&-1\\1&1&-1&-1\end{pmatrix}\begin{pmatrix}9\\3\\5\\3\end{pmatrix}
=\frac14\begin{pmatrix}-18+0+10+12\\9-3+5-3\\9+3-5-3\end{pmatrix}=\begin{pmatrix}1\\2\\1\end{pmatrix},
\]
and indeed, with \(X'X\) of Example~\ref{r1-exm-data},
\(X'X\hat\beta=(4+8+8,\ 4+16+8,\ 8+16+20)'=(20,28,44)'\), which is
\(X'y=(9+3+5+3,\ 18+10,\ 27+9+5+3)'\). The fitted values and
residuals are \(\hat y=Hy=(8,4,6,2)'\) and \(e=My=(1,-1,-1,1)'\), and the smallest value of \(S\) is
\(S(\hat\beta)=e'e=4\). As it is positive, no \(\beta\) gives \(X\beta=y\): the system \(X\beta=y\) has no
solution. As Proposition~\ref{r6-prop-residuals} says, \(X'e=0\)
(Example~\ref{r6-exm-hat}), \(\one'e=1-1-1+1=0\) and \(\hat y'e=8-4-6+2=0\). By the identity of
Theorem~\ref{r6-thm-least-squares}(a), for instance \(\beta=(1,2,2)'\) gives
\(S(\beta)=4+\lVert X(0,0,-1)'\rVert^2=4+(9+9+1+1)=24\).
\end{example}

\begin{corollary}\label{r6-cor-simple-regression}
Let \(k=1\) (\textbf{simple regression}): \(y_i=\beta_0+\beta_1x_i+\varepsilon_i\), where \(x_i=x_{i1}\) is a number,
and \(X=[\one_n\ x]\) with \(x=(x_1,\dots,x_n)'\). Put
\[
S_{xx}=\sum_{i=1}^n(x_i-\bar x)^2=\sum_{i=1}^nx_i^2-n\bar x^2,\qquad
S_{xy}=\sum_{i=1}^n(x_i-\bar x)(y_i-\bar y)=\sum_{i=1}^nx_iy_i-n\bar x\bar y
\]
(expand, using \(\sum_ix_i=n\bar x\) and \(\sum_iy_i=n\bar y\), as in Example~\ref{r1-exm-centering}).
\begin{enumerate}
\item \(X\) has full column rank if and only if \(x\) is not constant, that is, \(S_{xx}>0\).
\item If \(S_{xx}>0\), then \(\hat\beta_1=S_{xy}/S_{xx}\) and \(\hat\beta_0=\bar y-\hat\beta_1\bar x\): the fitted line
\(\hat\beta_0+\hat\beta_1x\) passes through the point \((\bar x,\bar y)\).
\end{enumerate}
\end{corollary}
\begin{proof}
(a) As \(\one\ne0\), the columns \(\one,x\) are dependent if and only if \(x=c\one\) for some \(c\)
(Proposition~\ref{r3-prop-dependence}(c)), that is, \(x\) is constant; and \(S_{xx}\), a sum of squares, is \(0\) exactly then.

(b) By (a), \(X\) has full column rank, so \(\hat\beta=(X'X)^{-1}X'y\)
(Theorem~\ref{r6-thm-normal-equations}(b)). Here
\(X'X=\left(\begin{smallmatrix}n&n\bar x\\n\bar x&\sum_ix_i^2\end{smallmatrix}\right)\), with determinant
\(n\sum_ix_i^2-n^2\bar x^2=nS_{xx}\), and \(X'y=(n\bar y,\ \sum_ix_iy_i)'\). By Theorem~\ref{r5d-thm-adjugate}
for \(2\times2\) matrices,
\[
\hat\beta=\frac1{nS_{xx}}\begin{pmatrix}\sum_ix_i^2&-n\bar x\\-n\bar x&n\end{pmatrix}\begin{pmatrix}n\bar y\\\sum_ix_iy_i\end{pmatrix},
\qquad\text{so}\qquad\hat\beta_1=\frac{n\sum_ix_iy_i-n^2\bar x\bar y}{nS_{xx}}=\frac{S_{xy}}{S_{xx}}
\]
The first normal equation, \(n\hat\beta_0+n\bar x\hat\beta_1=n\bar y\), gives \(\hat\beta_0+\hat\beta_1\bar x=\bar y\).
\end{proof}

\subsection*{Rank in regression}

\begin{proposition}\label{r6-prop-collinearity}
Let \(X\in\R^{n\times p}\) with \(p\geqslant2\). Then \(X\) has full column rank if and only if no column of
\(X\) is a linear combination of the other columns. If some column is such a combination
(\textbf{exact collinearity}), then \(Xz=0\) for some \(z\ne0\), the matrix \(X'X\) is singular, and \(\beta\) and
\(\beta+tz\) give the same \(X\beta\) for every \(t\in\R\).
\end{proposition}
\begin{proof}
By Theorem~\ref{thm:r4-rank-rules}(b), \(X\) has full column rank if and only if its columns are
independent, and \(p\geqslant2\) vectors are dependent if and only if one of them is a linear combination
of the others (Proposition~\ref{r3-prop-dependence}(c)). If they are dependent, \(Xz=0\) for some \(z\ne0\)
(Theorem~\ref{thm:r4-rank-rules}(b)), \(X'X\) is singular (Theorem~\ref{r6-thm-normal-equations}(b)
with \(p=k+1\)), and \(X(\beta+tz)=X\beta+tXz=X\beta\).
\end{proof}
For example, if a regression uses the scores \(x_{i1}\), \(x_{i2}\) of two tests and also their total
\(x_{i3}=x_{i1}+x_{i2}\), then column \(4\) of \(X\) is the sum of columns \(2\) and \(3\), and
\(X(0,1,1,-1)'=0\).

\begin{proposition}\label{r6-prop-reparametrization}
Let \(X\in\R^{n\times p}\), \(T\in\R^{p\times q}\) and \(X_1=XT\in\R^{n\times q}\), with \(\rank X_1=\rank X\).
Then \(\mathcal C(X_1)=\mathcal C(X)\), the least-squares fitted values (and residuals) of the designs \(X_1\)
and \(X\) are the same, and if \(\hat\beta_1\) is a least-squares estimate for \(X_1\), then \(T\hat\beta_1\) is one
for \(X\).
\end{proposition}
\begin{proof}
\(\mathcal C(X_1)\subseteq\mathcal C(X)\) (Theorem~\ref{thm:r4-rank-product}(a)), and both have dimension
\(\rank X\), so they are equal (Theorem~\ref{thm:r4-counting}(a)). For either design, the fitted values
are the projection \(\hat y\) of \(y\) onto this subspace (Theorem~\ref{r6-thm-least-squares}(b)). Finally
\(X(T\hat\beta_1)=X_1\hat\beta_1=\hat y\), so \(\lVert y-X(T\hat\beta_1)\rVert^2=\lVert y-\hat y\rVert^2\leqslant\lVert y-X\beta\rVert^2\) for
every \(\beta\) (Theorem~\ref{r6-thm-best-approximation}).
\end{proof}
For instance, adding to the design \(X\) of Example~\ref{r6-exm-least-squares} the total
\(X_{\cdot1}+X_{\cdot2}\) of its two variables gives \(\tilde X=[X\ \ X_{\cdot1}+X_{\cdot2}]\in\R^{4\times4}\), with
\(X=\tilde XT\) for \(T\) the first three columns of \(I_4\), and \(\rank\tilde X=\rank X=3\). So \(\tilde X\) has the
same fitted values \((8,4,6,2)'\), and \(T(1,2,1)'=(1,2,1,0)'\) is a least-squares estimate for \(\tilde X\).

\begin{theorem}\label{r6-thm-oneway}
Let \(n_1,\dots,n_g\geqslant1\) and \(n=n_1+\dots+n_g\). Arrange \(n\) observations in \(g\) groups (first the
\(n_1\) observations of group \(1\), then the \(n_2\) of group \(2\), and so on), and let the
\textbf{indicator} \(d_l\in\R^n\) of group \(l\) have entry \(1\) in the rows of group \(l\) and \(0\) elsewhere.
The model \(y_i=\mu+\alpha_l+\varepsilon_i\) for an observation \(i\) of group \(l\) reads \(y=X\beta+\varepsilon\) with
\(\beta=(\mu,\alpha_1,\dots,\alpha_g)'\) and
\[
X=[\one_n\ d_1\ \cdots\ d_g]\in\R^{n\times(g+1)}
\]
Let \(X_1=[d_1\ \cdots\ d_g]\) (intercept dropped) and \(X_2=[\one_n\ d_2\ \cdots\ d_g]\) (indicator \(d_1\)
dropped), both in \(\R^{n\times g}\).
\begin{enumerate}
\item \(X_1\) has full column rank, and \(X=X_1S\) with \(S=[\one_g\ I_g]\in\R^{g\times(g+1)}\) of full row rank.
\item \(\rank X=g\) and \(\mathcal N(X)=\Span\{(-1,1,\dots,1)'\}\): \(X\) does not have full column rank.
\item \(X_1=XT_1\) and \(X_2=XT_2\), where \(T_1,T_2\in\R^{(g+1)\times g}\) are \(I_{g+1}\) with its first,
respectively second, column deleted. Both \(X_1\) and \(X_2\) have full column rank \(g\), and
\(\mathcal C(X_1)=\mathcal C(X_2)=\mathcal C(X)\).
\end{enumerate}
\end{theorem}
\begin{proof}
(a)
\begin{itemize}
\item No observation lies in two groups, so \(d_l'd_m=0\) for \(l\ne m\), and \(d_l'd_l=n_l\). If
\(c_1d_1+\dots+c_gd_g=0\), multiplying by \(d_m'\) gives \(c_mn_m=0\), so \(c_m=0\): \(X_1\) has full column
rank (Theorem~\ref{thm:r4-rank-rules}(b)).
\item Every observation lies in exactly one group, so \(d_1+\dots+d_g=\one_n\), that is,
\(X_1\one_g=\one_n\). By the block product (Theorem~\ref{r5b-thm-block}),
\(X_1S=[X_1\one_g\ \ X_1I_g]=[\one_n\ d_1\ \cdots\ d_g]=X\).
\item The columns of \(S\) include those of \(I_g\), so \(\mathcal C(S)=\R^g\): \(S\) has full row rank
(Theorem~\ref{thm:r4-rank-rules}(c)).
\end{itemize}
(b) By (a) and Theorem~\ref{thm:r4-rank-product}(b), \(\rank X=\rank S=g\) and \(\mathcal N(X)=\mathcal N(S)\). Now
\(S(c_0,c_1,\dots,c_g)'=c_0\one_g+(c_1,\dots,c_g)'\) is \(0\) if and only if \(c_l=-c_0\) for every \(l\), that is,
\((c_0,c_1,\dots,c_g)'=-c_0(-1,1,\dots,1)'\).

(c)
\begin{itemize}
\item The columns of \(XT_1\) are the columns of \(X\) other than the first (Theorem~\ref{r1-thm-views}(a),(b)), so
\(XT_1=X_1\);
likewise \(XT_2=X_2\).
\item \(X_1\) has full column rank by (a). For \(X_2\): if \(c_1\one_n+c_2d_2+\dots+c_gd_g=0\), then
\(\one_n=d_1+\dots+d_g\) turns this into \(c_1d_1+(c_1+c_2)d_2+\dots+(c_1+c_g)d_g=0\), and (a) gives
\(c_1=0\) and then \(c_2=\dots=c_g=0\).
\item So \(\rank X_1=\rank X_2=g=\rank X\), and Proposition~\ref{r6-prop-reparametrization} gives
\(\mathcal C(X_1)=\mathcal C(X_2)=\mathcal C(X)\).\qedhere
\end{itemize}
\end{proof}
Using the intercept together with all \(g\) indicators (the \textbf{dummy variable trap}) thus makes \(X'X\)
singular (Theorem~\ref{r6-thm-normal-equations}(b)); dropping the intercept or one indicator gives a
design of full column rank with the same column space, hence the same fitted values
(Proposition~\ref{r6-prop-reparametrization}).

\begin{example}\label{r6-exm-oneway}
Six plants of three types have weights \(5,7,6\) (type \(1\)), \(3,5\) (type \(2\)) and \(2\) (type \(3\)). With
\(y=(5,7,6,3,5,2)'\), the model of Theorem~\ref{r6-thm-oneway} (\(g=3\), \(n_1=3\), \(n_2=2\), \(n_3=1\)) has
\[
X=\begin{pmatrix}1&1&0&0\\1&1&0&0\\1&1&0&0\\1&0&1&0\\1&0&1&0\\1&0&0&1\end{pmatrix},\qquad
X'X=\begin{pmatrix}6&3&2&1\\3&3&0&0\\2&0&2&0\\1&0&0&1\end{pmatrix},\qquad
X'y=\begin{pmatrix}28\\18\\8\\2\end{pmatrix}
\]
The \((j,l)\) entry of \(X'X\) is \(\col_j(X)'\col_l(X)\) (Theorems~\ref{r1-thm-views}(d)
and~\ref{r1-thm-transpose}), and \(X'y\) holds the total and the
three group totals.
\begin{itemize}
\item \(\rank X=3\) and \(\mathcal N(X)=\Span\{(-1,1,1,1)'\}\) (Theorem~\ref{r6-thm-oneway}), so \(X'X\) is
singular (Theorem~\ref{r6-thm-normal-equations}(b)).
\item \(\hat\beta=(0,6,4,2)'\) solves the normal equations: \(X'X\hat\beta=(18+8+2,\ 18,\ 8,\ 2)'=X'y\). By
Theorem~\ref{r6-thm-normal-equations}(c), all solutions are
\[
\begin{pmatrix}0\\6\\4\\2\end{pmatrix}+t\begin{pmatrix}-1\\1\\1\\1\end{pmatrix}
=\begin{pmatrix}-t\\6+t\\4+t\\2+t\end{pmatrix}\qquad(t\in\R),
\]
and all give \(X\hat\beta=(6,6,6,4,4,2)'\): the fitted value \(\hat\mu+\hat\alpha_l\) is the mean of group \(l\)
(\(18/3=6\), \(8/2=4\), \(2/1=2\)).
\item Intercept dropped: \(X_1'X_1=\diag(3,2,1)\) and \(X_1'y=(18,8,2)'\), so
\(\hat\beta=(X_1'X_1)^{-1}X_1'y=(6,4,2)'\), the three group means.
\item Indicator \(d_1\) dropped:
\[
X_2'X_2=\begin{pmatrix}6&2&1\\2&2&0\\1&0&1\end{pmatrix},\qquad
X_2'y=\begin{pmatrix}28\\8\\2\end{pmatrix},\qquad
X_2'X_2\begin{pmatrix}6\\-2\\-4\end{pmatrix}=\begin{pmatrix}36-4-4\\12-4\\6-4\end{pmatrix}=X_2'y
\]
so \(\hat\beta=(6,-2,-4)'\), the unique solution (Theorem~\ref{r6-thm-normal-equations}(b)): the mean \(6\) of
group \(1\) and the differences \(4-6\), \(2-6\) of the other group means from it.
\item In all three designs \(\hat y=(6,6,6,4,4,2)'\), as Proposition~\ref{r6-prop-reparametrization} says;
moreover \(T_1(6,4,2)'=(0,6,4,2)'\) and \(T_2(6,-2,-4)'=(6,0,-2,-4)'\) are the solutions above with \(t=0\)
and \(t=-6\).
\end{itemize}
\end{example}

\begin{proposition}\label{r6-prop-quadratic}
Let \(x=(x_1,\dots,x_n)'\in\R^n\), and let \(X\in\R^{n\times3}\) have \(i\)th row \((1,x_i,x_i^2)\): the design
matrix of the quadratic regression \(y_i=\beta_0+\beta_1x_i+\beta_2x_i^2+\varepsilon_i\) (\(k=2\), with
\(x_{i1}=x_i\) and \(x_{i2}=x_i^2\)). Then \(X\) has full column rank if and only if \(x_1,\dots,x_n\) take at
least three different values.
\end{proposition}
\begin{proof}
\begin{itemize}
\item If every \(x_i\) is one of two numbers \(a,b\) (possibly \(a=b\)), then
\((x_i-a)(x_i-b)=ab-(a+b)x_i+x_i^2=0\) for every \(i\); so \(X(ab,\,-(a+b),\,1)'=0\), and \(X\) does not
have full column rank (Theorem~\ref{thm:r4-rank-rules}(b)).
\item Let \(x_i=a\), \(x_j=b\), \(x_l=c\) be three different values. If \(Xv=0\), rows \(i,j,l\) give \(Mv=0\) for
\[
M=\begin{pmatrix}1&a&a^2\\1&b&b^2\\1&c&c^2\end{pmatrix}
\]
The operations \(R_2\leftarrow R_2-R_1\), \(R_3\leftarrow R_3-R_1\) give the rows \((0,\,b-a,\,b^2-a^2)\) and
\((0,\,c-a,\,c^2-a^2)\); dividing them by \(b-a\ne0\) and \(c-a\ne0\) gives \((0,1,a+b)\) and \((0,1,a+c)\);
then \(R_3\leftarrow R_3-R_2\) gives \((0,0,c-b)\), and dividing by \(c-b\ne0\) gives an echelon form with a
leading \(1\) in every column. So \(Mv=0\) only for \(v=0\) (Theorem~\ref{thm:r4-pivot-test} and
Remark~\ref{rem:r4-echelon}), and hence
\(Xv=0\) only for \(v=0\).\qedhere
\end{itemize}
\end{proof}

\begin{example}\label{r6-exm-quadratic}
(a) For \(x=(0,1,2,3)'\), four different values, \(X\) has full column rank by
Proposition~\ref{r6-prop-quadratic}. Directly:
\[
\begin{pmatrix}1&0&0\\1&1&1\\1&2&4\\1&3&9\end{pmatrix}
\xrightarrow[R_4\leftarrow R_4-R_1]{R_2\leftarrow R_2-R_1,\ R_3\leftarrow R_3-R_1}
\begin{pmatrix}1&0&0\\0&1&1\\0&2&4\\0&3&9\end{pmatrix}
\xrightarrow[R_4\leftarrow R_4-3R_2]{R_3\leftarrow R_3-2R_2}
\begin{pmatrix}1&0&0\\0&1&1\\0&0&2\\0&0&6\end{pmatrix}
\xrightarrow[R_3\leftarrow\frac12R_3]{R_4\leftarrow R_4-3R_3}
\begin{pmatrix}1&0&0\\0&1&1\\0&0&1\\0&0&0\end{pmatrix},
\]
an echelon form with three nonzero rows, so \(\rank X=3\).

(b) If a dose is given at only two levels, \(1\) and \(3\), say \(x=(1,3,3,1)'\), then \(x_i^2=4x_i-3\) for
every \(i\) (\(1=4-3\), \(9=12-3\)). So column \(3\) of \(X\) is \(4\) times column \(2\) minus \(3\) times
column \(1\), and \(X(3,-4,1)'=0\). Columns \(1\) and \(2\), \((1,1,1,1)'\) and \((1,3,3,1)'\), are not multiples
of each other, so \(\rank X=2\). Thus \(\beta\) and \(\beta+t(3,-4,1)'\) give the same \(X\beta\): two dose levels
cannot determine the curvature \(\beta_2\).
\end{example}

\subsection*{Sums of squares and \(R^2\)}

\begin{definition}\label{r6-def-sums-of-squares}
For a least-squares fit, with fitted values \(\hat y\) and residuals \(e=y-\hat y\), the \textbf{total},
\textbf{regression} and \textbf{error sums of squares} are
\[
\mathrm{SST}=\sum_{i=1}^n(y_i-\bar y)^2,\qquad\mathrm{SSR}=\sum_{i=1}^n(\hat y_i-\bar y)^2,\qquad
\mathrm{SSE}=\sum_{i=1}^n(y_i-\hat y_i)^2=e'e
\]
If \(\mathrm{SST}>0\), the \textbf{coefficient of determination} is \(R^2=\mathrm{SSR}/\mathrm{SST}\).
\end{definition}

\begin{theorem}\label{r6-thm-anova}
If \(\one_n\in\mathcal C(X)\) (for instance, if the model has an intercept), then
\(\mathrm{SST}=\mathrm{SSR}+\mathrm{SSE}\), where \(\mathrm{SST}=\lVert C_ny\rVert^2=y'C_ny\) and
\(\mathrm{SSR}=\lVert C_n\hat y\rVert^2\). Hence, if \(\mathrm{SST}>0\), then \(R^2=1-\mathrm{SSE}/\mathrm{SST}\) and
\(0\leqslant R^2\leqslant1\).
\end{theorem}
\begin{proof}
\begin{itemize}
\item For \(v\in\R^n\), \(C_nv=v-\bar v\one\), so \(\lVert C_nv\rVert^2=\sum_i(v_i-\bar v)^2\), and
\(\lVert C_nv\rVert^2=v'C_n'C_nv=v'C_nv\) (Proposition~\ref{r1-prop-centering}). With \(v=y\) this is SST; with
\(v=\hat y\), whose mean is \(\bar y\) (Proposition~\ref{r6-prop-residuals}(b)), it is SSR.
\item \(C_ne=e-\bar e\one=e\), as \(\bar e=0\) (Proposition~\ref{r6-prop-residuals}(b)). So
\(C_ny=C_n\hat y+C_ne=C_n\hat y+e\).
\item \((C_n\hat y)'e=\hat y'C_ne=\hat y'e=0\) (Proposition~\ref{r6-prop-residuals}(c)). By Pythagoras,
\(\lVert C_ny\rVert^2=\lVert C_n\hat y\rVert^2+\lVert e\rVert^2\), that is, \(\mathrm{SST}=\mathrm{SSR}+\mathrm{SSE}\).
\item Dividing by \(\mathrm{SST}>0\) gives \(R^2=1-\mathrm{SSE}/\mathrm{SST}\); and
\(0\leqslant\mathrm{SSR}\leqslant\mathrm{SST}\), as \(\mathrm{SSE}\geqslant0\).\qedhere
\end{itemize}
\end{proof}
In terms of projections: \(y=\bar y\one+C_ny\) along \(\Span\{\one\}\oplus\Span\{\one\}^\perp\)
(Example~\ref{r6-exm-centering}), and \(C_ny=C_n\hat y+e\) splits the centered response into
\(C_n\hat y=\hat y-\bar y\one\in\mathcal C(X)\) and \(e\in\mathcal C(X)^\perp\); \(\mathrm{SST}=\mathrm{SSR}+\mathrm{SSE}\) is
Pythagoras for this split. For the design \(X=\one_n\) alone, \(\hat y=\bar y\one\), so \(\mathrm{SSR}=0\) and
\(\mathrm{SSE}=\mathrm{SST}\); \(R^2\) is the fraction of SST that the other columns of \(X\) explain.

\begin{example}\label{r6-exm-anova}
In Example~\ref{r6-exm-least-squares}, \(\bar y=\frac14(9+3+5+3)=5\), so
\[
C_4y=\begin{pmatrix}4\\-2\\0\\-2\end{pmatrix},\qquad
C_4\hat y=\begin{pmatrix}3\\-1\\1\\-3\end{pmatrix},\qquad
e=\begin{pmatrix}1\\-1\\-1\\1\end{pmatrix},\qquad C_4y=C_4\hat y+e
\]
Hence \(\mathrm{SST}=16+4+0+4=24\), \(\mathrm{SSR}=9+1+1+9=20\) and \(\mathrm{SSE}=1+1+1+1=4\); indeed
\(24=20+4\), and \(R^2=\frac{20}{24}=\frac56\).
\end{example}

\needspace{14\baselineskip}
\subsection*{Exercises and solutions}

\begin{problem}\label{r6-ex-complements}
\begin{enumerate}
\item Find all vectors of \(\R^3\) that are orthogonal to \(x=(2,2,-1)'\).
\item What is the dimension of the subspace of \(\R^6\) consisting of the vectors orthogonal to both
\((1,1,-2,3,4,5)'\) and \((0,0,1,1,0,7)'\)?
\end{enumerate}
\end{problem}
\begin{solution}
(a) \(y=(y_1,y_2,y_3)'\) is orthogonal to \(x\) if and only if \(x'y=2y_1+2y_2-y_3=0\), that is,
\(y_3=2y_1+2y_2\). So these vectors are
\[
y=\begin{pmatrix}y_1\\y_2\\2y_1+2y_2\end{pmatrix}=y_1\begin{pmatrix}1\\0\\2\end{pmatrix}+y_2\begin{pmatrix}0\\1\\2\end{pmatrix}\qquad(y_1,y_2\in\R):
\]
\(\Span\{x\}^\perp=\mathcal N(x')=\Span\{(1,0,2)',(0,1,2)'\}\), of dimension \(2=3-1\)
(Theorem~\ref{r6-thm-complement}).

(b) Let \(a_1,a_2\) be the two vectors and \(A=[a_1\ a_2]\in\R^{6\times2}\). The subspace is
\(\Span\{a_1,a_2\}^\perp=\mathcal N(A')\) (Theorem~\ref{r6-thm-complement}(a)). The vectors are independent:
if \(c_1a_1+c_2a_2=0\), the first entry gives \(c_1=0\), and then \(c_2a_2=0\) with \(a_2\ne0\) gives \(c_2=0\).
So \(\dim\Span\{a_1,a_2\}=2\), and the answer is \(6-2=4\) (Theorem~\ref{r6-thm-complement}(c)).
\end{solution}

\begin{problem}\label{r6-ex-qr-regression}
Let \(X\) be the design matrix of Example~\ref{r6-exm-hat}.
\begin{enumerate}
\item Find the QR decomposition of \(X\) by the Gram--Schmidt process.
\item Show that \(Q'Q=I_3\) but \(QQ'\ne I_4\).
\end{enumerate}
\end{problem}
\begin{solution}
(a) The columns are \(a_1=\one_4\), \(a_2=(2,0,2,0)'\) and \(a_3=(3,3,1,1)'\).
\begin{itemize}
\item \(b_1=a_1\), \(\lVert b_1\rVert^2=4\).
\item \(a_2'b_1=4\), so \(t_{12}=1\) and \(b_2=a_2-b_1=(1,-1,1,-1)'\), with \(\lVert b_2\rVert^2=4\).
\item \(a_3'b_1=8\) and \(a_3'b_2=3-3+1-1=0\), so \(t_{13}=2\), \(t_{23}=0\) and \(b_3=a_3-2b_1=(1,1,-1,-1)'\),
with \(\lVert b_3\rVert^2=4\).
\item All three lengths are \(2\): \(q_j=\frac12b_j\), \(r_{jj}=2\), \(r_{12}=2t_{12}=2\), \(r_{13}=2t_{13}=4\) and
\(r_{23}=0\):
\[
Q=\frac12\begin{pmatrix}1&1&1\\1&-1&1\\1&1&-1\\1&-1&-1\end{pmatrix},\qquad
R=\begin{pmatrix}2&2&4\\0&2&0\\0&0&2\end{pmatrix}
\]
Check: the third column of \(QR\) is \(Q(4,0,2)'=\frac12(4+2,\ 4+2,\ 4-2,\ 4-2)'=(3,3,1,1)'=a_3\).
\end{itemize}
(b) The columns of \(Q\) are orthonormal, so \(Q'Q=I_3\). But
\[
QQ'=\frac14\begin{pmatrix}3&1&1&-1\\1&3&-1&1\\1&-1&3&1\\-1&1&1&3\end{pmatrix}\ne I_4;
\]
for instance, its \((1,1)\) entry is \(\frac14(1+1+1)=\frac34\). The rows of \(Q\) are not orthonormal.
\end{solution}

\begin{problem}\label{r6-ex-projection-drill}
Let \(X=\begin{pmatrix}1&1\\1&0\\0&1\end{pmatrix}\), \(W=\mathcal C(X)\) and \(y=(4,2,5)'\).
\begin{enumerate}
\item Find a basis of \(W^\perp\).
\item Compute the hat matrix \(H\) and the residual maker \(M\). Check that \(H^2=H\), \(HX=X\) and \(MX=0\),
and find \(\tr H\) and \(\tr M\).
\item Write \(y=Hy+My\), the decomposition along \(\R^3=W\oplus W^\perp\), and check that the two parts are
orthogonal.
\end{enumerate}
\end{problem}
\begin{solution}
(a) \(W^\perp=\mathcal N(X')\) (Theorem~\ref{r6-thm-complement}(a)). Row-reduce \(X'\):
\[
\begin{pmatrix}1&1&0\\1&0&1\end{pmatrix}
\xrightarrow{R_2\leftarrow R_2-R_1}
\begin{pmatrix}1&1&0\\0&-1&1\end{pmatrix}
\xrightarrow{R_2\leftarrow-R_2}
\begin{pmatrix}1&1&0\\0&1&-1\end{pmatrix}
\xrightarrow{R_1\leftarrow R_1-R_2}
\begin{pmatrix}1&0&1\\0&1&-1\end{pmatrix}
\]
With \(x_3=t\) free, \(x_1=-t\) and \(x_2=t\); so \((-1,1,1)'\) is a basis of \(W^\perp\), of dimension
\(3-2=1\).

(b) \(X'X=\left(\begin{smallmatrix}2&1\\1&2\end{smallmatrix}\right)\), with determinant \(3\), so
\((X'X)^{-1}=\frac13\left(\begin{smallmatrix}2&-1\\-1&2\end{smallmatrix}\right)\) (Theorem~\ref{r5d-thm-adjugate}),
\[
(X'X)^{-1}X'=\frac13\begin{pmatrix}2&-1\\-1&2\end{pmatrix}\begin{pmatrix}1&1&0\\1&0&1\end{pmatrix}
=\frac13\begin{pmatrix}1&2&-1\\1&-1&2\end{pmatrix},
\]
\[
H=X(X'X)^{-1}X'=\begin{pmatrix}1&1\\1&0\\0&1\end{pmatrix}\frac13\begin{pmatrix}1&2&-1\\1&-1&2\end{pmatrix}
=\frac13\begin{pmatrix}2&1&1\\1&2&-1\\1&-1&2\end{pmatrix},
\]
and \(M=I_3-H=\frac13\left(\begin{smallmatrix}1&-1&-1\\-1&1&1\\-1&1&1\end{smallmatrix}\right)\). Then
\[
9H^2=\begin{pmatrix}2&1&1\\1&2&-1\\1&-1&2\end{pmatrix}^2=\begin{pmatrix}6&3&3\\3&6&-3\\3&-3&6\end{pmatrix}=9H,\qquad
HX=\frac13\begin{pmatrix}3&3\\3&0\\0&3\end{pmatrix}=X,
\]
so \(H^2=H\), \(HX=X\) and \(MX=X-HX=0\). Finally \(\tr H=\frac13(2+2+2)=2\) and
\(\tr M=\frac13(1+1+1)=1\), the ranks of \(H\) and \(M\) (Corollary~\ref{r6-cor-trace-hat}).

(c) \(Hy=\frac13(8+2+5,\ 4+4-5,\ 4-2+10)'=(5,1,4)'\), which is \(X(1,4)'\), and
\(My=y-Hy=(-1,1,1)'\), a vector of \(W^\perp\) by (a). They are orthogonal: \((Hy)'(My)=-5+1+4=0\).
\end{solution}

\begin{problem}\label{r6-ex-idempotent-drill}
Let \(A=\begin{pmatrix}2&-1&-1\\1&0&-1\\1&-1&0\end{pmatrix}\) and \(v=(2,1,0)'\).
\begin{enumerate}
\item Show that \(A\) is idempotent but not symmetric.
\item Find \(\rank A\), \(\tr A\), \(\rank(I_3-A)\) and \(\tr(I_3-A)\).
\item Decompose \(v\) along \(\R^3=\mathcal N(A)\oplus\mathcal C(A)\). Are the two parts orthogonal?
\end{enumerate}
\end{problem}
\begin{solution}
(a) Row by row,
\(\row_1(A)A=(4-1-1,\ -2+0+1,\ -2+1+0)=(2,-1,-1)\), \(\row_2(A)A=(2+0-1,\ -1+0+1,\ -1+0+0)=(1,0,-1)\) and
\(\row_3(A)A=(2-1+0,\ -1+0+0,\ -1+1+0)=(1,-1,0)\); so \(A^2=A\). It is not symmetric: \(a_{12}=-1\ne1=a_{21}\).

(b) Row-reduce \(A\):
\[
A\xrightarrow{R_1\leftrightarrow R_2}
\begin{pmatrix}1&0&-1\\2&-1&-1\\1&-1&0\end{pmatrix}
\xrightarrow[R_3\leftarrow R_3-R_1]{R_2\leftarrow R_2-2R_1}
\begin{pmatrix}1&0&-1\\0&-1&1\\0&-1&1\end{pmatrix}
\xrightarrow[R_2\leftarrow-R_2]{R_3\leftarrow R_3-R_2}
\begin{pmatrix}1&0&-1\\0&1&-1\\0&0&0\end{pmatrix},
\]
an echelon form with two nonzero rows, so \(\rank A=2=2+0+0=\tr A\). Next \(I_3-A=\left(\begin{smallmatrix}-1&1&1\\-1&1&1\\-1&1&1\end{smallmatrix}\right)\) has three
equal nonzero rows, so \(\rank(I_3-A)=1=-1+1+1=\tr(I_3-A)\), as Theorem~\ref{r6-thm-idempotent}(b),(c) say.

(c) \(Av=(4-1-0,\ 2+0-0,\ 2-1+0)'=(3,2,1)'\in\mathcal C(A)\) and \(v-Av=(-1,-1,-1)'\), which lies in \(\mathcal N(A)\):
\(A(-1,-1,-1)'=(-2+1+1,\ -1+0+1,\ -1+1+0)'=0\). So \(v=(-1,-1,-1)'+(3,2,1)'\)
(Theorem~\ref{r6-thm-idempotent}(a)). The parts are not orthogonal: \((v-Av)'(Av)=-3-2-1=-6\). (By
Theorem~\ref{r6-thm-idempotent}(d), they would be if \(A\) were symmetric.)
\end{solution}

\begin{problem}\label{r6-ex-reflection}
Show that \(I_n-2A\) is an orthogonal matrix for every symmetric idempotent \(A\in\R^{n\times n}\). In
particular, if \(h\in\R^n\) and \(h'h=1\), then \(I_n-2hh'\) is orthogonal.
\end{problem}
\begin{solution}
Using \(A'=A\) and \(A^2=A\), \((I-2A)'(I-2A)=(I-2A)^2=I-4A+4A^2=I\); so \(I-2A\) is orthogonal
(Definition~\ref{r6-def-orthonormal}). The matrix \(hh'\) is symmetric and idempotent: \((hh')'=hh'\) and
\((hh')(hh')=h(h'h)h'=hh'\).
\end{solution}

\begin{problem}\label{r6-ex-simple-by-hand}
Five observations of a variable \(x\) and a response \(y\) are
\[
\begin{array}{c|ccccc}x_i&1&2&3&4&5\\\hline y_i&3&2&4&4&7\end{array}
\]
Fit the simple regression \(y_i=\beta_0+\beta_1x_i+\varepsilon_i\).
\begin{enumerate}
\item Compute \(X'X\), \(X'y\), \((X'X)^{-1}\) and \(\hat\beta\).
\item Compute \(\hat y\) and \(e\), and verify that \(\sum_ie_i=0\) and \(\sum_ix_ie_i=0\).
\item Compute SSE.
\item Compute SST and SSR, verify \(\mathrm{SST}=\mathrm{SSR}+\mathrm{SSE}\), and compute \(R^2\).
\end{enumerate}
\end{problem}
\begin{solution}
(a) \(X=[\one_5\ x]\), \(\sum_ix_i=15\), \(\sum_ix_i^2=55\), \(\sum_iy_i=20\) and \(\sum_ix_iy_i=3+4+12+16+35=70\):
\[
X'X=\begin{pmatrix}5&15\\15&55\end{pmatrix},\qquad X'y=\begin{pmatrix}20\\70\end{pmatrix},\qquad
(X'X)^{-1}=\frac1{50}\begin{pmatrix}55&-15\\-15&5\end{pmatrix}=\frac1{10}\begin{pmatrix}11&-3\\-3&1\end{pmatrix}
\]
(\(\det X'X=275-225=50\); Theorem~\ref{r5d-thm-adjugate}), and
\(\hat\beta=\frac1{10}(220-210,\ -60+70)'=(1,1)'\). Check with Corollary~\ref{r6-cor-simple-regression}:
\(\bar x=3\), \(\bar y=4\), \(S_{xx}=55-45=10\) and \(S_{xy}=70-60=10\), so \(\hat\beta_1=1\) and
\(\hat\beta_0=4-3=1\).

(b) \(\hat y_i=1+x_i\): \(\hat y=(2,3,4,5,6)'\) and \(e=y-\hat y=(1,-1,0,-1,1)'\). Then
\(\sum_ie_i=1-1+0-1+1=0\) and \(\sum_ix_ie_i=1-2+0-4+5=0\) (Proposition~\ref{r6-prop-residuals}).

(c) \(\mathrm{SSE}=1+1+0+1+1=4\).

(d) \(\mathrm{SST}=\sum_i(y_i-4)^2=1+4+0+0+9=14\) and \(\mathrm{SSR}=\sum_i(\hat y_i-4)^2=4+1+0+1+4=10\);
indeed \(10+4=14\) (Theorem~\ref{r6-thm-anova}), and \(R^2=\frac{10}{14}=\frac57\).
\end{solution}

\begin{problem}\label{r6-ex-rank-deficient}
Add to the design matrix \(X\) of Example~\ref{r6-exm-hat} the sum of its two variables as a fourth column:
\(\tilde X=[\one_4\ X_{\cdot1}\ X_{\cdot2}\ X_{\cdot1}+X_{\cdot2}]\in\R^{4\times4}\), and let \(y=(9,3,5,3)'\).
\begin{enumerate}
\item Find \(\rank\tilde X\) and a basis of \(\mathcal N(\tilde X)\).
\item Solve the normal equations \(\tilde X'\tilde X\beta=\tilde X'y\) by row reduction.
\item Show that all least-squares estimates give the fitted values \(\hat y=(8,4,6,2)'\) of
Example~\ref{r6-exm-least-squares}, and find SSE.
\end{enumerate}
\end{problem}
\begin{solution}
(a) The columns are \(\one_4\), \((2,0,2,0)'\), \((3,3,1,1)'\) and \((5,3,3,1)'\). The fourth is the sum of the
second and the third, so \(\tilde X(0,1,1,-1)'=0\); the first three are the columns of \(X\), which are
independent (Example~\ref{r6-exm-least-squares}). So \(\rank\tilde X=3\), \(\dim\mathcal N(\tilde X)=4-3=1\)
(Theorem~\ref{thm:r4-rank-nullity}), and \((0,1,1,-1)'\) is a basis of \(\mathcal N(\tilde X)\).

(b) The entries of \(\tilde X'\tilde X\) and \(\tilde X'y\) are inner products of columns (Theorems~\ref{r1-thm-views}(d)
and~\ref{r1-thm-transpose}):
\[
\left(\begin{array}{cccc|c}4&4&8&12&20\\4&8&8&16&28\\8&8&20&28&44\\12&16&28&44&72\end{array}\right)
\xrightarrow[R_4\leftarrow R_4-3R_1]{R_2\leftarrow R_2-R_1,\ R_3\leftarrow R_3-2R_1}
\left(\begin{array}{cccc|c}4&4&8&12&20\\0&4&0&4&8\\0&0&4&4&4\\0&4&4&8&12\end{array}\right)
\]
\[
\xrightarrow{R_4\leftarrow R_4-R_2}
\left(\begin{array}{cccc|c}4&4&8&12&20\\0&4&0&4&8\\0&0&4&4&4\\0&0&4&4&4\end{array}\right)
\xrightarrow[R_1\leftarrow\frac14R_1,\ R_2\leftarrow\frac14R_2,\ R_3\leftarrow\frac14R_3]{R_4\leftarrow R_4-R_3}
\left(\begin{array}{cccc|c}1&1&2&3&5\\0&1&0&1&2\\0&0&1&1&1\\0&0&0&0&0\end{array}\right)
\]
The last matrix is an echelon form with pivots in columns \(1,2,3\) and none in the last column, so the
system is consistent and \(\beta_3=t\) is free (Theorem~\ref{thm:r2-consistency}). Back substitution:
\(\beta_2=1-t\), \(\beta_1=2-t\) and \(\beta_0=5-\beta_1-2\beta_2-3\beta_3=5-(2-t)-2(1-t)-3t=1\). So the solutions are
\[
\beta=\begin{pmatrix}1\\2-t\\1-t\\t\end{pmatrix}=\begin{pmatrix}1\\2\\1\\0\end{pmatrix}+t\begin{pmatrix}0\\-1\\-1\\1\end{pmatrix}\qquad(t\in\R),
\]
a particular solution plus \(\mathcal N(\tilde X'\tilde X)=\mathcal N(\tilde X)\), as
Theorem~\ref{r6-thm-normal-equations}(c) says.

(c) These solutions are the least-squares estimates (Theorem~\ref{r6-thm-least-squares}(a)), and each
gives \(\tilde X\beta=X(1,2,1)'+t\tilde X(0,-1,-1,1)'=(8,4,6,2)'\). So \(\hat y\) and \(e=(1,-1,-1,1)'\) are those
of Example~\ref{r6-exm-least-squares}, and \(\mathrm{SSE}=4\).
\end{solution}

\begin{problem}\label{r6-ex-hat-identities}
Let \(X\in\R^{n\times(k+1)}\) have full column rank, with hat matrix \(H\) and residual maker \(M\).
\begin{enumerate}
\item In the model \(y=X\beta+\varepsilon\), show that \(e=M\varepsilon\) and \(\mathrm{SSE}=\varepsilon'M\varepsilon\).
\item Let \(\one_n\in\mathcal C(X)\). Show that \(H\one_n=\one_n\), \(H\bar J_n=\bar J_nH=\bar J_n\) and
\(MC_n=C_nM=M\).
\item Under (b), show that \(H-\bar J_n\) is symmetric and idempotent of rank \(k\), that
\((H-\bar J_n)M=0\), and that \(\mathrm{SSR}=y'(H-\bar J_n)y\).
\end{enumerate}
\end{problem}
\begin{solution}
(a) \(MX=0\) (Theorem~\ref{r6-thm-hat}(c)), so \(e=My=M(X\beta+\varepsilon)=MX\beta+M\varepsilon=M\varepsilon\)
(Theorem~\ref{r6-thm-least-squares}(c)). As \(M\) is symmetric and idempotent,
\(\mathrm{SSE}=e'e=\varepsilon'M'M\varepsilon=\varepsilon'M\varepsilon\).

(b)
\begin{itemize}
\item \(\one_n=Xc\) for some \(c\), so \(H\one_n=HXc=Xc=\one_n\) (Theorem~\ref{r6-thm-hat}(c)).
\item \(H\bar J_n=\frac1nH\one\one'=\frac1n\one\one'=\bar J_n\), and, as \(H\) and \(\bar J_n\) are symmetric,
\(\bar J_nH=(H\bar J_n)'=\bar J_n\).
\item \(M\bar J_n=\bar J_n-H\bar J_n=0\), so \(MC_n=M(I_n-\bar J_n)=M\), and \(C_nM=(MC_n)'=M\).
\end{itemize}
(c)
\begin{itemize}
\item \(H-\bar J_n\) is symmetric, and by (b) and \(\bar J_n^2=\bar J_n\) (Proposition~\ref{r1-prop-centering}),
\((H-\bar J_n)^2=H^2-H\bar J_n-\bar J_nH+\bar J_n^2=H-\bar J_n-\bar J_n+\bar J_n=H-\bar J_n\).
\item Its rank is its trace (Theorem~\ref{r6-thm-idempotent}(c)):
\(\tr H-\tr\bar J_n=(k+1)-1=k\) (Corollary~\ref{r6-cor-trace-hat}; the diagonal entries of \(\bar J_n\) are \(\frac1n\)).
\item \((H-\bar J_n)M=HM-\bar J_nM=0-(M\bar J_n)'=0\), by Theorem~\ref{r6-thm-hat}(c) and (b).
\item \(\hat y-\bar y\one=Hy-\bar J_ny=(H-\bar J_n)y\) (Example~\ref{r6-exm-centering}), so
\(\mathrm{SSR}=\lVert(H-\bar J_n)y\rVert^2=y'(H-\bar J_n)'(H-\bar J_n)y=y'(H-\bar J_n)y\).
\end{itemize}
\end{solution}

\begin{problem}\label{r6-ex-correlation}
Let \(x,y\in\R^n\) with \(S_{xx}>0\) and \(S_{yy}=\sum_i(y_i-\bar y)^2>0\) (notation of
Corollary~\ref{r6-cor-simple-regression}), and put \(r=S_{xy}/\sqrt{S_{xx}S_{yy}}\).
\begin{enumerate}
\item Show that \(\lvert r\rvert\leqslant1\).
\item In the simple regression of \(y\) on \(x\), show that \(\mathrm{SSR}=\hat\beta_1^2S_{xx}\) and \(R^2=r^2\).
\end{enumerate}
\end{problem}
\begin{solution}
(a) Let \(u=C_nx=x-\bar x\one\) and \(v=C_ny=y-\bar y\one\) (Example~\ref{r1-exm-centering}). Then
\(S_{xy}=u'v\), \(S_{xx}=\lVert u\rVert^2\) and \(S_{yy}=\lVert v\rVert^2\). By the Cauchy--Schwarz inequality
(Theorem~\ref{r6-thm-cauchy-schwarz}(a)), \(\lvert S_{xy}\rvert\leqslant\sqrt{S_{xx}}\sqrt{S_{yy}}\), that is,
\(\lvert r\rvert\leqslant1\).

(b) By Corollary~\ref{r6-cor-simple-regression}(b),
\(\hat y_i-\bar y=\hat\beta_0+\hat\beta_1x_i-(\hat\beta_0+\hat\beta_1\bar x)=\hat\beta_1(x_i-\bar x)\), so
\[
\mathrm{SSR}=\sum_i(\hat y_i-\bar y)^2=\hat\beta_1^2S_{xx}=\frac{S_{xy}^2}{S_{xx}^2}\,S_{xx}=\frac{S_{xy}^2}{S_{xx}},\qquad
R^2=\frac{\mathrm{SSR}}{\mathrm{SST}}=\frac{S_{xy}^2}{S_{xx}S_{yy}}=r^2,
\]
since \(\mathrm{SST}=S_{yy}\).
\end{solution}

\begin{problem}\label{r6-ex-r2-monotone}
Let \(X\in\R^{n\times p}\) with \(\one_n\in\mathcal C(X)\), let \(x\in\R^n\), let \(X_+=[X\ x]\) be \(X\) with the column
\(x\) added, and let \(y\in\R^n\) with \(\mathrm{SST}>0\). Let SSE and \(R^2\) belong to the least-squares fit with
\(X\), and \(\mathrm{SSE}_+\) and \(R^2_+\) to the fit with \(X_+\). Show that \(\mathrm{SSE}_+\leqslant\mathrm{SSE}\) and
\(R^2_+\geqslant R^2\): adding a column cannot decrease \(R^2\).
\end{problem}
\begin{solution}
\begin{itemize}
\item Every column of \(X\) is a column of \(X_+\), so \(\mathcal C(X)\subseteq\mathcal C(X_+)\), and both contain \(\one_n\).
\item The fitted values \(\hat y_+\) of \(X_+\) form the vector of \(\mathcal C(X_+)\) closest to \(y\)
(Theorems~\ref{r6-thm-least-squares}(b) and~\ref{r6-thm-best-approximation}), and the fitted values
\(\hat y\) of \(X\) lie in \(\mathcal C(X)\subseteq\mathcal C(X_+)\). Hence
\(\mathrm{SSE}_+=\lVert y-\hat y_+\rVert^2\leqslant\lVert y-\hat y\rVert^2=\mathrm{SSE}\).
\item SST does not depend on the design, so by Theorem~\ref{r6-thm-anova},
\(R^2_+=1-\mathrm{SSE}_+/\mathrm{SST}\geqslant1-\mathrm{SSE}/\mathrm{SST}=R^2\).
\end{itemize}
\end{solution}

\begin{problem}\label{r6-ex-four-regions}
Wages \(y_i\) of \(n\) workers are regressed on years of schooling \(x_i\), an intercept, and the indicators
\(d_1,d_2,d_3,d_4\in\R^n\) of all four regions of a country (each worker lives in exactly one region):
\(X=[\one_n\ x\ d_1\ d_2\ d_3\ d_4]\in\R^{n\times6}\). (a) Show that \(X\) does not have full column rank, so
that \(\hat\beta=(X'X)^{-1}X'y\) is not defined. (b) Let \(X_2=[\one_n\ x\ d_2\ d_3\ d_4]\) (indicator \(d_1\)
dropped). Show that \(\mathcal C(X_2)=\mathcal C(X)\), so that \(\rank X_2=\rank X\).
\end{problem}
\begin{solution}
(a) Each worker lives in exactly one region, so \(d_1+d_2+d_3+d_4=\one_n\), that is,
\(X(-1,0,1,1,1,1)'=0\). So \(X\) does not have full column rank (Theorem~\ref{thm:r4-rank-rules}(b)), and
\(X'X\) is singular (Theorem~\ref{r6-thm-normal-equations}(b)): the dummy variable trap.

(b) Every column of \(X_2\) is a column of \(X\), so \(\mathcal C(X_2)\subseteq\mathcal C(X)\). Conversely
\(d_1=\one_n-d_2-d_3-d_4\in\mathcal C(X_2)\), so every column of \(X\) lies in the subspace \(\mathcal C(X_2)\), and
therefore \(\mathcal C(X)\subseteq\mathcal C(X_2)\) (Theorem~\ref{r3-thm-span}). Hence \(\mathcal C(X_2)=\mathcal C(X)\), and
\(\rank X_2=\rank X\); by Proposition~\ref{r6-prop-reparametrization} (with \(X_2=XT\), \(T\) the columns
\(1,2,4,5,6\) of \(I_6\)), the two designs give the same fitted values.
\end{solution}

\begin{problem}\label{r6-ex-two-way}
Six plants of three varieties are grown under two fertilizer treatments, one plant for each
combination. The model \(y_{ij}=\mu+\alpha_i+\gamma_j+\varepsilon_{ij}\) (variety \(i\), treatment \(j\)) has
\(\beta=(\mu,\alpha_1,\alpha_2,\alpha_3,\gamma_1,\gamma_2)'\) and, with the observations ordered
\(y_{11},y_{12},y_{21},y_{22},y_{31},y_{32}\),
\[
X=\begin{pmatrix}1&1&0&0&1&0\\1&1&0&0&0&1\\1&0&1&0&1&0\\1&0&1&0&0&1\\1&0&0&1&1&0\\1&0&0&1&0&1\end{pmatrix}
\]
Show that \(\rank X=4\), and find a basis of \(\mathcal N(X)\).
\end{problem}
\begin{solution}
Let \(c_1,\dots,c_6\) be the columns of \(X\).
\begin{itemize}
\item Each plant has one variety and one treatment, so \(c_2+c_3+c_4=c_1\) and \(c_5+c_6=c_1\); that is,
\(z_1=(-1,1,1,1,0,0)'\) and \(z_2=(-1,0,0,0,1,1)'\) lie in \(\mathcal N(X)\).
\item \(c_2,c_3,c_4,c_5\) are independent: if \(a_2c_2+a_3c_3+a_4c_4+a_5c_5=0\), rows \(2\), \(4\), \(6\) give
\(a_2=a_3=a_4=0\), and then row \(1\) gives \(a_5=0\).
\item \(c_1=c_2+c_3+c_4\) and \(c_6=c_1-c_5\) are combinations of \(c_2,c_3,c_4,c_5\), so these four columns
span \(\mathcal C(X)\) and form a basis of it: \(\rank X=4\).
\item \(\dim\mathcal N(X)=6-4=2\) (Theorem~\ref{thm:r4-rank-nullity}), and \(z_1,z_2\) are independent (entries
\(2\) and \(5\) of \(t_1z_1+t_2z_2\) are \(t_1\) and \(t_2\)); so they form a basis of \(\mathcal N(X)\)
(Theorem~\ref{thm:r4-counting}(d)).
\end{itemize}
\end{solution}
\end{document}
